\chapter{Brave-New Lie Integration and Comparison}
\label{ch:brave-new-lie-integration}

\section{Purpose, continuity, and the native Lie III problem}
\label{sec:bnli-purpose}

The preceding three chapters have completed the nonlinear, differentiated, and
operation-theoretic stages of brave-new Lie theory.  The nonlinear theory
associates to a brave-new formal Lie groupoid
\[
    \mathfrak G_\bullet\rightrightarrows X
\]
its effective formal quotient
\[
    q_{\mathfrak G}:X\twoheadrightarrow Q_{\mathfrak G},
\]
its canonical nonlinear anchor
\[
    \phi_{\mathfrak G}:\mathfrak G_\bullet\longrightarrow R_{X/S,\bullet},
\]
formal isotropy, adjoint transport, disk--jet calculus, Cartan descent, and the
nonlinear formal-isotropy commutator.  Spectral differentiation associates to
the same groupoid the global formal leaf
\[
    \mathcal L^{\mathrm{fm,sp}}_{\mathfrak G}
    \in \operatorname{FMP}^{\mathrm{sp}}_{X/S}
\]
and its spectral partition-Lie algebroid
\[
    \mathbb A^{\pi,E_\infty}_{\mathfrak G}
    \in \operatorname{LieAlgd}^{\pi,E_\infty}_{X/S},
\]
while retaining the intrinsic stable arrow, actual isotropy, and Hopf-derived
primitive constructions as separate objects.  The Atiyah--Spencer chapter then
constructs the geometric Cech and Cartan cobars, normalized Spencer objects,
tangent representations and tangent enveloping monads, multi-filtered Artinian
operation systems, and the integral equivalence
\[
    \Xi^{\mathrm{Art}}_{\mathfrak G}:
    \mathbb A^{\pi,E_\infty}_{\mathfrak G}
    \xrightarrow{\sim}
    \mathbb A^{\mathrm{Sp,Art}}_{\mathfrak G}.
\]

Formal integration on the spectral Artinian side is already constructed.
On every coherent affine--etale pair
\[
    (U,V)=(\operatorname{Spec}B,\operatorname{Spec}R)
\]
the differentiation chapter constructs an inverse functor
\[
    \operatorname{Int}^{\pi,E_\infty}_{B/R}:
    \operatorname{LieAlgd}^{\pi,E_\infty}_{B/R}
    \xrightarrow{\sim}
    \operatorname{FMP}^{\mathrm{sp}}_{B/R},
\]
and proves its coherent one-variable and two-variable spectral-etale
transport.  These affine inverse functors therefore assemble by the same
coherent coefficient descent to a global inverse
\[
    \operatorname{Int}^{\mathrm{fm}}_{X/S}:
    \operatorname{LieAlgd}^{\pi,E_\infty}_{X/S}
    \xrightarrow{\sim}
    \operatorname{FMP}^{\mathrm{sp}}_{X/S}.
\]
The remaining problem is geometric: starting from the resulting global
spectral formal-moduli problem, construct a native fixed-object formal groupoid
in the big spectral Nisnevich $\infty$-topos and compare its nonlinear geometry
with the deformation-theoretic materialization.

\paragraph{The native materialization problem.}
For
\[
    \mathfrak a\in
    \operatorname{LieAlgd}^{\pi,E_\infty}_{X/S},
\]
construct functorially a native formal Lie groupoid
\[
    \operatorname{Int}^{\mathrm{gpd}}_{X/S}(\mathfrak a)
    \rightrightarrows X
\]
whose spectral differentiation is $\mathfrak a$.  The construction must not
choose a prorepresenting complete algebra, an Artinian square-zero chain, a
nilpotent filtration, an affine presentation, a connection, a splitting, a
smooth atlas, or a global smooth integration.

The construction below has two geometric realizations which are deliberately
not identified.  The first is the deformation-local spectral formal-moduli
stack
\[
    \operatorname{Mat}^{\mathrm{sp}}_{X/S}(F)
\]
attached to a global formal-moduli problem $F$.  The second is a native
fixed-object quotient
\[
    X\twoheadrightarrow Q_F^{\mathrm{nat}}
\]
built from effective images of Artinian cells after inducing every local cell
to the common object space $X$.  The bridge between them is a theorem,
\[
    L_S^{\mathrm{Def}}Q_F^{\mathrm{nat}}
    \xrightarrow{\sim}
    \operatorname{Mat}^{\mathrm{sp}}_{X/S}(F),
\]
not an identification by definition.

There is a second, logically independent question.  If $\mathfrak G$ is an
arbitrary native formal Lie groupoid, its controller determines the canonical
reconstructed groupoid
\[
    \mathfrak G^{\mathrm{Lie}}
    :=
    \operatorname{Int}^{\mathrm{gpd}}
    \operatorname{Diff}^{\pi,E_\infty}(\mathfrak G).
\]
There need not be a canonical direct raw morphism
$\mathfrak G\to\mathfrak G^{\mathrm{Lie}}$.  What is constructed
unconditionally is a two-sided comparison through the formalized relation of
$\mathfrak G$.  Moreover, the failure of native integration to be known fully
faithful is measured by explicit Artinian-cell mapping discrepancies.  No
vanishing condition on those discrepancies is promoted to an input category
or additional admissibility condition.

\paragraph{Non-conflation convention.}
The following objects remain distinct unless a theorem identifies them:
\begin{enumerate}
    \item a spectral formal-moduli problem
    $F\in\operatorname{FMP}^{\mathrm{sp}}_{X/S}$;
    \item its spectral formal-moduli-stack materialization;
    \item the minimal affine native Artinian-cell quotient on a chosen affine
    coefficient chart;
    \item the global fixed-object native Artinian-cell quotient obtained by
    inducing all local cells to $X$;
    \item the Cech relation of either native quotient;
    \item the corresponding native formal Lie groupoid;
    \item the deformation formalization of the native quotient;
    \item the spectral partition-Lie controller;
    \item actual formal isotropy;
    \item the homotopy anchor fibre;
    \item the independent Hopf-derived-primitives spectral Lie algebra;
    \item the native Artinian-cell mapping discrepancies.
\end{enumerate}

\paragraph{No hidden integration hypotheses.}
No prorepresenting algebra, chosen nilpotent filtration, Artinian presentation,
finite square-zero chain, pronilpotence condition, materializability condition,
Artinian-density condition, native-rigidity condition, deformation-completeness
condition, admissibility structure, connection, splitting, monodromy choice,
smooth atlas, source-simply-connected integration, representation-recognition
structure, or Hopf presentation is part of the input.  Nilpotent factorizations and coherent
approximations occur only inside proofs and are not retained.  Spectral
pro-reduction, left-exact deformation localization, finite square-zero
generation, two-stage materialization, formal-completion transition maps,
Artinian Cech effectivity, effective images, global Artinian Grothendieck
constructions, induced native cells, native quotient colimits, Cech nerves,
comparison maps, and discrepancy spaces are functorial outputs.

\paragraph{Coefficient-domain convention.}
The spectral formal-moduli and spectral partition-Lie categories are the
locally coherent global coefficient categories constructed in the
differentiation chapter.  Fixed-object native constructions take place in the
big spectral Nisnevich $\infty$-topos.  Coherent spectral-etale transport is
used on the formal-moduli and materialized formal-stack sides with the
formal-completion variance constructed below.  The native integration functor
itself is asserted functorially for fixed $X/S$.  No raw changing-object etale
exchange theorem for native groupoids is asserted: Chapter~13 already shows
that fixed-object formality is not invariant under arbitrary atlas
replacement.  Rationality enters only in the dg, PBW, CE, and classical BCH
comparison statements.  The spectral pro-reduction construction below is made
directly from $\pi_0$ and Eilenberg--Mac Lane spectral rings and does not use
an integral spectral-to-animated equivalence.

\paragraph{Variance notation.}
Two coefficient-change symbols occur below and have different sources.  For a
morphism of coherent affine--etale coefficient pairs,
\(f_{\mathrm{FMP},!}\) denotes the cocartesian transport functor in the
Grothendieck coefficient system constructed in the differentiation chapter.
Later, for an actual change of the global coefficient pair, the symbol
\(f_{\mathrm{FMP}}^*\) denotes restriction along the already constructed base
and object variables.  These are different variances and are not silently
identified or treated as inverse functors.

\paragraph{Chapter-level output.}
The chapter constructs:
\begin{enumerate}
    \item spectral formal integration $\operatorname{Int}^{\mathrm{fm}}$, assembled
    from the constructed affine inverse functors and their coherent
    spectral-etale transport;
    \item the spectral pro-reduction endofunctor
    $\mathsf P^{\mathrm{sp}}$, its left-exact idempotent localization theorem,
    and the intrinsic infinitesimal prestack
    $U^{\mathrm{inf,sp}}=\mathsf P^{\mathrm{sp}}U$;
    \item nil-centred deformation-local targets and spectrally laft
    formal-moduli stacks, together with reduced-centre detection on the laft
    locus;
    \item left exactness of the global spectral deformation reflector
    $L_S^{\mathrm{Def}}$;
    \item finite square-zero generation without a finite-presentation input
    and spectral Artinian Cech effectivity in the deformation-local topos;
    \item bounded coherent approximation of unrestricted square-zero
    coefficients, filtered representative categories, the correctly varianced
    relative obstruction map for backward Artinian edges, filtered contractible
    coherent-lift categories,
    finite coherent localization resolutions in the groupoid completion, unrestricted
    square-zero excision, relative Postnikov reconstruction, explicit Artinian
    factorization categories, the Cartesian-colimit fibre lemma, relative
    split-tangent continuity, and the first nilpotent left Kan extension with
    its recognition theorem;
    \item the bounded-test model for convergent spectral prestacks, filtered
    representative categories, the pointed filtered-fibre and pointed
    pro-reduction fibre formulae, simultaneous compatible representatives for
    finite squares, bounded pro-pushouts and pointed pro-pushouts, Postnikov
    comparison for nilpotent pullbacks, the second-stage
    dotted-lift materialization on unrestricted
    nilpotent neighbourhoods, pro-infinitesimal completion of arbitrary
    centred targets, and the theorem that ambient deformation locality together with spectral laftness
    implies nil-laftness on bounded nilpotent diagrams;
    \item spectral Nisnevich descent for laft nil-isomorphic deformation
    prestacks, proved through the raw split square-zero tangent spectrum and
    its stabilization extension before pro-coherent approximation;
    \item affine presentation-free materialization, spectral-etale
    cocartesianity of centred Artinian restriction, and inverse equivalence
    between Artinian formal-moduli problems and spectral formal-moduli stacks;
    \item reduced-centred completion stability, native relative reduced targets
    $W_{U/V}\to W_{X/S}$, representable open reduced-chart morphisms on finite
    affine open bases, the formal-completion Cartesian-section category,
    finite-subobject descent, recovery of every coherent affine-etale pair
    from a finite reduced-open basis, formal-completion rather than raw
    target-base-change transition maps, and global materialization by ordinary
    effective Cech/basis descent in the left-exact deformation-local subtopos,
    with spectral laftness and nil-centredness recovered from the local affine
    materializations;
    \item native Artinian effective-image cells and their formal Cech
    groupoids;
    \item minimal affine native Artinian-cell quotients and their local mapping
    discrepancy systems;
    \item the global Artinian Grothendieck category, the functorial induced
    fixed-object native cells, and the global native quotient
    $X\twoheadrightarrow Q_F^{\mathrm{nat}}$;
    \item the canonical deformation-formalization comparison
    \[
        L_S^{\mathrm{Def}}Q_F^{\mathrm{nat}}
        \xrightarrow{\sim}
        \operatorname{Mat}^{\mathrm{sp}}_{X/S}(F);
    \]
    \item the native groupoid integration functor
    $\operatorname{Int}^{\mathrm{gpd}}_{X/S}$, formal-leaf recovery, and
    \[
        \operatorname{Diff}^{\pi,E_\infty}
        \operatorname{Int}^{\mathrm{gpd}}
        \simeq \operatorname{id};
    \]
    \item local and global native-cell mapping discrepancies, without assuming
    their vanishing;
    \item for an arbitrary formal groupoid, the canonical two-sided Lie
    comparison span through its formalized relation, together with the
    reconstructed-span compatibility theorem;
    \item quotient, relation, stable, isotropy, adjoint, support-marked
    Spencer, and native mapping discrepancy systems;
    \item invariance of the Artinian Spencer/partition-Lie operation system
    under reconstructed integration;
    \item finite-limit stability of the relative formal-stack slice, stable
    formal representation materialization, and a canonically typed,
    contravariantly functorial native Cartan realization;
    \item the unconditional actual-isotropy-to-anchor-fibre comparison;
    \item native realization of the formalized adjoint tangent representation,
    together with the separate nonlinear stable
    formalized-adjoint-to-stable-actual-isotropy comparison;
    \item one-object integration and the parallel Hopf-derived-primitives
    output, without manufacturing an unproved Hopf/partition comparison map;
    \item materialized formal-spectrum construction, comparison with independently
    constructed adic formal spectra, rational Artinian CE materialization, and
    the correctly based-looped classical BCH comparison on its separate theorem domain;
    \item comparison spans for already constructed nonlinear actions, gauge
    groupoids, ordinary smooth Lie algebroids, and pre-existing smooth
    integrations;
    \item fixed-object functoriality, deformation-formalized zero recovery, and
    categorical finite-limit comparison after deformation formalization.
\end{enumerate}

The logical progression is
\[
\boxed{
\begin{aligned}
\text{spectral partition-Lie algebroid}
&\longrightarrow
\text{spectral formal-moduli problem}
\\
&\longrightarrow
\text{spectral formal-moduli stack}
\\
&\longrightarrow
\text{global fixed-object native Artinian-cell quotient}
\\
&\longrightarrow
\text{native formal Cech groupoid}
\\
&\longrightarrow
\text{differentiation recovery and comparison discrepancies}.
\end{aligned}}
\]

\section{Formal integration is already complete}
\label{sec:bnli-formal-integration}

Retain the affine differentiation equivalences and their coherent
spectral-etale transport from the differentiation chapter.  For every coherent
affine-etale pair
\[
    (U,V)=(\operatorname{Spec}B,\operatorname{Spec}R)
\]
write
\[
    \operatorname{Diff}^{\pi,E_\infty}_{B/R}:
    \operatorname{FMP}^{\mathrm{sp}}_{B/R}
    \xrightarrow{\sim}
    \operatorname{LieAlgd}^{\pi,E_\infty}_{B/R}
\]
for the affine equivalence and
\[
    \operatorname{Int}^{\pi,E_\infty}_{B/R}:
    \operatorname{LieAlgd}^{\pi,E_\infty}_{B/R}
    \xrightarrow{\sim}
    \operatorname{FMP}^{\mathrm{sp}}_{B/R}
\]
for the inverse functor already constructed there by spectral
cotangent--square-zero duality and cellular recognition.

\begin{definition}[Global spectral formal integration]
\label{def:bnli-formal-integration}
Let
\[
    \mathfrak a\in\operatorname{LieAlgd}^{\pi,E_\infty}_{X/S}.
\]
For every coherent affine-etale pair $(U,V)$ define
\[
    \left(
      \operatorname{Int}^{\mathrm{fm}}_{X/S}(\mathfrak a)
    \right)_{U/V}
    :=
    \operatorname{Int}^{\pi,E_\infty}_{B/R}
    \left(
      \mathfrak a|_{U/V}
    \right).
\]
The coherent one-variable and two-variable spectral-etale transport of the
affine inverse functors supplies the cocartesian transition equivalences on
morphisms of coherent affine-etale pairs.  Spectral-etale descent therefore
defines a global functor
\[
\boxed{
    \operatorname{Int}^{\mathrm{fm}}_{X/S}:
    \operatorname{LieAlgd}^{\pi,E_\infty}_{X/S}
    \longrightarrow
    \operatorname{FMP}^{\mathrm{sp}}_{X/S}.}
\]
For $\mathfrak a$ write
\[
    \mathcal F_{\mathfrak a}
    :=
    \operatorname{Int}^{\mathrm{fm}}_{X/S}(\mathfrak a).
\]
\end{definition}

\begin{theorem}[Spectral formal Lie III equivalence]
\label{thm:bnli-formal-lie-three}
The preceding functor is the inverse of global spectral differentiation.
There are canonical coherent equivalences
\[
\boxed{
    \operatorname{Diff}^{\pi,E_\infty}_{X/S}
    \operatorname{Int}^{\mathrm{fm}}_{X/S}
    \xrightarrow{\sim}
    \operatorname{id}}
\]
and
\[
\boxed{
    \operatorname{Int}^{\mathrm{fm}}_{X/S}
    \operatorname{Diff}^{\pi,E_\infty}_{X/S}
    \xrightarrow{\sim}
    \operatorname{id}.}
\]
\end{theorem}

\begin{proof}
On every coherent affine-etale pair the two composites are the unit and
counit equivalences of the already constructed mutually inverse affine
functors
\[
    \operatorname{Diff}^{\pi,E_\infty}_{B/R}
    \quad\text{and}\quad
    \operatorname{Int}^{\pi,E_\infty}_{B/R}.
\]
The differentiation chapter proves that these affine equivalences commute
with the coherent one-variable and two-variable spectral-etale transport.
Consequently their units and counits form morphisms of the corresponding
cocartesian coefficient sections.  Spectral-etale descent glues them to the
displayed global natural equivalences.  No new choice of an abstract inverse
equivalence is made in the present chapter.
\end{proof}

\section{Spectral formal-moduli stacks under a centre}
\label{sec:bnli-defcent}

The target of geometric materialization must retain two logically distinct
properties: deformation locality and equality of the entire nilpotent support
of the centre.  Reduced-classical equality detects the latter on the
spectrally laft locus, but is too weak for the larger deformation-local target
category on which the materialization adjunction is defined.  We therefore
construct spectral pro-reduction first and recover the reduced-classical
formulation as a consequence.  We also prove here the left exactness of the
global deformation reflector.  That geometric input will replace any separate
bar-length convergence argument in the Artinian Cech descent theorem.

\subsection{Spectral convergence and local almost finite presentation}

\begin{definition}[Spectrally convergent and locally almost finitely presented maps]
\label{def:bnli-spectral-laft}
Let $f:Y\to Z$ be a morphism of spectral prestacks over $S$.  It is
\emph{spectrally convergent} if, for every connective spectral affine test
$T=\operatorname{Spec}C\to Z$, the canonical square
\[
\begin{tikzcd}
Y(C) \arrow[r] \arrow[d] & \displaystyle\lim_n Y(\tau_{\le n}C) \arrow[d]\\
Z(C) \arrow[r] & \displaystyle\lim_n Z(\tau_{\le n}C)
\end{tikzcd}
\]
is Cartesian.

It is \emph{spectrally locally almost finitely presented}, abbreviated
\emph{spectrally laft}, if it is spectrally convergent and, for every $n\ge0$
and every filtered diagram $C_\alpha$ of $n$-truncated connective spectral
algebras over an affine test of $Z$, with colimit $C$, the square
\[
\begin{tikzcd}
\operatorname*{colim}_\alpha Y(C_\alpha) \arrow[r] \arrow[d] &
Y(C) \arrow[d]\\
\operatorname*{colim}_\alpha Z(C_\alpha) \arrow[r] & Z(C)
\end{tikzcd}
\]
is Cartesian.
\end{definition}

\begin{remark}
Spectral laft is a property of a morphism.  No compact presentation, affine
atlas, or chosen filtered system is retained.
\end{remark}

\subsection{Spectral pro-reduction and nil-centred targets}

\begin{definition}[Spectral pro-reduction of an affine]
\label{def:bnli-prored-affine}
Let $C$ be a connective commutative spectral ring.  Let
$\mathcal I_{\mathrm{nil}}(C)$ be the filtered poset of nilpotent ideals
$I\subseteq\pi_0(C)$, ordered by inclusion.  Thus
\[
    I\leq J \quad\Longleftrightarrow\quad I\subseteq J,
\]
and the transition morphism is the quotient map
\[
    H(\pi_0C/I)\longrightarrow H(\pi_0C/J).
\]
The poset is filtered because $I+J$ is nilpotent whenever $I$ and $J$ are
nilpotent.  Define the formal pro-reduction of $C$ by the ind-object
\[
\boxed{
    C^{\mathrm{prored}}
    :=
    ``\operatorname*{colim}_{I\in\mathcal I_{\mathrm{nil}}(C)}''
    H(\pi_0C/I).}
\]
For a spectral prestack $Y$ put
\[
\boxed{
    Y(C^{\mathrm{prored}})
    :=
    \operatorname*{colim}_{I\in\mathcal I_{\mathrm{nil}}(C)}
    Y\!\left(H(\pi_0C/I)\right).}
\]
Geometrically the affine spectra of these quotients form the opposite
pro-system; the displayed algebraic colimit notation records the covariant
quotient maps.
\end{definition}

\begin{lemma}[Functoriality of spectral pro-reduction]
\label{lem:bnli-prored-functoriality}
A morphism $f:C\to D$ of connective commutative spectral rings induces
canonically a morphism of ind-rings
\[
    f^{\mathrm{prored}}:
    C^{\mathrm{prored}}\longrightarrow D^{\mathrm{prored}}.
\]
These maps preserve identities and composition.  Equivalently, affine
pro-reduction is contravariantly functorial as a pro-affine construction.
\end{lemma}

\begin{proof}
For a nilpotent ideal $I\subseteq\pi_0(C)$, let
\[
    J_I:=(\pi_0(f)(I))\pi_0(D)\subseteq\pi_0(D).
\]
If $I^m=0$, then $J_I^m=0$, so $J_I$ is nilpotent.  The map $\pi_0(f)$
induces
\[
    \pi_0(C)/I\longrightarrow\pi_0(D)/J_I
\]
and hence
\[
    H(\pi_0(C)/I)\longrightarrow H(\pi_0(D)/J_I).
\]
If $I\subseteq I'$, then $J_I\subseteq J_{I'}$, and the resulting square of
quotient maps commutes.  These maps define a morphism of the filtered
ind-systems.  For a composite $C\to D\to E$, the ideal generated by the image
of $I$ in $\pi_0(E)$ agrees with the ideal generated after the two successive
steps, so identities and composition agree coherently in the ind-category.
\end{proof}

\begin{construction}[The spectral pro-reduction endofunctor]
\label{con:bnli-prored-endofunctor}
For a spectral prestack $Y$ define
\[
\boxed{
    (\mathsf P^{\mathrm{sp}}Y)(C)
    :=
    Y(C^{\mathrm{prored}})
    =
    \operatorname*{colim}_{I\in\mathcal I_{\mathrm{nil}}(C)}
    Y\!\left(H(\pi_0C/I)\right).}
\]
Functoriality of Lemma~\ref{lem:bnli-prored-functoriality} makes
$\mathsf P^{\mathrm{sp}}$ an endofunctor of spectral prestacks.  The canonical
maps
\[
    C\longrightarrow H\pi_0(C)\longrightarrow H(\pi_0(C)/I)
\]
induce a natural coaugmentation
\[
\boxed{
    \eta_Y^{\mathrm{prored}}:
    Y\longrightarrow \mathsf P^{\mathrm{sp}}Y.}
\]
\end{construction}

\begin{theorem}[Spectral pro-reduction is a left-exact localization]
\label{thm:bnli-prored-localization}
The endofunctor $\mathsf P^{\mathrm{sp}}$ preserves colimits and finite
limits, and the natural maps
\[
    \mathsf P^{\mathrm{sp}}Y
    \longrightarrow
    (\mathsf P^{\mathrm{sp}})^2Y
\]
induced by the coaugmentation are equivalences.  Hence the full subcategory
of \emph{pro-reduced-local} prestacks
\[
    \operatorname{PStk}^{\mathrm{sp,prored}}
    :=
    \{Y:\eta_Y^{\mathrm{prored}}\text{ is an equivalence}\}
\]
is reflective, with reflector $\mathsf P^{\mathrm{sp}}$.  In particular, for
every pro-reduced-local $W$ there is a canonical equivalence
\[
\boxed{
    \operatorname{Map}(\mathsf P^{\mathrm{sp}}Y,W)
    \xrightarrow{\sim}
    \operatorname{Map}(Y,W).}
\]
\end{theorem}

\begin{proof}
Colimits of prestacks are computed objectwise, so colimit preservation follows
by interchanging the two colimits.  Let $K$ be a finite indexing
$\infty$-category and $\{Y_k\}_{k\in K}$ a finite diagram.  Filtered colimits
of spaces commute with finite limits, hence
\[
\begin{aligned}
    \mathsf P^{\mathrm{sp}}\!\left(\lim_{k\in K}Y_k\right)(C)
    &\simeq
    \operatorname*{colim}_{I}
    \lim_{k\in K}Y_k(H(\pi_0C/I))\\
    &\simeq
    \lim_{k\in K}
    \operatorname*{colim}_{I}
    Y_k(H(\pi_0C/I))\\
    &\simeq
    \lim_{k\in K}
    \mathsf P^{\mathrm{sp}}Y_k(C).
\end{aligned}
\]
Thus $\mathsf P^{\mathrm{sp}}$ is left exact.

We prove idempotence together with the assertion that the actual maps induced
by the coaugmentation are equivalences.  Put $R:=\pi_0C$ and let
\[
    \mathcal N(R)
    :=
    \{I\subseteq R:I\text{ is nilpotent}\}.
\]
Expanding the definition twice gives
\[
    (\mathsf P^{\mathrm{sp}})^2Y(C)
    \simeq
    \operatorname*{colim}_{I\in\mathcal N(R)}
    \operatorname*{colim}_{J\in\mathcal N(R/I)}
    Y\!\left(H((R/I)/J)\right).
\]
Let
\[
    \operatorname{Flag}_{\mathrm{nil}}(R)
    :=
    \{(I\subseteq K):I,K\in\mathcal N(R)\}.
\]
A nilpotent ideal $J\subseteq R/I$ has inverse image $K\supseteq I$.  If
$I^m=0$ and $(K/I)^n=0$, then $K^n\subseteq I$ and hence $K^{nm}=0$.
Conversely every flag of nilpotent ideals determines the nilpotent ideal
$K/I$ of $R/I$.  Therefore
\[
\boxed{
    (\mathsf P^{\mathrm{sp}})^2Y(C)
    \simeq
    \operatorname*{colim}_{(I\subseteq K)\in
      \operatorname{Flag}_{\mathrm{nil}}(R)}
    Y(H(R/K)).}
\]

Consider
\[
    p:\operatorname{Flag}_{\mathrm{nil}}(R)\longrightarrow\mathcal N(R),
    \qquad p(I\subseteq K)=K,
\]
and
\[
    s:\mathcal N(R)\longrightarrow\operatorname{Flag}_{\mathrm{nil}}(R),
    \qquad s(K)=(0\subseteq K).
\]
There is an adjunction $s\dashv p$: a morphism
$s(K)\to(I\subseteq K')$ is precisely an inclusion $K\subseteq K'$.
There is also a diagonal section
\[
    d:\mathcal N(R)\longrightarrow\operatorname{Flag}_{\mathrm{nil}}(R),
    \qquad d(K)=(K\subseteq K),
\]
with $p\dashv d$.  Thus $p$ has both a left and a right adjoint.  In
particular, because $p$ is a right adjoint to $s$, it is final and induces a
canonical equivalence
\[
\boxed{
    (\mathsf P^{\mathrm{sp}})^2Y(C)
    \xrightarrow{\sim}
    \mathsf P^{\mathrm{sp}}Y(C).}
\]

It remains to identify the coaugmentation maps with this finality
calculation.  The morphism
\[
    \mathsf P^{\mathrm{sp}}\eta_Y:
    \mathsf P^{\mathrm{sp}}Y
    \longrightarrow
    (\mathsf P^{\mathrm{sp}})^2Y
\]
is induced, at the $I$-stage, by the zero ideal of $R/I$.  Under the flag description this is the diagonal section $d$ defined above.
Since $p\circ d=\operatorname{id}$, the composite of
$\mathsf P^{\mathrm{sp}}\eta_Y(C)$ with the finality equivalence induced by
$p$ is the identity of $\mathsf P^{\mathrm{sp}}Y(C)$.  Thus
$\mathsf P^{\mathrm{sp}}\eta_Y$ is an equivalence.

For the other coaugmentation first observe directly from the definition that
\[
    (\mathsf P^{\mathrm{sp}}Y)(C)
    \xrightarrow{\sim}
    (\mathsf P^{\mathrm{sp}}Y)(H\pi_0C),
\]
since both sides are the same filtered colimit
\[
    \operatorname*{colim}_{K\in\mathcal N(R)}Y(H(R/K)).
\]
Under this identification the outer coaugmentation
\[
    \eta_{\mathsf P^{\mathrm{sp}}Y}:
    \mathsf P^{\mathrm{sp}}Y
    \longrightarrow
    (\mathsf P^{\mathrm{sp}})^2Y
\]
is represented by the outer ideal $I=0$.  In the flag description it is
therefore the section $s(K)=(0\subseteq K)$.  Since
$p\circ s=\operatorname{id}$, composition with the finality equivalence is the
identity.  Hence $\eta_{\mathsf P^{\mathrm{sp}}Y}$ is an equivalence.
The two identifications are natural in $C$ and $Y$, and therefore
\[
\boxed{
    \mathsf P^{\mathrm{sp}}\eta
    \simeq
    \eta\mathsf P^{\mathrm{sp}}
    \simeq
    \operatorname{id}.}
\]

Finally let $W$ be pro-reduced-local.  Precomposition with the coaugmentation
gives
\[
    \operatorname{Map}(\mathsf P^{\mathrm{sp}}Y,W)
    \longrightarrow
    \operatorname{Map}(Y,W).
\]
An inverse sends $f:Y\to W$ to
\[
    \mathsf P^{\mathrm{sp}}Y
    \xrightarrow{\mathsf P^{\mathrm{sp}}f}
    \mathsf P^{\mathrm{sp}}W
    \xrightarrow{\sim}
    W.
\]
The two composites are identities by naturality and the two coaugmentation
idempotence identities just proved.  This is the reflector mapping property.
\end{proof}

\begin{remark}[No animated comparison is used]
The construction of $\mathsf P^{\mathrm{sp}}$ uses only $\pi_0$, nilpotent
ideals, Eilenberg--Mac Lane spectral rings, filtered colimits of spaces, and
the functoriality just proved.  It is therefore intrinsic to the spectral
coefficient theory.  No integral spectral-to-animated Artinian equivalence is
invoked.
\end{remark}

\begin{definition}[Spectral nil-isomorphism]
\label{def:bnli-nil-isomorphism}
A morphism $f:Y\to Z$ of spectral prestacks is a \emph{spectral
nil-isomorphism} if any, hence all, of the following equivalent conditions
hold:
\begin{enumerate}
    \item for every connective affine test $C$, the induced map
    \[
        Y(C^{\mathrm{prored}})
        \xrightarrow{\sim}
        Z(C^{\mathrm{prored}})
    \]
    is an equivalence;
    \item
    \[
        \mathsf P^{\mathrm{sp}}f:
        \mathsf P^{\mathrm{sp}}Y
        \xrightarrow{\sim}
        \mathsf P^{\mathrm{sp}}Z
    \]
    is an equivalence;
    \item for every pro-reduced-local prestack $W$, precomposition induces
    \[
        \operatorname{Map}(Z,W)
        \xrightarrow{\sim}
        \operatorname{Map}(Y,W).
    \]
\end{enumerate}
\end{definition}

\begin{proof}
The equivalence of the first two conditions is the definition of
$\mathsf P^{\mathrm{sp}}$.  The equivalence of the second and third follows
from Theorem~\ref{thm:bnli-prored-localization} by applying the reflector
mapping property to both source and target.
\end{proof}

\begin{lemma}[Nilpotent quotients have the same pro-reduction]
\label{lem:bnli-prored-thickening}
Let $A\to B$ be connective and suppose
$\pi_0(A)\twoheadrightarrow\pi_0(B)$ has nilpotent kernel.  Then the formal
pro-reductions are canonically equivalent,
\[
    A^{\mathrm{prored}}
    \xrightarrow{\sim}
    B^{\mathrm{prored}}.
\]
Equivalently, for every spectral prestack $Y$,
\[
    Y(A^{\mathrm{prored}})
    \xrightarrow{\sim}
    Y(B^{\mathrm{prored}}).
\]
\end{lemma}

\begin{proof}
Put
\[
    K:=\ker(\pi_0A\to\pi_0B).
\]
Nilpotent ideals of $\pi_0(B)$ identify with nilpotent ideals of $\pi_0(A)$
containing $K$.  This subsystem is cofinal: if $I\subseteq\pi_0(A)$ is
nilpotent, then $I+K$ is nilpotent and contains $K$.  Quotienting by an ideal
containing $K$ gives the same discrete quotient whether it is formed from
$\pi_0(A)$ or $\pi_0(B)$.  Hence the two filtered systems of discrete
nilpotent quotients are canonically cofinal and define the same pro-reduction.
Higher homotopy groups do not enter Definition~\ref{def:bnli-prored-affine}.
\end{proof}

\begin{lemma}[Representable nilpotent thickenings are pro-reduction equivalences]
\label{lem:bnli-representable-prored-thickening}
Let $A\to B$ be connective and suppose
$\pi_0(A)\twoheadrightarrow\pi_0(B)$ has nilpotent kernel.  Then the induced
morphism of representable spectral prestacks
\[
    \operatorname{Spec}B\longrightarrow\operatorname{Spec}A
\]
is a spectral nil-isomorphism.  Equivalently,
\[
\boxed{
    \mathsf P^{\mathrm{sp}}(\operatorname{Spec}B)
    \xrightarrow{\sim}
    \mathsf P^{\mathrm{sp}}(\operatorname{Spec}A).}
\]
These equivalences are natural in commutative diagrams of nilpotent
thickenings.  Consequently a functorial family of such centre maps gives a
natural equivalence after applying $\mathsf P^{\mathrm{sp}}$, and its inverse
is canonically coherent in the corresponding functor category.
\end{lemma}

\begin{proof}
Fix a connective spectral test $C$ and put
\[
    R:=\pi_0C,
    \qquad
    K:=\ker(\pi_0A\to\pi_0B),
    \qquad
    K^m=0.
\]
Since every ring $H(R/I)$ occurring in the pro-reduction system is discrete,
the classical truncation adjunction gives
\[
\begin{aligned}
    \bigl(\mathsf P^{\mathrm{sp}}\operatorname{Spec}A\bigr)(C)
    &\simeq
    \operatorname*{colim}_{I\subseteq R,\ I\ {\rm nil}}
    \operatorname{Hom}_{\operatorname{CRing}}
    (\pi_0A,R/I),\\
    \bigl(\mathsf P^{\mathrm{sp}}\operatorname{Spec}B\bigr)(C)
    &\simeq
    \operatorname*{colim}_{I\subseteq R,\ I\ {\rm nil}}
    \operatorname{Hom}_{\operatorname{CRing}}
    (\pi_0B,R/I).
\end{aligned}
\]
The comparison is induced by precomposition with the surjection
$\pi_0A\twoheadrightarrow\pi_0B$.

Let
\[
    \alpha:\pi_0A\longrightarrow R/I
\]
represent a point on the first line.  Let
\[
    L:=\alpha(K)(R/I)\subseteq R/I
\]
be the ideal generated by the image of $K$.  Since $K^m=0$, one has
$L^m=0$.  Let $J\subseteq R$ be the inverse image of $L$.  Then
$J/I=L$ is nilpotent.  Since $I$ is nilpotent as well, $J$ is nilpotent: if
$I^r=0$, then $J^m\subseteq I$ and therefore $J^{mr}=0$.  After the
refinement
\[
    R/I\longrightarrow R/J,
\]
the map $\alpha$ kills $K$ and hence factors uniquely through
$\pi_0B=\pi_0A/K$.  Thus every point of
$\mathsf P^{\mathrm{sp}}(\operatorname{Spec}A)(C)$ is represented by a point
coming from $\mathsf P^{\mathrm{sp}}(\operatorname{Spec}B)(C)$.

Conversely, suppose two maps
\[
    \beta_1,\beta_2:\pi_0B\longrightarrow R/I
\]
have images which agree in the filtered colimit after precomposition with
$\pi_0A\twoheadrightarrow\pi_0B$.  After passing to a common nilpotent
refinement $R/I\to R/J$, their two composites from $\pi_0A$ agree.  The
surjectivity of $\pi_0A\to\pi_0B$ then implies that the refined maps
$\beta_1,\beta_2:\pi_0B\to R/J$ agree.  Hence the comparison is injective on
the filtered colimit as well.

The displayed comparison is therefore an equivalence for every connective
$C$.  Since it is the natural morphism induced by $A\to B$, the equivalence
is natural both in $C$ and in commutative diagrams of nilpotent thickenings.
This proves the spectral nil-isomorphism statement.  A pointwise equivalence
of functors is an equivalence in the functor category, so its inverse carries
all identity, composition, and higher coherence automatically.
\end{proof}

\begin{definition}[Centred deformation-local targets]
\label{def:bnli-defcent-target}
Let
\[
    \operatorname{DefTgt}(X/S)
\]
be the full subcategory of $(\mathcal X_{/S})_{X/}$ spanned by centre maps
\[
    q:X\longrightarrow Y
\]
for which $Y$ is $L_S^{\mathrm{Def}}$-local.  Let
\[
    \operatorname{DefCent}^{\mathrm{nil}}(X/S)
    \subseteq
    \operatorname{DefTgt}(X/S)
\]
be the full subcategory on those centres which are spectral nil-isomorphisms.
Both are properties of the already specified centre map; no further field of
structure is added.
\end{definition}

\begin{lemma}[Connective recognition of pullbacks along a degree-zero surjection]
\label{lem:bnli-connective-pullback-recognition}
Let
\[
    P^{\mathrm{u}}:=A_0\times_{A_{01}}A_1
\]
be a pullback formed in unrestricted commutative ring spectra, where
$A_0,A_1,A_{01}$ are connective.  If one of the maps
\[
    \pi_0(A_i)\longrightarrow\pi_0(A_{01})
\]
is surjective, then $P^{\mathrm{u}}$ is connective.  Consequently it also
computes the pullback in connective commutative ring spectra.
\end{lemma}

\begin{proof}
The forgetful functor from commutative ring spectra to spectra creates limits,
so there is a fibre sequence of underlying spectra
\[
    P^{\mathrm{u}}
    \longrightarrow
    A_0\oplus A_1
    \longrightarrow
    A_{01}.
\]
All three displayed input rings are connective.  Hence the only possible
negative homotopy group of the pullback is
\[
    \pi_{-1}(P^{\mathrm{u}})
    \simeq
    \operatorname{coker}
    \left(
      \pi_0(A_0)\oplus\pi_0(A_1)
      \longrightarrow
      \pi_0(A_{01})
    \right).
\]
The assumed surjectivity of one component makes this cokernel vanish, and all
lower negative homotopy groups vanish by connectivity of the inputs.  Thus
$P^{\mathrm{u}}$ is connective.

The connective-cover theorem for structured ring spectra computes a limit in
connective commutative ring spectra by taking the unrestricted limit and then
its connective cover.  Since the unrestricted pullback is already connective,
the connective cover changes nothing.
\end{proof}

\begin{theorem}[Left exactness of spectral deformation localization]
\label{thm:bnli-def-left-exact}
The accessible deformation localization of the differentiation chapter
\[
    L_S^{\mathrm{Def}}:
    \mathcal X_{/S}\longrightarrow\mathcal X^{\mathrm{Def}}_{/S}
\]
is left exact.  Consequently $\mathcal X^{\mathrm{Def}}_{/S}$ is an
$\infty$-topos, the fully faithful inclusion
$\mathcal X^{\mathrm{Def}}_{/S}\hookrightarrow\mathcal X_{/S}$ preserves all
limits, and $L_S^{\mathrm{Def}}$ preserves finite limits.
\end{theorem}

\begin{proof}
Let $W_S^{\mathrm{Def}}$ be the generating set of the differentiation chapter.
It consists of the Postnikov-convergence maps
\[
    c_A:
    \operatorname*{colim}_{n\geq0}h_{\tau_{\leq n}A}
    \longrightarrow h_A
\]
and the nilpotent-excision maps
\[
    w_\square:
    h_{A_0}\amalg_{h_{A_{01}}}h_{A_1}
    \longrightarrow h_C,
    \qquad
    C\simeq A_0\times_{A_{01}}A_1,
\]
where the displayed pullback is connective and one
$A_i\to A_{01}$ is surjective on $\pi_0$ with nilpotent kernel.  Write $W$
for the strongly saturated class generated by $W_S^{\mathrm{Def}}$.  The
left-exactness criterion for an accessible localization of an $\infty$-topos
reduces the theorem to showing that $W$ is stable under pullback.  It is enough
to prove that every generator is universally a local equivalence.

First consider $w_\square$ and pull it back along an affine map
$h_E\to h_C$.  Put
\[
    E_i:=E\otimes_CA_i,
    \qquad
    E_{01}:=E\otimes_CA_{01}.
\]
The original connective pullback, viewed in unrestricted commutative ring
spectra, has underlying fibre sequence
\[
    C\longrightarrow A_0\oplus A_1\longrightarrow A_{01}.
\]
Extension of scalars $E\otimes_C-$ is exact in the stable category of
$C$-modules.  Therefore
\[
    E\longrightarrow E_0\oplus E_1\longrightarrow E_{01}
\]
is a fibre sequence.  Since limits of unrestricted commutative ring spectra
are created on underlying spectra, the canonical map
\[
    E\longrightarrow E_0\times_{E_{01}}E_1
\]
is an equivalence in unrestricted commutative ring spectra.

Suppose, for definiteness, that
$\pi_0(A_0)\twoheadrightarrow\pi_0(A_{01})$ is the nilpotent leg.  Connectivity
of relative tensor products gives
\[
    \pi_0(E_i)
    \simeq
    \pi_0(E)\otimes_{\pi_0(C)}\pi_0(A_i).
\]
Hence
\[
    \pi_0(E_0)\twoheadrightarrow\pi_0(E_{01})
\]
is surjective, and its kernel is generated by the image of the original
nilpotent kernel, hence is nilpotent.  Lemma~\ref{lem:bnli-connective-pullback-recognition}
now shows that the unrestricted pullback is connective.  It therefore computes
the pullback in connective commutative ring spectra, and the affine pullback
of $w_\square$ is again a nilpotent-excision generator.

Now consider $c_A$ and an affine map $h_E\to h_A$, corresponding to a
connective algebra map $A\to E$.  Its pullback is
\[
    \operatorname*{colim}_{n\geq0}h_{E_n}
    \longrightarrow h_E,
    \qquad
    E_n:=E\otimes_A\tau_{\leq n}A.
\]
Let
\[
    K_n:=\operatorname{fib}(A\to\tau_{\leq n}A).
\]
Then $K_n$ is $(n+1)$-connective, and exactness of extension of scalars gives
\[
    \operatorname{fib}(E\to E_n)
    \simeq E\otimes_AK_n.
\]
Since $E$ is connective, the tensor-connectivity theorem makes the latter
$(n+1)$-connective.  Consequently, for every fixed $d$,
\[
    \tau_{\leq d}E_n\xrightarrow{\sim}\tau_{\leq d}E
    \qquad(n\geq d).
\]
If $Z$ is $L_S^{\mathrm{Def}}$-local, Postnikov convergence and commutation of
limits in spaces give
\[
\begin{aligned}
    \lim_n Z(E_n)
    &\simeq
    \lim_n\lim_d Z(\tau_{\leq d}E_n)\\
    &\simeq
    \lim_d\lim_n Z(\tau_{\leq d}E_n)\\
    &\simeq
    \lim_d Z(\tau_{\leq d}E)\\
    &\simeq Z(E).
\end{aligned}
\]
Thus every affine pullback of $c_A$ is an
$L_S^{\mathrm{Def}}$-equivalence.

Finally let $P$ be an arbitrary object over the target of one of the two kinds
of generators.  By affine density in slices, $P$ is a colimit of affine
representables over that target.  Pullback in an $\infty$-topos preserves
colimits, so the pullback of the generator along $P$ is the colimit in the
arrow category of its affine pullbacks.  Applying the left adjoint
$L_S^{\mathrm{Def}}$ turns this into a colimit of equivalences and hence an
equivalence.  Therefore every generator is universally in $W$.

The class of morphisms all of whose pullbacks lie in $W$ is strongly
saturated.  It contains $W_S^{\mathrm{Def}}$, hence contains the strongly
saturated class $W$ generated by that set.  Therefore $W$ is pullback-stable,
and the accessible-localization left-exactness criterion proves that
$L_S^{\mathrm{Def}}$ is left exact.
\end{proof}

\begin{lemma}[Deformation locality under represented change of ambient base]
\label{lem:bnli-def-local-base-composition}
Let $p:V\to S$ be a represented spectral scheme over $S$, and let
\[
    Y\longrightarrow V
\]
be an $L_V^{\mathrm{Def}}$-local native sheaf.  Then the composite
\[
    Y\longrightarrow V\longrightarrow S
\]
is $L_S^{\mathrm{Def}}$-local.
\end{lemma}

\begin{proof}
The differentiation chapter proves that every represented spectral scheme is
$L_S^{\mathrm{Def}}$-local.  We verify the two defining deformation equations
for $Y$ after forgetting the specified map to $V$.

First let $C$ be a connective affine test over $S$.  The map
\[
    Y(C)\longrightarrow V(C)
\]
exhibits $Y(C)$ as the total space of the family whose fibre over
$v:C\to V$ is
\[
    \operatorname{Map}_{V}(\operatorname{Spec}C_v,Y),
\]
where $C_v$ denotes the same affine equipped with the chosen $V$-structure.
Represented Postnikov convergence gives
\[
    V(C)\xrightarrow{\sim}\lim_nV(\tau_{\le n}C).
\]
For a compatible point $v=(v_n)$, deformation locality of $Y$ over $V$ gives
\[
    \operatorname{Map}_{V}(\operatorname{Spec}C_v,Y)
    \xrightarrow{\sim}
    \lim_n
    \operatorname{Map}_{V}
    (\operatorname{Spec}(\tau_{\le n}C)_{v_n},Y).
\]
Thus the canonical map
\[
    Y(C)\longrightarrow\lim_nY(\tau_{\le n}C)
\]
is an equivalence on the base $V(C)$ and on every homotopy fibre, hence is an
equivalence.

Now let
\[
    C\simeq A_0\times_{A_{01}}A_1
\]
be a connective nilpotent-excision pullback over $S$.  Because $V$ is
represented,
\[
    V(C)
    \xrightarrow{\sim}
    V(A_0)\times_{V(A_{01})}V(A_1).
\]
Fix a compatible point of this base pullback.  It equips the four affine
vertices with compatible $V$-structures without changing the underlying
nilpotent-excision square.  Deformation locality of $Y$ over $V$ gives a
Cartesian square on the corresponding homotopy fibres of
$Y(-)\to V(-)$.  A square of total spaces over a Cartesian base is Cartesian
when all its compatible homotopy-fibre squares are Cartesian.  Therefore
\[
    Y(C)
    \xrightarrow{\sim}
    Y(A_0)\times_{Y(A_{01})}Y(A_1).
\]
The two deformation equations characterize $L_S^{\mathrm{Def}}$-local
objects, proving the claim.
\end{proof}

\begin{lemma}[Ambient-base preservation of deformation and de Rham equivalences]
\label{lem:bnli-local-equivalence-base-composition}
Let $p:V\to S$ be any morphism in the native spectral Nisnevich topos.  The
postcomposition functor
\[
    \Sigma_p:\mathcal X_{/V}\longrightarrow\mathcal X_{/S}
\]
carries relative de Rham equivalences over $V$ to relative de Rham
equivalences over $S$.  It likewise carries relative deformation equivalences
over $V$ to relative deformation equivalences over $S$.
\end{lemma}

\begin{proof}
The relative de Rham local-equivalence class over a base is the strongly
saturated class generated by affine spectral square-zero immersions in the
corresponding slice.  Postcomposition does not change the underlying affine
square-zero immersion, so it carries every generator over $V$ to a generator
over $S$.  Since $\Sigma_p$ is a left adjoint, it preserves all colimits and
therefore pushouts and transfinite composites; every functor preserves retract
diagrams.  It consequently carries the strongly saturated class generated over
$V$ into the strongly saturated class generated over $S$.

For deformation equivalences use the generating presentation by Postnikov
convergence maps and nilpotent-excision maps.  Again postcomposition changes
only the structural map to the ambient base and leaves the underlying affine
generator unchanged.  The same strong-saturation argument applies.
\end{proof}

\begin{definition}[Spectral formal-moduli stack]
\label{def:bnli-formal-moduli-stack}
A \emph{spectral formal-moduli stack on $X/S$} is a nil-centred
deformation-local target
\[
    X\longrightarrow Y
\]
whose centre map is spectrally laft.  Write
\[
\boxed{
    \operatorname{ModuliStk}^{\mathrm{sp}}(X/S)
    \subset
    \operatorname{DefCent}^{\mathrm{nil}}(X/S)}
\]
for the full subcategory spanned by this property.
\end{definition}

\begin{proposition}[Reduced-centre detection on the laft locus]
\label{prop:bnli-formal-stack-reduced-centre}
Every spectral formal-moduli stack has reduced-classically equivalent centre:
\[
    D_{\mathrm{red}}X
    \xrightarrow{\sim}
    D_{\mathrm{red}}Y.
\]
Conversely, for a spectrally laft deformation-local centre, the spectral
nil-isomorphism condition is equivalent to reduced-classical equivalence of
the centre.
\end{proposition}

\begin{proof}
If $X\to Y$ is a spectral nil-isomorphism, evaluate it on a discrete reduced
ring $H(R)$.  Such a ring has no nonzero nilpotent ideal, so its pro-reduction
system is the singleton $H(R)$ itself.  Hence the centre is an equivalence on
all reduced classical affine tests, giving
$D_{\mathrm{red}}X\simeq D_{\mathrm{red}}Y$.

Conversely suppose $X\to Y$ is spectrally laft and reduced-classically an
equivalence.  Let $D$ be an ordinary ring and write $D_{\mathrm{red}}$ for its
reduction.  The canonical map
\[
    \operatorname*{colim}_{I\in\mathcal I_{\mathrm{nil}}(D)}D/I
    \xrightarrow{\sim}
    D_{\mathrm{red}}
\]
is an isomorphism.  Indeed, finite sums of nilpotent ideals are nilpotent, and
the union of all nilpotent ideals is the nilradical.  Applying the relative
filtered-colimit condition in the definition of spectral laft gives a
Cartesian square
\[
\begin{tikzcd}
\operatorname*{colim}_I X(H(D/I)) \arrow[r] \arrow[d] &
X(H(D_{\mathrm{red}})) \arrow[d]\\
\operatorname*{colim}_I Y(H(D/I)) \arrow[r] &
Y(H(D_{\mathrm{red}})).
\end{tikzcd}
\]
The right vertical map is an equivalence by the reduced-centre hypothesis, so
the left vertical map is an equivalence.  This is exactly the spectral
nil-isomorphism condition on the pro-reduction associated to $D$.  Applying
the same argument to $D=\pi_0C$ for an arbitrary connective spectral test $C$
proves the claim.
\end{proof}

\subsection{Artinian restriction}

Fix a coherent affine-etale pair
\[
    V=\operatorname{Spec}R\longrightarrow S,
    \qquad
    U=\operatorname{Spec}B\longrightarrow X\times_SV.
\]
For a centred target $X\to Y$, let $u:U\to Y$ be the restricted centre.

\begin{definition}[Affine Artinian leaf]
\label{def:bnli-artinian-leaf}
Define
\[
    \operatorname{Leaf}^{\mathrm{Art}}_{U/V}(Y)(A\to B)
    :=
    \operatorname{fib}_{u}
    \left(
       \operatorname{Map}_S(\operatorname{Spec}A,Y)
       \longrightarrow
       \operatorname{Map}_S(U,Y)
    \right).
\]
\end{definition}

\begin{proposition}[Artinian restriction is formal-moduli valued]
\label{prop:bnli-artinian-leaf-fmp}
For every $L_S^{\mathrm{Def}}$-local centred target $Y$, the preceding functor
belongs to $\operatorname{FMP}^{\mathrm{sp}}_{B/R}$.  It is natural in
morphisms of centred targets and in morphisms of coherent affine-etale pairs.
\end{proposition}

\begin{proof}
At the identity object $B\to B$ the lifting fibre is contractible.  A
Schlessinger square is one of the nilpotent-excision squares inverted by
$L_S^{\mathrm{Def}}$.  Evaluation of the local target carries that square to a
pullback of spaces, and taking the homotopy fibre over the compatible centre
preserves the pullback.  Naturality follows from restriction of mapping spaces
and the coherent Artinian transport constructed in the differentiation
chapter.  At this stage only the canonical transport morphisms are used; their
invertibility under spectral-etale change of centre is the content of the next
proposition.
\end{proof}

\begin{proposition}[Etale cocartesianity of centred Artinian restriction]
\label{prop:bnli-artinian-leaf-etale}
Let
\[
    f:(U',V')\longrightarrow(U,V)
\]
be a morphism of coherent affine-etale pairs, and let $Y$ be
$L_S^{\mathrm{Def}}$-local with centre $u:U\to Y$.  Write $u':U'\to Y$ for
the restricted centre.  Then the canonical centre-change morphism
\[
\boxed{
    \lambda_{f,Y}:
    f_{\mathrm{FMP},!}
    \operatorname{Leaf}^{\mathrm{Art}}_{U/V}(Y)
    \xrightarrow{\sim}
    \operatorname{Leaf}^{\mathrm{Art}}_{U'/V'}(Y)}
\]
is an equivalence.  These equivalences are coherent for identities,
composition, and higher spectral-etale pasting.  The same statement holds
after the evident base change of $Y$ to the new affine base.
\end{proposition}

\begin{proof}
Put
\[
    F:=\operatorname{Leaf}^{\mathrm{Art}}_{U/V}(Y),
    \qquad
    F':=\operatorname{Leaf}^{\mathrm{Art}}_{U'/V'}(Y).
\]
Composition of a centred Artinian lift with the canonical spectral-etale
transport of its test gives the natural transformation whose Schlessinger
localization is $\lambda_{f,Y}$.

Apply affine spectral differentiation.  Coherent two-variable Artinian
transport gives a canonical equivalence
\[
    \operatorname{Diff}^{\pi,E_\infty}_{B'/R'}
    \left(
      f_{\mathrm{FMP},!}F
    \right)
    \simeq
    f^*
    \operatorname{Diff}^{\pi,E_\infty}_{B/R}(F).
\]
The pre-tangent identification theorem of the differentiation chapter
identifies the underlying anchored pro-coherent objects of the two affine
differentiations with
\[
    \mathbb T^{\mathrm{pc}}_{U/Y;V}[1]
    \longrightarrow
    \mathbb T^{\mathrm{pc}}_{B/R}[1]
\]
and
\[
    \mathbb T^{\mathrm{pc}}_{U'/Y;V'}[1]
    \longrightarrow
    \mathbb T^{\mathrm{pc}}_{B'/R'}[1].
\]
Because $Y$ is $L_S^{\mathrm{Def}}$-local, coherent spectral-etale tangent
transport identifies the second anchored object with the $f^\sharp$-transport
of the first.  On underlying pro-coherent sources, $f^\sharp$ is ordinary
pro-coherent pullback.  Under these identifications the underlying anchored
morphism of
$\operatorname{Diff}^{\pi,E_\infty}(\lambda_{f,Y})$
is exactly that base-change equivalence.

The forgetful functor from spectral partition-Lie algebroids to anchored
pro-coherent objects is conservative by monadicity.  Hence
$\operatorname{Diff}^{\pi,E_\infty}(\lambda_{f,Y})$ is an equivalence.
Affine spectral differentiation is itself an equivalence, so
$\lambda_{f,Y}$ is an equivalence.  The identity, composition, and
higher-pasting coherences are inherited from the coherent anchored
spectral-etale transport and the uniqueness of Schlessinger localization.
\end{proof}

\begin{definition}[Global Artinian restriction]
\label{def:bnli-global-artinian-restriction}
By Proposition~\ref{prop:bnli-artinian-leaf-etale}, the affine Artinian leaf
functors form a coherent spectral-etale cocartesian section of the global
formal-moduli coefficient system.  Taking global sections defines
\[
\boxed{
    \operatorname{Leaf}^{\mathrm{Art}}_{X/S}:
    \operatorname{DefTgt}(X/S)
    \longrightarrow
    \operatorname{FMP}^{\mathrm{sp}}_{X/S}.}
\]
\end{definition}

\section{Finite nilpotent factorization and spectral Artinian Cech descent}
\label{sec:bnli-artinian-cech-descent}

The native quotient construction uses effective images of Artinian centres.
The bridge from an Artinian thickening to its native effective image is now a
geometric consequence of Theorem~\ref{thm:bnli-def-left-exact}.  In the
left-exact deformation-local topos, one structured square-zero thickening is
an effective epimorphism; finite square-zero generation then gives effectivity
for every spectral Artinian thickening.  No adic bar-length convergence
argument is required.

\begin{proposition}[Finite square-zero generation without finite-presentation input]
\label{prop:bnli-unrestricted-finite-square-zero}
Let $A\to B$ be a connective augmented spectral algebra such that
$\pi_0(A)\to\pi_0(B)$ is surjective with nilpotent kernel and
$\operatorname{fib}(A\to B)$ is eventually coconnective.  Then $A\to B$
admits a finite factorization by structured square-zero extensions whose
coefficients are restrictions of connective, eventually coconnective
$B$-modules.  No such factorization is functorially retained.
\end{proposition}

\begin{proof}
Put
\[
    I:=\ker(\pi_0A\to\pi_0B),
    \qquad I^m=0.
\]

First separate the discrete nilpotent part.  Form the pullback in unrestricted
commutative ring spectra
\[
    A^{(1)}
    :=
    H\pi_0(A)\times_{H\pi_0(B)}B.
\]
The map $H\pi_0(A)\to H\pi_0(B)$ is surjective on $\pi_0$, so
Lemma~\ref{lem:bnli-connective-pullback-recognition} makes $A^{(1)}$
connective; hence this is also the pullback in connective commutative ring
spectra.  For $1\leq k\leq m$ put
\[
    D_k:=H(\pi_0(A)/I^k),
    \qquad
    \widetilde D_k:=D_k\times_{D_1}B.
\]
Every $\widetilde D_k$ is connective by
Lemma~\ref{lem:bnli-connective-pullback-recognition}, because
$D_k\to D_1$ is surjective on $\pi_0$.  Then
\[
    \widetilde D_1\simeq B,
    \qquad
    \widetilde D_m\simeq A^{(1)},
\]
and pullback pasting gives
\[
    \widetilde D_{k+1}
    \simeq
    D_{k+1}\times_{D_k}\widetilde D_k.
\]
The map $D_{k+1}\to D_k$ is a discrete structured square-zero extension with
coefficient $H(I^k/I^{k+1})$.  The ideal $I$ annihilates
$I^k/I^{k+1}$, so the coefficient action factors through
$\pi_0(B)=\pi_0(A)/I$ and hence is the restriction of the discrete, therefore
eventually coconnective, connective $B$-module $H(I^k/I^{k+1})$.
Structured square-zero extensions are stable under affine base change, so
\[
    A^{(1)}=\widetilde D_m\longrightarrow\cdots
    \longrightarrow\widetilde D_1=B
\]
is a finite structured square-zero factorization with coefficients restricted
from $B$.

It remains to factor $A\to A^{(1)}$.  The map $A\to A^{(1)}$ is an
isomorphism on $\pi_0$.  We first produce a uniform upper Postnikov bound for
its cofiber.  Since
\[
    \operatorname{fib}(A^{(1)}\to B)
    \simeq H(I)
\]
is bounded and $\operatorname{fib}(A\to B)$ is eventually coconnective, fibre
transitivity for
\[
    A\longrightarrow A^{(1)}\longrightarrow B
\]
shows that $\operatorname{fib}(A\to A^{(1)})$ is eventually coconnective.
Therefore
\[
    C_1:=\operatorname{cofib}(A\to A^{(1)})
\]
is $1$-connective and $T$-coconnective for some finite $T$.

Inductively suppose that maps
\[
    A\longrightarrow A^{(i)}\longrightarrow A^{(1)}
\]
have been constructed with
\[
    C_i:=\operatorname{cofib}(A\to A^{(i)})
\]
$i$-connective and $T$-coconnective.  The canonical relative
cofiber--cotangent morphism is the composite
\[
    h_i:
    C_i
    \longrightarrow
    A^{(i)}\otimes_A C_i
    \longrightarrow
    L_{A^{(i)}/A}.
\]
Because $\pi_0(A)\xrightarrow{\sim}\pi_0(A^{(i)})$, the first arrow induces
an isomorphism on $\pi_i$, and the relative spectral Hurewicz theorem gives
an isomorphism on $\pi_i$ for the second.  Hence
\[
\boxed{
    \pi_i(h_i):\pi_iC_i\xrightarrow{\sim}\pi_iL_{A^{(i)}/A}.}
\]
Put
\[
    N_i:=\pi_iC_i,
    \qquad
    M_i:=HN_i[i-1].
\]
The relative cotangent complex is $i$-connective in this situation.  Its first
Postnikov projection, followed by the inverse of $\pi_i(h_i)$, therefore
defines
\[
    \eta_i:
    L_{A^{(i)}/A}
    \longrightarrow
    H\pi_iL_{A^{(i)}/A}[i]
    \xrightarrow{\ H(\pi_i(h_i)^{-1})[i]\ }
    HN_i[i]
    =
    M_i[1].
\]
The composite
\[
    C_i\xrightarrow{h_i}L_{A^{(i)}/A}
    \xrightarrow{\eta_i}HN_i[i]
\]
induces the identity on $N_i=\pi_iC_i$.  Since $C_i$ is $i$-connective and
$HN_i[i]$ is $i$-truncated, the canonical truncation equivalence
\[
    \operatorname{Map}(C_i,HN_i[i])
    \simeq
    \operatorname{Hom}(\pi_iC_i,N_i)
\]
shows that this composite is exactly the canonical $i$-th Postnikov
projection
\[
    C_i\longrightarrow H\pi_i(C_i)[i].
\]

Let
\[
    A^{(i+1)}\longrightarrow A^{(i)}
\]
be the structured square-zero extension over $A$ classified by $\eta_i$.
Its underlying square-zero fibre sequence is
\[
    M_i\longrightarrow A^{(i+1)}\longrightarrow A^{(i)}.
\]
Octahedrality for
\[
    A\longrightarrow A^{(i+1)}\longrightarrow A^{(i)}
\]
then gives a cofiber sequence
\[
    C_{i+1}\longrightarrow C_i\longrightarrow M_i[1].
\]
The last arrow is the canonical Postnikov projection just identified, and
therefore
\[
\boxed{
    C_{i+1}
    \simeq
    \tau_{\geq i+1}C_i.}
\]
Thus connectivity rises by one while the upper coconnective bound remains
$T$.

The coefficient $M_i$ need not itself descend from $B$, but it admits a
finite $I$-adic filtration
\[
    M_i\supseteq IM_i\supseteq\cdots\supseteq I^mM_i=0.
\]
Its successive cofibres are
\[
    H(I^kN_i/I^{k+1}N_i)[i-1].
\]
The displayed homotopy modules are annihilated by $I$, so their
$\pi_0(A^{(i)})$-actions factor through $\pi_0(B)$.  Consequently every
successive cofiber is the restriction of a connective, eventually
coconnective $B$-module; in fact it is concentrated in the single homotopy
degree $i-1$.  The coefficient-filtration refinement lemma of the differentiation chapter,
which requires only a finite filtration of the coefficient and not coherence
of its successive quotients, refines
\[
    A^{(i+1)}\to A^{(i)}
\]
into finitely many structured square-zero extensions with coefficients
restricted from those connective $B$-modules.  The refinement changes neither
the composite nor the Postnikov calculation above.

For $i>T$ one has $C_i\simeq0$, so $A\simeq A^{(i)}$.  Concatenating the
finitely many refined Postnikov-killing stages with the discrete nilpotent
tower from $A^{(1)}$ to $B$ gives the desired finite factorization.  All
choices occur only inside the proof and no factorization is retained on the
arrow $A\to B$.
\end{proof}

\subsection{Artinian effectivity in the deformation-local topos}

\begin{lemma}[Structured square-zero centres are deformation-effective]
\label{lem:bnli-square-zero-def-effective}
Let $B_\alpha\to B$ be a structured square-zero extension over a coherent
centre, classified by a connective coefficient.  Then the geometric centre
map
\[
    \operatorname{Spec}B\longrightarrow\operatorname{Spec}B_\alpha
\]
is an effective epimorphism in the deformation-local $\infty$-topos
$\mathcal X^{\mathrm{Def}}_{/S}$.
\end{lemma}

\begin{proof}
Write the structured square-zero extension as the pullback
\[
    B_\alpha
    \simeq
    B\times_{B\oplus M[1]}B,
\]
where the two maps $B\to B\oplus M[1]$ are the classifying derivation and the
zero derivation.  The corresponding nilpotent-excision generator becomes an
equivalence after applying $L_S^{\mathrm{Def}}$; all four affine spectral
schemes involved are already deformation-local.  Hence in
$\mathcal X^{\mathrm{Def}}_{/S}$ there is a canonical pushout square
\[
\begin{tikzcd}
\operatorname{Spec}(B\oplus M[1]) \arrow[r] \arrow[d] &
\operatorname{Spec}B \arrow[d]\\
\operatorname{Spec}B \arrow[r] &
\operatorname{Spec}B_\alpha.
\end{tikzcd}
\]
Each of the two maps out of $\operatorname{Spec}(B\oplus M[1])$ is split by
the section induced from the augmentation $B\oplus M[1]\to B$, and hence is
an effective epimorphism.  Effective epimorphisms are stable under pushout in
an $\infty$-topos.  Therefore either coprojection
$\operatorname{Spec}B\to\operatorname{Spec}B_\alpha$ is an effective
epimorphism.
\end{proof}

\begin{theorem}[Spectral Artinian Cech effectivity and descent]
\label{thm:bnli-artinian-cech-descent}
Let $A\to B$ be spectral Artinian.  Then
\[
    \operatorname{Spec}B\longrightarrow\operatorname{Spec}A
\]
is an effective epimorphism in $\mathcal X^{\mathrm{Def}}_{/S}$.  Consequently,
for every $L_S^{\mathrm{Def}}$-local target $Z$, the canonical augmented Cech
comparison is an equivalence
\[
\boxed{
    \operatorname{Map}_S(\operatorname{Spec}A,Z)
    \xrightarrow{\sim}
    \operatorname*{Tot}_{[n]\in\Delta}
    \operatorname{Map}_S
    \left(
      \operatorname{Spec}(B^{\otimes_A(n+1)}),Z
    \right).}
\]
The equivalence is natural in the Artinian arrow, in $Z$, and under coherent
spectral-etale transport inside the constructed Artinian coefficient system.
\end{theorem}

\begin{proof}
The spectral Artinian factorization theorem of the differentiation chapter
provides, only inside this proof, a finite factorization
\[
    A=A_m\longrightarrow A_{m-1}\longrightarrow\cdots
    \longrightarrow A_0=B
\]
by structured square-zero extensions with connective coherent coefficients.
By Lemma~\ref{lem:bnli-square-zero-def-effective}, every geometric map
\[
    \operatorname{Spec}A_{i-1}
    \longrightarrow
    \operatorname{Spec}A_i
\]
is an effective epimorphism in $\mathcal X^{\mathrm{Def}}_{/S}$.  Effective
epimorphisms compose, so
$\operatorname{Spec}B\to\operatorname{Spec}A$ is effective there.  The
conclusion is independent of the temporary factorization because effectivity
is a property of the canonical geometric centre map.

Theorem~\ref{thm:bnli-def-left-exact} implies that the inclusion of the
local topos preserves the finite limits defining the Cech nerve.  Since every
represented spectral scheme is already deformation-local, the Cech nerve of
the preceding effective epimorphism in $\mathcal X^{\mathrm{Def}}_{/S}$ is
literally
\[
    [n]\longmapsto
    \operatorname{Spec}(B^{\otimes_A(n+1)}).
\]
Effectivity gives
\[
    \operatorname{Spec}A
    \simeq
    \left|
      \check C_\bullet(\operatorname{Spec}B/\operatorname{Spec}A)
    \right|
\]
in the deformation-local topos.  Mapping this geometric realization into a
local object $Z$ gives the displayed totalization.  Naturality is the
naturality of the canonical Cech augmentation; coherent spectral-etale
transport preserves the Artinian square-zero towers, relative tensor products,
and the deformation localization generators.
\end{proof}

\subsection{Native effective images of Artinian thickenings}

For a spectral Artinian arrow $A\to B$, put
\[
    T_A:=\operatorname{Spec}A,
    \qquad
    U:=\operatorname{Spec}B,
\]
and write
\[
    U\xrightarrow{q_A}I_A\lhook\joinrel\longrightarrow T_A
\]
for the effective-epimorphism--monomorphism factorization of the centre in the
native big spectral Nisnevich topos.

\begin{corollary}[Artinian effective image is a deformation equivalence]
\label{cor:bnli-artinian-effective-image-def}
For every spectral Artinian arrow $A\to B$, the canonical monomorphism
\[
    i_A:I_A\longrightarrow T_A
\]
is an $L_S^{\mathrm{Def}}$-equivalence.  Since $T_A$ is represented and hence
$L_S^{\mathrm{Def}}$-local,
\[
\boxed{
    L_S^{\mathrm{Def}}(I_A)
    \xrightarrow{\sim}
    T_A.}
\]
\end{corollary}

\begin{proof}
Since $q_A:U\twoheadrightarrow I_A$ is an effective epimorphism in the native
topos,
\[
    I_A\simeq|\check C_\bullet(q_A)|.
\]
Because $i_A:I_A\hookrightarrow T_A$ is a monomorphism, it identifies the
kernel pairs and all iterated Cech terms:
\[
    \check C_\bullet(q_A)
    \simeq
    \check C_\bullet(U\to T_A).
\]
Apply $L_S^{\mathrm{Def}}$.  It preserves geometric realizations because it is
a left adjoint and preserves the finite limits defining every Cech term by
Theorem~\ref{thm:bnli-def-left-exact}.  Therefore
\[
\begin{aligned}
    L_S^{\mathrm{Def}}I_A
    &\simeq
    \left|
      \check C_\bullet^{\mathrm{Def}}(U/T_A)
    \right|\\
    &\simeq T_A,
\end{aligned}
\]
where the second equivalence is Theorem~\ref{thm:bnli-artinian-cech-descent}.
This equivalence is induced by $i_A$.
\end{proof}

\begin{proposition}[Artinian cell quotient is already formal]
\label{prop:bnli-artinian-cell-formal}
The effective quotient $q_A:U\twoheadrightarrow I_A$ is a formal effective
quotient.  Equivalently,
\[
    L_V^{\mathrm{dR}}(U)
    \xrightarrow{\sim}
    L_V^{\mathrm{dR}}(I_A),
\]
and $\check C_\bullet(q_A)$ is a brave-new formal Lie groupoid on $U/V$.
\end{proposition}

\begin{proof}
The Artinian map $U\to T_A$ is a finite composite of structured square-zero
thickenings and is therefore a relative de Rham equivalence by Chapter~9.
Factor it as
\[
    U\xrightarrow{q_A}I_A\xrightarrow{i_A}T_A.
\]
Relative de Rham localization carries effective epimorphisms to effective
epimorphisms and is left exact, hence preserves monomorphisms.  Thus after
relative de Rham localization the two arrows are respectively an effective
epimorphism and a monomorphism and their composite is an equivalence.
Uniqueness of effective-epimorphism--monomorphism factorization forces both
localized arrows to be equivalences.  The formal quotient recognition theorem
of the formal-groupoid chapter gives the final assertion.
\end{proof}

\section{Affine spectral materialization}
\label{sec:bnli-affine-materialization}

Fix a coherent affine coefficient map $R\to B$ and write
\[
    V=\operatorname{Spec}R,
    \qquad
    U=\operatorname{Spec}B.
\]
We construct the geometric materialization of
$F\in\operatorname{FMP}^{\mathrm{sp}}_{B/R}$ without selecting a
prorepresenting complete algebra.

\subsection{Finite and unrestricted nilpotent neighbourhoods}

\begin{definition}[Nilpotent affine neighbourhood categories]
\label{def:bnli-nil-neighbourhoods}
Let
$\operatorname{Aff}^{\mathrm{sp,afp-nil,+}}_{U/V}$ be the $\infty$-category of
affine spectral maps $U\to T$ over $V$ whose coordinate augmentation is
spectral Artinian.  Thus
\[
    \operatorname{Aff}^{\mathrm{sp,afp-nil,+}}_{U/V}
    \simeq
    \left(\operatorname{Art}^{\mathrm{sp}}_{R/B}\right)^{\mathrm{op}}.
\]
Let $\operatorname{Aff}^{\mathrm{sp,nil,+}}_{U/V}$ be the full category of
affine centred maps
\[
    U\longrightarrow T=\operatorname{Spec}C
\]
for which $\pi_0(C)\to\pi_0(B)$ is surjective with nilpotent kernel and
$\operatorname{fib}(C\to B)$ is eventually coconnective.  No
almost-finite-presentation condition is imposed.  Write
\[
    u:
    \operatorname{Aff}^{\mathrm{sp,afp-nil,+}}_{U/V}
    \hookrightarrow
    \operatorname{Aff}^{\mathrm{sp,nil,+}}_{U/V}
\]
for the inclusion.
\end{definition}

\begin{definition}[Centred Yoneda presheaf and dotted-lift restriction]
\label{def:bnli-centred-yoneda}
On $\operatorname{Aff}^{\mathrm{sp,nil,+}}_{U/V}$ let
\[
\boxed{
    h_U^{\mathrm{cent}}(T)
    :=
    \operatorname{Map}_{U/}(T,U).}
\]
Thus $h_U^{\mathrm{cent}}$ is the Yoneda presheaf of the identity
neighbourhood $U\to U$ in the centred neighbourhood category.  If
$u:U\to Y$ is an ambient centred prestack over $V$, define its dotted-lift
restriction by
\[
\boxed{
    \operatorname{Lift}^{\mathrm{nil}}_U(Y)(T)
    :=
    \operatorname{Map}_{U/}(T,Y).}
\]
The centre map is the natural transformation
\[
    h_U^{\mathrm{cent}}
    \longrightarrow
    \operatorname{Lift}^{\mathrm{nil}}_U(Y)
\]
obtained by composition with $u$.  This notation is kept separate from the
ambient represented prestack $U(T)=\operatorname{Map}_V(T,U)$.
\end{definition}


Regard a formal-moduli problem $F$ as a presheaf on the finite geometric
category and let $u_!F$ be its left Kan extension to all nilpotent affine
neighbourhoods.

\begin{lemma}[Nilpotent descent of almost-perfect modules]
\label{lem:bnli-nilpotent-descent-aperf}
Let $C\to B$ be a finite composite of structured square-zero extensions whose
coefficients are restrictions of connective coherent $B$-modules.  Let $M$ be
a connective $C$-module.  If
\[
    B\otimes_C M
\]
is almost perfect over $B$, then $M$ is almost perfect over $C$.
\end{lemma}

\begin{proof}
It is enough to treat one square-zero stage and then iterate.  Suppose
$C\to D$ is a structured square-zero extension with coefficient $N$, where
$N$ is connective and almost perfect over $D$, and suppose
$D\otimes_C M$ is almost perfect.  The square-zero fibre sequence
\[
    N\longrightarrow C\longrightarrow D
\]
remains a fibre sequence after tensoring over $C$ with $M$:
\[
    N\otimes_C M
    \longrightarrow M
    \longrightarrow D\otimes_C M.
\]
Because the $C$-action on $N$ factors through $D$,
\[
    N\otimes_C M
    \simeq
    N\otimes_D(D\otimes_C M).
\]
The right-hand side is almost perfect over $D$, and hence after restriction it
is almost perfect over $C$ because $D$ is almost perfect as a $C$-module by
spectral square-zero finiteness.  The same is true of $D\otimes_CM$.  Almost
perfect modules form a stable subcategory, so the fibre sequence makes $M$
almost perfect over $C$.  Induction through the finite square-zero tower proves
the lemma.
\end{proof}

\begin{proposition}[Cotangent criterion for almost finite presentation without a surjectivity hypothesis]
\label{prop:bnli-general-afp-cotangent}
Let $A\to C$ be a morphism of connective commutative spectral rings.  Suppose
that $\pi_0(C)$ is finitely presented as a $\pi_0(A)$-algebra and that
$L_{C/A}$ is almost perfect as a $C$-module.  Then $C$ is almost finitely
presented over $A$.
\end{proposition}

\begin{proof}
Choose finitely many algebra generators of $\pi_0(C)$ over $\pi_0(A)$.  Let
\[
    P:=\operatorname{Sym}_A(A^{\oplus r})
\]
be the corresponding finite free connective commutative $A$-algebra and choose
a morphism $P\to C$ lifting the selected generators.  Then
\[
    \pi_0(P)\twoheadrightarrow\pi_0(C),
\]
and $\pi_0(C)$ is finitely presented as a $\pi_0(P)$-algebra.

The algebra $P$ is finite cell over $A$, so $L_{P/A}$ is perfect.  Cotangent
transitivity gives a cofiber sequence of $C$-modules
\[
    C\otimes_PL_{P/A}
    \longrightarrow
    L_{C/A}
    \longrightarrow
    L_{C/P}.
\]
The first term is perfect and the second is almost perfect by hypothesis;
hence $L_{C/P}$ is almost perfect.  The morphism $P\to C$ is surjective on
$\pi_0$, so the surjective spectral cell criterion of the differentiation
chapter applies and makes $C$ almost finitely presented over $P$.  The finite
free algebra $P$ is finitely presented over $A$, and almost finite presentation
is stable under composition.  Therefore $C$ is almost finitely presented over
$A$.
\end{proof}

\begin{lemma}[Morphisms between Artinian neighbourhoods are almost finitely presented]
\label{lem:bnli-artinian-morphism-afp}
Let
\[
    A\longrightarrow A'\longrightarrow B
\]
be a morphism between spectral Artinian extensions of the coherent centre
$B/R$.  Then $A'$ is almost finitely presented as a connective commutative
$A$-algebra.  Equivalently, the corresponding affine morphism
$\operatorname{Spec}A'\to\operatorname{Spec}A$ is spectrally almost finitely
presented.
\end{lemma}

\begin{proof}
First we prove module finiteness.  The spectral Artinian factorization theorem
writes $A'\to B$, only inside this proof, as a finite composite of structured
square-zero extensions by connective coherent $B$-modules.  Since $A\to B$ is
spectral Artinian, $B$ is almost perfect as an $A$-module.  Every coherent
$B$-module occurring in the tower is almost perfect over $B$ and therefore,
by restriction through the almost-perfect $A$-module $B$, is almost perfect
over $A$.  Induction through the square-zero fibre sequences makes $A'$ almost
perfect as an $A$-module.  In particular $\pi_0(A')$ is finitely presented as
a $\pi_0(A)$-module.

A commutative algebra which is finitely presented as a module is finitely
presented as an algebra.  Indeed choose finitely many module generators,
including the unit, finitely many module relations, and for every ordered pair
of generators record their product as a finite linear combination of the same
generators.  The module relations together with these finitely many
multiplication relations present the algebra.  Hence
$\pi_0(A')$ is finitely presented as a $\pi_0(A)$-algebra.

Cotangent transitivity for $A\to A'\to B$ gives a cofiber sequence
\[
    B\otimes_{A'}L_{A'/A}
    \longrightarrow
    L_{B/A}
    \longrightarrow
    L_{B/A'}.
\]
The last two terms are almost perfect because $A\to B$ and $A'\to B$ are
spectral Artinian.  Thus $B\otimes_{A'}L_{A'/A}$ is almost perfect.  The
relative cotangent complex $L_{A'/A}$ is connective.  Apply
Lemma~\ref{lem:bnli-nilpotent-descent-aperf} to the finite Artinian
square-zero tower $A'\to B$; it follows that $L_{A'/A}$ itself is almost
perfect over $A'$.

Proposition~\ref{prop:bnli-general-afp-cotangent} now applies to $A\to A'$:
its degree-zero algebra is finitely presented and its relative cotangent
complex is almost perfect.  Therefore $A'$ is almost finitely presented over
$A$.
\end{proof}

\begin{lemma}[Almost-finitely-presented affine maps are spectrally laft]
\label{lem:bnli-affine-afp-laft}
Let $A\to C$ be an almost finitely presented morphism of connective spectral
rings.  Then the represented affine morphism
\[
    \operatorname{Spec}C
    \longrightarrow
    \operatorname{Spec}A
\]
is spectrally laft in the sense of
Definition~\ref{def:bnli-spectral-laft}.
\end{lemma}

\begin{proof}
Represented spectral prestacks are Postnikov convergent.  Let
$\{D_\alpha\}$ be a filtered diagram of $n$-truncated connective
$A$-algebras with colimit $D$.  Almost finite presentation of $C$ over $A$
is exactly the equivalence
\[
    \operatorname*{colim}_\alpha
    \operatorname{Map}_{\operatorname{CAlg}_A}(C,D_\alpha)
    \xrightarrow{\sim}
    \operatorname{Map}_{\operatorname{CAlg}_A}(C,D).
\]
After allowing the structural map to $\operatorname{Spec}A$ to vary, this is
precisely the assertion that the relative filtered-colimit square in
Definition~\ref{def:bnli-spectral-laft} is Cartesian.
\end{proof}

\begin{lemma}[Bounded coherent approximation of square-zero coefficients]
\label{lem:bnli-bounded-coherent-approximation}
Let $C$ be a coherent connective spectral ring and let $M$ be a connective
$C$-module which is $n$-coconnective for some finite $n$.  Then there is a
filtered diagram
\[
    \{M_\alpha\}_{\alpha\in A}
    \subset
    \operatorname{Coh}(C)\cap\operatorname{Mod}_{C,[0,n]}
\]
and a canonical colimit equivalence
\[
\boxed{
    \operatorname*{colim}_{\alpha}M_\alpha
    \xrightarrow{\sim}
    M.}
\]
If $L$ is an almost-perfect $C$-module, then
\[
\boxed{
    \operatorname{Map}_C(L,M)
    \xrightarrow{\sim}
    \operatorname*{colim}_{\alpha}
    \operatorname{Map}_C(L,M_\alpha).}
\]
The approximation diagram is used only inside proofs and is not retained as
structure on $M$.
\end{lemma}

\begin{proof}
The stable module category is compactly generated by its perfect objects, so
there is a filtered presentation
\[
    M\simeq\operatorname*{colim}_{\alpha}P_\alpha,
    \qquad
    P_\alpha\in\operatorname{Perf}(C).
\]
The standard module $t$-structure is compatible with filtered colimits.  Since
$M\in\operatorname{Mod}_{C,[0,n]}$, applying the bounded truncation functor
termwise gives
\[
\begin{aligned}
    M
    &\simeq \tau_{[0,n]}M\\
    &\simeq
    \tau_{[0,n]}
    \operatorname*{colim}_{\alpha}P_\alpha\\
    &\simeq
    \operatorname*{colim}_{\alpha}\tau_{[0,n]}P_\alpha.
\end{aligned}
\]
Put
\[
    M_\alpha:=\tau_{[0,n]}P_\alpha.
\]

We check coherence.  A perfect $C$-module is a retract of a finite cell
$C$-module.  Since $C$ is coherent, every homotopy group of a finite cell
$C$-module is finitely presented over $\pi_0(C)$: this follows by induction on
the finite cell filtration from finite presentation of every $\pi_i(C)$ and
coherence of $\pi_0(C)$.  Retracts preserve finite presentation.  Therefore
for every perfect $P_\alpha$ and every integer $i$, the module
$\pi_i(P_\alpha)$ is finitely presented.  The bounded truncation
$M_\alpha$ is consequently eventually coconnective with finitely presented
homotopy groups, hence belongs to
$\operatorname{Coh}(C)\cap\operatorname{Mod}_{C,[0,n]}$.
This proves the first assertion.

An almost-perfect $C$-module is compact on every fixed bounded Postnikov range.
All $M_\alpha$ and $M$ lie in the same range $[0,n]$, so
\[
    \operatorname{Map}_C(L,M)
    \simeq
    \operatorname*{colim}_{\alpha}
    \operatorname{Map}_C(L,M_\alpha).
\]
No approximation diagram is retained after this argument.
\end{proof}

\begin{lemma}[Filtered representatives in a filtered colimit of spaces]
\label{lem:bnli-filtered-representatives}
Let $I$ be a filtered $\infty$-category, let $D:I\to\mathcal S$, and let
$x\in\operatorname*{colim}_I D$.  Define $\operatorname{Rep}_D(x)$ to have
objects $(i,y,\gamma)$, where $y\in D(i)$ and $\gamma$ is a path from the
image of $y$ in the filtered colimit to $x$; morphisms are refinements in $I$
together with the compatible higher homotopies.  Then:
\begin{enumerate}
    \item $\operatorname{Rep}_D(x)$ is nonempty and filtered;
    \item every finite compatible family of representatives, paths, and higher
    homotopies admits a common refinement;
    \item its groupoid completion is contractible.
\end{enumerate}
No representative is canonically retained.
\end{lemma}

\begin{proof}
The terminal space is compact with respect to filtered colimits of spaces.
Hence the point
\[
    *\xrightarrow{x}\operatorname*{colim}_I D
\]
is represented at some stage, proving nonemptiness.

Let $K\to\operatorname{Rep}_D(x)$ be a finite simplicial diagram.  Its
projection to $I$ admits a cone after passage to some stage $i$, because $I$
is filtered.  Transport every chosen representative to $i$.  What remains is
a finite collection of points of $D(i)$, paths between their images in the
filtered colimit, paths from those images to $x$, and finitely many higher
coherence simplices.

Filtered colimits of spaces commute with mapping out of finite simplicial
sets.  Consequently, after passage to one further stage $j\geq i$, all of
these finitely many paths and higher simplices are represented simultaneously
in $D(j)$.  They give a cone
\[
    K^\triangleright\longrightarrow\operatorname{Rep}_D(x).
\]
Thus $\operatorname{Rep}_D(x)$ is filtered, and the same argument is precisely
the asserted simultaneous-refinement property.

A nonempty filtered $\infty$-category has contractible groupoid completion.
Therefore
\[
\boxed{
    \left|\operatorname{Rep}_D(x)\right|\simeq *.}
\]
\end{proof}


\begin{definition}[Artinian factorization category]
\label{def:bnli-artinian-factorization-category}
Let $F\in\operatorname{FMP}^{\mathrm{sp}}_{B/R}$ and let
$S\in\operatorname{Aff}^{\mathrm{sp,nil,+}}_{U/V}$.  Define
\[
    \operatorname{Fact}_{\mathrm{Art}}(S;F)
\]
to be the Grothendieck construction whose objects are triples
\[
    (T,\xi,\sigma),
\]
where $U\to T$ is an Artinian neighbourhood, $\xi\in F(T)$, and
$\sigma:S\to T$ is a morphism of centred geometric neighbourhoods.  Morphisms
are the evident compatible morphisms of Artinian neighbourhoods carrying the
chosen point to the chosen point.  The pointwise left-Kan formula gives
\[
\boxed{
    u_!F(S)
    \simeq
    \left|
      \operatorname{Fact}_{\mathrm{Art}}(S;F)
    \right|.}
\]
\end{definition}

\begin{lemma}[Finite Artinian excision for formal-moduli problems]
\label{lem:bnli-fmp-finite-artinian-excision}
Let $F\in\operatorname{FMP}^{\mathrm{sp}}_{B/R}$.  A pushout square in the
geometric spectral Artinian neighbourhood category whose one leg is a
structured square-zero thickening is carried by $F$ to a pullback square of
spaces.  The same statement holds when that leg is a finite composite of
structured square-zero thickenings.
\end{lemma}

\begin{proof}
For one structured square-zero leg this is exactly the Schlessinger equation
in Definition~14.14.2 of the differentiation chapter, written in geometric
variance.  For a finite composite, factor the geometric leg into its finitely
many square-zero stages.  Successively base change the remaining vertex along
those stages.  Every intermediate square is again an Artinian square-zero
pushout, so $F$ carries it to a pullback.  Pasting finitely many pullback
squares gives the asserted pullback for the composite.  The factorization is
used only in this proof and is not retained.
\end{proof}

\begin{lemma}[Coherent backward lifting across an unrestricted square-zero corner]
\label{lem:bnli-coherent-backward-lift}
Let
\[
\begin{tikzcd}
S_0 \arrow[r] \arrow[d] & T \arrow[d]\\
S \arrow[r] & P
\end{tikzcd}
\]
be a commutative square of centred affine nilpotent neighbourhoods in which
$T$ and $P$ are spectral Artinian and $S_0\to S$ is a structured square-zero
thickening.  Write
\[
S_0=\operatorname{Spec}C_0,
\qquad
S=\operatorname{Spec}C,
\qquad
T=\operatorname{Spec}A_1,
\qquad
P=\operatorname{Spec}A,
\]
so the coordinate-ring diagram is
\[
A\longrightarrow A_1\longrightarrow C_0,
\qquad
A\longrightarrow C\longrightarrow C_0.
\]
Let $M$ be the connective eventually coconnective $C_0$-module classifying the
structured square-zero extension $C\to C_0$.  Then there is a filtered weakly
contractible $\infty$-category
\[
\operatorname{Lift}^{\mathrm{coh}}_{A_1/A}(C/C_0)
\]
of coherent relative square-zero refinements.  An object of this category
canonically determines a diagram
\[
\begin{tikzcd}
S_0 \arrow[r] \arrow[d] & T \arrow[d]\\
S \arrow[r] & \widetilde T \arrow[r] & P
\end{tikzcd}
\]
in which $T\to\widetilde T$ is a structured square-zero thickening with
coherent coefficient.  In particular $\widetilde T$ is spectral Artinian.
The construction is functorial under refinements, and no coherent
approximation is retained after passage to the resulting contractible
category.
\end{lemma}

\begin{proof}
The structured square-zero extension $C\to C_0$, regarded in commutative
$A$-algebras, has a relative class
\[
    \delta_A:L_{C_0/A}\longrightarrow M[1].
\]
The map $A_1\to C_0$ has the formal-lifting obstruction
\[
\begin{aligned}
    o:
    L_{A_1/A}\otimes_{A_1}C_0
    &\longrightarrow L_{C_0/A}
    \xrightarrow{\delta_A}M[1].
\end{aligned}
\]
Equivalently, by restriction of scalars and adjunction, it is a morphism
\[
\boxed{
    \bar o:L_{A_1/A}
    \longrightarrow
    \operatorname{Res}^{C_0}_{A_1}M[1].}
\]
This is the obstruction supplied by the fixed-extension lifting theorem of
Chapter~2; in particular the target cotangent complex is not
$L_{C_0/A_1}$ and no variance reversal is being used.

The spectral Artinian source $A_1$ is coherent.  Moreover
Lemma~\ref{lem:bnli-artinian-morphism-afp} makes $A\to A_1$ almost finitely
presented, hence $L_{A_1/A}$ is almost perfect.  The restricted coefficient
$\operatorname{Res}^{C_0}_{A_1}M$ is connective and bounded above.  Apply
Lemma~\ref{lem:bnli-bounded-coherent-approximation} over $A_1$ to obtain a
filtered diagram
\[
    \operatorname{Res}^{C_0}_{A_1}M
    \simeq
    \operatorname*{colim}_{\alpha}N_\alpha,
    \qquad
    N_\alpha\in\operatorname{Coh}(A_1)
\]
in one common bounded Postnikov range.  Almost-perfect compactness on that
bounded range gives
\[
\boxed{
    \operatorname{Map}_{A_1}
    \left(L_{A_1/A},\operatorname{Res}^{C_0}_{A_1}M[1]\right)
    \simeq
    \operatorname*{colim}_{\alpha}
    \operatorname{Map}_{A_1}(L_{A_1/A},N_\alpha[1]).}
\]
Let $\operatorname{Lift}^{\mathrm{coh}}_{A_1/A}(C/C_0)$ be the category of
representatives of the point $\bar o$ in this filtered colimit: an object is
$(\alpha,o_\alpha,h_\alpha)$ with
\[
    o_\alpha:L_{A_1/A}\longrightarrow N_\alpha[1]
\]
and a specified path from its image to $\bar o$; morphisms are refinements
with the compatible higher homotopies.  This is exactly the representative
category of Lemma~\ref{lem:bnli-filtered-representatives}.  It is therefore
filtered, every finite family of obstruction representatives and coherence
simplices admits a simultaneous refinement, and
\[
\boxed{
    \left|
      \operatorname{Lift}^{\mathrm{coh}}_{A_1/A}(C/C_0)
    \right|\simeq *.}
\]

For such a representative, let
\[
    A_{1,\alpha}\longrightarrow A_1
\]
be the $A$-relative structured square-zero extension classified by
$o_\alpha$.  Pulling $C\to C_0$ back along $A_1\to C_0$ gives the
$A$-relative structured square-zero extension
\[
    A_1\times_{C_0}C\longrightarrow A_1.
\]
Naturality of structured square-zero classification identifies its class with
$\bar o$.  The coefficient map $N_\alpha\to M$, together with the specified
homotopy from $o_\alpha$ to $\bar o$, therefore induces a canonical morphism
of structured square-zero extensions
\[
    A_{1,\alpha}
    \longrightarrow
    A_1\times_{C_0}C.
\]
Composing with the second projection gives $A_{1,\alpha}\to C$.  On affine
spectra this is exactly
\[
    S\longrightarrow\widetilde T_\alpha:=\operatorname{Spec}A_{1,\alpha},
    \qquad
    T\longrightarrow\widetilde T_\alpha\longrightarrow P.
\]
The coefficient $N_\alpha$ is connective coherent, so
$A_{1,\alpha}\to A_1$ is an Artinian square-zero stage.  Since $A_1\to B$ is
Artinian, finite composition makes $A_{1,\alpha}\to B$ Artinian.  All maps are
functorial under refinement of $(\alpha,o_\alpha,h_\alpha)$, which proves the
claim.
\end{proof}

\begin{lemma}[Finite coherent localization resolutions]
\label{lem:bnli-finite-zigzag-resolution}
Let $\mathcal C$ be a small $\infty$-category and let
\[
    \ell:\mathcal C\longrightarrow\mathcal C^{\mathrm{gpd}}
\]
be its groupoid completion.  For objects $c_0,c_1\in\mathcal C$ and a path
\[
    \gamma:\ell(c_0)\longrightarrow\ell(c_1)
\]
there is a nonempty space
\[
    \operatorname{LocRes}_{\mathcal C}(\gamma)
\]
of finite coherent localization resolutions of $\gamma$.  An element may be
represented by a finite zigzag of actual morphisms of $\mathcal C$
\[
    c_0
    \longleftrightarrow c_1'
    \longleftrightarrow\cdots
    \longleftrightarrow c_m=c_1
\]
together with a finite hammock of higher simplices recording its coherent
identification with $\gamma$.  The resolution space is contractible:
\[
\boxed{
    \operatorname{LocRes}_{\mathcal C}(\gamma)\simeq *.}
\]
Likewise, for every point $x\in\mathcal C^{\mathrm{gpd}}$, the space of an
object $c\in\mathcal C$ together with a finite coherent localization path
$\ell(c)\to x$ is contractible.  The same assertion holds for every finite
simplicial family of points, paths, and higher homotopies.
\end{lemma}

\begin{proof}
Localize $\mathcal C$ at the class of all of its morphisms.  Use the standard
simplicial localization model in which mapping simplices are represented by
finite hammocks.  Denote the resulting simplicial category by
\[
    L^{\mathrm{fin}}(\mathcal C).
\]
Its objects are those of $\mathcal C$; a $0$-simplex in a mapping space is a
finite word of actual morphisms, traversed either with or against their
categorical orientation, and an $n$-simplex is a finite hammock containing the
finitely many higher coherence simplices relating such words.

Because every morphism of $\mathcal C$ has been inverted, the canonical
functor
\[
    \mathcal C\longrightarrow L^{\mathrm{fin}}(\mathcal C)
\]
sends every arrow to an equivalence.  The simplicial localization universal
property therefore identifies its homotopy-coherent nerve with the
$\infty$-categorical localization
\[
    \mathcal C[\operatorname{Mor}(\mathcal C)^{-1}]
    \simeq
    \mathcal C^{\mathrm{gpd}}.
\]
Consequently, for every $c_0,c_1$, the realization of the finite-hammock
mapping object is canonically equivalent to
\[
    \operatorname{Map}_{\mathcal C^{\mathrm{gpd}}}
    (\ell(c_0),\ell(c_1)).
\]
Define $\operatorname{LocRes}_{\mathcal C}(\gamma)$ to be the homotopy fibre
of this equivalence over the fixed path $\gamma$.  It is therefore
contractible.  By construction every point of the fibre uses only finitely
many arrows and finitely many higher simplices, which gives the displayed
finite coherent zigzag description.

For the point-representative statement, use the same localization model before
fixing the source object.  Its realization is
$\mathcal C^{\mathrm{gpd}}$, so the homotopy fibre over a fixed point $x$ is
again contractible and consists exactly of an object of $\mathcal C$ with a
finite coherent localization path to $x$.  Replacing $\Delta^0$ or
$\Delta^1$ by any finite simplicial parameter gives the simultaneous finite
family statement, because the finite-hammock localization is simplicially
functorial in that parameter.
\end{proof}

\begin{theorem}[Unrestricted square-zero excision for the first Kan extension]
\label{thm:bnli-unrestricted-square-zero-kan}
Let $F\in\operatorname{FMP}^{\mathrm{sp}}_{B/R}$ and let
\[
\begin{tikzcd}
S_0 \arrow[r] \arrow[d] & S'_0 \arrow[d]\\
S \arrow[r] & S'
\end{tikzcd}
\]
be an arbitrary pushout square in
$\operatorname{Aff}^{\mathrm{sp,nil,+}}_{U/V}$ in which $S_0\to S$ is a
structured square-zero thickening.  Then the canonical comparison
\[
\boxed{
    u_!F(S')
    \xrightarrow{\sim}
    u_!F(S)\times_{u_!F(S_0)}u_!F(S'_0)}
\]
is an equivalence.
\end{theorem}

\begin{proof}
Write
\[
    \mathcal C_T:=\operatorname{Fact}_{\mathrm{Art}}(T;F),
    \qquad
    X_T:=|\mathcal C_T|\simeq u_!F(T),
\]
and put
\[
    R:=X_S\times_{X_{S_0}}X_{S'_0}.
\]
Let
\[
    \Phi:X_{S'}\longrightarrow R
\]
be the canonical restriction map.

We first resolve every point of $R$ by finite Artinian data.  A point of $R$
consists of
\[
    x\in X_S,
    \qquad
    y\in X_{S'_0},
\]
and a specified path
\[
    \gamma:x|_{S_0}\simeq y|_{S_0}
\]
in $X_{S_0}=|\mathcal C_{S_0}|$.  Let $R^{\mathrm{res}}$ be the space of the
following data over such a triple: an actual representative
\[
    (P_0,\eta_0,S\to P_0)\in\mathcal C_S
\]
of $x$, an actual representative
\[
    (T_m,\xi_m,S'_0\to T_m)\in\mathcal C_{S'_0}
\]
of $y$, and a finite coherent localization resolution of $\gamma$ between
their restrictions to $\mathcal C_{S_0}$.  By the point and path clauses of
Lemma~\ref{lem:bnli-finite-zigzag-resolution}, every homotopy fibre of
\[
    R^{\mathrm{res}}\longrightarrow R
\]
is contractible.  Hence
\[
\boxed{
    R^{\mathrm{res}}\xrightarrow{\sim}R.}
\]

Choose a representative finite localization word
\[
    T_0=P_0
    \longleftrightarrow T_1
    \longleftrightarrow\cdots
    \longleftrightarrow T_m
\]
inside one such resolution.  We propagate the Artinian lift of $S$ along this
word.  Inductively maintain an Artinian $P_i$, a map $S\to P_i$, a finite
composite of coherent square-zero thickenings
\[
    T_i\longrightarrow P_i,
\]
and a point $\eta_i\in F(P_i)$ whose restriction is the prescribed point
$\xi_i\in F(T_i)$.

If the next actual edge is oriented
\[
    T_i\longrightarrow T_{i+1},
\]
form the geometric pushout
\[
\boxed{
    P_{i+1}:=P_i\amalg_{T_i}T_{i+1}.}
\]
Base change preserves the finite coherent square-zero tower, so
$T_{i+1}\to P_{i+1}$ is again such a tower and $P_{i+1}$ is Artinian.  The
points $\eta_i$ and $\xi_{i+1}$ agree over $T_i$.  Finite Artinian excision,
Lemma~\ref{lem:bnli-fmp-finite-artinian-excision}, therefore gives a
contractible space of extensions
\[
    \eta_{i+1}\in F(P_{i+1}).
\]

If instead the actual edge is traversed backwards,
\[
    T_{i+1}\longrightarrow T_i,
\]
compose it with $T_i\to P_i$.  The resulting square
\[
\begin{tikzcd}
S_0 \arrow[r] \arrow[d] & T_{i+1} \arrow[d]\\
S \arrow[r] & P_i
\end{tikzcd}
\]
satisfies Lemma~\ref{lem:bnli-coherent-backward-lift}.  Its filtered coherent
lift category has contractible groupoid completion and gives Artinian objects
$P_{i+1}$ with
\[
    T_{i+1}\longrightarrow P_{i+1}\longrightarrow P_i,
    \qquad
    S\longrightarrow P_{i+1}.
\]
Pulling $\eta_i$ back to $P_{i+1}$ gives $\eta_{i+1}$, and its restriction to
$T_{i+1}$ is $\xi_{i+1}$ by the compatibility encoded in the localization
edge.

The localization resolution contains only finitely many arrows and finitely
many higher coherence simplices.  Forward propagation is functorial by
pushout and Artinian excision.  At backward arrows,
Lemma~\ref{lem:bnli-filtered-representatives} applied inside the coherent-lift
category shows that every finite family of obstruction representatives,
comparison paths, and higher coherence simplices admits one simultaneous
refinement.  Consequently the space of complete propagation data over a fixed
finite localization resolution is contractible.

Let
\[
    \mathscr R\longrightarrow R^{\mathrm{res}}
\]
be the space of finite localization resolutions together with all such
propagation data.  The preceding paragraph shows that this map has
contractible homotopy fibres, and hence the composite
\[
\boxed{
    \pi:\mathscr R\xrightarrow{\sim}R^{\mathrm{res}}
    \xrightarrow{\sim}R}
\]
is an equivalence.

At the final stage the maps
\[
    S\longrightarrow P_m,
    \qquad
    S'_0\longrightarrow T_m\longrightarrow P_m
\]
agree on $S_0$.  Therefore the pushout universal property gives
\[
    S'=S\amalg_{S_0}S'_0\longrightarrow P_m.
\]
Together with $\eta_m\in F(P_m)$ this defines a canonical map
\[
    P:\mathscr R\longrightarrow X_{S'}.
\]
By construction
\[
\boxed{
    \Phi P\simeq\pi.}
\]
Choose an inverse $s:R\xrightarrow{\sim}\mathscr R$ to $\pi$ and put
\[
    \Psi:=Ps:R\longrightarrow X_{S'}.
\]
Then
\[
    \Phi\Psi\simeq\pi s\simeq\operatorname{id}_R.
\]

For the other composite there is a canonical constant-resolution map
\[
    j:X_{S'}\longrightarrow\mathscr R.
\]
Indeed, an actual Artinian factorization $S'\to P$ restricts to factorizations
of $S$ and $S'_0$ with the same Artinian target $P$; their restrictions to
$S_0$ agree strictly.  Use the length-zero localization resolution and no
nontrivial propagation.  Thus
\[
    \pi j=\Phi,
    \qquad
    Pj=\operatorname{id}_{X_{S'}}.
\]
The two maps
\[
    s\Phi,j:X_{S'}\rightrightarrows\mathscr R
\]
have the same composite with the equivalence $\pi$.  Hence the space of
homotopies between them over $R$ is contractible, so
\[
    s\Phi\simeq j.
\]
Applying $P$ gives
\[
    \Psi\Phi
    =Ps\Phi
    \simeq
    Pj
    =
    \operatorname{id}_{X_{S'}}.
\]
Therefore $\Phi$ is an equivalence, proving the theorem with all higher
coherence included in the resolution spaces.
\end{proof}

\begin{theorem}[First nilpotent left-Kan extension]
\label{thm:bnli-nil-kan-deformation}
For every $F\in\operatorname{FMP}^{\mathrm{sp}}_{B/R}$, the left Kan
extension $u_!F$ has deformation theory on
$\operatorname{Aff}^{\mathrm{sp,nil,+}}_{U/V}$: its value at $U$ is
contractible and it carries pushouts along nilpotent closed immersions to
pullbacks of spaces.  Moreover
\[
\boxed{
    F\xrightarrow{\sim}u^*u_!F.}
\]
\end{theorem}

\begin{proof}
The unit is an equivalence because $u$ is fully faithful.  By
Theorem~\ref{thm:bnli-unrestricted-square-zero-kan}, $u_!F$ satisfies
excision for every structured square-zero leg.  Proposition~\ref{prop:bnli-unrestricted-finite-square-zero}
expresses every nilpotent embedding in the unrestricted neighbourhood category
as a finite composite of structured square-zero extensions.  Pasting the
resulting pullback squares proves nilpotent excision.  Finally
$u_!F(U)\simeq F(B)\simeq *$.
\end{proof}


\begin{lemma}[Relative Postnikov reconstruction for a bounded fibre]
\label{lem:bnli-relative-postnikov-reconstruction}
Let
\[
    f:C\longrightarrow D
\]
be a morphism of connective commutative spectral rings and suppose
\[
    K:=\operatorname{fib}(C\to D)
\]
is $(n-1)$-coconnective.  Put
\[
    C_n:=\tau_{\leq n}C,
    \qquad
    D_n:=\tau_{\leq n}D.
\]
Then the canonical map
\[
\boxed{
    K
    \xrightarrow{\sim}
    \operatorname{fib}(C_n\to D_n)}
\]
is an equivalence.  Consequently the natural square
\[
\begin{tikzcd}
C \arrow[r] \arrow[d] & D \arrow[d]\\
C_n \arrow[r] & D_n
\end{tikzcd}
\]
is Cartesian and
\[
\boxed{
    C
    \xrightarrow{\sim}
    D\times_{D_n}C_n.}
\]
Equivalently, on affine spectra,
\[
\boxed{
    \operatorname{Spec}C
    \simeq
    \operatorname{Spec}D
    \amalg_{\operatorname{Spec}D_n}
    \operatorname{Spec}C_n.}
\]
\end{lemma}

\begin{proof}
Let
\[
    K_n:=\operatorname{fib}(C_n\to D_n).
\]
For $i<n$, truncation does not change the homotopy groups of $C$ or $D$.
Comparison of the long exact homotopy sequences therefore gives
\[
    \pi_iK\xrightarrow{\sim}\pi_iK_n.
\]
In degree $n$, the hypothesis gives $\pi_nK=0$.  Hence the long exact
sequence for $K$ makes
\[
    \pi_nC\longrightarrow\pi_nD
\]
injective.  Since $\pi_nC_n=\pi_nC$ and $\pi_nD_n=\pi_nD$, the long exact
sequence for $K_n$ gives
\[
    \pi_nK_n
    =
    \ker(\pi_nC\to\pi_nD)
    =0.
\]
Both $K$ and $K_n$ are $(n-1)$-coconnective, so the canonical map
$K\to K_n$ is an equivalence on every homotopy group and hence an
equivalence.

The induced map on horizontal fibres is an equivalence, so the square of
underlying spectra is Cartesian.  Limits of commutative spectral rings are
created on underlying spectra, so the square is Cartesian in commutative
spectral rings.  The two displayed reconstruction formulas follow.
\end{proof}

\begin{definition}[Relative nilpotent laft property]
\label{def:bnli-nil-laft}
Let $G\to G'$ be a morphism of presheaves on
$\operatorname{Aff}^{\mathrm{sp,nil,+}}_{U/V}$.  We say that it is
\emph{nil-laft} if, for every integer $n\geq0$ and every cofiltered diagram
$S_\alpha$ of nilpotent affine neighbourhoods with limit $S$ such that all
coordinate fibres over $B$ are $(n-1)$-coconnective, the canonical map
\[
\boxed{
    \operatorname*{colim}_\alpha G(S_\alpha)
    \longrightarrow
    G(S)\times_{G'(S)}
    \operatorname*{colim}_\alpha G'(S_\alpha)}
\]
is an equivalence.  For a centred functor with centre map
$h_U^{\mathrm{cent}}\to G$, the phrase ``the centre is nil-laft'' refers to
this relative condition.  No cofiltered presentation is retained as structure.
\end{definition}


\begin{construction}[Relative split tangent spectrum of a nilpotent centre]
\label{con:bnli-relative-split-tangent}
Let
\[
    h_U^{\mathrm{cent}}\longrightarrow G
\]
be a centred presheaf on
$\operatorname{Aff}^{\mathrm{sp,nil,+}}_{U/V}$ with deformation theory, so
$G(U)\simeq *$.  For a connective $B$-module $M$ and $r\geq0$, put
\[
    S_B(M[r])
    :=
    \operatorname{Spec}(B\oplus M[r]).
\]
The split retraction
\[
    S_B(M[r])\longrightarrow U
\]
supplies compatible basepoints.  Define
\[
\begin{aligned}
    \Theta_G^r(M)
    &:=
    \operatorname{fib}_{*}
    \left(
      G(S_B(M[r]))\longrightarrow G(U)
    \right)
    \simeq G(S_B(M[r])),\\
    \Theta_U^r(M)
    &:=
    \operatorname{fib}_{\operatorname{id}_U}
    \left(
      h_U^{\mathrm{cent}}(S_B(M[r]))
      \longrightarrow h_U^{\mathrm{cent}}(U)
    \right).
\end{aligned}
\]
The centre induces a morphism
\[
    \Theta_U^r(M)\longrightarrow\Theta_G^r(M).
\]

The split square-zero identity
\[
    B\oplus M[r]
    \simeq
    B\times_{B\oplus M[r+1]}B
\]
is carried to a pullback by the centred Yoneda presheaf
$h_U^{\mathrm{cent}}$ and by deformation theory of $G$.  After taking the indicated fibres at the split basepoints it
gives canonical loop equivalences in $r$.  Thus the two sequences define
spectra
\[
    \Theta_U^{\mathrm{sp}}(M),
    \qquad
    \Theta_G^{\mathrm{sp}}(M),
\]
and we define the \emph{relative split tangent spectrum} by
\[
\boxed{
    \Theta_{U/G}^{\mathrm{sp}}(M)
    :=
    \operatorname{fib}
    \left(
      \Theta_U^{\mathrm{sp}}(M)
      \longrightarrow
      \Theta_G^{\mathrm{sp}}(M)
    \right).}
\]
This construction is functorial in centred morphisms of deformation
presheaves.
\end{construction}

\begin{lemma}[Nil-laft continuity of the relative split tangent]
\label{lem:bnli-relative-split-tangent-continuity}
Let $U\to G$ have deformation theory and nil-laft centre.  Let $M$ be a
connective eventually coconnective $B$-module and let
\[
    M\simeq\operatorname*{colim}_{\alpha}M_\alpha
\]
be a bounded coherent approximation as in
Lemma~\ref{lem:bnli-bounded-coherent-approximation}, with every
$M_\alpha$ in one common Postnikov interval.  Then
\[
\boxed{
    \operatorname*{colim}_{\alpha}
    \Theta_{U/G}^{\mathrm{sp}}(M_\alpha)
    \xrightarrow{\sim}
    \Theta_{U/G}^{\mathrm{sp}}(M).}
\]
\end{lemma}

\begin{proof}
Fix $r\geq0$.  The affine split square-zero neighbourhoods
\[
    S_\alpha:=S_B(M_\alpha[r])
\]
form a cofiltered geometric diagram with limit
\[
    S:=S_B(M[r]).
\]
Their coordinate fibres over $B$ lie in one common bounded Postnikov range,
so nil-laftness gives a Cartesian square
\[
\begin{tikzcd}
\operatorname*{colim}_\alpha h_U^{\mathrm{cent}}(S_\alpha)
    \arrow[r] \arrow[d] &
h_U^{\mathrm{cent}}(S) \arrow[d]\\
\operatorname*{colim}_\alpha G(S_\alpha)
    \arrow[r] &
G(S).
\end{tikzcd}
\]
Pull this square back along the identity point in
$h_U^{\mathrm{cent}}(U)$ and the compatible
split basepoints in $G(U)\simeq *$.  Filtered colimits of spaces commute with
finite limits, so the resulting square of tangent fibres is again Cartesian:
\[
\begin{tikzcd}
\operatorname*{colim}_\alpha \Theta_U^r(M_\alpha)
    \arrow[r] \arrow[d] &
\Theta_U^r(M) \arrow[d]\\
\operatorname*{colim}_\alpha \Theta_G^r(M_\alpha)
    \arrow[r] &
\Theta_G^r(M).
\end{tikzcd}
\]
Taking the homotopy fibre of the vertical maps and again commuting the
filtered colimit with that finite limit gives
\[
    \operatorname*{colim}_\alpha
    \operatorname{fib}
    \left(
      \Theta_U^r(M_\alpha)\to\Theta_G^r(M_\alpha)
    \right)
    \xrightarrow{\sim}
    \operatorname{fib}
    \left(
      \Theta_U^r(M)\to\Theta_G^r(M)
    \right).
\]
The construction is compatible with the split suspension maps in $r$.
Passing to the associated spectra gives the displayed equivalence.
\end{proof}

\begin{lemma}[Split coefficient base change over an arbitrary neighbourhood]
\label{lem:bnli-split-coefficient-base-change}
Let $A\to B$ be any connective augmented spectral algebra and let $M$ be a
$B$-module, regarded as an $A$-module by restriction.  For every integer
$r\geq0$, let $A\oplus_BM[r]$ denote the split square-zero $A$-algebra with
this restricted action.  Then
\[
\boxed{
    A\oplus_BM[r]
    \simeq
    A\times_B(B\oplus M[r]).}
\]
Consequently, if $P$ is a deformation presheaf on unrestricted nilpotent
neighbourhoods, then
\[
\boxed{
    P(\operatorname{Spec}(A\oplus_BM[r]))
    \simeq
    P(\operatorname{Spec}A)
    \times_{P(U)}
    P(S_B(M[r])).}
\]
In particular, since $P(U)\simeq *$, equivalences for
$\operatorname{Spec}A$ and for $S_B(M[r])$ imply the corresponding
equivalence for $\operatorname{Spec}(A\oplus_BM[r])$.
\end{lemma}

\begin{proof}
The pullback on the first line has underlying spectrum $A\oplus M[r]$.
The multiplication on the square-zero summand is zero and its $A$-action is
the restriction of the $B$-action along $A\to B$, so the pullback is exactly
$A\oplus_BM[r]$ as a commutative spectral algebra.  Passing to affine spectra
turns the ring pullback into a pushout
\[
    \operatorname{Spec}(A\oplus_BM[r])
    \simeq
    \operatorname{Spec}A
    \amalg_U
    S_B(M[r]).
\]
The nilpotent closed immersion in this pushout is
\[
    U\longrightarrow S_B(M[r]),
\]
induced by the augmentation $B\oplus M[r]\to B$.  The reverse map
$S_B(M[r])\to U$ is its split retraction, induced by the section
$B\to B\oplus M[r]$.  Deformation theory of $P$ applied to the nilpotent leg
$U\to S_B(M[r])$ carries the pushout to the displayed pullback of spaces.
\end{proof}

\begin{lemma}[Cartesian colimit fibres for space-valued diagrams]
\label{lem:bnli-cartesian-colimit-fibres}
Let $J$ be a small $\infty$-category and let
\[
    \eta:H\longrightarrow H'
\]
be a natural transformation in $\operatorname{Fun}(J,\mathcal S)$.  Suppose
that for every morphism $j\to k$ in $J$ the naturality square
\[
\begin{tikzcd}
H(j) \arrow[r] \arrow[d] & H(k) \arrow[d]\\
H'(j) \arrow[r] & H'(k)
\end{tikzcd}
\]
is Cartesian.  Then for every $j\in J$ the canonical square
\[
\begin{tikzcd}
H(j) \arrow[r] \arrow[d] & H'(j) \arrow[d]\\
\operatorname*{colim}_{J}H \arrow[r] &
\operatorname*{colim}_{J}H'
\end{tikzcd}
\]
is Cartesian.
\end{lemma}

\begin{proof}
Unstraighten $H$ and $H'$ to left fibrations
\[
    E\longrightarrow J,
    \qquad
    E'\longrightarrow J,
\]
and let $f:E\to E'$ be the map over $J$ induced by $\eta$.  For a point
\[
    e'=(j,b)\in E',
    \qquad b\in H'(j),
\]
put
\[
    K(e')
    :=
    \operatorname{fib}_b(H(j)\to H'(j)).
\]
Transport in the two left fibrations makes $K$ a functor
\[
    K:E'\longrightarrow\mathcal S.
\]
The Cartesian naturality hypothesis says exactly that every transport map of
$K$ is an equivalence.  Hence $K$ factors canonically through the groupoid
completion
\[
    E'\longrightarrow(E')^{\mathrm{gpd}}.
\]
Moreover, $E\to E'$ is the Grothendieck construction of this fibre functor.
After groupoid completion we therefore obtain the Grothendieck construction of
the same locally constant family over the space $(E')^{\mathrm{gpd}}$.  Its
homotopy fibre over the image of $e'=(j,b)$ is precisely $K(e')$.

Groupoid completion of a left fibration computes the colimit of its
straightening, so
\[
    E^{\mathrm{gpd}}
    \simeq
    \operatorname*{colim}_{J}H,
    \qquad
    (E')^{\mathrm{gpd}}
    \simeq
    \operatorname*{colim}_{J}H'.
\]
Thus the homotopy fibre over the image of every $b\in H'(j)$ of
\[
    \operatorname*{colim}_{J}H
    \longrightarrow
    \operatorname*{colim}_{J}H'
\]
is canonically
\[
    \operatorname{fib}_b(H(j)\to H'(j)).
\]
Equivalently, the map
\[
    H(j)
    \longrightarrow
    H'(j)\times_{\operatorname*{colim}H'}\operatorname*{colim}H
\]
is an equivalence.  This is the required Cartesian square.
\end{proof}

\begin{theorem}[First-stage nilpotent recognition]
\label{thm:bnli-first-stage-recognition}
Left Kan extension restricts to an equivalence between
$\operatorname{FMP}^{\mathrm{sp}}_{B/R}$ and the full subcategory of
presheaves on $\operatorname{Aff}^{\mathrm{sp,nil,+}}_{U/V}$ which have
deformation theory and whose centre map is nil-laft.  Its inverse is
restriction along $u$.
\end{theorem}

\begin{proof}
Theorem~\ref{thm:bnli-nil-kan-deformation} gives deformation theory and the
unit equivalence.  It remains first to prove that the centred Yoneda map
\[
    h_U^{\mathrm{cent}}\longrightarrow u_!F
\]
is nil-laft.

Let $\{S_\alpha\}_{\alpha\in K}$ be a cofiltered diagram as in
Definition~\ref{def:bnli-nil-laft}, with limit $S$.  Write $C_\alpha$, $C$,
and $B$ for the corresponding coordinate rings and choose $n$ such that
\[
    K_\alpha:=\operatorname{fib}(C_\alpha\to B)
\]
is $(n-1)$-coconnective for every $\alpha$.  Exactness of filtered colimits in
spectra gives
\[
    K:=\operatorname{fib}(C\to B)
    \simeq
    \operatorname*{colim}_\alpha K_\alpha,
\]
so $K$ is also $(n-1)$-coconnective.  Because every map $C_\alpha\to C$ lies
over the same target $B$,
\[
    \operatorname{fib}(C_\alpha\to C)
    \simeq
    \operatorname{fib}(K_\alpha\to K)
\]
is $(n-1)$-coconnective.  Lemma~\ref{lem:bnli-relative-postnikov-reconstruction}
therefore gives
\[
\boxed{
    S_\alpha
    \simeq
    (S_\alpha)_{\le n}
    \amalg_{S_{\le n}}S.}
\]

Let
\[
    J:=(\!\int F)^{\mathrm{op}}.
\]
For $j=(T,\xi)\in J$ define
\[
\begin{aligned}
    H(j)
    &:=
    \operatorname*{colim}_\alpha
    \operatorname{Map}_{U/}(S_\alpha,T),\\
    H'(j)
    &:=
    \operatorname{Map}_{U/}(S,T).
\end{aligned}
\]
The natural maps $S\to S_\alpha$ define a natural transformation
$H\to H'$.  By the Artinian factorization formula and Fubini for colimits,
\[
\boxed{
\begin{aligned}
    \operatorname*{colim}_{j\in J}H(j)
    &\simeq
    \operatorname*{colim}_\alpha u_!F(S_\alpha),\\
    \operatorname*{colim}_{j\in J}H'(j)
    &\simeq
    u_!F(S).
\end{aligned}}
\]

We verify the hypothesis of
Lemma~\ref{lem:bnli-cartesian-colimit-fibres}.  Let $j\to j'$ correspond to a
morphism of Artinian intermediaries
\[
    T\longrightarrow T'.
\]
Lemma~\ref{lem:bnli-artinian-morphism-afp} and
Lemma~\ref{lem:bnli-affine-afp-laft} show that $T\to T'$ is spectrally laft.
The relative Postnikov reconstruction above is a pushout in the ambient
geometric category,
\[
    S_\alpha
    \simeq
    (S_\alpha)_{\leq n}\amalg_{S_{\leq n}}S.
\]
Therefore, without introducing a nonexistent map $S\to T$, its mapping
universal property gives
\[
\begin{aligned}
    \operatorname{Map}_V(S_\alpha,T)
    &\simeq
    \operatorname{Map}_V(S,T)
    \times_{\operatorname{Map}_V(S_{\leq n},T)}
    \operatorname{Map}_V((S_\alpha)_{\leq n},T),\\
    \operatorname{Map}_V(S_\alpha,T')
    &\simeq
    \operatorname{Map}_V(S,T')
    \times_{\operatorname{Map}_V(S_{\leq n},T')}
    \operatorname{Map}_V((S_\alpha)_{\leq n},T').
\end{aligned}
\]
The coordinate algebras of $\{(S_\alpha)_{\leq n}\}$ form a filtered diagram
of $n$-truncated algebras with colimit the coordinate algebra of
$S_{\leq n}$.  Spectral laftness of $T\to T'$ therefore says exactly that
\[
\begin{tikzcd}
\operatorname*{colim}_\alpha
\operatorname{Map}_V((S_\alpha)_{\leq n},T)
    \arrow[r] \arrow[d] &
\operatorname{Map}_V(S_{\leq n},T) \arrow[d]\\
\operatorname*{colim}_\alpha
\operatorname{Map}_V((S_\alpha)_{\leq n},T')
    \arrow[r] &
\operatorname{Map}_V(S_{\leq n},T')
\end{tikzcd}
\]
is Cartesian.  Filtered colimits of spaces commute with the finite pullbacks
in the two mapping formulas.  Substitution consequently gives a Cartesian
square
\[
\begin{tikzcd}
\operatorname*{colim}_\alpha\operatorname{Map}_V(S_\alpha,T)
    \arrow[r] \arrow[d] &
\operatorname{Map}_V(S,T) \arrow[d]\\
\operatorname*{colim}_\alpha\operatorname{Map}_V(S_\alpha,T')
    \arrow[r] &
\operatorname{Map}_V(S,T').
\end{tikzcd}
\]
All Artinian intermediaries and all transition maps are under the fixed centre
$U$.  Taking the compatible homotopy fibres over the prescribed maps
$U\to T$ and $U\to T'$ therefore produces the Cartesian naturality square
\[
\begin{tikzcd}
H(j) \arrow[r] \arrow[d] & H(j') \arrow[d]\\
H'(j) \arrow[r] & H'(j').
\end{tikzcd}
\]
required by Lemma~\ref{lem:bnli-cartesian-colimit-fibres}.

Apply Lemma~\ref{lem:bnli-cartesian-colimit-fibres} to the distinguished
Artinian identity element $(U,*)\in J$.  Since
\[
    H(U,*)=
    \operatorname*{colim}_\alpha h_U^{\mathrm{cent}}(S_\alpha),
    \qquad
    H'(U,*)=h_U^{\mathrm{cent}}(S),
\]
we obtain a Cartesian square
\[
\begin{tikzcd}
\operatorname*{colim}_\alpha h_U^{\mathrm{cent}}(S_\alpha)
    \arrow[r] \arrow[d] &
h_U^{\mathrm{cent}}(S) \arrow[d]\\
\operatorname*{colim}_\alpha u_!F(S_\alpha)
    \arrow[r] &
u_!F(S).
\end{tikzcd}
\]
This is exactly nil-laftness of the centred Yoneda map
$h_U^{\mathrm{cent}}\to u_!F$.

We turn to the converse.  Let $G$ have deformation theory and nil-laft
centre and put
\[
    H:=u_!u^*G.
\]
The counit
\[
    e:H\longrightarrow G
\]
is an equivalence on the Artinian subcategory.  By the first half of the
proof, both centred Yoneda maps
\[
    h_U^{\mathrm{cent}}\longrightarrow H,
    \qquad
    h_U^{\mathrm{cent}}\longrightarrow G
\]
are nil-laft.

Let $M$ be a connective eventually coconnective $B$-module.  Choose a
bounded coherent approximation
\[
    M\simeq\operatorname*{colim}_\alpha M_\alpha
\]
as in Lemma~\ref{lem:bnli-bounded-coherent-approximation}.  For every
$\alpha$ and every $r\geq0$, the split square-zero augmentation
\[
    B\oplus M_\alpha[r]\longrightarrow B
\]
is spectral Artinian.  Hence the Artinian counit gives an equivalence
\[
    \Theta_H^{\mathrm{sp}}(M_\alpha)
    \xrightarrow{\sim}
    \Theta_G^{\mathrm{sp}}(M_\alpha)
\]
and therefore an equivalence of the relative tangent fibres
\[
    \Theta_{U/H}^{\mathrm{sp}}(M_\alpha)
    \xrightarrow{\sim}
    \Theta_{U/G}^{\mathrm{sp}}(M_\alpha).
\]
Lemma~\ref{lem:bnli-relative-split-tangent-continuity}, applied to the two
nil-laft centres, gives
\[
\begin{aligned}
    \Theta_{U/H}^{\mathrm{sp}}(M)
    &\simeq
    \operatorname*{colim}_\alpha
    \Theta_{U/H}^{\mathrm{sp}}(M_\alpha)\\
    &\xrightarrow{\sim}
    \operatorname*{colim}_\alpha
    \Theta_{U/G}^{\mathrm{sp}}(M_\alpha)
    \simeq
    \Theta_{U/G}^{\mathrm{sp}}(M).
\end{aligned}
\]
Compare the two fibre sequences of spectra
\[
\begin{tikzcd}
\Theta_{U/H}^{\mathrm{sp}}(M)
    \arrow[r] \arrow[d,"\sim"] &
\Theta_U^{\mathrm{sp}}(M)
    \arrow[r] \arrow[d,equal] &
\Theta_H^{\mathrm{sp}}(M)
    \arrow[d]\\
\Theta_{U/G}^{\mathrm{sp}}(M)
    \arrow[r] &
\Theta_U^{\mathrm{sp}}(M)
    \arrow[r] &
\Theta_G^{\mathrm{sp}}(M).
\end{tikzcd}
\]
Stability gives
\[
\boxed{
    \Theta_H^{\mathrm{sp}}(M)
    \xrightarrow{\sim}
    \Theta_G^{\mathrm{sp}}(M).}
\]
In particular, for every $r\geq0$,
\[
\boxed{
    H(S_B(M[r]))
    \xrightarrow{\sim}
    G(S_B(M[r])).}
\]

We propagate this split comparison to structured square-zero stages over an
arbitrary intermediate neighbourhood.  Suppose
\[
    A'\longrightarrow A\longrightarrow B
\]
is a structured square-zero extension whose coefficient is the restriction
of a connective eventually coconnective $B$-module $M$, and suppose
$H(\operatorname{Spec}A)\to G(\operatorname{Spec}A)$ is already an
equivalence.  The classifying derivation expresses the extension as
\[
    A'
    \simeq
    A\times_{A\oplus_BM[1]}A.
\]
By Lemma~\ref{lem:bnli-split-coefficient-base-change},
\[
\begin{aligned}
H(\operatorname{Spec}(A\oplus_BM[1]))
&\simeq
H(\operatorname{Spec}A)
\times_{H(U)}
H(S_B(M[1])),\\
G(\operatorname{Spec}(A\oplus_BM[1]))
&\simeq
G(\operatorname{Spec}A)
\times_{G(U)}
G(S_B(M[1])).
\end{aligned}
\]
The right-hand terms correspond by equivalences, so the middle comparison
is an equivalence.  Applying deformation theory to the structured
square-zero pullback for $A'$ now gives
\[
    H(\operatorname{Spec}A')
    \xrightarrow{\sim}
    G(\operatorname{Spec}A').
\]

Finally let $A\to B$ be any unrestricted nilpotent neighbourhood.
Proposition~\ref{prop:bnli-unrestricted-finite-square-zero}, in its bounded
coefficient form, gives a finite factorization
\[
    A=A_m\longrightarrow A_{m-1}\longrightarrow\cdots
    \longrightarrow A_0=B
\]
whose square-zero coefficients are restrictions of connective eventually
coconnective $B$-modules.  Starting from
$H(U)\simeq G(U)\simeq *$ and applying the preceding propagation step
inductively gives
\[
    H(\operatorname{Spec}A_i)
    \xrightarrow{\sim}
    G(\operatorname{Spec}A_i)
\]
for every $i$, in particular for $A_m=A$.  Thus the counit is an
equivalence on every unrestricted nilpotent neighbourhood.  Restriction and
left Kan extension are inverse on the stated full subcategories.
\end{proof}

\begin{proposition}[Ambient deformation locality and spectral laftness imply nil-laftness]
\label{prop:bnli-spectral-laft-implies-nil-laft}
Let
\[
    u:U\longrightarrow Y
\]
be a spectrally laft centred prestack over $V$.  Assume that $Y$ is
spectrally convergent and satisfies ambient nilpotent excision on connective
affine tests; in particular it is enough that $Y$ be
$L_V^{\mathrm{Def}}$-local.  Then the dotted-lift restriction
\[
    h_U^{\mathrm{cent}}
    \longrightarrow
    \operatorname{Lift}^{\mathrm{nil}}_U(Y)
\]
is nil-laft in the sense of Definition~\ref{def:bnli-nil-laft}.
\end{proposition}

\begin{proof}
Let $\{S_\alpha=\operatorname{Spec}C_\alpha\}$ be a cofiltered diagram of
unrestricted nilpotent neighbourhoods with limit
$S=\operatorname{Spec}C$, and suppose
\[
    K_\alpha:=\operatorname{fib}(C_\alpha\to B)
\]
is $(n-1)$-coconnective for every $\alpha$.  Exactness of filtered colimits
in spectra gives
\[
    K:=\operatorname{fib}(C\to B)
    \simeq
    \operatorname*{colim}_\alpha K_\alpha,
\]
so $K$ is $(n-1)$-coconnective.  Because all maps lie over $B$,
\[
    \operatorname{fib}(C_\alpha\to C)
    \simeq
    \operatorname{fib}(K_\alpha\to K)
\]
is also $(n-1)$-coconnective.  Put
\[
    C_n:=\tau_{\leq n}C,
    \qquad
    C_{\alpha,n}:=\tau_{\leq n}C_\alpha.
\]
Lemma~\ref{lem:bnli-relative-postnikov-reconstruction} gives
\[
\boxed{
    C_\alpha
    \simeq
    C\times_{C_n}C_{\alpha,n}.}
\]

The map $C\to C_n$ is an isomorphism on $\pi_0$, hence is an allowed
nilpotent-excision leg.  Ambient nilpotent excision for $Y$ gives
\[
\boxed{
    Y(C_\alpha)
    \simeq
    Y(C)\times_{Y(C_n)}Y(C_{\alpha,n}).}
\]
The represented prestack $U$ satisfies the same pullback equation.  Since
Postnikov truncation commutes with filtered colimits of connective spectral
algebras,
\[
    C_n\simeq\operatorname*{colim}_\alpha C_{\alpha,n}.
\]
Spectral laftness of the centre $U\to Y$, applied to this uniformly
$n$-truncated diagram, gives
\[
\boxed{
    \operatorname*{colim}_\alpha U(C_{\alpha,n})
    \simeq
    U(C_n)\times_{Y(C_n)}
    \operatorname*{colim}_\alpha Y(C_{\alpha,n}).}
\]
Filtered colimits of spaces commute with finite limits.  Substituting the last
equivalence into the pullback formulas yields
\[
\boxed{
    \operatorname*{colim}_\alpha U(C_\alpha)
    \simeq
    U(C)\times_{Y(C)}
    \operatorname*{colim}_\alpha Y(C_\alpha).}
\]

It remains to pass from ambient mapping spaces to the centred presheaves.  The
fixed neighbourhood structures $U\to S_\alpha$ and $U\to S$ identify
\[
\begin{aligned}
    h_U^{\mathrm{cent}}(S_\alpha)
    &\simeq
    \operatorname{fib}_{\operatorname{id}_U}
    \left(U(C_\alpha)\to U(B)\right),\\
    \operatorname{Lift}^{\mathrm{nil}}_U(Y)(S_\alpha)
    &\simeq
    \operatorname{fib}_{u}
    \left(Y(C_\alpha)\to Y(B)\right),
\end{aligned}
\]
and similarly at $S$.  Taking these compatible homotopy fibres in the
preceding Cartesian equivalence, and commuting the filtered colimit with the
finite fibre, gives
\[
\boxed{
\begin{aligned}
    \operatorname*{colim}_\alpha h_U^{\mathrm{cent}}(S_\alpha)
    \simeq{}&
    h_U^{\mathrm{cent}}(S)
    \times_{\operatorname{Lift}^{\mathrm{nil}}_U(Y)(S)}\\
    &\operatorname*{colim}_\alpha
    \operatorname{Lift}^{\mathrm{nil}}_U(Y)(S_\alpha).
\end{aligned}}
\]
This is exactly nil-laftness of the dotted-lift centre.
\end{proof}

\subsection{Second-stage geometric materialization}

\subsubsection{Convergent prestacks from eventually coconnective tests}

Let
\[
    j:
    \operatorname{CAlg}^{\mathrm{sp,cn},<\infty}_{/R}
    \hookrightarrow
    \operatorname{CAlg}^{\mathrm{sp,cn}}_{/R}
\]
be the inclusion of eventually coconnective connective spectral $R$-algebras.

\begin{proposition}[Bounded-test model for convergent spectral prestacks]
\label{prop:bnli-convergent-bounded-tests}
Restriction along $j$ has a fully faithful right adjoint
\[
    j_*:
    \operatorname{Fun}
    \left(
      \operatorname{CAlg}^{\mathrm{sp,cn},<\infty}_{/R},
      \mathcal S
    \right)
    \longrightarrow
    \operatorname{PStk}^{\mathrm{sp}}_{/V},
\]
and for every functor $H$ on eventually coconnective tests there is a
canonical formula
\[
\boxed{
    (j_*H)(C)
    \simeq
    \lim_{n\geq0}H(\tau_{\leq n}C).}
\]
The essential image of $j_*$ is precisely the full subcategory of spectrally
convergent prestacks.  Consequently
$\operatorname{PStk}^{\mathrm{sp,conv}}_{/V}$ is presentable, and its
colimits are computed pointwise on eventually coconnective affine tests,
followed by the displayed right Kan extension.
\end{proposition}

\begin{proof}
For a connective $C$, the right Kan extension formula is the limit over
bounded targets $C\to D$.  If $D$ is $n$-truncated, every map $C\to D$
factors through the universal truncation
\[
    C\longrightarrow\tau_{\leq n}C.
\]
Thus restriction to the Postnikov truncations is initial for the limit in the
right-Kan indexing category: for every bounded target, factorization through
the appropriate truncation is contractibly unique.  The right-Kan limit is
therefore computed by the Postnikov tower, giving the displayed formula.  If $C$ is itself
eventually coconnective, the Postnikov tower is eventually constant, so the
unit
\[
    H\longrightarrow j^*j_*H
\]
is an equivalence.  Hence $j_*$ is fully faithful.

A prestack $Y$ lies in the essential image if and only if
\[
    Y(C)\xrightarrow{\sim}\lim_nY(\tau_{\leq n}C)
\]
for every connective $C$, which is exactly spectral convergence.  Presentability
and the colimit statement now follow from the equivalence with the functor
category on bounded tests.
\end{proof}

\begin{definition}[Spectral infinitesimal prestack of the centre]
\label{def:bnli-inf-prestack}
Define
\[
\boxed{
    U^{\mathrm{inf,sp}}
    :=
    \mathsf P^{\mathrm{sp}}U.}
\]
Equivalently,
\[
    U^{\mathrm{inf,sp}}(C)
    =
    U(C^{\mathrm{prored}}).
\]
This prestack is spectrally convergent because its value depends only on
$\pi_0(C)$ and the filtered nilpotent-quotient system.  The canonical map
\[
    U\longrightarrow U^{\mathrm{inf,sp}}
\]
is the pro-reduction coaugmentation and is a spectral nil-isomorphism.
\end{definition}

Let
\[
    \operatorname{PStk}^{\mathrm{sp,conv}}_{/V,U^{\mathrm{inf,sp}}}
\]
denote the $\infty$-category of spectrally convergent prestacks over
$U^{\mathrm{inf,sp}}$, and let
\[
    \left(
      \operatorname{PStk}^{\mathrm{sp,conv}}_{/V,U^{\mathrm{inf,sp}}}
    \right)_{U/}
\]
be the category of diagrams
\[
    U\longrightarrow Y\longrightarrow U^{\mathrm{inf,sp}}
\]
whose composite is the pro-reduction coaugmentation.

\begin{definition}[Bounded pro-pushout associated to a nilpotent neighbourhood]
\label{def:bnli-pro-pushout}
Let
\[
    U=\operatorname{Spec}B
    \longrightarrow
    T=\operatorname{Spec}C
\]
be an unrestricted nilpotent neighbourhood with $C$ eventually
coconnective, and put
\[
    K:=\ker(\pi_0C\to\pi_0B).
\]
Let $\mathcal I_{\mathrm{nil}}(C;K)$ be the cofinal filtered subposet of
nilpotent ideals $I\subseteq\pi_0C$ containing $K$, and set
\[
    T_I:=\operatorname{Spec}H(\pi_0C/I).
\]
The inclusion $K\subseteq I$ gives canonical maps
\[
    T_I\longrightarrow U,
    \qquad
    T_I\longrightarrow T.
\]
Define the \emph{pro-pushout neighbourhood} by the filtered pro-system
\[
\boxed{
    U\amalg_{T^{\mathrm{prored}}}T
    :=
    \left\{
      U\amalg_{T_I}T
    \right\}_{I\in\mathcal I_{\mathrm{nil}}(C;K)}.}
\]
For a presheaf $G$ on unrestricted nilpotent neighbourhoods, write
\[
\boxed{
    G(U\amalg_{T^{\mathrm{prored}}}T)
    :=
    \operatorname*{colim}_{I\in\mathcal I_{\mathrm{nil}}(C;K)}
    G(U\amalg_{T_I}T).}
\]
Each pushout is an affine nilpotent neighbourhood of $U$: on coordinate
rings it is the pullback
\[
    P_I
    :=
    B\times_{H(\pi_0C/I)}C.
\]
It is connective because $C\to H(\pi_0C/I)$ is surjective on $\pi_0$.
Moreover, writing $R=\pi_0C$ and using
$\pi_0B\simeq R/K$, one has
\[
    \pi_0P_I
    \simeq
    (R/K)\times_{R/I}R
\]
and therefore
\[
\boxed{
    \ker(\pi_0P_I\to\pi_0B)
    \simeq I.}
\]
Thus the degree-zero kernel is nilpotent, and pullback of the underlying
spectra gives
\[
    \operatorname{fib}(P_I\to B)
    \simeq
    \operatorname{fib}
    \left(
      C\to H(\pi_0C/I)
    \right).
\]
This fibre is eventually coconnective because $C$ is.  The absolute
eventual-coconnectivity hypothesis on $C$ is used only for this pro-pushout
construction; an unrestricted nilpotent neighbourhood need not satisfy it.
\end{definition}



\begin{lemma}[Pointed filtered-fibre formula]
\label{lem:bnli-pointed-filtered-fibre}
Let $I$ be filtered and let $E_i\to B_i$ be a natural transformation of
$I$-diagrams of spaces.  For
\[
    b\in\operatorname*{colim}_{i\in I}B_i
\]
there is a canonical equivalence
\[
\boxed{
    \operatorname*{colim}_{(i,b_i,\gamma)\in\operatorname{Rep}_B(b)}
    \operatorname{fib}_{b_i}(E_i\to B_i)
    \xrightarrow{\sim}
    \operatorname{fib}_{b}
    \left(
      \operatorname*{colim}_iE_i
      \longrightarrow
      \operatorname*{colim}_iB_i
    \right).}
\]
The equivalence is natural in the diagram and in $b$.
\end{lemma}

\begin{proof}
A point of the right-hand fibre is represented by a stage $i$, a point
$e_i\in E_i$, and a path from the image of its base point in the filtered
colimit of the $B_i$ to $b$.  This is exactly an object of the Grothendieck
construction of the left-hand fibre diagram over
$\operatorname{Rep}_B(b)$.  A morphism or higher simplex on either side is
represented after a common later stage because the indexing category is
filtered and the relevant finite simplicial set is compact.  Hence the two
Grothendieck constructions have equivalent groupoid completions.  Groupoid
completion computes the displayed colimit of fibres on the left and the
homotopy fibre on the right, proving the formula.
\end{proof}

\begin{lemma}[Pointed pro-reduction fibre formula]
\label{lem:bnli-pointed-prored-fibre}
Let $U$ be a spectral prestack, let
\[
    Y\longrightarrow\mathsf P^{\mathrm{sp}}U
\]
be a morphism, let $C$ be eventually coconnective, and let
\[
    x\in(\mathsf P^{\mathrm{sp}}U)(C).
\]
For a nilpotent ideal $I\subseteq\pi_0C$ put
$T_I=\operatorname{Spec}H(\pi_0C/I)$ and define the filtered diagram
$D_U(I):=U(T_I)$.  Write
\[
    \operatorname{Rep}_{U,C}(x)
    :=
    \operatorname{Rep}_{D_U}(x).
\]
For $r=(I,x_I)\in\operatorname{Rep}_{U,C}(x)$ let
$\bar x_I\in(\mathsf P^{\mathrm{sp}}U)(T_I)$ be the pro-reduction image of
$x_I$.  Then
\[
\boxed{
\operatorname*{colim}_{r\in\operatorname{Rep}_{U,C}(x)}
\operatorname{fib}_{\bar x_I}
\left(
  Y(T_I)\longrightarrow(\mathsf P^{\mathrm{sp}}U)(T_I)
\right)
\simeq
\operatorname{fib}_{x}
\left(
  (\mathsf P^{\mathrm{sp}}Y)(C)
  \longrightarrow
  (\mathsf P^{\mathrm{sp}}U)(C)
\right),}
\]
where the map on the right is obtained by applying
$\mathsf P^{\mathrm{sp}}$ and then the idempotence equivalence
$(\mathsf P^{\mathrm{sp}})^2U\simeq\mathsf P^{\mathrm{sp}}U$.
\end{lemma}

\begin{proof}
Expand the target of the right-hand morphism before using idempotence:
\[
    (\mathsf P^{\mathrm{sp}})^2U(C)
    \simeq
    \operatorname*{colim}_{(I\subseteq K)\in
      \operatorname{Flag}_{\mathrm{nil}}(\pi_0C)}
    U(H(\pi_0C/K)).
\]
An object representing $x$ in the iterated target may first be represented at
an explicit flag $(I\subseteq K)$ and a point of $U(H(\pi_0C/K))$.  Its
category of choices of such an inner representative is filtered and
contractible by Lemma~\ref{lem:bnli-filtered-representatives}.  Refining the
outer ideal from $I$ to $K$ sends it to the diagonal flag $(K\subseteq K)$.
Define the diagonal representative functor from ordinary representatives of
$x$ in $\operatorname*{colim}_K U(H(\pi_0C/K))$ to representatives in the
iterated colimit by this diagonal flag.  Every object of the latter category
admits a morphism to an object in the image, and every finite diagram in the
corresponding comma category admits a common diagonal refinement.  The comma
categories are therefore nonempty and filtered, hence have contractible
groupoid completion.  The diagonal representative functor is final.  Under the idempotence equivalence the diagonal
representative is precisely the pro-reduction image of a representative
$x_K\in U(T_K)$.  Thus the representative category for $x$ in the target
colimit may be replaced cofinally by $\operatorname{Rep}_{U,C}(x)$.

Now apply Lemma~\ref{lem:bnli-pointed-filtered-fibre} to the natural
transformation
\[
    I\longmapsto Y(T_I)
    \longrightarrow
    I\longmapsto(\mathsf P^{\mathrm{sp}}U)(T_I).
\]
The colimit of the source is $(\mathsf P^{\mathrm{sp}}Y)(C)$; the colimit of
the target is $(\mathsf P^{\mathrm{sp}})^2U(C)$, which is canonically
$(\mathsf P^{\mathrm{sp}}U)(C)$.  The pointed filtered-fibre formula is
therefore exactly the displayed equivalence.
\end{proof}

\begin{lemma}[Simultaneous representatives for a bounded nilpotent-excision square]
\label{lem:bnli-simultaneous-square-representatives}
Let
\[
\begin{tikzcd}
C \arrow[r] \arrow[d] & A_0 \arrow[d]\\
A_1 \arrow[r] & A_{01}
\end{tikzcd}
\]
be a pullback of connective spectral rings in which
$\pi_0(A_0)\twoheadrightarrow\pi_0(A_{01})$ has nilpotent kernel, and let
$x_\bullet$ be a compatible point of the corresponding pullback square of
spaces $(\mathsf P^{\mathrm{sp}}U)(-)$.  There is a filtered weakly
contractible $\infty$-category
\[
    \operatorname{Rep}^{\square}_U(x_\bullet)
\]
of compatible four-vertex nilpotent representatives with the following
additional property.  If
\[
    D_C,
    \quad D_0,
    \quad D_1,
    \quad D_{01}
\]
are their discrete quotient rings, then the two maps along the nilpotent legs
are equivalences
\[
\boxed{
    D_0\xrightarrow{\sim}D_{01},
    \qquad
    D_C\xrightarrow{\sim}D_1,}
\]
and the four $D$'s form a Cartesian ring square.  The projection from
$\operatorname{Rep}^{\square}_U(x_\bullet)$ to the ordinary representative
category at each vertex is cofinal.
\end{lemma}

\begin{proof}
Write
\[
    R_C=\pi_0C,
    \quad R_i=\pi_0A_i,
    \quad R_{01}=\pi_0A_{01},
\]
and let
\[
    K_0=\ker(R_0\to R_{01}),
    \qquad
    K_C=\ker(R_C\to R_1).
\]
The second kernel is nilpotent because $C\to A_1$ is the base change of the
nilpotent surjection $A_0\to A_{01}$.  Start with finitely many chosen representatives of the four points.  Let
$J_1$ be the finite sum in $R_1$ of all prescribed ideals at the $A_1$-vertex
and the ideals generated by the images of all prescribed ideals at the
$C$-vertex.  Let $J_{01}$ be the finite sum in $R_{01}$ of all prescribed
ideals at the terminal vertex, the ideals generated by the images of all
prescribed ideals at the $A_0$-vertex, and the image of $J_1$.  Finite sums
and images of nilpotent ideals are nilpotent, so $J_1$ and $J_{01}$ are
nilpotent.

Now put
\[
    J_0:=(R_0\to R_{01})^{-1}(J_{01}),
    \qquad
    J_C:=(R_C\to R_1)^{-1}(J_1).
\]
The kernels $K_0$ and $K_C$ are nilpotent.  An extension of a nilpotent ideal
by a nilpotent kernel is nilpotent, so $J_0$ and $J_C$ are nilpotent.  By
construction $J_0$ and $J_C$ refine every prescribed ideal at their
respective vertices.  If $c\in J_C$, its image in $R_1$ lies in $J_1$, hence
its image in $R_{01}$ lies in $J_{01}$.  The pullback relation then implies
that its image in $R_0$ belongs to $J_0$.  Thus all four quotient maps are
defined compatibly.

The resulting discrete quotients satisfy
\[
    R_0/J_0\xrightarrow{\sim}R_{01}/J_{01},
    \qquad
    R_C/J_C\xrightarrow{\sim}R_1/J_1.
\]
The horizontal maps commute, so the square of quotient rings is Cartesian:
with both vertical maps equivalences it is identified with the same
horizontal map on the top and bottom.  Passing to Eilenberg--Mac Lane rings
gives the asserted square of the $D$'s.

The compatible point $x_\bullet$ supplies finitely many paths and higher
homotopies asserting compatibility of the represented points after passage
to the four filtered colimits.  Compactness of their finite simplicial domain
moves all of them to one further system of nilpotent refinements of the form
just constructed.  Hence compatible representatives with the displayed
vertical equivalences exist.  Applying the construction to any finite diagram
of such representatives gives a common cone, so
$\operatorname{Rep}^{\square}_U(x_\bullet)$ is filtered and weakly
contractible.

For cofinality at a chosen vertex, begin with an arbitrary representative at
that vertex and include its ideal and finite path data in the prescribed
finite family above.  The construction produces a simultaneous
nilpotent-leg-saturated refinement.  The same procedure for a finite diagram
in the comma category gives a cone, so each comma category is filtered and
nonempty.  The projection to every individual representative category is
therefore cofinal.
\end{proof}

\begin{definition}[Bounded pointed pro-pushout]
\label{def:bnli-pointed-pro-pushout}
Let
\[
    x:T=\operatorname{Spec}C
    \longrightarrow
    U^{\mathrm{inf,sp}}
\]
be a point with $C$ eventually coconnective.  Let
$\operatorname{Rep}_{U}(x)$ denote the representative category
$\operatorname{Rep}_{U,C}(x)$ of
Lemma~\ref{lem:bnli-pointed-prored-fibre}, equivalently the $\infty$-category
of representatives of $x$ in the filtered-colimit presentation
\[
    U^{\mathrm{inf,sp}}(C)
    =
    \operatorname*{colim}_{I\in\mathcal I_{\mathrm{nil}}(C)}
    U\!\left(H(\pi_0C/I)\right).
\]
Thus an object consists of a nilpotent ideal
$I\subseteq\pi_0C$, a map
\[
    x_I:
    T_I:=\operatorname{Spec}H(\pi_0C/I)
    \longrightarrow U,
\]
and a path from the image of $x_I$ in the displayed filtered colimit to
$x$.  Morphisms are refinements of the nilpotent quotient together with the
compatible homotopies.

Lemma~\ref{lem:bnli-filtered-representatives} makes
$\operatorname{Rep}_{U}(x)$ filtered, hence weakly contractible.  In
particular every finite compatible collection of representative points,
paths, and higher homotopies admits a common refinement.  After any one
representative has been chosen, the full subsystem of its refinements is
cofinal.  For
$r=(I,x_I)\in\operatorname{Rep}_{U}(x)$ define
\[
\boxed{
    P_{x,r}(T)
    :=
    U\amalg_{T_I}T.}
\]
For a presheaf $G$ on unrestricted nilpotent neighbourhoods set
\[
\boxed{
    G(P_x(T))
    :=
    \operatorname*{colim}_{r\in\operatorname{Rep}_{U}(x)}
    G(P_{x,r}(T)).}
\]
Every $P_{x,r}(T)$ is an unrestricted nilpotent neighbourhood of $U$.  On
coordinate rings it is
\[
    B\times_{H(\pi_0C/I)}C.
\]
The map to $B$ is surjective on $\pi_0$ with nilpotent kernel, and its fibre
is eventually coconnective because $C$ is eventually coconnective.

If $U\to T$ is itself an eventually coconnective unrestricted nilpotent
neighbourhood and $x$ is the canonical composite
\[
    T\longrightarrow\mathsf P^{\mathrm{sp}}T
    \xrightarrow{\sim}U^{\mathrm{inf,sp}},
\]
where the second arrow is the coherent inverse supplied by
Lemma~\ref{lem:bnli-representable-prored-thickening}, then the representatives coming from ideals containing
$\ker(\pi_0C\to\pi_0B)$ form a cofinal subsystem and
\[
    G(P_x(T))
    \simeq
    G(U\amalg_{T^{\mathrm{prored}}}T)
\]
in the notation of Definition~\ref{def:bnli-pro-pushout}.
\end{definition}

\begin{construction}[Second-stage materialization adjunction]
\label{con:bnli-second-stage-adjunction}
For an unrestricted nilpotent neighbourhood $U\to T$, the coaugmentation
\[
    T\longrightarrow\mathsf P^{\mathrm{sp}}T
\]
followed by the coherent inverse of the natural equivalence of
Lemma~\ref{lem:bnli-representable-prored-thickening},
\[
    \mathsf P^{\mathrm{sp}}T
    \xrightarrow{\sim}
    \mathsf P^{\mathrm{sp}}U
    =
    U^{\mathrm{inf,sp}},
\]
defines a centred diagram
\[
    U\longrightarrow T\longrightarrow U^{\mathrm{inf,sp}}.
\]
Represented spectral prestacks are convergent.  Hence the preceding assignment
defines a functor from
$\operatorname{Aff}^{\mathrm{sp,nil,+}}_{U/V}$ to the presentable category
\[
    \left(
      \operatorname{PStk}^{\mathrm{sp,conv}}_{/V,U^{\mathrm{inf,sp}}}
    \right)_{U/}.
\]
By the free-cocompletion universal property of the presheaf category, it
extends uniquely to a colimit-preserving functor
\[
\boxed{
    \widetilde v_!:
    \operatorname{Fun}
    \left(
      (\operatorname{Aff}^{\mathrm{sp,nil,+}}_{U/V})^{\mathrm{op}},
      \mathcal S
    \right)
    \longrightarrow
    \left(
      \operatorname{PStk}^{\mathrm{sp,conv}}_{/V,U^{\mathrm{inf,sp}}}
    \right)_{U/}.}
\]
Let $\widetilde v^*$ be its right adjoint.
\end{construction}

\begin{proposition}[Dotted-lift and bounded pointed pro-pushout formulas]
\label{prop:bnli-dotted-lift-formula}
For a diagram
\[
    U\longrightarrow Y\longrightarrow U^{\mathrm{inf,sp}}
\]
and an unrestricted nilpotent neighbourhood $U\to T$, the value
$\widetilde v^*Y(T)$ is canonically the space of dotted lifts
\[
\begin{tikzcd}
U \arrow[r] \arrow[d] & Y \arrow[d]\\
T \arrow[r] & U^{\mathrm{inf,sp}}.
\end{tikzcd}
\]

More generally, let
\[
    x:T=\operatorname{Spec}C
    \longrightarrow U^{\mathrm{inf,sp}}
\]
be any eventually coconnective affine test.  For every presheaf $G$ on
unrestricted nilpotent neighbourhoods there is a canonical equivalence
\[
\boxed{
    \operatorname{Map}_{/U^{\mathrm{inf,sp}}}
    (T,\widetilde v_!G)_x
    \xrightarrow{\sim}
    G(P_x(T)),}
\]
where the left-hand side denotes maps whose composite to
$U^{\mathrm{inf,sp}}$ is the fixed point $x$.  In particular, for a bounded
nilpotent neighbourhood equipped with its canonical infinitesimal point,
\[
\boxed{
    \operatorname{Map}_{/U^{\mathrm{inf,sp}}}
    (T,\widetilde v_!G)
    \simeq
    G(U\amalg_{T^{\mathrm{prored}}}T).}
\]
\end{proposition}

\begin{proof}
The dotted-lift description is the adjunction formula evaluated on the
representable nilpotent neighbourhood $T$.

Fix now the bounded test $x:T\to U^{\mathrm{inf,sp}}$.  Both sides of the
pointed formula preserve colimits in $G$.  On the left, evaluation on an
eventually coconnective test is pointwise in the bounded-test model of
Proposition~\ref{prop:bnli-convergent-bounded-tests}.  Taking the fibre over
the fixed point $x$ is pullback in the slice of spaces over
$U^{\mathrm{inf,sp}}(T)$; such pullback preserves colimits.  On the right,
colimit preservation is immediate from the defining filtered colimit.  It therefore suffices to take $G$ represented by an
unrestricted nilpotent neighbourhood $U\to T'$.

The left-hand side then classifies maps
\[
    g:T\longrightarrow T'
\]
whose composite with $T'\to U^{\mathrm{inf,sp}}$ is $x$.  The represented
nilpotent thickening $U\to T'$ has a canonical pro-reduction equivalence
\[
\boxed{
    \mathsf P^{\mathrm{sp}}T'
    \xrightarrow{\sim}
    \mathsf P^{\mathrm{sp}}U
    =U^{\mathrm{inf,sp}}}
\]
by Lemma~\ref{lem:bnli-representable-prored-thickening}.  Transport the fixed
point $x$ through the inverse equivalence and denote the resulting point of
$(\mathsf P^{\mathrm{sp}}T')(C)$ by $x'$.  The condition on $g$ is now the
identity
\[
    \eta_{T'}(g)=x'
\]
in $(\mathsf P^{\mathrm{sp}}T')(C)$.

Choose a representative
\[
    r=(I,x_I)\in\operatorname{Rep}_{U}(x).
\]
Composing $x_I:T_I\to U$ with $U\to T'$ gives a second point of $T'(T_I)$.
Its image in $(\mathsf P^{\mathrm{sp}}T')(C)$ is also $x'$, by naturality of
the displayed pro-reduction equivalence.  Thus the two actual maps
\[
    T_I\rightrightarrows T',
\]
one obtained by restricting $g$ and the other by $T_I\to U\to T'$, represent
the same point of the filtered colimit defining
$(\mathsf P^{\mathrm{sp}}T')(C)$.  Lemma~\ref{lem:bnli-filtered-representatives}
therefore gives, after a common refinement in $\operatorname{Rep}_U(x)$, a
specified homotopy between these two maps, together with every finite higher
coherence required in a family.

By the pushout universal property, such data are exactly maps under $U$
\[
    U\amalg_{T_I}T
    \longrightarrow T'.
\]
Taking the filtered colimit over the category of representatives gives
\[
    \operatorname{Map}_{/U^{\mathrm{inf,sp}}}(T,T')_x
    \simeq
    \operatorname*{colim}_{r\in\operatorname{Rep}_{U}(x)}
    \operatorname{Map}_{U/}(P_{x,r}(T),T'),
\]
which is the asserted formula for the represented presheaf.  Colimit
preservation gives the formula for arbitrary $G$.  The final specialization is
the cofinality statement in
Definition~\ref{def:bnli-pointed-pro-pushout}.
\end{proof}

\begin{lemma}[Postnikov comparison for nilpotent pullbacks]
\label{lem:bnli-postnikov-nilpotent-pullback}
Let
\[
    C\simeq A_0\times_{A_{01}}A_1
\]
be a pullback of connective commutative spectral rings and suppose
$A_0\to A_{01}$ is surjective on $\pi_0$ with nilpotent kernel.  For
$n\geq0$ put
\[
\boxed{
    C^{\square}_{\leq n}
    :=
    \tau_{\leq n}A_0
    \times_{\tau_{\leq n}A_{01}}
    \tau_{\leq n}A_1.}
\]
Then $C^{\square}_{\leq n}$ is connective and $n$-truncated.  Moreover, for
every fixed $d$ the canonical map
\[
\boxed{
    \tau_{\leq d}C
    \xrightarrow{\sim}
    \tau_{\leq d}C^{\square}_{\leq n}}
\]
is an equivalence for all $n>d$.
\end{lemma}

\begin{proof}
The degree-zero map
\[
    \pi_0(\tau_{\leq n}A_0)
    \longrightarrow
    \pi_0(\tau_{\leq n}A_{01})
\]
is the original surjection.  Hence the underlying pullback in commutative ring
spectra is connective: in the fibre sequence
\[
    C^{\square}_{\leq n}
    \longrightarrow
    \tau_{\leq n}A_0\oplus\tau_{\leq n}A_1
    \longrightarrow
    \tau_{\leq n}A_{01}
\]
the only possible negative homotopy group is killed by that surjection.
Because $n$-truncated spectra are closed under limits, the pullback is also
$n$-truncated.

Compare this fibre sequence with
\[
    C\longrightarrow A_0\oplus A_1\longrightarrow A_{01}.
\]
For every degree strictly below $n$, the maps from the second sequence to the
first are isomorphisms on the homotopy groups of the two right-hand terms.
The long exact homotopy sequences and the five-lemma therefore identify
\[
    \pi_i(C)\xrightarrow{\sim}
    \pi_i(C^{\square}_{\leq n})
    \qquad(i\leq d<n).
\]
This is exactly the displayed truncation equivalence.
\end{proof}

\begin{lemma}[Convergent extension of bounded nilpotent excision]
\label{lem:bnli-convergent-excision-extension}
Let $Y$ be a spectrally convergent prestack.  Suppose $Y$ carries every
nilpotent-excision pullback square whose four affine rings are eventually
coconnective to a pullback of spaces.  Then $Y$ satisfies nilpotent excision
on all connective affine tests.
\end{lemma}

\begin{proof}
Use the notation of
Lemma~\ref{lem:bnli-postnikov-nilpotent-pullback}.  For every $n$ the square
\[
\begin{tikzcd}
C^{\square}_{\leq n} \arrow[r] \arrow[d] &
\tau_{\leq n}A_0 \arrow[d]\\
\tau_{\leq n}A_1 \arrow[r] &
\tau_{\leq n}A_{01}
\end{tikzcd}
\]
is a bounded nilpotent-excision square.  Hence
\[
    Y(C^{\square}_{\leq n})
    \simeq
    Y(\tau_{\leq n}A_0)
    \times_{Y(\tau_{\leq n}A_{01})}
    Y(\tau_{\leq n}A_1).
\]
The preceding lemma and spectral convergence give
\[
\begin{aligned}
    \lim_nY(C^{\square}_{\leq n})
    &\simeq
    \lim_n\lim_d
    Y(\tau_{\leq d}C^{\square}_{\leq n})\\
    &\simeq
    \lim_dY(\tau_{\leq d}C)\\
    &\simeq Y(C).
\end{aligned}
\]
Taking the inverse limit of the bounded pullback equations, and using that
limits of spaces commute with finite limits, gives
\[
    Y(C)
    \simeq
    Y(A_0)\times_{Y(A_{01})}Y(A_1).
\]
\end{proof}

\begin{definition}[Nil-isomorphic deformation diagram]
\label{def:bnli-nil-deformation-diagram}
Let $\operatorname{DefDiag}^{\mathrm{nil}}(U/V)$ be the full subcategory of
diagrams
\[
    U\longrightarrow Y\longrightarrow U^{\mathrm{inf,sp}}
\]
for which $Y$ is spectrally convergent, satisfies nilpotent excision, and
$U\to Y$ is a spectral nil-isomorphism.
\end{definition}

\begin{theorem}[Second-stage materialization equivalence]
\label{thm:bnli-second-stage-equivalence}
The adjunction $\widetilde v_!\dashv\widetilde v^*$ restricts to an
equivalence between:
\begin{enumerate}
    \item functors on unrestricted nilpotent neighbourhoods which have
    deformation theory;
    \item nil-isomorphic deformation diagrams in
    $\operatorname{DefDiag}^{\mathrm{nil}}(U/V)$.
\end{enumerate}
\end{theorem}

\begin{proof}
Let $G$ have deformation theory and put
\[
    Y:=\widetilde v_!G.
\]
By construction $Y$ is spectrally convergent.

We first prove nilpotent excision on bounded tests.  Consider a pullback
square of eventually coconnective connective spectral rings
\[
\begin{tikzcd}
C \arrow[r] \arrow[d] &
A_0 \arrow[d]\\
A_1 \arrow[r] &
A_{01}
\end{tikzcd}
\]
in which $A_0\to A_{01}$ is surjective on $\pi_0$ with nilpotent kernel.
The base-changed map $C\to A_1$ has the same property.  Hence
Lemma~\ref{lem:bnli-prored-thickening} identifies the pro-reductions of
$A_0$ and $A_{01}$ and those of $C$ and $A_1$.  Consequently
$U^{\mathrm{inf,sp}}$ carries the displayed square to a pullback of spaces.

It therefore suffices to check the $Y$-square on homotopy fibres over a
compatible point $x_\bullet$ of this base pullback.  Apply
Lemma~\ref{lem:bnli-simultaneous-square-representatives}.  Its filtered
weakly contractible category
$\operatorname{Rep}^{\square}_U(x_\bullet)$ is cofinal in each of the four
pointed representative systems.  Reindex all four pointed-pro-pushout colimits by this single canonical
refinement category.  At every stage the four discrete representative rings
form a Cartesian square and the two maps along the nilpotent legs are
isomorphisms.  On affine spectra the representative square is therefore a
pushout.  The original bounded affine square is also a pushout geometrically.
Forming
\[
    P_x(T)=U\amalg_{T_I}T
\]
is a finite colimit of the span $U\leftarrow T_I\to T$, and finite colimits
commute with finite colimits.  Pushout pasting therefore identifies the four
stagewise pointed pro-pushouts with a pushout square of unrestricted
nilpotent neighbourhoods of $U$.  One leg remains nilpotent, being obtained by
base change from the original nilpotent leg.  Deformation theory of $G$
carries every such stagewise pushout to a pullback.  Filtered colimits of
spaces commute with finite limits, so
Proposition~\ref{prop:bnli-dotted-lift-formula} shows that the four homotopy
fibres of the $Y$-square form a pullback.  Thus $Y$ satisfies nilpotent
excision on bounded tests.  Lemma~\ref{lem:bnli-convergent-excision-extension}
upgrades this to nilpotent excision on all connective tests.

Next we prove that the centre $U\to Y$ is a spectral nil-isomorphism.  The
category-of-elements co-Yoneda presentation writes $G$ as a colimit of
representable nilpotent neighbourhoods.  Let $\int G$ denote the category of
elements of the covariant functor on the opposite geometric neighbourhood
category.  The identity neighbourhood $U\to U$ is initial geometrically,
hence terminal after passage to the opposite category.  Since $G(U)\simeq *$,
the object $(U,*)$ is terminal in $\int G$.  Therefore the opposite indexing
category used in the co-Yoneda colimit is weakly contractible.

Evaluate pro-reduction of the corresponding colimit presentation of $Y$ on a
connective ring $C$.  Every test
$H(\pi_0C/I)$ occurring in the definition of $\mathsf P^{\mathrm{sp}}$ is
discrete, so colimits of convergent prestacks are pointwise on all of these
tests.  Hence
\[
\begin{aligned}
    (\mathsf P^{\mathrm{sp}}Y)(C)
    &\simeq
    \operatorname*{colim}_{e}
    (\mathsf P^{\mathrm{sp}}T_e)(C)\\
    &\simeq
    \operatorname*{colim}_{e}
    (\mathsf P^{\mathrm{sp}}U)(C)\\
    &\simeq
    (\mathsf P^{\mathrm{sp}}U)(C).
\end{aligned}
\]
The middle equivalence is
Lemma~\ref{lem:bnli-representable-prored-thickening}, because every represented cell
$U\to T_e$ is a nilpotent neighbourhood; the last equivalence uses weak
contractibility of the indexing category.  Thus
\[
\boxed{
    \mathsf P^{\mathrm{sp}}U
    \xrightarrow{\sim}
    \mathsf P^{\mathrm{sp}}Y,}
\]
so the centre is a spectral nil-isomorphism.

We now verify the adjunction unit on every unrestricted nilpotent
neighbourhood.  Let
\[
    U=\operatorname{Spec}B
    \longrightarrow
    T=\operatorname{Spec}C
\]
be such a neighbourhood and put
\[
    K:=\operatorname{fib}(C\to B).
\]
Choose $n$ such that $K$ is $(n-1)$-coconnective, and write
\[
    C_n:=\tau_{\leq n}C,
    \qquad
    B_n:=\tau_{\leq n}B,
    \qquad
    T_n:=\operatorname{Spec}C_n,
    \qquad
    U_n:=\operatorname{Spec}B_n.
\]
Lemma~\ref{lem:bnli-relative-postnikov-reconstruction}, applied to
$C\to B$, gives
\[
    K
    \xrightarrow{\sim}
    \operatorname{fib}(C_n\to B_n)
\]
and the Cartesian reconstruction
\[
\boxed{
    C
    \xrightarrow{\sim}
    B\times_{B_n}C_n.}
\]
Passing to affine spectra gives the relative Postnikov pushout
\[
\boxed{
    T
    \simeq
    U\amalg_{U_n}T_n.}
\]

Put
\[
    K_0:=\ker(\pi_0C\to\pi_0B)
\]
and, for every nilpotent ideal
$I\supseteq K_0$, write
\[
    D_I:=\operatorname{Spec}H(\pi_0C/I).
\]
The bounded points of $T_n$ and $U_n$ over $U^{\mathrm{inf,sp}}$ admit these
common representatives.  Define
\[
    P^T_{n,I}:=U\amalg_{D_I}T_n,
    \qquad
    P^U_{n,I}:=U\amalg_{D_I}U_n.
\]
Mapping the relative Postnikov pushout into $Y=\widetilde v_!G$ under $U$
and over $U^{\mathrm{inf,sp}}$, followed by the bounded pointed
pro-pushout formula, gives
\[
\widetilde v^*Y(T)
\simeq
\left(
  \operatorname*{colim}_{I\supseteq K_0}G(P^T_{n,I})
\right)
\times_{
  \operatorname*{colim}_{I\supseteq K_0}G(P^U_{n,I})
}
G(U).
\]
For each $I$, pushout pasting gives a pushout square
\[
\begin{tikzcd}
P^U_{n,I} \arrow[r] \arrow[d] &
U \arrow[d]\\
P^T_{n,I} \arrow[r] &
T.
\end{tikzcd}
\]
The right-hand leg is the original nilpotent neighbourhood
$U\to T$.  Hence deformation theory of $G$ gives
\[
    G(T)
    \xrightarrow{\sim}
    G(P^T_{n,I})
    \times_{G(P^U_{n,I})}
    G(U).
\]
Taking the filtered colimit over $I$ and commuting it with the finite
pullback identifies the right-hand side with the preceding expression for
$\widetilde v^*Y(T)$.  Thus
\[
\boxed{
    G(T)
    \xrightarrow{\sim}
    \widetilde v^*\widetilde v_!G(T)}
\]
for every unrestricted nilpotent neighbourhood.  This proves the unit
equivalence without imposing any absolute Postnikov bound on the coherent
centre $B$.

Conversely let
\[
    U\xrightarrow{y_0}Y\longrightarrow U^{\mathrm{inf,sp}}
\]
be a nil-isomorphic deformation diagram.  The dotted-lift functor
$G:=\widetilde v^*Y$ has deformation theory: its value at the identity
neighbourhood is contractible, and mapping a pushout of nilpotent
neighbourhoods into the nilpotent-excisive target $Y$, with the centre fixed,
gives a pullback of dotted-lift spaces.

It remains to prove the counit
\[
    \widetilde v_!G\longrightarrow Y.
\]
Both sides are convergent, so it is enough to evaluate on an eventually
coconnective affine test
\[
    x:T=\operatorname{Spec}C
    \longrightarrow U^{\mathrm{inf,sp}}.
\]
For
$r=(I,x_I)\in\operatorname{Rep}_{U}(x)$, put
\[
    P_{x,r}(T)=U\amalg_{T_I}T.
\]
The pointed pro-pushout formula gives
\[
    \left(
      \widetilde v_!G
    \right)(T)_x
    \simeq
    \operatorname*{colim}_{r}
    \operatorname{Map}_{U/,/U^{\mathrm{inf,sp}}}
    (P_{x,r}(T),Y).
\]
Nilpotent excision for $Y$ identifies the term indexed by $r$ with
\[
    Y(T)_x
    \times_{Y(T_I)_{x_I}}
    \{y_0|_{T_I}\},
\]
where the subscripts denote fibres over the displayed infinitesimal points.
Filtered colimits commute with finite limits, so
\[
    \left(
      \widetilde v_!G
    \right)(T)_x
    \simeq
    Y(T)_x
    \times_{
      \operatorname*{colim}_{r}Y(T_I)_{x_I}
    }
    *.
\]
Apply $\mathsf P^{\mathrm{sp}}$ to the centred diagram
\[
    U\longrightarrow Y\longrightarrow\mathsf P^{\mathrm{sp}}U.
\]
The first arrow becomes an equivalence because $U\to Y$ is a spectral
nil-isomorphism, while the composite becomes the idempotence equivalence of
$\mathsf P^{\mathrm{sp}}U$.  Hence
\[
\boxed{
    \mathsf P^{\mathrm{sp}}Y
    \xrightarrow{\sim}
    \mathsf P^{\mathrm{sp}}U}
\]
for the map induced by the displayed structure morphism.  The pointed
pro-reduction fibre formula,
Lemma~\ref{lem:bnli-pointed-prored-fibre}, therefore gives
\[
\boxed{
    \operatorname*{colim}_{r\in\operatorname{Rep}_U(x)}
    Y(T_I)_{x_I}
    \simeq
    \operatorname{fib}_{x}
    \left(
      (\mathsf P^{\mathrm{sp}}Y)(C)
      \to
      (\mathsf P^{\mathrm{sp}}U)(C)
    \right)
    \simeq *.}
\]
Here $Y(T_I)_{x_I}$ denotes the fibre of
$Y(T_I)\to(\mathsf P^{\mathrm{sp}}U)(T_I)$ over the pro-reduction image of
$x_I$, exactly as in the preceding pointed pro-pushout calculation.  Therefore
\[
    \left(
      \widetilde v_!G
    \right)(T)_x
    \xrightarrow{\sim}
    Y(T)_x.
\]
This is natural in $x$, hence the counit is an equivalence on every bounded
test.  Spectral convergence extends it to every connective affine test.

The unit and counit are therefore equivalences on the stated full
subcategories, proving the theorem.
\end{proof}

\begin{proposition}[Forgetting the infinitesimal target]
\label{prop:bnli-forget-inf-target}
For every spectral nil-isomorphism $U\to Y$, the mapping space
\[
\boxed{
    \operatorname{Map}_{U/}
    (Y,U^{\mathrm{inf,sp}})
    \simeq *.}
\]
Consequently forgetful passage
\[
    (U\to Y\to U^{\mathrm{inf,sp}})
    \longmapsto
    (U\to Y)
\]
is fully faithful on nil-isomorphic diagrams.  Restricted to convergent
deformation diagrams it gives an equivalence onto convergent nil-centred
deformation prestacks under $U$.  The same statement therefore holds on the
spectrally laft locus used below.
\end{proposition}

\begin{proof}
By Definition~\ref{def:bnli-inf-prestack},
$U^{\mathrm{inf,sp}}=\mathsf P^{\mathrm{sp}}U$ is pro-reduced-local.  Since
$U\to Y$ is a spectral nil-isomorphism,
Definition~\ref{def:bnli-nil-isomorphism} gives an equivalence
\[
    \operatorname{Map}(Y,U^{\mathrm{inf,sp}})
    \xrightarrow{\sim}
    \operatorname{Map}(U,U^{\mathrm{inf,sp}}).
\]
Taking the homotopy fibre over the pro-reduction coaugmentation
$U\to U^{\mathrm{inf,sp}}$ gives the displayed contractible mapping space.
Thus the map to the infinitesimal target is contractibly determined and is
not additional structure.
\end{proof}

\begin{construction}[Pro-infinitesimal completion of an arbitrary centred target]
\label{con:bnli-pro-infinitesimal-completion}
Let $u:U\to Y$ be any centred spectral prestack.  Define
\[
\boxed{
    Y_U^{\mathrm{pinf}}
    :=
    Y\times_{\mathsf P^{\mathrm{sp}}Y}
    \mathsf P^{\mathrm{sp}}U.}
\]
The naturality square for the pro-reduction coaugmentations gives a canonical
centre
\[
    U\longrightarrow Y_U^{\mathrm{pinf}}
\]
and a canonical map
\[
    Y_U^{\mathrm{pinf}}\longrightarrow U^{\mathrm{inf,sp}}
\]
under $U$.
\end{construction}

\begin{proposition}[Universal property of pro-infinitesimal completion]
\label{prop:bnli-pro-infinitesimal-universal}
The centre
\[
    U\longrightarrow Y_U^{\mathrm{pinf}}
\]
is a spectral nil-isomorphism.  For every nil-isomorphic centred prestack
$U\to Z$, projection induces a canonical equivalence
\[
\boxed{
    \operatorname{Map}_{U/}
    (Z,Y_U^{\mathrm{pinf}})
    \xrightarrow{\sim}
    \operatorname{Map}_{U/}(Z,Y).}
\]
If $Y$ is spectrally convergent, then $Y_U^{\mathrm{pinf}}$ is spectrally
convergent and therefore belongs to the target of
$\widetilde v_!$.  In that case, for every unrestricted nilpotent affine
neighbourhood $U\to T$,
\[
\boxed{
    \widetilde v^*
    \left(
      Y_U^{\mathrm{pinf}}
    \right)(T)
    \simeq
    \operatorname{Map}_{U/}(T,Y).}
\]
\end{proposition}

\begin{proof}
Apply the left-exact functor $\mathsf P^{\mathrm{sp}}$ to the defining
pullback.  Idempotence gives
\[
\begin{aligned}
    \mathsf P^{\mathrm{sp}}
    \left(
      Y_U^{\mathrm{pinf}}
    \right)
    &\simeq
    \mathsf P^{\mathrm{sp}}Y
    \times_{\mathsf P^{\mathrm{sp}}Y}
    \mathsf P^{\mathrm{sp}}U\\
    &\simeq
    \mathsf P^{\mathrm{sp}}U.
\end{aligned}
\]
Thus the centre is a nil-isomorphism.

Let $Z$ be nil-isomorphic to $U$.  A map $Z\to Y$ under $U$ induces after
pro-reduction a map
\[
    \mathsf P^{\mathrm{sp}}Z
    \longrightarrow
    \mathsf P^{\mathrm{sp}}Y.
\]
Using the canonical equivalence
$\mathsf P^{\mathrm{sp}}Z\simeq\mathsf P^{\mathrm{sp}}U$, naturality forces
this pro-reduced component to be the one induced by the centre $U\to Y$.
Thus the lift to the pullback $Y_U^{\mathrm{pinf}}$ exists through a
contractible space.  This proves the mapping equivalence.

The prestacks $\mathsf P^{\mathrm{sp}}Y$ and
$\mathsf P^{\mathrm{sp}}U$ are convergent because they depend only on
$\pi_0$ and its nilpotent quotient system.  If $Y$ is convergent, the defining
finite pullback therefore makes $Y_U^{\mathrm{pinf}}$ convergent as well.
An unrestricted nilpotent neighbourhood $T$ is nil-isomorphic to $U$ by
Lemma~\ref{lem:bnli-representable-prored-thickening}.  The preceding mapping equivalence,
together with the dotted-lift description of $\widetilde v^*$, then gives the
final formula.
\end{proof}

\begin{lemma}[Base change and right cancellation for nil-laft maps]
\label{lem:bnli-nil-laft-base-change-cancellation}
Let
\[
    G\longrightarrow G'
\]
be a nil-laft morphism of presheaves on
$\operatorname{Aff}^{\mathrm{sp,nil,+}}_{U/V}$.
\begin{enumerate}
    \item For every morphism $H\to G'$, the base change
    \[
        G\times_{G'}H\longrightarrow H
    \]
    is nil-laft.
    \item For composable morphisms
    \[
        G\longrightarrow H\longrightarrow K,
    \]
    if $H\to K$ and $G\to K$ are nil-laft, then $G\to H$ is nil-laft.
\end{enumerate}
\end{lemma}

\begin{proof}
Let $\{S_\alpha\}$ be a cofiltered diagram occurring in
Definition~\ref{def:bnli-nil-laft}, with limit $S$.  Filtered colimits of
spaces commute with finite limits.  For the first assertion,
\[
\begin{aligned}
\operatorname*{colim}_\alpha
(G\times_{G'}H)(S_\alpha)
&\simeq
\left(\operatorname*{colim}_\alpha G(S_\alpha)\right)
\times_{\operatorname*{colim}_\alpha G'(S_\alpha)}
\left(\operatorname*{colim}_\alpha H(S_\alpha)\right)\\
&\simeq
G(S)\times_{G'(S)}
\operatorname*{colim}_\alpha H(S_\alpha)\\
&\simeq
(G\times_{G'}H)(S)
\times_{H(S)}
\operatorname*{colim}_\alpha H(S_\alpha).
\end{aligned}
\]
This is the nil-laft equation for the base-changed map.

For the second assertion, use nil-laftness of $H\to K$ and then of
$G\to K$:
\[
\begin{aligned}
G(S)\times_{H(S)}
\operatorname*{colim}_\alpha H(S_\alpha)
&\simeq
G(S)\times_{H(S)}
\left(
H(S)\times_{K(S)}
\operatorname*{colim}_\alpha K(S_\alpha)
\right)\\
&\simeq
G(S)\times_{K(S)}
\operatorname*{colim}_\alpha K(S_\alpha)\\
&\simeq
\operatorname*{colim}_\alpha G(S_\alpha).
\end{aligned}
\]
Hence $G\to H$ is nil-laft.
\end{proof}

\begin{proposition}[Local almost finite presentation of two-stage materialization]
\label{prop:bnli-two-stage-laft}
For every $F\in\operatorname{FMP}^{\mathrm{sp}}_{B/R}$, the centre map
\[
    U\longrightarrow
    \operatorname{oblv}(\widetilde v_!u_!F)
\]
is spectrally laft.
\end{proposition}

\begin{proof}
Write
\[
    Y:=
    \operatorname{oblv}(\widetilde v_!u_!F).
\]
Both $U$ and $Y$ are spectrally convergent, so the convergence square in the
definition of a spectrally laft morphism is already Cartesian.  It remains to
prove the relative filtered-colimit condition on uniformly truncated tests.

The affine co-Yoneda formula expresses $F$ as the colimit of its Artinian
corepresentable cells.  Since both realization functors are
colimit-preserving, $Y$ is the corresponding colimit of the represented
Artinian neighbourhoods
\[
    U\longrightarrow S=\operatorname{Spec}A
\]
indexed by Artinian elements of $F$.  On every $n$-truncated affine test this
colimit is computed pointwise by
Proposition~\ref{prop:bnli-convergent-bounded-tests}.

Let $T\to Y$ be an $n$-truncated affine test.  In the pointwise
category-of-elements presentation choose an Artinian element
\[
    \xi:h_A^{\mathrm{Art}}\longrightarrow F,
\]
a factorization
\[
    T\longrightarrow S=\operatorname{Spec}A\longrightarrow Y,
\]
and a path in the mapping space $\operatorname{Map}(T,Y)$ from the displayed
composite to the given map $T\to Y$.  Homotopic maps into $Y$ have equivalent
homotopy pullbacks.  Since the relative laft condition is stable under affine
base change, it is therefore enough to show that
\[
\boxed{
    U\times_YS\longrightarrow S}
\]
is spectrally laft.

Let $F_U$ denote the formal-moduli problem obtained by applying the inverse
first- and second-stage recognition equivalences to the represented centred
object
\[
    U\longrightarrow U\longrightarrow U^{\mathrm{inf,sp}}.
\]
Under the same equivalences, the centre $U\to Y$ corresponds to a canonical
morphism
\[
    F_U\longrightarrow F,
\]
while the represented Artinian neighbourhood $S=\operatorname{Spec}A$
corresponds to $h_A^{\mathrm{Art}}$.  Put
\[
\boxed{
    H_{A,\xi}
    :=
    F_U\times_Fh_A^{\mathrm{Art}}}
\]
and
\[
    Z:=U\times_YS.
\]
We now identify this particular ambient pullback with the two-stage
realization of $H_{A,\xi}$ without asserting that materialization preserves
arbitrary ambient finite limits.

Regard $U$, $S$, $Y$, and $Z$ as diagrams over $U^{\mathrm{inf,sp}}$ through
the canonical infinitesimal maps.  Since $\widetilde v^*$ is a right adjoint,
it preserves the displayed pullback.  The first- and second-stage recognition
unit equivalences therefore give
\[
\begin{aligned}
    G_Z
    &:=\widetilde v^*Z\\
    &\simeq
    \widetilde v^*U
    \times_{\widetilde v^*Y}
    \widetilde v^*S\\
    &\simeq
    u_!F_U
    \times_{u_!F}
    u_!h_A^{\mathrm{Art}}.
\end{aligned}
\]

The functor $G_Z$ has deformation theory because that property is preserved
by finite limits.  By construction $F_U$ is the formal-moduli problem
corresponding under the first-stage recognition equivalence to the represented
centred functor $U$ on unrestricted nilpotent neighbourhoods.  Hence there is
a canonical equivalence
\[
    u_!F_U\xrightarrow{\sim}U,
\]
and under this equivalence the morphism
\[
    u_!F_U\longrightarrow u_!F
\]
is exactly the centre $U\to u_!F$.  The first half of
Theorem~\ref{thm:bnli-first-stage-recognition} proves that this centre is
nil-laft.

The projection
\[
    G_Z\longrightarrow u_!h_A^{\mathrm{Art}}
\]
is the base change of that centre morphism.  It is therefore nil-laft by
Lemma~\ref{lem:bnli-nil-laft-base-change-cancellation}.  Its composite with
the centre
\[
    U\longrightarrow G_Z
\]
is the centre
\[
    U\longrightarrow u_!h_A^{\mathrm{Art}},
\]
which is nil-laft by the same first-stage theorem applied to
$h_A^{\mathrm{Art}}$.  Right cancellation in
Lemma~\ref{lem:bnli-nil-laft-base-change-cancellation} therefore shows that
$U\to G_Z$ is nil-laft.

Restriction to the Artinian subcategory preserves the pullback, so
\[
\begin{aligned}
    u^*G_Z
    &\simeq
    F_U\times_Fh_A^{\mathrm{Art}}\\
    &=H_{A,\xi}.
\end{aligned}
\]
The first-stage recognition theorem now gives a canonical equivalence
\[
\boxed{
    G_Z\xrightarrow{\sim}u_!H_{A,\xi}.}
\]

It remains to identify the second-stage realization.  The prestack $Z$ is
spectrally convergent and nilpotent-excisive because both properties are
preserved by finite limits.  Moreover, left exactness of
$\mathsf P^{\mathrm{sp}}$, the nil-isomorphism $U\to Y$, and
Lemma~\ref{lem:bnli-representable-prored-thickening} for the Artinian centre
$U\to S$ give
\[
\begin{aligned}
    \mathsf P^{\mathrm{sp}}Z
    &\simeq
    \mathsf P^{\mathrm{sp}}U
    \times_{\mathsf P^{\mathrm{sp}}Y}
    \mathsf P^{\mathrm{sp}}S\\
    &\simeq
    \mathsf P^{\mathrm{sp}}U.
\end{aligned}
\]
Thus $U\to Z$ is a spectral nil-isomorphism.  The second-stage recognition
counit therefore gives
\[
\boxed{
    Z
    \xrightarrow{\sim}
    \widetilde v_!G_Z
    \xrightarrow{\sim}
    \widetilde v_!u_!H_{A,\xi}.}
\]
Under this equivalence the projection $Z\to S$ is the two-stage realization
of $H_{A,\xi}\to h_A^{\mathrm{Art}}$ by naturality of the recognition units
and counits.

Apply co-Yoneda in the slice over $h_A^{\mathrm{Art}}$.  It presents
$H_{A,\xi}$ as a colimit of corepresentable cells
\[
    h_{A'}^{\mathrm{Art}}
    \longrightarrow
    h_A^{\mathrm{Art}}.
\]
By Yoneda, such a morphism is the same as a morphism of spectral Artinian
extensions
\[
    A\longrightarrow A'\longrightarrow B,
\]
and geometrically gives
\[
    S'=\operatorname{Spec}A'
    \longrightarrow
    S=\operatorname{Spec}A.
\]
Lemma~\ref{lem:bnli-artinian-morphism-afp} shows that every $A'$ is almost
finitely presented over $A$.  Thus every represented map $S'\to S$ is
spectrally laft by
Lemma~\ref{lem:bnli-affine-afp-laft}.

Now let $\{D_\alpha\}$ be a filtered diagram of $m$-truncated connective
$A$-algebras with colimit $D$.  Evaluation of every represented cell
$S'\to S$ commutes with this filtered colimit.  The realization of
$H_{A,\xi}$ is a colimit of these cells, so the two colimits commute:
\[
\begin{aligned}
    \operatorname*{colim}_{\alpha}
    (U\times_YS)(D_\alpha)
    &\simeq
    \operatorname*{colim}_{(A',\eta)}
    \operatorname*{colim}_{\alpha}
    \operatorname{Map}_{\operatorname{CAlg}_A}(A',D_\alpha)\\
    &\simeq
    \operatorname*{colim}_{(A',\eta)}
    \operatorname{Map}_{\operatorname{CAlg}_A}(A',D)\\
    &\simeq
    (U\times_YS)(D).
\end{aligned}
\]
This is precisely the relative laft equation over the represented base $S$.
Therefore $U\times_YS\to S$ is spectrally laft.  Pulling back along
$T\to S$ proves the relative laft equation over the original truncated test
$T\to Y$.  Since $n$ and the filtered diagram were arbitrary,
$U\to Y$ is spectrally laft.
\end{proof}

\begin{lemma}[Raw split square-zero tangent spectrum]
\label{lem:bnli-raw-square-zero-tangent}
Let $D$ be any connective spectral $R$-algebra.  Let $Y$ be a spectral
prestack which satisfies Postnikov convergence and nilpotent excision on
connective affine tests, and fix a point
\[
    x\in Y(D).
\]
For a connective $D$-module $M$ and $n\geq0$ put
\[
    \Theta^n_{Y,D,x}(M)
    :=
    \operatorname{fib}_{x}
    \left(
      Y(D\oplus M[n])
      \longrightarrow
      Y(D)
    \right).
\]
Then the split square-zero suspension identities make
$n\mapsto\Theta^n_{Y,D,x}(M)$ an $\Omega$-spectrum.  These spectra extend
uniquely to an exact functor
\[
\boxed{
    \Theta^{\mathrm{sp}}_{Y,D,x}:
    \operatorname{Mod}_D
    \longrightarrow
    \operatorname{Sp}.}
\]
If $D\to D'$ is spectral etale, $x'$ is the restricted point, and $N$ is a
$D'$-module, then restriction of scalars induces a canonical equivalence
\[
\boxed{
    \Theta^{\mathrm{sp}}_{Y,D,x}
    \left(
      \operatorname{Res}^{D'}_D N
    \right)
    \xrightarrow{\sim}
    \Theta^{\mathrm{sp}}_{Y,D',x'}(N).}
\]
No sheaf condition, coherence condition, or pro-coherent approximation is
required for this construction.
\end{lemma}

\begin{proof}
For connective $M$ there is a canonical pullback of connective spectral
algebras
\[
    D\oplus M[n]
    \simeq
    D\times_{D\oplus M[n+1]}D.
\]
For $n\ge0$ the map $D\to D\oplus M[n+1]$ is an isomorphism on $\pi_0$ and
therefore is a nilpotent-excision leg.  Nilpotent excision for $Y$, followed
by the homotopy fibre over $x$, gives
\[
    \Theta^n_{Y,D,x}(M)
    \xrightarrow{\sim}
    \Omega\Theta^{n+1}_{Y,D,x}(M).
\]
Thus the spaces $\Theta^n_{Y,D,x}(M)$ assemble functorially into a spectrum,
which we denote $\Theta^{\mathrm{cn}}_{Y,D,x}(M)$.

Let
\[
    M'\longrightarrow M\longrightarrow M''
\]
be a fibre sequence of connective $D$-modules.  The split square-zero algebras
form a pullback square
\[
\begin{tikzcd}
D\oplus M' \arrow[r] \arrow[d] &
D \arrow[d]\\
D\oplus M \arrow[r] &
D\oplus M''.
\end{tikzcd}
\]
Because $M'$ is connective, the map
$\pi_0(M)\to\pi_0(M'')$ is surjective.  Hence the lower horizontal algebra
map is surjective on $\pi_0$ with square-zero kernel, so this is a
nilpotent-excision square.  Evaluating $Y$ and passing to fibres over $x$
gives a fibre sequence of spectra
\[
    \Theta^{\mathrm{cn}}_{Y,D,x}(M')
    \longrightarrow
    \Theta^{\mathrm{cn}}_{Y,D,x}(M)
    \longrightarrow
    \Theta^{\mathrm{cn}}_{Y,D,x}(M'').
\]
Therefore
\[
    \Theta^{\mathrm{cn}}_{Y,D,x}:
    \operatorname{Mod}_{D,\ge0}
    \longrightarrow
    \operatorname{Sp}
\]
is reduced and excisive.

The connective module category has stabilization
\[
\boxed{
    \operatorname{Sp}
    \left(
      \operatorname{Mod}_{D,\ge0}
    \right)
    \simeq
    \operatorname{Mod}_D.}
\]
By the universal property of stabilization, every reduced excisive
spectrum-valued functor on $\operatorname{Mod}_{D,\ge0}$ extends uniquely to
an exact functor on its stabilization.  This gives
\[
    \Theta^{\mathrm{sp}}_{Y,D,x}:
    \operatorname{Mod}_D\to\operatorname{Sp}.
\]

Now let $D\to D'$ be spectral etale.  For every connective $D'$-module $N$
there is a canonical pullback
\[
\begin{tikzcd}
D\oplus \operatorname{Res}^{D'}_D N \arrow[r] \arrow[d] &
D \arrow[d]\\
D'\oplus N \arrow[r] &
D'.
\end{tikzcd}
\]
The lower horizontal map is split square-zero and hence a
nilpotent-excision leg.  Evaluating $Y$ and taking the compatible fibres over
$x$ and $x'$ gives a natural equivalence of reduced excisive functors on
connective $D'$-modules
\[
    \Theta^{\mathrm{cn}}_{Y,D,x}
    \left(
      \operatorname{Res}^{D'}_D(-)
    \right)
    \xrightarrow{\sim}
    \Theta^{\mathrm{cn}}_{Y,D',x'}(-).
\]
The universal property of stabilization extends this equivalence uniquely to
all $D'$-modules.  Uniqueness supplies coherence for identities, composition,
and spectral-etale pullback.  Postnikov convergence is not needed in this
local tangent construction itself; it is retained in the lemma's hypotheses
because it is the form in which the construction is applied in the
materialization and Nisnevich-descent arguments below.
\end{proof}

\begin{lemma}[Relative structured square-zero obstruction fibre]
\label{lem:bnli-relative-square-zero-obstruction-fibre}
Let $P$ be a deformation presheaf on unrestricted nilpotent neighbourhoods.
Let
\[
    A\longrightarrow D
\]
be a connective augmentation, fix $x\in P(D)$, and let
\[
    A'\longrightarrow A
\]
be a structured square-zero extension whose coefficient is the restriction of
a connective $D$-module $M$.  Put
\[
    F_A
    :=
    P(\operatorname{Spec}A)
    \times_{P(\operatorname{Spec}D)}
    \{x\}
\]
and
\[
    R_{D,x}(M[1])
    :=
    P(\operatorname{Spec}(D\oplus M[1]))
    \times_{P(\operatorname{Spec}D)}
    \{x\}.
\]
Then the classifying derivation induces a canonical obstruction morphism
\[
    o_{A'/A}:F_A\longrightarrow R_{D,x}(M[1]),
\]
and there is a canonical Cartesian square
\[
\begin{tikzcd}
F_{A'} \arrow[r] \arrow[d] &
F_A \arrow[d,"o_{A'/A}"]\\
* \arrow[r,"0_x"'] &
R_{D,x}(M[1]).
\end{tikzcd}
\]
The construction is natural under base change of the augmented structured
square-zero extension.
\end{lemma}

\begin{proof}
The structured square-zero extension is the pullback
\[
    A'
    \simeq
    A\times_{A\oplus_DM[1]}A,
\]
where the two maps from $A$ are the classifying derivation and the zero
derivation.  By
Lemma~\ref{lem:bnli-split-coefficient-base-change},
\[
    A\oplus_DM[1]
    \simeq
    A\times_D(D\oplus M[1]),
\]
and deformation theory of $P$ therefore gives
\[
    P(\operatorname{Spec}(A\oplus_DM[1]))
    \simeq
    P(\operatorname{Spec}A)
    \times_{P(\operatorname{Spec}D)}
    P(\operatorname{Spec}(D\oplus M[1])).
\]
After taking the fibre over $x$, the target of the two maps from
$P(\operatorname{Spec}A)$ identifies with
\[
    F_A\times R_{D,x}(M[1]).
\]
Under this identification the zero derivation is
$y\mapsto(y,0_x)$ and the classifying derivation is
$y\mapsto(y,o_{A'/A}(y))$.  Applying deformation theory to the pullback
which defines $A'$ consequently gives
\[
\begin{aligned}
    F_{A'}
    &\simeq
    F_A
    \times_{F_A\times R_{D,x}(M[1])}
    F_A\\
    &\simeq
    F_A\times_{R_{D,x}(M[1])}\{0_x\}.
\end{aligned}
\]
This is the displayed Cartesian square.  Naturality follows from functoriality
of structured square-zero pullbacks, split coefficient base change, and the
classifying derivation.
\end{proof}

\begin{theorem}[Spectral Nisnevich descent for formal-moduli prestacks]
\label{thm:bnli-formal-stack-nisnevich}
Let $U$ be a spectral affine scheme and let
\[
    U\longrightarrow Y
\]
be a spectrally laft nil-isomorphism of convergent spectral prestacks.  Assume
that $Y$ has nilpotent deformation theory.  Then $Y$ satisfies spectral
Nisnevich descent.  Consequently it defines an object of the native spectral
Nisnevich $\infty$-topos; as a native sheaf it is moreover
$L_S^{\mathrm{Def}}$-local.
\end{theorem}

\begin{proof}
The affine finite-basis criterion reduces spectral Nisnevich descent to
Nisnevich distinguished squares.  The morphisms in such a square are flat, so
Postnikov truncation preserves the distinguished-square property.  Spectral
convergence therefore reduces the proof to an eventually coconnective
distinguished square
\[
\begin{tikzcd}
A \arrow[r] \arrow[d] &
A_1 \arrow[d]\\
A_2 \arrow[r] &
A_{12}.
\end{tikzcd}
\]

Choose the filtered system of nilpotent ideals
$I\subseteq\pi_0(A)$ and put
\[
    I_i:=I\pi_0(A_i),
    \qquad
    i\in\{1,2,12\}.
\]
Flat spectral-etale and open base change gives
\[
    \operatorname*{colim}_I
    \pi_0(A_i)/I_i
    \xrightarrow{\sim}
    \pi_0(A_i)_{\mathrm{red}}.
\]
For every vertex
$i\in\{\varnothing,1,2,12\}$, spectral laftness of the centre gives a
Cartesian square
\[
\begin{tikzcd}
\operatorname*{colim}_I U(H(\pi_0(A_i)/I_i))
    \arrow[r] \arrow[d] &
\operatorname*{colim}_I Y(H(\pi_0(A_i)/I_i))
    \arrow[d]\\
U(H(\pi_0(A_i)_{\mathrm{red}}))
    \arrow[r] &
Y(H(\pi_0(A_i)_{\mathrm{red}})).
\end{tikzcd}
\]
The \emph{bottom horizontal} map is an equivalence because the reduced
discrete test is pro-reduced and $U\to Y$ is a nil-isomorphism.  Since the
square is Cartesian, the top horizontal map is therefore an equivalence:
\[
\boxed{
\operatorname*{colim}_I U(H(\pi_0(A_i)/I_i))
\xrightarrow{\sim}
\operatorname*{colim}_I Y(H(\pi_0(A_i)/I_i)).}
\]
We now prove a fibrewise Nisnevich statement at every fixed nilpotent stage.

Fix $I$ and a point
\[
    x\in U(D),
    \qquad
    D:=H(\pi_0(A)/I).
\]
Put
\[
    D_i:=H(\pi_0(A_i)/I_i)
\]
and let $x_i$ be the induced points.  These discrete coefficient rings are
literally the spectral base changes of $D$ along the Nisnevich vertices.  For
every $i\in\{1,2,12\}$, flat spectral-etale or open base change gives
\[
\begin{aligned}
    \pi_k(D\otimes_AA_i)
    &\simeq
    \pi_k(D)\otimes_{\pi_0(A)}\pi_0(A_i).
\end{aligned}
\]
The right-hand side vanishes for $k>0$, while in degree zero it is
\[
    (\pi_0(A)/I)\otimes_{\pi_0(A)}\pi_0(A_i)
    \simeq
    \pi_0(A_i)/I_i.
\]
Consequently
\[
\boxed{
    D\otimes_AA_i\xrightarrow{\sim}D_i.}
\]
Thus the square of the $D_i$ is exactly the base change of the original
Nisnevich distinguished square along $D\to A$, and quasi-coherent module
Nisnevich descent applies to the coefficient modules used below.  We claim
that
\[
\begin{tikzcd}
Y(A)\times_{Y(D)}\{x\} \arrow[r] \arrow[d] &
Y(A_1)\times_{Y(D_1)}\{x_1\} \arrow[d]\\
Y(A_2)\times_{Y(D_2)}\{x_2\} \arrow[r] &
Y(A_{12})\times_{Y(D_{12})}\{x_{12}\}
\end{tikzcd}
\]
is Cartesian.

Apply Proposition~\ref{prop:bnli-unrestricted-finite-square-zero} to the
nilpotent eventually coconnective arrow $A\to D$.  Choose, only inside this
proof, a finite factorization
\[
    A=A^{(r)}
    \longrightarrow\cdots\longrightarrow
    A^{(0)}=D
\]
in which
\[
    A^{(k)}\longrightarrow A^{(k-1)}
\]
is a structured square-zero extension by the restriction of a connective
eventually coconnective $D$-module $J^{(k)}$.  Flat base change along the
Nisnevich square gives
\[
    A_i^{(k)}
    :=
    A^{(k)}\otimes_AA_i,
    \qquad
    J_i^{(k)}
    :=
    J^{(k)}\otimes_DD_i.
\]
Every
$A_i^{(k)}\to A_i^{(k-1)}$
is the corresponding structured square-zero extension and
$A_i^{(0)}=D_i$.

Define
\[
    F_i^{(k)}
    :=
    Y(A_i^{(k)})
    \times_{Y(D_i)}
    \{x_i\}.
\]
We prove by induction on $k$ that the four $F_i^{(k)}$ form a Cartesian
square.  For $k=0$ all four spaces are contractible.  At the inductive step,
Lemma~\ref{lem:bnli-relative-square-zero-obstruction-fibre}, applied over
$D_i$, gives a natural Cartesian square
\[
\begin{tikzcd}
F_i^{(k)} \arrow[r] \arrow[d] &
F_i^{(k-1)} \arrow[d,"o_i^{(k)}"]\\
* \arrow[r,"0"'] &
R_i^{(k)},
\end{tikzcd}
\]
where
\[
    R_i^{(k)}
    :=
    Y(D_i\oplus J_i^{(k)}[1])
    \times_{Y(D_i)}
    \{x_i\}.
\]
The obstruction morphisms $o_i^{(k)}$ are compatible under the flat base
changes of the Nisnevich square.  By
Lemma~\ref{lem:bnli-raw-square-zero-tangent},
\[
    R_i^{(k)}
    \simeq
    \Omega^\infty
    \Theta^{\mathrm{sp}}_{Y,D,x}
    \left(
      \operatorname{Res}^{D_i}_D
      J_i^{(k)}[1]
    \right).
\]

Quasi-coherent spectral modules satisfy Nisnevich descent.  Hence the
restriction of the coefficient square is a pullback in $\operatorname{Mod}_D$:
\[
\boxed{
    J^{(k)}
    \xrightarrow{\sim}
    \operatorname{Res}^{D_1}_D J_1^{(k)}
    \times_{\operatorname{Res}^{D_{12}}_D J_{12}^{(k)}}
    \operatorname{Res}^{D_2}_D J_2^{(k)}.}
\]
The functor
$\Theta^{\mathrm{sp}}_{Y,D,x}$
is exact and therefore preserves finite limits, and $\Omega^\infty$ preserves
limits.  Thus the four $R_i^{(k)}$ form a Cartesian square.  The induction
hypothesis gives a Cartesian square of the four $F_i^{(k-1)}$, and the
base-change naturality of the obstruction morphisms gives a morphism from
that Cartesian square to the Cartesian square of the $R_i^{(k)}$.  Taking the
pointwise homotopy fibres over the compatible zero sections is a finite-limit
construction.  Finite limits commute with finite limits, so the four
$F_i^{(k)}$ again form a Cartesian square.  At $k=r$ this proves the
fibrewise claim.

For this fixed $I$, put
\[
    Z_{i,I}
    :=
    Y(A_i)\times_{Y(D_i)}U(D_i),
    \qquad
    i\in\{\varnothing,1,2,12\}.
\]
The four represented spaces $U(D_i)$ form a Cartesian Nisnevich square.
The preceding fibrewise claim says that, over every compatible point of this
base square, the four homotopy fibres of $Z_{i,I}\to U(D_i)$ form a
Cartesian square.  The total-space criterion for a square over a Cartesian
base therefore gives a Cartesian square
\[
\begin{tikzcd}
Z_{\varnothing,I} \arrow[r] \arrow[d] &
Z_{1,I} \arrow[d]\\
Z_{2,I} \arrow[r] &
Z_{12,I}.
\end{tikzcd}
\]

Filtered colimits of spaces commute with finite limits, so taking the colimit
over $I$ preserves this Cartesian square.  For each vertex $i$,
\[
\begin{aligned}
\operatorname*{colim}_I Z_{i,I}
&\simeq
Y(A_i)
\times_{\operatorname*{colim}_IY(D_i)}
\operatorname*{colim}_I U(D_i)\\
&\simeq
Y(A_i),
\end{aligned}
\]
where the second equivalence is the top horizontal equivalence proved at the
beginning of the argument.  Thus the original square
\[
\begin{tikzcd}
Y(A) \arrow[r] \arrow[d] &
Y(A_1) \arrow[d]\\
Y(A_2) \arrow[r] &
Y(A_{12})
\end{tikzcd}
\]
is Cartesian.  This proves Nisnevich descent on eventually coconnective
distinguished squares.  Spectral convergence extends the result to arbitrary
connective Nisnevich distinguished squares.  Hence $Y$ is a spectral
Nisnevich sheaf.  It already satisfies Postnikov convergence and nilpotent
excision by hypothesis.  The characterization of the global spectral
deformation reflector from the differentiation chapter therefore identifies
this native sheaf as $L_S^{\mathrm{Def}}$-local.
\end{proof}

\begin{definition}[Affine spectral materialization]
\label{def:bnli-affine-materialization}
For $F\in\operatorname{FMP}^{\mathrm{sp}}_{B/R}$ define
\[
\boxed{
    \operatorname{Mat}^{\mathrm{sp}}_{B/R}(F)
    :=
    \operatorname{oblv}
    \left(
      \widetilde v_!u_!F
    \right),}
\]
where $\operatorname{oblv}$ is the forgetful equivalence of
Proposition~\ref{prop:bnli-forget-inf-target}.  Proposition~\ref{prop:bnli-two-stage-laft}
shows that the centre is spectrally laft, and
Theorem~\ref{thm:bnli-formal-stack-nisnevich} shows that the target is already
a spectral Nisnevich sheaf and, once regarded as a native sheaf,
$L_V^{\mathrm{Def}}$-local.  No additional sheafification is inserted into the
definition.
\end{definition}

\begin{theorem}[Affine materialization adjunction and recognition]
\label{thm:bnli-affine-materialization-equivalence}
Artinian restriction admits affine materialization as a fully faithful left
adjoint on centred deformation-local targets,
\[
\boxed{
    \operatorname{Mat}^{\mathrm{sp}}_{B/R}:
    \operatorname{FMP}^{\mathrm{sp}}_{B/R}
    \rightleftarrows
    \operatorname{DefTgt}(U/V):
    \operatorname{Leaf}^{\mathrm{Art}}_{B/R}.}
\]
Its essential image is precisely
$\operatorname{ModuliStk}^{\mathrm{sp}}(U/V)$.  Hence restriction gives inverse
equivalences
\[
\boxed{
    \operatorname{ModuliStk}^{\mathrm{sp}}(U/V)
    \simeq
    \operatorname{FMP}^{\mathrm{sp}}_{B/R}.}
\]
For every $Y\in\operatorname{DefTgt}(U/V)$ there is a canonical evaluation
counit
\[
\boxed{
    \epsilon_Y^{\mathrm{mat}}:
    \operatorname{Mat}^{\mathrm{sp}}_{B/R}
    \operatorname{Leaf}^{\mathrm{Art}}_{B/R}(Y)
    \longrightarrow Y.}
\]
It is an equivalence exactly when the centre belongs to the spectrally laft
nil-centred essential image.
\end{theorem}

\begin{proof}
Let $Y\in\operatorname{DefTgt}(U/V)$ and form its pro-infinitesimal
completion
\[
    Y_U^{\mathrm{pinf}}
    =
    Y\times_{\mathsf P^{\mathrm{sp}}Y}
    \mathsf P^{\mathrm{sp}}U.
\]
Because $Y$ is deformation-local it is spectrally convergent, and
$\mathsf P^{\mathrm{sp}}Y$ and $\mathsf P^{\mathrm{sp}}U$ are convergent by
construction.  Hence $Y_U^{\mathrm{pinf}}$ is a convergent prestack over
$U^{\mathrm{inf,sp}}$.

For $F\in\operatorname{FMP}^{\mathrm{sp}}_{B/R}$, the materialized source
$\operatorname{Mat}^{\mathrm{sp}}(F)$ is nil-isomorphic to $U$.  Therefore
Proposition~\ref{prop:bnli-pro-infinitesimal-universal} and
Proposition~\ref{prop:bnli-forget-inf-target} give
\[
\begin{aligned}
\operatorname{Map}_{U/}
\left(
  \operatorname{Mat}^{\mathrm{sp}}(F),Y
\right)
&\simeq
\operatorname{Map}_{U/,/U^{\mathrm{inf,sp}}}
\left(
  \widetilde v_!u_!F,
  Y_U^{\mathrm{pinf}}
\right)\\
&\simeq
\operatorname{Map}
\left(
  u_!F,
  \widetilde v^*Y_U^{\mathrm{pinf}}
\right)\\
&\simeq
\operatorname{Map}
\left(
  F,
  u^*\widetilde v^*Y_U^{\mathrm{pinf}}
\right).
\end{aligned}
\]
For an Artinian neighbourhood $U\to T$, the last right-adjoint value is,
by Proposition~\ref{prop:bnli-pro-infinitesimal-universal},
\[
    \widetilde v^*Y_U^{\mathrm{pinf}}(T)
    \simeq
    \operatorname{Map}_{U/}(T,Y).
\]
After restriction to the Artinian subcategory this is exactly
$\operatorname{Leaf}^{\mathrm{Art}}_{B/R}(Y)$.  Hence
\[
\boxed{
\operatorname{Map}_{U/}
\left(
  \operatorname{Mat}^{\mathrm{sp}}_{B/R}(F),Y
\right)
\simeq
\operatorname{Map}_{\operatorname{FMP}^{\mathrm{sp}}_{B/R}}
\left(
  F,
  \operatorname{Leaf}^{\mathrm{Art}}_{B/R}(Y)
\right).}
\]
This is the required adjunction.

The unit is the composite of the unit equivalences in
Theorem~\ref{thm:bnli-first-stage-recognition} and
Theorem~\ref{thm:bnli-second-stage-equivalence}, so
$\operatorname{Mat}^{\mathrm{sp}}$ is fully faithful.
Proposition~\ref{prop:bnli-two-stage-laft} and
Theorem~\ref{thm:bnli-formal-stack-nisnevich} show that its image consists of
spectrally laft nil-centred deformation-local targets.

Conversely suppose $U\to Y$ is deformation-local, spectrally laft, and
nil-centred.  Deformation locality supplies ambient Postnikov convergence and
ambient nilpotent excision, hence restricts to deformation theory on
unrestricted nilpotent neighbourhoods.  The hypotheses of
Proposition~\ref{prop:bnli-spectral-laft-implies-nil-laft} are therefore
satisfied, and the restricted dotted-lift centre is nil-laft.  Since $U\to Y$ is a spectral
nil-isomorphism,
\[
    \mathsf P^{\mathrm{sp}}U
    \xrightarrow{\sim}
    \mathsf P^{\mathrm{sp}}Y.
\]
Thus the projection
\[
    Y_U^{\mathrm{pinf}}
    =
    Y\times_{\mathsf P^{\mathrm{sp}}Y}\mathsf P^{\mathrm{sp}}U
    \longrightarrow Y
\]
is an equivalence.  The first-stage recognition theorem therefore identifies
the restriction of $Y$ to unrestricted nilpotent neighbourhoods with the left
Kan extension of its Artinian leaf, and the second-stage counit identifies the
resulting convergent diagram with $Y_U^{\mathrm{pinf}}\simeq Y$.  Hence the
materialization counit
\[
    \operatorname{Mat}^{\mathrm{sp}}
    \operatorname{Leaf}^{\mathrm{Art}}(Y)
    \longrightarrow Y
\]
is an equivalence.

Therefore the essential image is precisely
$\operatorname{ModuliStk}^{\mathrm{sp}}(U/V)$, and Artinian restriction and
materialization restrict to inverse equivalences on that subcategory.  The
final counit criterion is the usual essential-image criterion for a fully
faithful left adjoint.
\end{proof}

\begin{corollary}[Affine Artinian density]
\label{thm:bnli-affine-density}
For every $F\in\operatorname{FMP}^{\mathrm{sp}}_{B/R}$ the adjunction unit is a
canonical equivalence
\[
\boxed{
    F
    \xrightarrow{\sim}
    \operatorname{Leaf}^{\mathrm{Art}}_{B/R}
    \left(
      \operatorname{Mat}^{\mathrm{sp}}_{B/R}(F)
    \right).}
\]
\end{corollary}

\begin{remark}[No density or infinitesimal-target input]
There is no separate ``Artinian-dense'' structure and no retained choice of a
map to $U^{\mathrm{inf,sp}}$ on a formal-moduli stack.  Density is the proved
unit of the two-stage materialization equivalence, and the infinitesimal-target
map is contractibly determined on the laft nil-isomorphic locus.
\end{remark}

\section{Spectral-etale formal-completion descent and global materialization}
\label{sec:bnli-materialization-descent}

Assume throughout this section that $X/S$ lies in the locally coherent qcqs
geometric range of the coherent affine-etale coefficient system constructed in
the differentiation chapter.  Let
$\operatorname{AffEt}^{\mathrm{coh}}(X/S)$ be its coherent affine-etale
indexing category.  Under a change of centre, a formal-moduli stack is
transported by formal completion along the new centre, not by raw target base
change.  To globalize this transition without mistyping the target over $S$,
we first construct the relative reduced chart through which every coherent
affine-etale centre factors.

\begin{definition}[Reduced-centred formal completion]
\label{def:bnli-reduced-centred-completion}
Let $U'\to Y'$ be a centred target.  Define its reduced-centred completion by
\[
\boxed{
    \widehat Y'^{\,\mathrm{sp}}_{U'}
    :=
    Y'\times_{D_{\mathrm{red}}Y'}D_{\mathrm{red}}U'.}
\]
When $U'\to Y'$ is already a spectral formal-moduli stack,
Proposition~\ref{prop:bnli-formal-stack-reduced-centre} identifies the reduced
classical shadows, so the projection
$\widehat Y'^{\,\mathrm{sp}}_{U'}\to Y'$ is an equivalence.
\end{definition}

\begin{proposition}[Stability and etale base change of reduced-centred completion]
\label{prop:bnli-completion-stability}
Let $u:U\to Y$ be a centred spectral prestack with $Y$
$L_S^{\mathrm{Def}}$-local, and define
\[
    \widehat Y^{\,\mathrm{sp}}_U
    =
    Y\times_{D_{\mathrm{red}}Y}D_{\mathrm{red}}U.
\]
Then:
\begin{enumerate}
    \item $\widehat Y^{\,\mathrm{sp}}_U$ is
    $L_S^{\mathrm{Def}}$-local;
    \item if $U\to Y$ is spectrally laft, then both
    \[
        U\longrightarrow\widehat Y^{\,\mathrm{sp}}_U,
        \qquad
        \widehat Y^{\,\mathrm{sp}}_U\longrightarrow Y
    \]
    are spectrally laft;
    \item if $U\to Y$ is spectrally laft, then the completed centre
    \[
        U\longrightarrow\widehat Y^{\,\mathrm{sp}}_U
    \]
    is a spectral nil-isomorphism and hence a spectral formal-moduli stack;
    \item centred completion does not change the Artinian leaf:
    \[
        \operatorname{Leaf}^{\mathrm{Art}}_U
        \left(
          \widehat Y^{\,\mathrm{sp}}_U
        \right)
        \xrightarrow{\sim}
        \operatorname{Leaf}^{\mathrm{Art}}_U(Y);
    \]
    \item if $U'\to U$ is spectral etale and $U\to Y$ is a spectrally
    laft nil-isomorphism, then after the evident base change to the common
    spectral base the canonical square
    \[
    \begin{tikzcd}
    U' \arrow[r] \arrow[d] &
    U \arrow[d]\\
    \widehat Y^{\,\mathrm{sp}}_{U'} \arrow[r] &
    Y
    \end{tikzcd}
    \]
    is Cartesian.
\end{enumerate}
All comparisons are natural and coherent under spectral-etale composition and
pullback.
\end{proposition}

\begin{proof}
For a connective affine test $C$ one has
\[
\boxed{
    \widehat Y^{\,\mathrm{sp}}_U(C)
    \simeq
    Y(C)\times_{Y(H(\pi_0C)_{\mathrm{red}})}
    U(H(\pi_0C)_{\mathrm{red}}).}
\]
Postnikov truncations of $C$ have the same reduced degree-zero ring.  Hence
Postnikov convergence of $Y$ and of the represented centre $U$ gives
convergence of the completed target.

For nilpotent excision, let
\[
    C\simeq A_0\times_{A_{01}}A_1
\]
be a pullback of connective spectral rings in which
$A_0\to A_{01}$ is surjective on $\pi_0$ with nilpotent kernel.  Put
\[
    A_{i,\mathrm{red}}
    :=H((\pi_0A_i)_{\mathrm{red}}),
    \qquad
    C_{\mathrm{red}}
    :=H((\pi_0C)_{\mathrm{red}}).
\]
Then
\[
    A_{0,\mathrm{red}}\xrightarrow{\sim}A_{01,\mathrm{red}},
    \qquad
    C_{\mathrm{red}}\xrightarrow{\sim}A_{1,\mathrm{red}},
\]
the second equivalence being obtained by base change of the nilpotent leg.
Substituting the evaluation formula at the four vertices and using pullback
pasting, nilpotent excision for $Y$, and the represented pullback equation for
$U$ gives
\[
\begin{aligned}
    \widehat Y_U^{\,\mathrm{sp}}(A_0)
    \times_{\widehat Y_U^{\,\mathrm{sp}}(A_{01})}
    \widehat Y_U^{\,\mathrm{sp}}(A_1)
    &\simeq
    Y(C)\times_{Y(C_{\mathrm{red}})}U(C_{\mathrm{red}})\\
    &\simeq
    \widehat Y_U^{\,\mathrm{sp}}(C).
\end{aligned}
\]
Thus the completed target is deformation-local.

Let $\{C_\alpha\}$ be a filtered diagram of uniformly $n$-truncated
connective rings with colimit $C$.  Reduction commutes with this filtered
colimit:
\[
\boxed{
    \operatorname*{colim}_\alpha
    (\pi_0C_\alpha)_{\mathrm{red}}
    \xrightarrow{\sim}
    (\pi_0C)_{\mathrm{red}}.}
\]
Indeed, an element represented at a finite stage becomes nilpotent in the
colimit precisely when some finite power vanishes at a later finite stage.
Apply the laft Cartesian equation for $U\to Y$ to $\{C_\alpha\}$ and to the
reduced discrete diagram
$\{H((\pi_0C_\alpha)_{\mathrm{red}})\}$.  Filtered colimits of spaces commute
with finite limits, so the evaluation formula proves the relative laft
equations for both
\[
    U\to\widehat Y_U^{\,\mathrm{sp}}
    \quad\text{and}\quad
    \widehat Y_U^{\,\mathrm{sp}}\to Y.
\]

Apply the left-exact idempotent reduced-classical reflector to the defining
pullback.  It gives
\[
    D_{\mathrm{red}}\widehat Y_U^{\,\mathrm{sp}}
    \simeq
    D_{\mathrm{red}}U.
\]
Thus the completed centre is a reduced-classical equivalence.  Since its
target is deformation-local and its centre is spectrally laft,
Proposition~\ref{prop:bnli-formal-stack-reduced-centre} makes it a spectral
nil-isomorphism.

For an Artinian thickening the source and centre have the same reduced
classical affine.  Hence in the homotopy fibre defining the Artinian leaf the
extra reduced component is uniquely forced by the prescribed centre.  This
proves the fourth assertion.

Finally suppose $U'\to U$ is spectral etale and $U\to Y$ is laft and
nil-centred.  Then
$D_{\mathrm{red}}U\simeq D_{\mathrm{red}}Y$, while reduced-classical
Cartesianity of spectral-etale maps gives
\[
    U'
    \simeq
    U\times_{D_{\mathrm{red}}U}D_{\mathrm{red}}U'.
\]
Therefore
\[
\begin{aligned}
    U\times_Y\widehat Y^{\,\mathrm{sp}}_{U'}
    &\simeq
    U\times_Y
    \left(
      Y\times_{D_{\mathrm{red}}Y}D_{\mathrm{red}}U'
    \right)\\
    &\simeq
    U\times_{D_{\mathrm{red}}U}D_{\mathrm{red}}U'\\
    &\simeq U'.
\end{aligned}
\]
This is the displayed Cartesian square.  Coherence follows from functoriality
of finite pullbacks and the coherent reduced-classical spectral-etale
comparison.
\end{proof}

\begin{proposition}[Affine materialization exchange by formal completion]
\label{prop:bnli-affine-materialization-etale}
Let
\[
    f:(U',V')\longrightarrow(U,V)
\]
be a coherent spectral-etale change of affine pair and let
\[
    F':=f_{\mathrm{FMP},!}F.
\]
Put
\[
    Y:=\operatorname{Mat}^{\mathrm{sp}}_{U/V}(F),
    \qquad
    Y_{V'}:=Y\times_VV'.
\]
Then there is a canonical equivalence of centred formal-moduli stacks over
$V'$,
\[
\boxed{
    \operatorname{Mat}^{\mathrm{sp}}_{U'/V'}(F')
    \xrightarrow{\sim}
    \widehat{Y_{V'}}^{\,\mathrm{sp}}_{U'}.}
\]
These equivalences are coherent for identities, composition, and higher
spectral-etale pasting.
\end{proposition}

\begin{proof}
Coherent two-variable Artinian transport gives
\[
    f_{\mathrm{FMP},!}
    \operatorname{Leaf}^{\mathrm{Art}}_{U/V}(Y)
    \xrightarrow{\sim}
    \operatorname{Leaf}^{\mathrm{Art}}_{U'/V'}(Y_{V'}).
\]
The left side is $F'$.  By
Proposition~\ref{prop:bnli-completion-stability}, completion along $U'$ keeps
the Artinian leaf and produces a spectral formal-moduli stack.  Hence
\[
    \operatorname{Leaf}^{\mathrm{Art}}_{U'/V'}
    \left(
      \widehat{Y_{V'}}^{\,\mathrm{sp}}_{U'}
    \right)
    \simeq F'.
\]
Full faithfulness of affine materialization therefore identifies the completed
target with
$\operatorname{Mat}^{\mathrm{sp}}_{U'/V'}(F')$.  Coherence is inherited from
two-variable Artinian transport, reduced-classical Cartesianity, and the unit
and counit coherences of the affine materialization equivalence.
\end{proof}

\subsection{Relative reduced chart objects}

The relative reduced objects used for global descent are native objects of the
spectral Nisnevich topos.  They are not asserted to be represented spectral
schemes.  What is needed for descent is instead deformation locality of these
objects and representability of the open chart morphisms used in a chosen
finite cover.

\begin{definition}[Relative reduced chart]
\label{def:bnli-relative-reduced-chart}
Define the relative reduced target of $X/S$ by
\[
\boxed{
    W_{X/S}
    :=
    S\times_{D_{\mathrm{red}}S}D_{\mathrm{red}}X.}
\]
For a coherent affine-etale pair
\[
    V\longrightarrow S,
    \qquad
    U\longrightarrow X\times_SV,
\]
define
\[
\boxed{
    W_{U/V}
    :=
    V\times_{D_{\mathrm{red}}V}D_{\mathrm{red}}U.}
\]
The maps $V\to S$ and $D_{\mathrm{red}}U\to D_{\mathrm{red}}X$ induce a
canonical map
\[
    W_{U/V}\longrightarrow W_{X/S}.
\]
These are objects of the native spectral Nisnevich $\infty$-topos.  No
representability assertion is part of the definition.
\end{definition}

\begin{lemma}[Reduced-classical locality implies deformation locality]
\label{lem:bnli-dred-implies-def}
Every $D_{\mathrm{red}}$-local native sheaf is
$L_S^{\mathrm{Def}}$-local.  Consequently $W_{X/S}$ and every $W_{U/V}$ are
$L_S^{\mathrm{Def}}$-local when regarded over $S$.
\end{lemma}

\begin{proof}
Let $Z$ be $D_{\mathrm{red}}$-local.  For every connective affine ring $A$,
\[
    Z(A)\simeq Z(H(\pi_0A)_{\mathrm{red}}).
\]
The rings $A$ and $\tau_{\leq n}A$ have the same reduced classical
truncation.  Hence the Postnikov evaluation diagram for $Z$ is naturally
constant and
\[
    Z(A)\xrightarrow{\sim}\lim_nZ(\tau_{\leq n}A).
\]

Now let
\[
    C\simeq A_0\times_{A_{01}}A_1
\]
be a nilpotent-excision pullback with
$\pi_0(A_0)\twoheadrightarrow\pi_0(A_{01})$ nilpotent.  Reduction identifies
\[
    (A_0)_{\mathrm{red}}\simeq(A_{01})_{\mathrm{red}},
    \qquad
    C_{\mathrm{red}}\simeq(A_1)_{\mathrm{red}}.
\]
Therefore
\[
\begin{aligned}
    Z(A_0)\times_{Z(A_{01})}Z(A_1)
    &\simeq
    Z((A_0)_{\mathrm{red}})
    \times_{Z((A_{01})_{\mathrm{red}})}
    Z((A_1)_{\mathrm{red}})\\
    &\simeq Z((A_1)_{\mathrm{red}})\\
    &\simeq Z(C).
\end{aligned}
\]
Thus $Z$ satisfies the two deformation-locality equations.

It remains to prove the asserted consequence for the relative reduced chart
objects without regarding $D_{\mathrm{red}}X$ as an object over $S$.  Let
$t:T=\operatorname{Spec}A\to S$ be a connective affine test, put
\[
    T_{\mathrm{red}}
    :=
    \operatorname{Spec}H((\pi_0A)_{\mathrm{red}}),
\]
and let $t_{\mathrm{red}}\in S(T_{\mathrm{red}})$ be the restriction of the
structural point.  The defining pullback for $W_{X/S}$ gives the canonical
evaluation formula
\[
\boxed{
    \operatorname{Map}_S(T,W_{X/S})
    \simeq
    \operatorname{fib}_{t_{\mathrm{red}}}
    \left(
      X(T_{\mathrm{red}})
      \longrightarrow
      S(T_{\mathrm{red}})
    \right).}
\]
Every Postnikov truncation of $A$ has the same reduced degree-zero ring.
Hence the preceding formula identifies the Postnikov evaluation diagram of
$W_{X/S}$ over the compatible structural points with a constant diagram, and
therefore
\[
    \operatorname{Map}_S(T,W_{X/S})
    \xrightarrow{\sim}
    \lim_n
    \operatorname{Map}_S
    (\operatorname{Spec}\tau_{\leq n}A,W_{X/S}).
\]

Now let
\[
    C\simeq A_0\times_{A_{01}}A_1
\]
be a nilpotent-excision pullback over $S$, with
$\pi_0(A_0)\twoheadrightarrow\pi_0(A_{01})$ having nilpotent kernel.  As
above,
\[
    (A_0)_{\mathrm{red}}\simeq(A_{01})_{\mathrm{red}},
    \qquad
    C_{\mathrm{red}}\simeq(A_1)_{\mathrm{red}}.
\]
The structural points induced from the fixed map $\operatorname{Spec}C\to S$
are compatible with these equivalences.  Applying the boxed fibre formula at
the four vertices therefore gives canonical equivalences
\[
    \operatorname{Map}_S(\operatorname{Spec}A_0,W_{X/S})
    \xrightarrow{\sim}
    \operatorname{Map}_S(\operatorname{Spec}A_{01},W_{X/S})
\]
and
\[
    \operatorname{Map}_S(\operatorname{Spec}C,W_{X/S})
    \xrightarrow{\sim}
    \operatorname{Map}_S(\operatorname{Spec}A_1,W_{X/S}).
\]
Consequently
\[
    \operatorname{Map}_S(\operatorname{Spec}C,W_{X/S})
    \xrightarrow{\sim}
    \operatorname{Map}_S(\operatorname{Spec}A_0,W_{X/S})
    \times_{\operatorname{Map}_S(\operatorname{Spec}A_{01},W_{X/S})}
    \operatorname{Map}_S(\operatorname{Spec}A_1,W_{X/S}).
\]
Thus $W_{X/S}$ satisfies both deformation-locality equations and is
$L_S^{\mathrm{Def}}$-local.

Applying the same argument over $V$ to the morphism $U\to V$ shows that
$W_{U/V}$ is $L_V^{\mathrm{Def}}$-local.  Since $V\to S$ is represented,
Lemma~\ref{lem:bnli-def-local-base-composition} then makes the composite
$W_{U/V}\to V\to S$ $L_S^{\mathrm{Def}}$-local.
\end{proof}

\begin{proposition}[Relative reduced Cartesianity of coherent affine-etale charts]
\label{prop:bnli-relative-reduced-cartesianity}
For every coherent affine-etale pair $(U,V)$ there is a canonical Cartesian
square in the native spectral Nisnevich topos
\[
\begin{tikzcd}
U \arrow[r] \arrow[d] & X \arrow[d]\\
W_{U/V} \arrow[r] & W_{X/S}.
\end{tikzcd}
\]
Equivalently,
\[
\boxed{
    U
    \xrightarrow{\sim}
    X\times_{W_{X/S}}W_{U/V}.}
\]
The equivalence is coherent under morphisms of coherent affine-etale pairs.
\end{proposition}

\begin{proof}
Left exactness of $D_{\mathrm{red}}$ gives
\[
    D_{\mathrm{red}}(X\times_SV)
    \simeq
    D_{\mathrm{red}}X
    \times_{D_{\mathrm{red}}S}
    D_{\mathrm{red}}V.
\]
Consequently
\[
\begin{aligned}
    X\times_{W_{X/S}}W_{U/V}
    &\simeq
    (X\times_SV)
    \times_{D_{\mathrm{red}}(X\times_SV)}
    D_{\mathrm{red}}U.
\end{aligned}
\]
The map $U\to X\times_SV$ is spectral etale.  Reduced-classical Cartesianity
of spectral-etale maps identifies the last pullback with $U$.  Naturality and
higher coherence are inherited from finite pullback and the
reduced-classical Cartesianity theorem.
\end{proof}

\begin{lemma}[Representable open reduced charts]
\label{lem:bnli-reduced-open-charts}
Suppose $V\hookrightarrow S$ is a quasi-compact open immersion and
$U\hookrightarrow X\times_SV$ is a quasi-compact open immersion.  Then
\[
    W_{U/V}\longrightarrow W_{X/S}
\]
is a representable quasi-compact open immersion.

If $\{V_i\to S\}_{i=1}^r$ is a finite affine open cover and, for each $i$,
$\{U_{ij}\to X\times_SV_i\}_{j=1}^{r_i}$ is a finite affine open cover, then
\[
    \coprod_{i,j}W_{U_{ij}/V_i}
    \longrightarrow W_{X/S}
\]
is a finite open effective epimorphism.  Every finite intersection in its
Cech nerve admits a finite cover by objects $W_{U'/V'}$ coming from coherent
affine open refinements of the participating pairs.
\end{lemma}

\begin{proof}
Let $T=\operatorname{Spec}C\to W_{X/S}$ be an affine spectral test and put
$T_{\mathrm{red}}:=\operatorname{Spec}H(\pi_0C)_{\mathrm{red}}$.  By the
definition of $D_{\mathrm{red}}$, a point of $W_{X/S}(T)$ is a compatible
pair
\[
    s:T\to S,
    \qquad
    x_0:T_{\mathrm{red}}\to X
\]
whose composites to $D_{\mathrm{red}}S$ agree.  A lift of this point to
$W_{U/V}$ is precisely the conjunction of two open conditions:
\begin{enumerate}
    \item the map $s$ factors through the open $V\subseteq S$;
    \item the induced reduced map to $X\times_SV$ factors through the open
    $U\subseteq X\times_SV$.
\end{enumerate}
The first condition is represented by an open spectral subscheme of $T$.  The
second determines an open subscheme of $T_{\mathrm{red}}$; the equivalence
between open spectral subschemes of $T$ and open subschemes of its classical
truncation lifts it uniquely to an open spectral subscheme of $T$.  Their
intersection represents
\[
    T\times_{W_{X/S}}W_{U/V}.
\]
Quasi-compactness follows from quasi-compactness of the two given opens.  This
proves representability and the quasi-compact open assertion.

For the chosen finite covers, every point of $W_{X/S}$ is locally in some
$V_i$ and, after that restriction, its reduced $X$-component is locally in
some $U_{ij}$.  Thus the displayed finite family is an open covering family
and hence an effective epimorphism in the native topos.

A finite intersection imposes only finitely many open conditions on the base
and on the reduced $X$-component.  In the locally coherent qcqs range those
intersections are qcqs and admit finite coherent affine open refinements.
Applying the same construction to such a refinement gives a finite cover by
relative reduced objects $W_{U'/V'}$ refining every participating chart.
\end{proof}

\begin{proposition}[Formal completion as pullback on relative reduced charts]
\label{prop:bnli-completion-relative-reduced-pullback}
Let $U\to Y$ be a spectral formal-moduli stack over $V$, and let
\[
    f:(U',V')\longrightarrow(U,V)
\]
be a coherent spectral-etale change of affine pair.  Then $Y$ has a canonical
map $Y\to W_{U/V}$ and there is a canonical equivalence
\[
\boxed{
    \widehat{(Y\times_VV')}^{\,\mathrm{sp}}_{U'}
    \xrightarrow{\sim}
    Y\times_{W_{U/V}}W_{U'/V'}.}
\]
Thus formal-completion transport is ordinary Cartesian transport after
passing to the relative reduced chart system.
\end{proposition}

\begin{proof}
Reduced-centre detection gives
$D_{\mathrm{red}}Y\simeq D_{\mathrm{red}}U$.  Together with the structural
map $Y\to V$, this gives
\[
    Y\longrightarrow
    V\times_{D_{\mathrm{red}}V}D_{\mathrm{red}}U
    =W_{U/V}.
\]
Now expand the completed target:
\[
\begin{aligned}
    \widehat{(Y\times_VV')}^{\,\mathrm{sp}}_{U'}
    &\simeq
    (Y\times_VV')
    \times_{D_{\mathrm{red}}(Y\times_VV')}
    D_{\mathrm{red}}U'\\
    &\simeq
    (Y\times_VV')
    \times_{D_{\mathrm{red}}U
      \times_{D_{\mathrm{red}}V}D_{\mathrm{red}}V'}
    D_{\mathrm{red}}U'.
\end{aligned}
\]
On the other hand, expanding
$Y\times_{W_{U/V}}W_{U'/V'}$ gives the same finite pullback.  The
identification is therefore canonical and coherent.
\end{proof}

\begin{definition}[Formal-completion Cartesian sections]
\label{def:bnli-global-formal-moduli-stack-descent}
Let
\[
    \operatorname{CartFC}^{\mathrm{sp}}(X/S)
\]
be the $\infty$-category of Cartesian sections of the coherent affine-etale
formal-stack family whose fibre at $(U,V)$ is
$\operatorname{ModuliStk}^{\mathrm{sp}}(U/V)$ and whose transition functor
along
\[
    f:(U',V')\longrightarrow(U,V)
\]
is reduced-centred formal completion
\[
    Y\longmapsto
    \widehat{(Y\times_VV')}^{\,\mathrm{sp}}_{U'}.
\]
Equivalently, an object is a coherent family $(U,V)\mapsto Y_{U/V}$ for which
the canonical transition morphisms are equivalences
\[
    Y_{U'/V'}
    \xrightarrow{\sim}
    \widehat{(Y_{U/V}\times_VV')}^{\,\mathrm{sp}}_{U'}.
\]
By Proposition~\ref{prop:bnli-completion-relative-reduced-pullback}, the same
category is the category of Cartesian sections on the relative reduced chart
system:
\[
    Y_{U'/V'}
    \xrightarrow{\sim}
    Y_{U/V}\times_{W_{U/V}}W_{U'/V'}.
\]
This is a section category constructed from the transition functors; no
independently retained descent field is attached to a spectral partition-Lie
algebroid or to a global materialized target.
\end{definition}

\begin{lemma}[Finite subobject descent]
\label{lem:bnli-finite-subobject-descent}
Let $\mathcal X$ be an $\infty$-topos and let
\[
    Y_1,\ldots,Y_r\hookrightarrow Y
\]
be a finite family of subobjects whose union is $Y$.  Let $\mathcal P_r^\times$ be the finite poset of nonempty subsets of
$\{1,\ldots,r\}$ ordered by inclusion.  For $J\in\mathcal P_r^\times$ put
\[
    Y_J:=\bigcap_{j\in J}Y_j.
\]
The intersection assignment is a functor
$(\mathcal P_r^\times)^{\mathrm{op}}\to\mathcal X$.  The canonical finite
colimit map
\[
\boxed{
    \operatorname*{colim}_{J\in(\mathcal P_r^\times)^{\mathrm{op}}}
    Y_J
    \xrightarrow{\sim}Y}
\]
is an equivalence.  Consequently, for every $Z\in\mathcal X$,
\[
\boxed{
    \operatorname{Map}_{\mathcal X}(Y,Z)
    \xrightarrow{\sim}
    \operatorname*{lim}_{J\in\mathcal P_r^\times}
    \operatorname{Map}_{\mathcal X}(Y_J,Z).}
\]
Both indexing categories are finite.
\end{lemma}

\begin{proof}
For two subobjects the union is, by the universal calculus of subobjects in an
$\infty$-topos, the pushout
\[
    Y_1\amalg_{Y_1\cap Y_2}Y_2.
\]
Induct on $r$.  Writing $Y'=Y_1\cup\cdots\cup Y_{r-1}$, one has
\[
    Y\simeq Y'\amalg_{Y'\cap Y_r}Y_r.
\]
Distributivity of finite intersections over finite unions in the subobject
lattice identifies $Y'\cap Y_r$ with the union of the intersections
$Y_J\cap Y_r$ for nonempty $J\subseteq\{1,\ldots,r-1\}$.  Applying the
induction hypothesis to $Y'$ and to $Y'\cap Y_r$, then pasting the two finite
colimit diagrams, gives the stated finite colimit over all nonempty subsets.
Mapping a finite colimit into $Z$ gives the displayed finite limit.
\end{proof}

\begin{lemma}[Finite-open locality of spectral laftness]
\label{lem:bnli-finite-open-laft}
Let $f:Z\to Y$ be a morphism of spectrally convergent spectral Nisnevich
sheaves.  Suppose
\[
    \{Y_i\longrightarrow Y\}_{1\le i\le r}
\]
is a finite cover by quasi-compact open subobjects and every base change
\[
    Z\times_YY_i\longrightarrow Y_i
\]
is spectrally laft.  Then $f$ is spectrally laft.
\end{lemma}

\begin{proof}
Convergence is already assumed.  Fix $n\ge0$ and a filtered diagram of
$n$-truncated connective affine tests over $Y$,
\[
    C\simeq\operatorname*{colim}_\alpha C_\alpha.
\]
Pull the finite open cover back to $\operatorname{Spec}C$ and refine it by
finitely many principal opens, each subordinate to some $Y_i$.  The finitely
many defining elements and the finitely many relations expressing that they
cover the unit descend to one stage $\alpha_0$.  On the cofinal subsystem
$\alpha\ge\alpha_0$ the same principal opens give compatible finite covers of
all $\operatorname{Spec}C_\alpha$.

On each member of the cover and every nonempty finite intersection, the
structural map factors through one of the $Y_i$ after refinement.  Spectral
laftness is stable under affine base change, so the relative filtered-colimit
comparison is an equivalence on each such piece.  Apply
Lemma~\ref{lem:bnli-finite-subobject-descent} to the resulting finite family
of principal-open subobjects.  Sections are computed by the finite limit over
the nonempty intersection poset, not by an infinite Cech totalization.
Filtered colimits of spaces commute with this genuinely finite limit, and the
local equivalences therefore glue to
\[
    \operatorname*{colim}_\alpha Z(C_\alpha)
    \xrightarrow{\sim}
    Z(C)\times_{Y(C)}
    \operatorname*{colim}_\alpha Y(C_\alpha).
\]
This is the required relative laft equation.
\end{proof}

\begin{lemma}[Recovery of every coherent pair from a finite reduced-open basis]
\label{lem:bnli-basis-recovery}
Choose a finite coherent affine-open basis as in
Lemma~\ref{lem:bnli-reduced-open-charts}, and let
\[
    Y\longrightarrow W_{X/S}
\]
be the object obtained by effective descent from a formal-completion Cartesian
section on that basis.  Then for every coherent affine-etale pair $(U,V)$ the
canonical comparison
\[
\boxed{
    Y\times_{W_{X/S}}W_{U/V}
    \xrightarrow{\sim}
    Y_{U/V}}
\]
is an equivalence.  The comparison is independent, through a contractible
space, of the chosen finite basis.
\end{lemma}

\begin{proof}
Pull the finite open effective epimorphism
\[
    W=\coprod_aW_a\twoheadrightarrow W_{X/S}
\]
back along $W_{U/V}\to W_{X/S}$.  This gives a finite representable
quasi-compact open cover of $W_{U/V}$.  Every nonempty finite intersection in
that cover admits, by Lemma~\ref{lem:bnli-reduced-open-charts}, a finite cover
by relative reduced objects $W_{U'/V'}$ attached to coherent affine-open
refinements of $(U,V)$ and of the chosen basis pairs.

On every such refining object, the pullback of the descended object $Y$ is,
by its construction from the basis, the corresponding basis target.  The
formal-completion Cartesian-section equivalences identify that target with
the restriction of $Y_{U/V}$.  Hence the displayed comparison is an
equivalence after pullback to a finite Nisnevich open cover of every
intersection.  Both sides are spectral Nisnevich sheaves, so ordinary
effective descent makes the comparison an equivalence on each intersection
and then on $W_{U/V}$ itself.

If a second finite basis is chosen, the two bases admit a common finite
coherent affine-open refinement.  Pulling both descended objects to that
common refinement identifies them by the same Cartesian-section maps.
Effective descent gives a contractible space of comparison equivalences over
$W_{X/S}$.  Thus the recovery equivalence and the descended object are
basis-independent.
\end{proof}

\begin{proposition}[Descent of formal-completion Cartesian sections]
\label{prop:bnli-formal-completion-descent}
In the locally coherent qcqs range fixed above, restriction from global
centred targets induces a canonical equivalence
\[
\boxed{
    \operatorname{ModuliStk}^{\mathrm{sp}}(X/S)
    \xrightarrow{\sim}
    \operatorname{CartFC}^{\mathrm{sp}}(X/S).}
\]
Thus the global objects are recovered as the Cartesian sections of the
constructed formal-completion transition system.
\end{proposition}

\begin{proof}
Let $X\to Y$ be a global spectral formal-moduli stack.  Restrict it to every
coherent affine-etale pair $(U,V)$ and complete the restricted target along
$U$.  Proposition~\ref{prop:bnli-completion-stability} and
Proposition~\ref{prop:bnli-affine-materialization-etale} give the transition
equivalences, while
Proposition~\ref{prop:bnli-completion-relative-reduced-pullback} rewrites them
as Cartesian transition maps over the native relative reduced objects.  Their
cocycle coherences are inherited from finite pullback and the coherent
spectral-etale coefficient transport.  This constructs the restriction
functor in the statement.

Conversely, let a formal-completion Cartesian section be given.  Every
local target $Y_{U/V}$ is a spectral Nisnevich sheaf and is
$L_V^{\mathrm{Def}}$-local by affine materialization.  After composition with
$V\to S$, Lemma~\ref{lem:bnli-def-local-base-composition} makes it
$L_S^{\mathrm{Def}}$-local.

Choose a finite coherent affine open cover $\{V_i\to S\}$ and, for each $i$,
a finite coherent affine open cover
\[
    \{U_{ij}\to X\times_SV_i\}_j.
\]
Put
\[
    W_{ij}:=W_{U_{ij}/V_i},
    \qquad
    W:=\coprod_{i,j}W_{ij}.
\]
Lemma~\ref{lem:bnli-reduced-open-charts} gives a representable finite
quasi-compact open effective epimorphism
\[
\boxed{
    W\twoheadrightarrow W_{X/S}.}
\]
By Lemma~\ref{lem:bnli-dred-implies-def}, the target and all chart objects are
$L_S^{\mathrm{Def}}$-local.  Apply $L_S^{\mathrm{Def}}$ to the native Cech
presentation.  Left exactness preserves every finite fibre product in the
Cech nerve and the reflector, being a left adjoint, preserves its geometric
realization.  Since the target and chart objects are already local, the same
map is an effective epimorphism in the deformation-local $\infty$-topos
$\mathcal X^{\mathrm{Def}}_{/S}$.

By Proposition~\ref{prop:bnli-completion-relative-reduced-pullback}, the
supplied transition maps identify the restrictions of the local
$Y_{U_{ij}/V_i}$ by ordinary pullback along the $W$-charts.  Every finite
intersection in the Cech nerve admits, by
Lemma~\ref{lem:bnli-reduced-open-charts}, a finite cover by relative reduced
objects attached to coherent affine open refinements.  On each refining basis
piece all participating local targets are identified by the supplied
formal-completion transition equivalences.  Since the local targets are
Nisnevich sheaves, these basiswise identifications descend uniquely to the
whole intersection.  The higher transition coherences therefore give a
coherent simplicial descent diagram on
\[
    W_\bullet:=\check C_\bullet(W/W_{X/S}).
\]

Effective descent in $\mathcal X^{\mathrm{Def}}_{/S}$ produces a unique
object
\[
\boxed{
    Y\longrightarrow W_{X/S}.}
\]
Lemma~\ref{lem:bnli-basis-recovery} then supplies, for every coherent
affine-etale pair, the canonical Cartesian recovery equivalence
\[
\boxed{
    Y_{U/V}
    \xrightarrow{\sim}
    Y\times_{W_{X/S}}W_{U/V}.}
\]
Thus recovery on arbitrary coherent pairs is a theorem, not an implicit
extension of the chosen basis.  This is ordinary Cech descent in the
non-hypercomplete native topos; no hyperdescent or hypercompletion is used.

The local centres descend simultaneously.  Relative reduced Cartesianity
gives
\[
    U_{ij}
    \xrightarrow{\sim}
    X\times_{W_{X/S}}W_{ij}.
\]
Hence the maps $U_{ij}\to Y_{U_{ij}/V_i}$ are the pullbacks of a unique
global centre
\[
\boxed{
    X\longrightarrow Y.}
\]
Two choices of finite open basis admit a common finite refinement, and
effective descent gives a contractible space of comparison equivalences, so
the centred target is canonical.

It remains to verify spectral laftness and nil-centredness.  The maps
\[
    Y_{ij}:=Y\times_{W_{X/S}}W_{ij}\longrightarrow Y
\]
form a finite cover by representable quasi-compact open subobjects, because
$W_{ij}\to W_{X/S}$ has that property.  The base change of the global centre
to $Y_{ij}$ is the affine materialized centre
\[
    U_{ij}\longrightarrow Y_{ij}\simeq Y_{U_{ij}/V_i},
\]
which is spectrally laft.  Both $X$ and $Y$ are convergent, the former because
it is represented and the latter because it is deformation-local.  Therefore
Lemma~\ref{lem:bnli-finite-open-laft} makes $X\to Y$ spectrally laft.

Apply $D_{\mathrm{red}}$.  Left exactness and idempotence give
\[
    D_{\mathrm{red}}W_{X/S}\simeq D_{\mathrm{red}}X,
    \qquad
    D_{\mathrm{red}}W_{ij}\simeq D_{\mathrm{red}}U_{ij}.
\]
Thus
\[
    D_{\mathrm{red}}Y_{ij}
    \simeq
    D_{\mathrm{red}}Y
    \times_{D_{\mathrm{red}}X}
    D_{\mathrm{red}}U_{ij}.
\]
The local formal-moduli-stack condition identifies
$D_{\mathrm{red}}U_{ij}\to D_{\mathrm{red}}Y_{ij}$ as an equivalence.  Let
\[
    p_{\mathrm{red}}:
    D_{\mathrm{red}}Y
    \longrightarrow
    D_{\mathrm{red}}W_{X/S}
    \simeq
    D_{\mathrm{red}}X
\]
be the reduced structural map induced by $Y\to W_{X/S}$.  Its base change
along the effective open cover
$\coprod_{i,j}D_{\mathrm{red}}U_{ij}\to D_{\mathrm{red}}X$ is
\[
    D_{\mathrm{red}}Y_{ij}
    \longrightarrow
    D_{\mathrm{red}}U_{ij},
\]
which is an equivalence because the local reduced centre is an equivalence.
Effective descent therefore makes $p_{\mathrm{red}}$ an equivalence.  The
global reduced centre
\[
    D_{\mathrm{red}}X\longrightarrow D_{\mathrm{red}}Y
\]
is a section of $p_{\mathrm{red}}$, so it is the inverse equivalence.  By
Proposition~\ref{prop:bnli-formal-stack-reduced-centre}, the spectrally laft
deformation-local centre $X\to Y$ is a spectral nil-isomorphism.  Thus it is a
global spectral formal-moduli stack.

Finally, the Cartesian formulas and
Proposition~\ref{prop:bnli-completion-relative-reduced-pullback} show that
completion of the descended target along every coherent affine-etale centre
recovers the prescribed local target.  Conversely, starting with a global
formal-moduli stack, its local completions are precisely these Cartesian
pullbacks.  Restriction and descent are inverse equivalences.
\end{proof}

\begin{definition}[Global materialization]
\label{def:bnli-global-materialization}
For
\[
    F\in\operatorname{FMP}^{\mathrm{sp}}_{X/S},
\]
materialize every coherent affine-etale restriction $F_{U/V}$.  The
formal-completion exchange equivalences of
Proposition~\ref{prop:bnli-affine-materialization-etale} make the resulting
family an object of $\operatorname{CartFC}^{\mathrm{sp}}(X/S)$.  Define
\[
\boxed{
    \operatorname{Mat}^{\mathrm{sp}}_{X/S}(F)
    \in
    \operatorname{ModuliStk}^{\mathrm{sp}}(X/S)}
\]
to be the global centred target supplied by
Proposition~\ref{prop:bnli-formal-completion-descent}.
\end{definition}

\begin{theorem}[Global Materialization Theorem]
\label{thm:bnli-global-materialization}
There is a canonical equivalence
\[
\boxed{
    \operatorname{ModuliStk}^{\mathrm{sp}}(X/S)
    \simeq
    \operatorname{FMP}^{\mathrm{sp}}_{X/S}.}
\]
The inverse is global materialization.  More generally,
\[
    \operatorname{Mat}^{\mathrm{sp}}_{X/S}
    \dashv
    \operatorname{Leaf}^{\mathrm{Art}}_{X/S}
\]
on $\operatorname{DefTgt}(X/S)$, the unit on formal-moduli problems is an
equivalence, and the essential image of materialization is precisely the
spectrally laft nil-centred deformation-local targets.
\end{theorem}

\begin{proof}
The affine materialization adjunctions and their unit equivalences are
compatible with coherent spectral-etale transport of formal-moduli problems
and, by Proposition~\ref{prop:bnli-affine-materialization-etale}, with
formal-completion transport of materialized targets.  The relative reduced
chart calculation turns that transition law into ordinary Cartesian transport,
and Proposition~\ref{prop:bnli-formal-completion-descent} therefore glues the
affine materializations to the global functor
$\operatorname{Mat}^{\mathrm{sp}}_{X/S}$.

We prove the adjunction on the larger category of deformation-local centred
targets.  Let
\[
    M:=\operatorname{Mat}^{\mathrm{sp}}_{X/S}(F)
\]
and let $X\to Y$ be $L_S^{\mathrm{Def}}$-local.  Choose a finite open reduced
chart cover
\[
    W=\coprod_{i,j}W_{U_{ij}/V_i}\twoheadrightarrow W_{X/S}
\]
as in Lemma~\ref{lem:bnli-reduced-open-charts}, and pull it back to $M$.
Because the cover is effective, its pullback $M_\bullet\to M$ is an
effective Cech cover.  Put
\[
    X_n:=X\times_{W_{X/S}}W_n,
    \qquad
    M_n:=M\times_{W_{X/S}}W_n.
\]
A map $M\to Y$ under $X$ is equivalently a compatible family of maps
$M_n\to Y$ whose restriction to $X_n$ is the prescribed centre.  Hence
\[
    \operatorname{Map}_{X/}(M,Y)
    \simeq
    \operatorname*{Tot}_{[n]\in\Delta}
    \operatorname{Map}_{X_n/}(M_n,Y).
\]
Refining each finite intersection by coherent affine open pairs,
Proposition~\ref{prop:bnli-formal-completion-descent} identifies every local
source with the corresponding affine materialization
\[
    \operatorname{Mat}^{\mathrm{sp}}_{U/V}(F_{U/V}).
\]

For such a pair, a map from the local source to $Y$ under $U$ factors through
the base change $Y\times_SV$.  Spectral-etale base change of the deformation
reflector makes this target deformation-local over $V$.  The affine
materialization adjunction therefore gives a canonical equivalence
\[
\begin{aligned}
    \operatorname{Map}_{U/}
    \left(
      \operatorname{Mat}^{\mathrm{sp}}_{U/V}(F_{U/V}),
      Y\times_SV
    \right)
    \simeq
    \operatorname{Map}_{\operatorname{FMP}^{\mathrm{sp}}_{U/V}}
    \left(
      F_{U/V},
      \operatorname{Leaf}^{\mathrm{Art}}_{U/V}(Y)
    \right).
\end{aligned}
\]
These equivalences are compatible on refinements and on all Cech
intersections.  On the left this is functoriality of restriction and affine
materialization; on the right it is the spectral-etale cocartesianity of
centred Artinian restriction.  Taking the Cech totalization therefore gives
\[
\boxed{
\operatorname{Map}_{X/}
\left(
  \operatorname{Mat}^{\mathrm{sp}}_{X/S}(F),Y
\right)
\simeq
\operatorname{Map}_{\operatorname{FMP}^{\mathrm{sp}}_{X/S}}
\left(
  F,\operatorname{Leaf}^{\mathrm{Art}}_{X/S}(Y)
\right).}
\]
This is the asserted global adjunction.

The adjunction unit restricts on every coherent affine-etale pair to the unit
of the affine materialization adjunction, hence is an equivalence.  Therefore
global materialization is fully faithful.  Proposition~\ref{prop:bnli-formal-completion-descent}
identifies its essential image with the global spectral formal-moduli stacks.
Consequently global Artinian restriction and global materialization are
inverse equivalences
\[
    \operatorname{ModuliStk}^{\mathrm{sp}}(X/S)
    \simeq
    \operatorname{FMP}^{\mathrm{sp}}_{X/S}.
\]
The counit of the larger adjunction is an equivalence exactly on this
spectrally laft nil-centred essential image.
\end{proof}

\begin{remark}[Why raw target base change is not the transition law]
Even when $V'=V$, a proper spectral-etale change of centre $U'\to U$ does not
turn a formal neighbourhood of $U$ into a formal neighbourhood of $U'$ by raw
target base change.  The correct transition is completion of the base-changed
target along the new centre.  Proposition~\ref{prop:bnli-completion-relative-reduced-pullback}
shows that this becomes ordinary pullback only after passing to the relative
reduced chart system.  This variance is retained below and is not replaced by
an atlas-invariance assertion for native formal groupoids.
\end{remark}

\section{Native Artinian cells and their geometric recognition}
\label{sec:bnli-native-cells}

The formal-moduli-stack materialization is not itself an effective quotient of
$U$ in the native topos.  Native groupoid integration is therefore constructed
from a second geometric realization: the effective images of the individual
Artinian cells.

Continue affinely with $U=\operatorname{Spec}B$ over
$V=\operatorname{Spec}R$.

\begin{definition}[Native Artinian quotient cell]
\label{def:bnli-native-artinian-cell}
For a spectral Artinian arrow $A\to B$, let
\[
    U\xrightarrow{q_A}I_A\lhook\joinrel\longrightarrow\operatorname{Spec}A
\]
be its effective-image factorization and put
\[
\boxed{
    \mathfrak G_A
    :=
    \check C_\bullet(q_A).}
\]
By Proposition~\ref{prop:bnli-artinian-cell-formal},
$\mathfrak G_A$ is a native formal Lie groupoid on $U/V$.
\end{definition}

\begin{proposition}[Maps between native Artinian cells]
\label{prop:bnli-native-cell-fully-faithful}
For spectral Artinian arrows $A\to B$ and $C\to B$, postcomposition with
$I_C\hookrightarrow\operatorname{Spec}C$ and restriction along
$I_A\hookrightarrow\operatorname{Spec}A$ give canonical equivalences
\[
    \operatorname{Map}_{U/}(I_A,I_C)
    \xrightarrow{\sim}
    \operatorname{Map}_{U/}(I_A,\operatorname{Spec}C)
\]
and
\[
    \operatorname{Map}_{U/}(\operatorname{Spec}A,\operatorname{Spec}C)
    \xrightarrow{\sim}
    \operatorname{Map}_{U/}(I_A,\operatorname{Spec}C).
\]
Consequently the first equivalence followed by the inverse of the second
determines a canonical equivalence
\[
\boxed{
    \operatorname{Map}_{U/}(I_A,I_C)
    \xrightarrow{\sim}
    \operatorname{Map}_{U/}
    (\operatorname{Spec}A,\operatorname{Spec}C).}
\]
Equivalently,
\[
\boxed{
    \operatorname{Map}_{\operatorname{FormLieGpd}(U/V)}
    (\mathfrak G_A,\mathfrak G_C)
    \simeq
    \operatorname{Map}_{\operatorname{Art}^{\mathrm{sp}}_{R/B}}(C,A).}
\]
\end{proposition}

\begin{proof}
We first identify maps into $I_C$.  A map
$f:I_A\to\operatorname{Spec}C$ under $U$ pulls back the monomorphism
$I_C\hookrightarrow\operatorname{Spec}C$ to a subobject
\[
    I_A\times_{\operatorname{Spec}C}I_C\hookrightarrow I_A.
\]
Its pullback along the effective epimorphism $q_A:U\twoheadrightarrow I_A$ is
all of $U$, because the prescribed map $U\to\operatorname{Spec}C$ factors
through $I_C$ by definition of the effective image.  Effective epimorphisms
detect equivalences of subobjects, so the displayed monomorphism is an
equivalence.  Thus every map $I_A\to\operatorname{Spec}C$ under $U$ factors
through $I_C$, and the space of such factorizations is contractible because
$I_C\to\operatorname{Spec}C$ is a monomorphism.  Hence
\[
    \operatorname{Map}_{U/}(I_A,I_C)
    \xrightarrow{\sim}
    \operatorname{Map}_{U/}(I_A,\operatorname{Spec}C).
\]

Corollary~\ref{cor:bnli-artinian-effective-image-def} says that
$I_A\to\operatorname{Spec}A$ is an $L_V^{\mathrm{Def}}$-equivalence, while
$\operatorname{Spec}C$ is deformation-local.  The localization mapping-space
property therefore gives the restriction equivalence
\[
    \operatorname{Map}_{U/}(\operatorname{Spec}A,\operatorname{Spec}C)
    \xrightarrow{\sim}
    \operatorname{Map}_{U/}(I_A,\operatorname{Spec}C).
\]
Composing the first equivalence with the inverse of this restriction
equivalence proves the first formula.  The second follows from contravariant
affine Yoneda and the fixed-object formal-groupoid/effective-quotient
equivalence of Chapter~13.
\end{proof}

\begin{remark}[No raw compactness theorem is inserted]
The preceding proposition controls maps between finite native cells.  We do
not infer from it that $I_A$ is compact for arbitrary filtered colimits in the
big spectral Nisnevich topos, nor that $\mathfrak G_A$ is compact for arbitrary
filtered colimits of raw formal groupoids.  Such a statement would be a
separate native coherence theorem and is not needed for object-level
integration below.
\end{remark}

\section{Minimal affine native Artinian-cell reconstruction}
\label{sec:bnli-native-quotient}

Continue on a coherent affine pair
\[
    U=\operatorname{Spec}B,
    \qquad
    V=\operatorname{Spec}R,
\]
and let
\[
    F\in\operatorname{FMP}^{\mathrm{sp}}_{B/R}.
\]
The affine construction in this section is the minimal native realization using
the identity affine centre $U$.  It is useful for the local raw mapping
problem.  The global integration functor will be constructed later by inducing
all local Artinian cells to the common fixed object space $X$; no raw
changing-atlas invariance is assumed between the two constructions.

Write $\int F$ for the Grothendieck category of elements of the covariant
Artinian functor $F$.  Its objects are pairs $(A\to B,\xi)$ with
$\xi\in F(A)$.  Since $B\to B$ is terminal in the Artinian category and
$F(B)\simeq *$, the object $(B,*)$ is terminal in $\int F$.

\begin{lemma}[Canonical affine Artinian co-Yoneda resolution]
\label{lem:bnli-artinian-coyoneda-resolution}
If
\[
    h_A^{\mathrm{Art}}(C)
    :=
    \operatorname{Map}_{\operatorname{Art}^{\mathrm{sp}}_{R/B}}(A,C),
\]
then there is a canonical colimit equivalence
\[
\boxed{
    \operatorname*{colim}_{(A,\xi)\in(\int F)^{\mathrm{op}}}
    h_A^{\mathrm{Art}}
    \xrightarrow{\sim}
    F.}
\]
No cellular presentation of $F$ is chosen.
\end{lemma}

\begin{proof}
This is the covariant co-Yoneda formula.  A point $\xi\in F(A)$ determines a
natural transformation $h_A^{\mathrm{Art}}\to F$, and a morphism in
$\int F$ induces the opposite morphism between the corresponding
corepresentables.  Evaluation on every Artinian test gives the ordinary
category-of-elements colimit formula.
\end{proof}

\begin{definition}[Minimal affine native Artinian-cell quotient]
\label{def:bnli-effective-leaf-image}
Define
\[
\boxed{
    Q_F^{\mathrm{nat,aff}}
    :=
    \operatorname*{colim}_{(A,\xi)\in(\int F)^{\mathrm{op}}}I_A}
\]
in the undercategory $(\mathcal X_{/V})_{U/}$, and let
\[
    q_F^{\mathrm{nat,aff}}:
    U\longrightarrow Q_F^{\mathrm{nat,aff}}
\]
be the induced centre.  Put
\[
    \mathfrak G_F^{\mathrm{nat,aff}}
    :=
    \check C_\bullet(q_F^{\mathrm{nat,aff}}).
\]
\end{definition}

\begin{proposition}[The minimal affine native centre is effective]
\label{prop:bnli-affine-native-centre-effective}
The morphism $q_F^{\mathrm{nat,aff}}$ is an effective epimorphism.
\end{proposition}

\begin{proof}
Factor it as
\[
    U\twoheadrightarrow J_F\lhook\joinrel\longrightarrow
    Q_F^{\mathrm{nat,aff}}.
\]
For every Artinian element, the composite
$U\twoheadrightarrow I_A\to Q_F^{\mathrm{nat,aff}}$ factors through $J_F$.
Since $U\to I_A$ is effective, the entire cocone map
$I_A\to Q_F^{\mathrm{nat,aff}}$ factors through the monomorphism $J_F$.
Hence the colimit cocone factors through $J_F$, so the identity of
$Q_F^{\mathrm{nat,aff}}$ factors through that monomorphism.  It is therefore an
equivalence.
\end{proof}

\begin{theorem}[The minimal affine native quotient is formal]
\label{thm:bnli-affine-native-quotient-formal}
There is a canonical equivalence
\[
    L_V^{\mathrm{dR}}U
    \xrightarrow{\sim}
    L_V^{\mathrm{dR}}Q_F^{\mathrm{nat,aff}}.
\]
Consequently $\mathfrak G_F^{\mathrm{nat,aff}}$ is already a brave-new formal
Lie groupoid on $U/V$.
\end{theorem}

\begin{proof}
Every cell quotient $U\twoheadrightarrow I_A$ is a relative de Rham
equivalence by Proposition~\ref{prop:bnli-artinian-cell-formal}.  Relative de
Rham localization preserves the colimit defining $Q_F^{\mathrm{nat,aff}}$.
Because $(B,*)$ is terminal in $\int F$, the opposite indexing category is
weakly contractible, and therefore
\[
\begin{aligned}
    L_V^{\mathrm{dR}}Q_F^{\mathrm{nat,aff}}
    &\simeq
    \operatorname*{colim}_{(\int F)^{\mathrm{op}}}L_V^{\mathrm{dR}}I_A\\
    &\simeq
    \operatorname*{colim}_{(\int F)^{\mathrm{op}}}L_V^{\mathrm{dR}}U\\
    &\simeq L_V^{\mathrm{dR}}U.
\end{aligned}
\]
Together with Proposition~\ref{prop:bnli-affine-native-centre-effective}, the
formal quotient recognition theorem gives the groupoid statement.
\end{proof}

\begin{theorem}[Affine deformation formalization of the minimal native quotient]
\label{thm:bnli-affine-native-quotient-materialization}
There is a canonical comparison
\[
    \kappa_F^{\mathrm{mat,aff}}:
    Q_F^{\mathrm{nat,aff}}
    \longrightarrow
    \operatorname{Mat}^{\mathrm{sp}}_{B/R}(F)
\]
and it becomes an equivalence after deformation localization:
\[
\boxed{
    L_V^{\mathrm{Def}}Q_F^{\mathrm{nat,aff}}
    \xrightarrow{\sim}
    \operatorname{Mat}^{\mathrm{sp}}_{B/R}(F).}
\]
\end{theorem}

\begin{proof}
For every Artinian element $(A,\xi)$, the monomorphism
$I_A\hookrightarrow\operatorname{Spec}A$ followed by the Artinian evaluation
map to $\operatorname{Mat}^{\mathrm{sp}}(F)$ gives a compatible cocone.  This
defines $\kappa_F^{\mathrm{mat,aff}}$.

Apply $L_V^{\mathrm{Def}}$ to the defining colimit.  Since each
$I_A\to\operatorname{Spec}A$ is a deformation equivalence,
\[
    L_V^{\mathrm{Def}}Q_F^{\mathrm{nat,aff}}
    \simeq
    \operatorname*{colim}_{(A,\xi)\in(\int F)^{\mathrm{op}}}^{\mathrm{Def}}
    \operatorname{Spec}A.
\]
By Lemma~\ref{lem:bnli-artinian-coyoneda-resolution} and preservation of
colimits by the materialization left adjoint,
\[
\begin{aligned}
    \operatorname{Mat}^{\mathrm{sp}}(F)
    &\simeq
    \operatorname*{colim}_{(A,\xi)}
    \operatorname{Mat}^{\mathrm{sp}}(h_A^{\mathrm{Art}})\\
    &\simeq
    \operatorname*{colim}_{(A,\xi)}^{\mathrm{Def}}
    \operatorname{Spec}A.
\end{aligned}
\]
The second equivalence uses affine materialization recognition for the
represented Artinian target.  Thus the two sides are the same colimit in the
deformation-local category, and the displayed comparison is induced termwise
by $I_A\to\operatorname{Spec}A$.
\end{proof}

\begin{proposition}[Functoriality of the minimal affine reconstruction]
\label{prop:bnli-affine-native-functoriality}
The assignment
\[
    F\longmapsto\mathfrak G_F^{\mathrm{nat,aff}}
\]
defines a functor
\[
    \operatorname{FMP}^{\mathrm{sp}}_{B/R}
    \longrightarrow
    \operatorname{FormLieGpd}(U/V).
\]
\end{proposition}

\begin{proof}
A natural transformation $F\to G$ induces a functor of categories of elements
and the compatible morphisms of native cells supplied by
Proposition~\ref{prop:bnli-native-cell-fully-faithful}.  Colimit functoriality
produces a map of minimal affine quotients under $U$, and taking Cech nerves
produces the formal-groupoid morphism.  Identity and composition are inherited
from the Grothendieck construction and effective-image functoriality.
\end{proof}

\subsection{Affine native mapping discrepancy}

Let $G\in\operatorname{FMP}^{\mathrm{sp}}_{B/R}$ and let $A\to B$ be spectral
Artinian.  The affine comparison above defines
\[
\boxed{
    \operatorname{ev}^{\mathrm{nat,aff}}_{A,G}:
    \operatorname{Map}_{U/}(I_A,Q_G^{\mathrm{nat,aff}})
    \longrightarrow
    G(A).}
\]

\begin{proposition}[Canonical section of affine native evaluation]
\label{prop:bnli-affine-native-evaluation-section}
The map $\operatorname{ev}^{\mathrm{nat,aff}}_{A,G}$ admits a canonical section
\[
    s^{\mathrm{aff}}_{A,G}:
    G(A)
    \longrightarrow
    \operatorname{Map}_{U/}(I_A,Q_G^{\mathrm{nat,aff}}).
\]
\end{proposition}

\begin{proof}
A point $x\in G(A)$ is the object $(A,x)$ of the category of elements defining
$Q_G^{\mathrm{nat,aff}}$.  Its colimit cocone supplies the required map
$I_A\to Q_G^{\mathrm{nat,aff}}$.  After deformation localization this is the
represented Artinian evaluation map, so the composite with
$\operatorname{ev}^{\mathrm{nat,aff}}_{A,G}$ is the identity.
\end{proof}

\begin{definition}[Affine native cell discrepancy]
\label{def:bnli-affine-native-cell-discrepancy}
For $x\in G(A)$ define
\[
\boxed{
    \mathscr E^{\mathrm{nat,aff}}_{A,G}(x)
    :=
    \operatorname{fib}_{x}
    \left(
      \operatorname{ev}^{\mathrm{nat,aff}}_{A,G}
    \right).}
\]
\end{definition}

\begin{lemma}[Corepresentable affine cell]
\label{lem:bnli-affine-corepresentable-cell}
For the Artinian corepresentable $h_A^{\mathrm{Art}}$ there is a canonical
equivalence
\[
    Q_{h_A^{\mathrm{Art}}}^{\mathrm{nat,aff}}
    \xrightarrow{\sim}
    I_A.
\]
\end{lemma}

\begin{proof}
The category of elements of $h_A^{\mathrm{Art}}$ has the distinguished object
$(A,\operatorname{id}_A)$ terminal after passage to the opposite indexing
category used by co-Yoneda.  The defining colimit of native cells therefore
reduces to $I_A$.
\end{proof}

\begin{theorem}[Affine mapping discrepancy formula]
\label{thm:bnli-affine-mapping-discrepancy}
Let $F,G\in\operatorname{FMP}^{\mathrm{sp}}_{B/R}$ and
$\alpha:F\to G$.  The natural map
\[
    \operatorname{Map}_{\operatorname{FormLieGpd}(U/V)}
    (\mathfrak G_F^{\mathrm{nat,aff}},\mathfrak G_G^{\mathrm{nat,aff}})
    \longrightarrow
    \operatorname{Map}_{\operatorname{FMP}^{\mathrm{sp}}_{B/R}}(F,G)
\]
is split by the affine native reconstruction functor, and the fibre over
$\alpha$ is
\[
\boxed{
    \operatorname*{lim}_{(A,\xi)\in\int F}
    \mathscr E^{\mathrm{nat,aff}}_{A,G}
    \bigl(\alpha_A(\xi)\bigr).}
\]
Consequently the minimal affine native reconstruction is fully faithful if and
only if every $\operatorname{ev}^{\mathrm{nat,aff}}_{A,G}$ is an equivalence.
\end{theorem}

\begin{proof}
The fixed-object quotient/groupoid equivalence identifies the raw groupoid
mapping space with
\[
    \operatorname{Map}_{U/}
    (Q_F^{\mathrm{nat,aff}},Q_G^{\mathrm{nat,aff}})
    \simeq
    \operatorname*{lim}_{(A,\xi)\in\int F}
    \operatorname{Map}_{U/}(I_A,Q_G^{\mathrm{nat,aff}}).
\]
Co-Yoneda gives
\[
    \operatorname{Map}(F,G)
    \simeq
    \operatorname*{lim}_{(A,\xi)\in\int F}G(A).
\]

For an element $(A,\xi)$, restriction of a raw quotient morphism along the
cocone map
\[
    I_A\longrightarrow Q_F^{\mathrm{nat,aff}}
\]
and subsequent affine formal-leaf differentiation is, by naturality of
$\kappa^{\mathrm{mat,aff}}$, precisely postcomposition with
\[
    Q_G^{\mathrm{nat,aff}}
    \longrightarrow
    L_V^{\mathrm{Def}}Q_G^{\mathrm{nat,aff}}
    \simeq
    \operatorname{Mat}^{\mathrm{sp}}(G)
\]
followed by the Artinian evaluation at $A$.  This is exactly
$\operatorname{ev}^{\mathrm{nat,aff}}_{A,G}$.  Hence the comparison of mapping
spaces is the limit of the affine native evaluation maps.  Its section is the
limit of their canonical sections.

Fibres commute with limits of spaces, so the fibre over $\alpha$ is the
displayed limit of the cell discrepancies.  If every cell evaluation is an
equivalence, full faithfulness follows.  Conversely apply full faithfulness
to the corepresentable source $h_A^{\mathrm{Art}}$ and use
Lemma~\ref{lem:bnli-affine-corepresentable-cell}.
\end{proof}

\begin{remark}[Local raw rigidity remains a theorem, not an input condition]
The preceding theorem isolates the local raw statement
\[
    \operatorname{Map}_{U/}(I_A,Q_G^{\mathrm{nat,aff}})
    \xrightarrow{\sim}
    G(A).
\]
It is not proved here and is not renamed as deformation-completeness,
admissibility, materializability, or Artinian-density structure.
\end{remark}

\section{Global fixed-object Artinian-cell quotient and groupoid integration}
\label{sec:bnli-global-integration}

The global native groupoid is not obtained by gluing the preceding affine raw
groupoids along changing-object etale maps.  Instead, all local Artinian cells
are induced to the common fixed object space $X$ inside the native topos and
then assembled there.

\subsection{The global Artinian test category and co-Yoneda formula}

\begin{definition}[Global spectral Artinian test category]
\label{def:bnli-global-artinian-category}
Let
\[
    \mathscr{Art}^{\mathrm{sp}}_{X/S}
    \longrightarrow
    \operatorname{AffEt}^{\mathrm{coh}}(X/S)
\]
be the Grothendieck construction of the cocartesian family
\[
    (U,V)=(\operatorname{Spec}B,\operatorname{Spec}R)
    \longmapsto
    \operatorname{Art}^{\mathrm{sp}}_{R/B}
\]
constructed in the differentiation chapter.  An object will be denoted
\[
    c=((U,V),A\to B).
\]
A global formal-moduli problem
$F\in\operatorname{FMP}^{\mathrm{sp}}_{X/S}$ determines, by evaluation on this
Grothendieck construction, a functor
\[
    F^{\mathrm{Art}}:
    \mathscr{Art}^{\mathrm{sp}}_{X/S}
    \longrightarrow\mathcal S
\]
which carries the cocartesian etale arrows to the canonical transport
equivalences.
\end{definition}

\begin{definition}[Global category of Artinian elements]
\label{def:bnli-global-category-elements}
For global $F$, let
\[
    \int_XF
\]
be the Grothendieck category of elements of $F^{\mathrm{Art}}$.  Thus an object
is a tuple
\[
    e=(c,\xi)
    =((U_e,V_e),A_e\to B_e,\xi_e),
    \qquad
    \xi_e\in F_{U_e/V_e}(A_e).
\]
\end{definition}

\begin{lemma}[Global Artinian co-Yoneda mapping formula]
\label{lem:bnli-global-artinian-coyoneda}
For $F,G\in\operatorname{FMP}^{\mathrm{sp}}_{X/S}$ there is a canonical
equivalence
\[
\boxed{
    \operatorname{Map}_{\operatorname{FMP}^{\mathrm{sp}}_{X/S}}(F,G)
    \xrightarrow{\sim}
    \operatorname*{lim}_{e=(c,\xi)\in\int_XF}
    G_c(A_c).}
\]
\end{lemma}

\begin{proof}
Straightening identifies a global formal-moduli problem with a cocartesian
section of the coherent affine-etale formal-moduli family, equivalently with a
functor on $\mathscr{Art}^{\mathrm{sp}}_{X/S}$ satisfying the Schlessinger and
cocartesian transport conditions.  The inclusion of these sections into the
ambient functor category is fully faithful.  The ordinary category-of-elements
co-Yoneda formula in that ambient functor category gives the displayed limit,
and the cocartesian compatibility is already built into the morphisms of
$\int_XF$.
\end{proof}

\subsection{Inducing local cells to the fixed object space}

For
\[
    c=((U,V),A\to B)
    \in\mathscr{Art}^{\mathrm{sp}}_{X/S},
\]
put
\[
    T_c:=\operatorname{Spec}A,
\]
and let
\[
    U\twoheadrightarrow I_c\hookrightarrow T_c
\]
be the local native Artinian effective-image factorization.

\begin{definition}[Induced global native and represented cells]
\label{def:bnli-induced-global-cells}
Form the pushouts in the native slice $\mathcal X_{/S}$
\[
\boxed{
    \overline I_c
    :=
    X\amalg_U I_c,
    \qquad
    \overline T_c
    :=
    X\amalg_U T_c.}
\]
Both are regarded as centred objects under $X$ through the first coprojection.
\end{definition}

\begin{proposition}[Functoriality of induced global Artinian cells]
\label{prop:bnli-induced-cell-functoriality}
The assignments
\[
    c\longmapsto \overline I_c,
    \qquad
    c\longmapsto \overline T_c
\]
assemble canonically into functors
\[
\boxed{
    \overline I,\overline T:
    \left(\mathscr{Art}^{\mathrm{sp}}_{X/S}\right)^{\mathrm{op}}
    \longrightarrow
    (\mathcal X_{/S})_{X/},}
\]
and the morphisms $I_c\to T_c$ assemble into a natural transformation
\[
    \overline I\longrightarrow\overline T.
\]
These functors use the already constructed coherent spectral-etale transport
of Artinian tests; they do not assert raw atlas invariance of a native formal
groupoid.
\end{proposition}

\begin{proof}
A morphism in the Grothendieck category
$\mathscr{Art}^{\mathrm{sp}}_{X/S}$ factors, through a contractible space of
choices, as a cocartesian spectral-etale transport followed by a morphism in a
single affine Artinian fibre.  We treat the two factors separately.

For a morphism $A\to A'$ in one fibre
$\operatorname{Art}^{\mathrm{sp}}_{R/B}$, the associated geometric morphism is
$T_{A'}\to T_A$ under $U$.  Functoriality of the
effective-epimorphism--monomorphism factorization gives a canonical map
\[
    I_{A'}\longrightarrow I_A
\]
under $U$, compatible with $I_{A'}\to T_{A'}$ and $I_A\to T_A$.  Pushing out
along $U\to X$ gives the required morphisms
\[
    \overline I_{A'}\longrightarrow\overline I_A,
    \qquad
    \overline T_{A'}\longrightarrow\overline T_A
\]
under $X$.

Now let
$f:(U',V')\to(U,V)$ be a coherent spectral-etale change of affine pair and
let $c'$ be the cocartesian transport of
$c=((U,V),A\to B)$.  The Artinian transport theorem of the differentiation
chapter supplies a canonical pushout square of coordinate rings; on affine
spectra this gives a Cartesian square
\[
\begin{tikzcd}
U' \arrow[r] \arrow[d] & U \arrow[d]\\
T_{c'} \arrow[r] & T_c.
\end{tikzcd}
\]
Effective-image factorization is stable under pullback in an $\infty$-topos.
Since $U'\simeq U\times_{T_c}T_{c'}$, the effective image of
$U'\to T_{c'}$ is therefore canonically
\[
    I_{c'}\simeq I_c\times_{T_c}T_{c'}.
\]
The projection supplies $I_{c'}\to I_c$, compatible with $U'\to U$.  Together
with the fixed maps of both centres to $X$, this induces, by the two pushout
universal properties, canonical maps
\[
    X\amalg_{U'}I_{c'}\longrightarrow X\amalg_UI_c,
    \qquad
    X\amalg_{U'}T_{c'}\longrightarrow X\amalg_UT_c
\]
under $X$.

The factorization system is functorial and the spectral-etale Artinian
transport was constructed with identity, composition, and higher pullback
pasting coherences.  The two constructions above therefore agree on mixed
factorizations and assemble to functors on the opposite Grothendieck category.
Compatibility of effective-image factorization with the monomorphisms
$I_c\hookrightarrow T_c$ gives the displayed natural transformation.
\end{proof}

\begin{proposition}[Induced native cells are formal effective quotients]
\label{prop:bnli-induced-cell-formal-effective}
For every global Artinian test $c$, the centre
\[
    X\longrightarrow\overline I_c
\]
is an effective epimorphism and a relative de Rham equivalence.
Furthermore the canonical morphism
\[
    \overline I_c\longrightarrow\overline T_c
\]
is an $L_S^{\mathrm{Def}}$-equivalence.
\end{proposition}

\begin{proof}
The first centre is the pushout of the effective epimorphism
$U\twoheadrightarrow I_c$.  Effective epimorphisms are stable under pushout in
an $\infty$-topos, so it is effective.

The local cell map $U\to I_c$ is a relative de Rham equivalence over its
affine base $V$.  By
Lemma~\ref{lem:bnli-local-equivalence-base-composition}, postcomposition
$V\to S$ carries it to a relative de Rham equivalence over $S$.  Relative de
Rham localization is a left adjoint, so
\[
\begin{aligned}
    L_S^{\mathrm{dR}}\overline I_c
    &\simeq
    L_S^{\mathrm{dR}}X
    \amalg_{L_S^{\mathrm{dR}}U}
    L_S^{\mathrm{dR}}I_c\\
    &\simeq
    L_S^{\mathrm{dR}}X.
\end{aligned}
\]
The comparison $I_c\to T_c$ is an $L_S^{\mathrm{Def}}$-equivalence by
Corollary~\ref{cor:bnli-artinian-effective-image-def}.  Applying the
colimit-preserving deformation reflector to the two pushouts therefore shows
that $\overline I_c\to\overline T_c$ is also a deformation equivalence.
\end{proof}

\subsection{The global fixed-object native quotient}

\begin{definition}[Global native and represented Artinian realizations]
\label{def:bnli-global-native-quotient}
For $F\in\operatorname{FMP}^{\mathrm{sp}}_{X/S}$, restrict the functors of
Proposition~\ref{prop:bnli-induced-cell-functoriality} along the projection
$(\int_XF)^{\mathrm{op}}\to
(\mathscr{Art}^{\mathrm{sp}}_{X/S})^{\mathrm{op}}$.  Define
\[
\boxed{
    Q_F^{\mathrm{nat}}
    :=
    \operatorname*{colim}_{e=(c,\xi)\in(\int_XF)^{\mathrm{op}}}
    \overline I_c}
\]
and
\[
\boxed{
    Q_F^{\mathrm{Art,rep}}
    :=
    \operatorname*{colim}_{e=(c,\xi)\in(\int_XF)^{\mathrm{op}}}
    \overline T_c}
\]
in the undercategory $(\mathcal X_{/S})_{X/}$.  Write
\[
    q_F^{\mathrm{nat}}:
    X\longrightarrow Q_F^{\mathrm{nat}}
\]
for the induced centre.
\end{definition}

\begin{proposition}[The global native centre is effective]
\label{prop:bnli-native-centre-effective}
The morphism $q_F^{\mathrm{nat}}$ is an effective epimorphism.
\end{proposition}

\begin{proof}
Factor it as
\[
    X\twoheadrightarrow J_F\lhook\joinrel\longrightarrow Q_F^{\mathrm{nat}}.
\]
For every element $e=(c,\xi)$, the composite
\[
    X\twoheadrightarrow\overline I_c\longrightarrow Q_F^{\mathrm{nat}}
\]
factors through $J_F$.  Since $X\to\overline I_c$ is effective, the pullback
of $J_F$ to $\overline I_c$ is all of $\overline I_c$; equivalently the cocone
map $\overline I_c\to Q_F^{\mathrm{nat}}$ factors through $J_F$.  Hence the
entire colimit cocone factors through $J_F$, so the identity of
$Q_F^{\mathrm{nat}}$ factors through the monomorphism $J_F$.  Therefore that
monomorphism is an equivalence.
\end{proof}

\begin{theorem}[The global native quotient is formal]
\label{thm:bnli-native-quotient-formal}
There is a canonical equivalence
\[
\boxed{
    L_S^{\mathrm{dR}}X
    \xrightarrow{\sim}
    L_S^{\mathrm{dR}}Q_F^{\mathrm{nat}}.}
\]
Consequently
\[
\boxed{
    \mathfrak G_F^{\mathrm{nat}}
    :=
    \check C_\bullet(q_F^{\mathrm{nat}})}
\]
is a brave-new formal Lie groupoid on the fixed object space $X$.
\end{theorem}

\begin{proof}
Apply relative de Rham localization to the colimit defining
$Q_F^{\mathrm{nat}}$.  By
Proposition~\ref{prop:bnli-induced-cell-formal-effective}, every centred object
\[
    L_S^{\mathrm{dR}}X
    \longrightarrow
    L_S^{\mathrm{dR}}\overline I_c
\]
is equivalent to the initial object
$\operatorname{id}_{L_S^{\mathrm{dR}}X}$ of the localized undercategory.  The
full subcategory of objects equivalent to an initial object is contractible,
so the localized diagram is canonically equivalent to the constant initial
diagram.  Its colimit is initial.  Hence
\[
    L_S^{\mathrm{dR}}Q_F^{\mathrm{nat}}
    \simeq
    L_S^{\mathrm{dR}}X.
\]
Together with Proposition~\ref{prop:bnli-native-centre-effective}, the formal
quotient recognition theorem gives the groupoid statement.
\end{proof}

\begin{definition}[Raw integrated relation]
\label{def:bnli-raw-integrated-relation}
Define
\[
\boxed{
    \mathcal R^{\mathrm{nat}}_{F,\bullet}
    :=
    \check C_\bullet(q_F^{\mathrm{nat}}).}
\]
Thus
\[
    \mathcal R^{\mathrm{nat}}_{F,n}
    \simeq
    \underbrace{
      X\times_{Q_F^{\mathrm{nat}}}\cdots
      \times_{Q_F^{\mathrm{nat}}}X
    }_{n+1}.
\]
\end{definition}

\subsection{Global comparison with formal-moduli-stack materialization}

Every global Artinian element $e=(c,\xi)$ supplies an evaluation map
$T_c\to\operatorname{Mat}^{\mathrm{sp}}_{X/S}(F)$ extending the centre on
$U_c$.  Together with the global centre $X\to\operatorname{Mat}^{\mathrm{sp}}(F)$
this gives, by the pushout universal property, a map
\[
    \overline T_c\longrightarrow\operatorname{Mat}^{\mathrm{sp}}(F),
\]
and hence a compatible map from $\overline I_c$.

\begin{definition}[Global native-to-materialized quotient comparison]
\label{def:bnli-native-materialized-comparison}
Define
\[
\boxed{
    \kappa_F^{\mathrm{mat}}:
    Q_F^{\mathrm{nat}}
    \longrightarrow
    \operatorname{Mat}^{\mathrm{sp}}_{X/S}(F)}
\]
by the global native-cell colimit universal property.
\end{definition}

\begin{theorem}[Deformation formalization of the global native quotient]
\label{thm:bnli-native-quotient-materialization}
The comparison $\kappa_F^{\mathrm{mat}}$ is an
$L_S^{\mathrm{Def}}$-equivalence.  More precisely,
\[
\boxed{
    L_S^{\mathrm{Def}}Q_F^{\mathrm{nat}}
    \xrightarrow{\sim}
    \operatorname{Mat}^{\mathrm{sp}}_{X/S}(F).}
\]
\end{theorem}

\begin{proof}
Proposition~\ref{prop:bnli-induced-cell-formal-effective} gives a termwise
deformation equivalence
\[
    \overline I_c\longrightarrow\overline T_c.
\]
Since $L_S^{\mathrm{Def}}$ preserves colimits,
\[
    L_S^{\mathrm{Def}}Q_F^{\mathrm{nat}}
    \simeq
    L_S^{\mathrm{Def}}Q_F^{\mathrm{Art,rep}}.
\]
We identify the latter object by its mapping universal property in the
reflective subcategory of deformation-local centred targets.

Let $X\to Y$ be $L_S^{\mathrm{Def}}$-local.  Mapping out of the global
represented Artinian realization gives
\[
\begin{aligned}
    \operatorname{Map}_{X/}(Q_F^{\mathrm{Art,rep}},Y)
    &\simeq
    \operatorname*{lim}_{e=(c,\xi)\in\int_XF}
    \operatorname{Map}_{X/}(\overline T_c,Y)\\
    &\simeq
    \operatorname*{lim}_{e=(c,\xi)\in\int_XF}
    \operatorname{Map}_{U_c/}(T_c,Y).
\end{aligned}
\]
The second equivalence is the pushout universal property.  By definition,
$\operatorname{Map}_{U_c/}(T_c,Y)$ is the value at $c$ of the global Artinian
leaf $\operatorname{Leaf}^{\mathrm{Art}}_{X/S}(Y)$.  Therefore
Lemma~\ref{lem:bnli-global-artinian-coyoneda} identifies the last limit with
\[
    \operatorname{Map}_{\operatorname{FMP}^{\mathrm{sp}}_{X/S}}
    \left(
      F,\operatorname{Leaf}^{\mathrm{Art}}_{X/S}(Y)
    \right).
\]
The global materialization adjunction identifies this with
\[
    \operatorname{Map}_{X/}
    \left(
      \operatorname{Mat}^{\mathrm{sp}}_{X/S}(F),Y
    \right).
\]
Thus $L_S^{\mathrm{Def}}Q_F^{\mathrm{Art,rep}}$ and
$\operatorname{Mat}^{\mathrm{sp}}_{X/S}(F)$ corepresent the same mapping functor
on the deformation-local subcategory.  Yoneda in that reflective subcategory
gives the displayed equivalence, and it is induced by the termwise Artinian
evaluation cocone.  Hence it is precisely the localization of
$\kappa_F^{\mathrm{mat}}$.
\end{proof}

\subsection{Global brave-new groupoid integration}

\begin{definition}[Brave-new groupoid integration]
\label{def:bnli-integration}
For
\[
    \mathfrak a\in\operatorname{LieAlgd}^{\pi,E_\infty}_{X/S}
\]
put
\[
    F_{\mathfrak a}
    :=
    \operatorname{Int}^{\mathrm{fm}}_{X/S}(\mathfrak a)
\]
and define
\[
\boxed{
    \operatorname{Int}^{\mathrm{gpd}}_{X/S}(\mathfrak a)
    :=
    \mathfrak G_{F_{\mathfrak a}}^{\mathrm{nat}}
    =
    \check C_\bullet
    \left(
      X\twoheadrightarrow Q_{F_{\mathfrak a}}^{\mathrm{nat}}
    \right).}
\]
No final unitwise completion is applied: formality of the quotient Cech nerve
is already Theorem~\ref{thm:bnli-native-quotient-formal}.
\end{definition}

\begin{theorem}[Groupoid Integration Functor]
\label{thm:bnli-integration-functor}
For fixed $X/S$, native groupoid integration is a functor
\[
\boxed{
    \operatorname{Int}^{\mathrm{gpd}}_{X/S}:
    \operatorname{LieAlgd}^{\pi,E_\infty}_{X/S}
    \longrightarrow
    \operatorname{FormLieGpd}(X/S).}
\]
No raw changing-object etale exchange is asserted.
\end{theorem}

\begin{proof}
Formal integration is functorial.  A natural transformation $F\to G$ induces
a functor of global Artinian categories of elements
\[
    \int_XF\longrightarrow\int_XG
\]
and compatible maps between the associated local effective-image cells.
Pushout along the fixed centre $X$, followed by colimit functoriality, gives a
map
\[
    Q_F^{\mathrm{nat}}\longrightarrow Q_G^{\mathrm{nat}}
\]
under $X$.  Taking Cech nerves produces the required formal-groupoid morphism.
Identity, composition, and all higher coherence are inherited from the
Grothendieck construction, effective-image functoriality, pushout, and colimit.
\end{proof}

\begin{theorem}[Global Formal-Leaf Recovery Theorem]
\label{thm:bnli-leaf-recovery}
For every
$\mathfrak a\in\operatorname{LieAlgd}^{\pi,E_\infty}_{X/S}$ there is a
canonical equivalence
\[
\boxed{
    \mathcal L^{\mathrm{fm,sp}}_{
      \operatorname{Int}^{\mathrm{gpd}}(\mathfrak a)}
    \xrightarrow{\sim}
    \operatorname{Int}^{\mathrm{fm}}(\mathfrak a).}
\]
\end{theorem}

\begin{proof}
The effective quotient of the integrated groupoid is
$Q_{F_{\mathfrak a}}^{\mathrm{nat}}$.  The differentiation chapter constructs
its global formal leaf by applying the spectral deformation reflector to the
quotient and then taking the reduced-centred model.  Theorem~\ref{thm:bnli-native-quotient-materialization} identifies the deformation-local
target with
$\operatorname{Mat}^{\mathrm{sp}}_{X/S}(F_{\mathfrak a})$.  This is already a
nil-centred spectral formal-moduli stack, so reduced-centred completion does
not change its Artinian leaf.  The global materialization theorem then gives
\[
    \operatorname{Leaf}^{\mathrm{Art}}
    \operatorname{Mat}^{\mathrm{sp}}(F_{\mathfrak a})
    \simeq F_{\mathfrak a}.
\]
These equivalences are natural on the coherent affine-etale coefficient
system and therefore identify the global formal leaf with
$F_{\mathfrak a}$.
\end{proof}

\begin{corollary}[Differentiation after groupoid integration]
\label{cor:bnli-diff-int-identity}
There is a canonical equivalence
\[
\boxed{
    \operatorname{Diff}^{\pi,E_\infty}_{X/S}
    \operatorname{Int}^{\mathrm{gpd}}_{X/S}
    \xrightarrow{\sim}
    \operatorname{id}_{
      \operatorname{LieAlgd}^{\pi,E_\infty}_{X/S}}.}
\]
\end{corollary}

\begin{proof}
Apply global spectral differentiation to
Theorem~\ref{thm:bnli-leaf-recovery} and use the formal integration
equivalence of Theorem~\ref{thm:bnli-formal-lie-three}.
\end{proof}

\section{Global native mapping discrepancy and the remaining rigidity problem}
\label{sec:bnli-native-density}

The identity
\[
    \operatorname{Diff}
    \operatorname{Int}^{\mathrm{gpd}}
    \simeq
    \operatorname{id}
\]
does not imply raw full faithfulness.  The global native quotient contains
nonlinear fixed-object mapping information which may disappear after
deformation localization.  We retain the exact cellwise discrepancy.

Let
\[
    c=((U,V),A\to B)
    \in\mathscr{Art}^{\mathrm{sp}}_{X/S}
\]
and write $\overline I_c=X\amalg_UI_c$.

\begin{definition}[Global native cell evaluation]
\label{def:bnli-global-native-evaluation}
For
$G\in\operatorname{FMP}^{\mathrm{sp}}_{X/S}$ define
\[
\boxed{
    \operatorname{ev}^{\mathrm{nat}}_{c,G}:
    \operatorname{Map}_{X/}(\overline I_c,Q_G^{\mathrm{nat}})
    \longrightarrow
    G_c(A).}
\]
It is the composite obtained by postcomposing with
$Q_G^{\mathrm{nat}}\to L_S^{\mathrm{Def}}Q_G^{\mathrm{nat}}\simeq
\operatorname{Mat}^{\mathrm{sp}}(G)$, using the deformation equivalence
$\overline I_c\to\overline T_c$, and then applying the Artinian-leaf
identification
\[
    \operatorname{Map}_{X/}
    (\overline T_c,\operatorname{Mat}^{\mathrm{sp}}(G))
    \simeq G_c(A).
\]
\end{definition}

\begin{proposition}[Canonical section of global native evaluation]
\label{prop:bnli-native-evaluation-section}
For every global Artinian test $c$ the map
$\operatorname{ev}^{\mathrm{nat}}_{c,G}$ admits a canonical section
\[
\boxed{
    s_{c,G}:
    G_c(A)
    \longrightarrow
    \operatorname{Map}_{X/}(\overline I_c,Q_G^{\mathrm{nat}}).}
\]
\end{proposition}

\begin{proof}
A point $x\in G_c(A)$ is exactly a global Artinian element $(c,x)$ of
$\int_XG$.  The defining colimit cocone for $Q_G^{\mathrm{nat}}$ therefore
supplies a map
\[
    \overline I_c\longrightarrow Q_G^{\mathrm{nat}}.
\]
After deformation localization this cocone map is the Artinian evaluation map
classified by $x$.  Hence the composite with
$\operatorname{ev}^{\mathrm{nat}}_{c,G}$ is the identity.
\end{proof}

\begin{definition}[Global native cell mapping discrepancy]
\label{def:bnli-native-cell-discrepancy}
For $x\in G_c(A)$ define
\[
\boxed{
    \mathscr E^{\mathrm{nat}}_{c,G}(x)
    :=
    \operatorname{fib}_{x}
    \left(
      \operatorname{ev}^{\mathrm{nat}}_{c,G}
    \right).}
\]
This is a functorial pointed space.  No vanishing condition is imposed as
input structure.
\end{definition}

\begin{theorem}[Global mapping discrepancy formula]
\label{thm:bnli-mapping-discrepancy}
Let $F,G\in\operatorname{FMP}^{\mathrm{sp}}_{X/S}$ and let
$\alpha:F\to G$.  Differentiation induces a split map
\[
\boxed{
    \operatorname{Map}_{\operatorname{FormLieGpd}(X/S)}
    (\mathfrak G_F^{\mathrm{nat}},\mathfrak G_G^{\mathrm{nat}})
    \longrightarrow
    \operatorname{Map}_{\operatorname{FMP}^{\mathrm{sp}}_{X/S}}(F,G),}
\]
whose section is native integration.  The fibre over $\alpha$ is canonically
\[
\boxed{
    \operatorname{fib}_{\alpha}
    \simeq
    \operatorname*{lim}_{e=(c,\xi)\in\int_XF}
    \mathscr E^{\mathrm{nat}}_{c,G}
    \bigl(\alpha_c(\xi)\bigr).}
\]
In particular, vanishing of all global native cell discrepancies implies raw
full faithfulness of native integration.
\end{theorem}

\begin{proof}
The fixed-object groupoid/effective-quotient equivalence gives
\[
\begin{aligned}
    \operatorname{Map}
    (\mathfrak G_F^{\mathrm{nat}},\mathfrak G_G^{\mathrm{nat}})
    &\simeq
    \operatorname{Map}_{X/}(Q_F^{\mathrm{nat}},Q_G^{\mathrm{nat}})\\
    &\simeq
    \operatorname*{lim}_{e=(c,\xi)\in\int_XF}
    \operatorname{Map}_{X/}(\overline I_c,Q_G^{\mathrm{nat}}).
\end{aligned}
\]
Lemma~\ref{lem:bnli-global-artinian-coyoneda} gives
\[
    \operatorname{Map}(F,G)
    \simeq
    \operatorname*{lim}_{e=(c,\xi)\in\int_XF}G_c(A_c).
\]

We identify the comparison map term by term.  Let
\[
    j_e:\overline I_c\longrightarrow Q_F^{\mathrm{nat}}
\]
be the cocone map corresponding to
$e=(c,\xi)\in\int_XF$, and let
\[
    g:Q_F^{\mathrm{nat}}\longrightarrow Q_G^{\mathrm{nat}}
\]
be a raw quotient morphism under $X$.  Functoriality of the global formal-leaf
construction sends $g$ to a morphism
\[
    \operatorname{Leaf}^{\mathrm{Art}}
    (Q_F^{\mathrm{nat}})
    \longrightarrow
    \operatorname{Leaf}^{\mathrm{Art}}
    (Q_G^{\mathrm{nat}}).
\]
Under formal-leaf recovery these are $F$ and $G$.

Evaluate this formal-leaf morphism on the Artinian element $e$.  By
naturality of
\[
    \kappa^{\mathrm{mat}}:
    Q^{\mathrm{nat}}
    \longrightarrow
    \operatorname{Mat}^{\mathrm{sp}},
\]
the resulting point of $G_c(A_c)$ is obtained by applying deformation
localization to
\[
    \overline I_c
    \xrightarrow{j_e}
    Q_F^{\mathrm{nat}}
    \xrightarrow{g}
    Q_G^{\mathrm{nat}},
\]
using the canonical equivalences
\[
    L_S^{\mathrm{Def}}\overline I_c
    \simeq
    \overline T_c,
    \qquad
    L_S^{\mathrm{Def}}Q_G^{\mathrm{nat}}
    \simeq
    \operatorname{Mat}^{\mathrm{sp}}(G),
\]
and then applying Artinian restriction.  This is exactly
\[
    \operatorname{ev}^{\mathrm{nat}}_{c,G}
    \left(
      g\circ j_e
    \right)
\]
by Definition~\ref{def:bnli-global-native-evaluation}.  Therefore, under the
two displayed limit descriptions, the differentiation map is literally
\[
\boxed{
    \operatorname*{lim}_{e=(c,\xi)\in\int_XF}
    \operatorname{ev}^{\mathrm{nat}}_{c,G}.}
\]

Native integration sends $\alpha:F\to G$ to the compatible family of native
cell maps supplied by the colimit cocones.  Hence the section of the
differentiation map is the limit of the canonical sections
$s_{c,G}$ from
Proposition~\ref{prop:bnli-native-evaluation-section}.

Finally, homotopy fibres commute with limits of spaces.  Taking the fibre of
the displayed limit map over the compatible family
$\{\alpha_c(\xi)\}_{e\in\int_XF}$ gives
\[
\begin{aligned}
    \operatorname{fib}_{\alpha}
    &\simeq
    \operatorname*{lim}_{e=(c,\xi)\in\int_XF}
    \operatorname{fib}_{\alpha_c(\xi)}
    \left(
      \operatorname{ev}^{\mathrm{nat}}_{c,G}
    \right)\\
    &=
    \operatorname*{lim}_{e=(c,\xi)\in\int_XF}
    \mathscr E^{\mathrm{nat}}_{c,G}
    \bigl(\alpha_c(\xi)\bigr),
\end{aligned}
\]
which proves the formula.
\end{proof}

\begin{remark}[The remaining global native rigidity problem]
The chapter does not prove that every
$\operatorname{ev}^{\mathrm{nat}}_{c,G}$ is an equivalence.  The exact fibre
formula records the obstruction without defining a ``native-rigid'',
``materializable'', ``Artinian-dense'', or ``admissible'' class.  The stronger
affine if-and-only-if criterion remains available in
Theorem~\ref{thm:bnli-affine-mapping-discrepancy}.
\end{remark}

\section{Native Brave-New Lie III and reconstructed objects}
\label{sec:bnli-def-complete}

The preceding constructions produce native integration on every spectral
partition-Lie algebroid and recover the controller after differentiation.  The
result is a one-sided native Lie III theorem.  A categorical equivalence with a
full raw subgroupoid category would additionally require the native rigidity
theorem isolated in Section~\ref{sec:bnli-native-density}.

\begin{theorem}[Native Brave-New Lie III]
\label{thm:bnli-lie-three}
For fixed $X/S$ there is a canonical functor
\[
\boxed{
    \operatorname{Int}^{\mathrm{gpd}}_{X/S}:
    \operatorname{LieAlgd}^{\pi,E_\infty}_{X/S}
    \longrightarrow
    \operatorname{FormLieGpd}(X/S)}
\]
and a canonical natural equivalence
\[
\boxed{
    \operatorname{Diff}^{\pi,E_\infty}_{X/S}
    \operatorname{Int}^{\mathrm{gpd}}_{X/S}
    \xrightarrow{\sim}
    \operatorname{id}_{\operatorname{LieAlgd}^{\pi,E_\infty}_{X/S}}.}
\]
For every $\mathfrak a$, applying integration to this recovery equivalence
gives a canonical equivalence between the two reconstructed objects
\[
\boxed{
    \operatorname{Int}^{\mathrm{gpd}}
    \operatorname{Diff}^{\pi,E_\infty}
    \operatorname{Int}^{\mathrm{gpd}}(\mathfrak a)
    \xrightarrow{\sim}
    \operatorname{Int}^{\mathrm{gpd}}(\mathfrak a).}
\]
No full faithfulness of $\operatorname{Int}^{\mathrm{gpd}}$ in the unrestricted
raw groupoid mapping spaces, no adjunction on the entire native groupoid
category, and no raw changing-object etale exchange or atlas-invariance theorem
is asserted.
\end{theorem}

\begin{proof}
The functor is Theorem~\ref{thm:bnli-integration-functor} and the recovery
equivalence is Corollary~\ref{cor:bnli-diff-int-identity}.  Applying the functor
$\operatorname{Int}^{\mathrm{gpd}}$ to that natural equivalence gives the
reconstructed-object equivalence.  The final sentence records exactly the
mapping-space scope established by Theorem~\ref{thm:bnli-mapping-discrepancy}.
\end{proof}

\begin{definition}[Reconstructed native formal groupoid]
\label{def:bnli-reconstructed-groupoid}
A \emph{reconstructed native formal groupoid} is an object presented as
\[
    \operatorname{Int}^{\mathrm{gpd}}(\mathfrak a)
\]
for a spectral partition-Lie algebroid $\mathfrak a$.  This terminology names
an output of the functor; it does not define a new input structure or assert
that the full essential image is categorically equivalent to the algebroid
category in raw mapping spaces.
\end{definition}

\section{The Lie comparison of an arbitrary formal groupoid}
\label{sec:bnli-completion}

Let
\[
    \mathfrak G\in\operatorname{FormLieGpd}(X/S),
    \qquad
    \mathfrak a_{\mathfrak G}
    :=
    \operatorname{Diff}^{\pi,E_\infty}(\mathfrak G).
\]

\begin{definition}[Controller integration of a formal groupoid]
\label{def:bnli-lie-completion}
Define
\[
\boxed{
    \mathfrak G^{\mathrm{Lie}}
    :=
    \operatorname{Int}^{\mathrm{gpd}}
    (\mathfrak a_{\mathfrak G}).}
\]
This object is defined for every formal groupoid.
\end{definition}

The differentiation chapter constructs the globally deformation-local target
\[
    q_{\mathfrak G}^{\mathrm{fm}}:
    X\longrightarrow Q_{\mathfrak G}^{\mathrm{fm}}
\]
and the Atiyah--Spencer chapter constructs the formalized relation
\[
    R_{\mathfrak G,\bullet}^{\mathrm{fm}}
    :=
    \check C_\bullet(q_{\mathfrak G}^{\mathrm{fm}})
\]
together with a simplicial map
\[
    \lambda^R_{\mathfrak G}:
    \mathfrak G_\bullet
    \longrightarrow
    R_{\mathfrak G,\bullet}^{\mathrm{fm}}.
\]

The materialization adjunction gives an evaluation counit
\[
    \epsilon_{\mathfrak G}^{\mathrm{mat}}:
    \operatorname{Mat}^{\mathrm{sp}}
    \left(
      \mathcal L_{\mathfrak G}^{\mathrm{fm,sp}}
    \right)
    \longrightarrow
    Q_{\mathfrak G}^{\mathrm{fm}}.
\]
The native quotient of $\mathfrak G^{\mathrm{Lie}}$ maps first to the
materialized target and then through this counit to
$Q_{\mathfrak G}^{\mathrm{fm}}$, hence induces a second simplicial relation map
\[
    \rho^R_{\mathfrak G}:
    \mathfrak G^{\mathrm{Lie}}_\bullet
    \longrightarrow
    R_{\mathfrak G,\bullet}^{\mathrm{fm}}.
\]

\begin{definition}[Formal Lie comparison object]
\label{def:bnli-lie-comparison-object}
Let
\[
\boxed{
    \mathfrak K_{\mathfrak G}^{\mathrm{Lie}}
    :=
    \widehat{
      R_{\mathfrak G,\bullet}^{\mathrm{fm}}
    }^{\,\mathrm{dR}}}
\]
be the unitwise formalization of the formalized relation.
\end{definition}

\begin{theorem}[Canonical Lie-comparison span]
\label{thm:bnli-completion-map}
For every native formal Lie groupoid there is a canonical functorial span in
$\operatorname{FormLieGpd}(X/S)$
\[
\boxed{
    \mathfrak G
    \xrightarrow{\ \widehat\lambda^{\mathrm{Lie}}_{\mathfrak G}\ }
    \mathfrak K_{\mathfrak G}^{\mathrm{Lie}}
    \xleftarrow{\ \widehat\rho^{\mathrm{Lie}}_{\mathfrak G}\ }
    \mathfrak G^{\mathrm{Lie}}.}
\]
It is natural in fixed-object formal-groupoid morphisms.  No raw
changing-object etale variance is asserted for this span.

If $\mathfrak G$ is itself a reconstructed object
$\operatorname{Int}^{\mathrm{gpd}}(\mathfrak a)$, the reconstructed-object
equivalence of Theorem~\ref{thm:bnli-lie-three} canonically identifies the two
outer sources.  For an arbitrary $\mathfrak G$, no direct morphism
$\mathfrak G\to\mathfrak G^{\mathrm{Lie}}$ is inferred from the span.
\end{theorem}

\begin{proof}
Both $\lambda^R_{\mathfrak G}$ and $\rho^R_{\mathfrak G}$ are maps from formal
fixed-object groupoids to the not-necessarily-formal groupoid
$R_{\mathfrak G}^{\mathrm{fm}}$.  Unitwise formalization is right adjoint to the
inclusion of formal fixed-object groupoids.  Each map therefore has a
contractibly unique lift through the formalization counit
\[
    \mathfrak K_{\mathfrak G}^{\mathrm{Lie}}
    \longrightarrow
    R_{\mathfrak G}^{\mathrm{fm}}.
\]
These are the two displayed arrows.  All constructions are functorial and the
lifts are unique by the coreflection mapping-space equivalence.  For a reconstructed source the direct comparison is supplied separately by
applying integration to the differentiation-recovery equivalence of
Theorem~\ref{thm:bnli-lie-three}.
\end{proof}

\begin{proposition}[Compatibility of the comparison span on reconstructed objects]
\label{prop:bnli-reconstructed-span-compatibility}
Let
\[
    \mathfrak G
    =
    \operatorname{Int}^{\mathrm{gpd}}(\mathfrak a)
\]
be reconstructed, and let
\[
    \varepsilon_{\mathfrak G}^{\mathrm{rec}}:
    \mathfrak G^{\mathrm{Lie}}
    \xrightarrow{\sim}
    \mathfrak G
\]
be the equivalence obtained by applying native integration to
$\operatorname{Diff}\operatorname{Int}\simeq\operatorname{id}$.  Then the two
arrows of the comparison span agree after this identification:
\[
\boxed{
    \widehat\rho^{\mathrm{Lie}}_{\mathfrak G}
    \simeq
    \widehat\lambda^{\mathrm{Lie}}_{\mathfrak G}
    \circ
    \varepsilon_{\mathfrak G}^{\mathrm{rec}}.}
\]
Equivalently, the corresponding quotient and simplicial relation maps agree
before unitwise formalization after identifying the two reconstructed outer
sources.
\end{proposition}

\begin{proof}
Put
\[
    F:=\operatorname{Int}^{\mathrm{fm}}(\mathfrak a).
\]
The effective quotient of $\mathfrak G$ is $Q_F^{\mathrm{nat}}$ and
Theorem~\ref{thm:bnli-native-quotient-materialization} gives
\[
    L_S^{\mathrm{Def}}Q_F^{\mathrm{nat}}
    \xrightarrow{\sim}
    \operatorname{Mat}^{\mathrm{sp}}(F),
\]
with the equivalence induced by
$\kappa_F^{\mathrm{mat}}$.

The formalized quotient of $\mathfrak G$ is obtained by applying
$L_S^{\mathrm{Def}}$ to $Q_F^{\mathrm{nat}}$ and then taking the
reduced-centred model.  Since $\operatorname{Mat}^{\mathrm{sp}}(F)$ is already
nil-centred, the projection from that reduced-centred model is an
equivalence.  Under the resulting canonical identification
\[
\boxed{
    Q_{\mathfrak G}^{\mathrm{fm}}
    \xrightarrow{\sim}
    \operatorname{Mat}^{\mathrm{sp}}(F),}
\]
the quotient map underlying
$\lambda^R_{\mathfrak G}$ is exactly
$\kappa_F^{\mathrm{mat}}$: both correspond, under the localization mapping
property into the same $L_S^{\mathrm{Def}}$-local target, to the equivalence
displayed above.

Next,
\[
    \mathcal L_{\mathfrak G}^{\mathrm{fm,sp}}
    \simeq F
\]
by formal-leaf recovery.  Hence the materialization counit occurring in the
definition of $\rho^R_{\mathfrak G}$ is
\[
    \operatorname{Mat}^{\mathrm{sp}}(F)
    \longrightarrow
    Q_{\mathfrak G}^{\mathrm{fm}}.
\]
Both source and target belong to the fully faithful materialization essential
image and have Artinian leaf $F$, so this counit is the inverse of the
preceding canonical equivalence.

Finally, the reconstructed equivalence
$\varepsilon_{\mathfrak G}^{\mathrm{rec}}$ is obtained by functoriality of the
same native quotient construction under the differentiation-recovery
equivalence.  Naturality of
$\kappa^{\mathrm{mat}}$ therefore identifies the native-to-materialized
quotient map for $\mathfrak G^{\mathrm{Lie}}$ with that for $\mathfrak G$
after $\varepsilon_{\mathfrak G}^{\mathrm{rec}}$.  The two quotient maps to
$Q_{\mathfrak G}^{\mathrm{fm}}$ consequently agree.  Taking Cech nerves gives
the equality of relation maps, and the uniqueness in the unitwise
formalization coreflection gives the displayed equality of formal comparison
arrows.
\end{proof}


\begin{remark}[Why the span is necessary]
The materialization counit points
\[
    \operatorname{Mat}
    \operatorname{Leaf}(Q_{\mathfrak G}^{\mathrm{fm}})
    \longrightarrow
    Q_{\mathfrak G}^{\mathrm{fm}},
\]
not in the reverse direction.  The span records this variance exactly.  A
reverse map is not manufactured from an adjunction counit.
\end{remark}

\section{The Lie-comparison discrepancy system}
\label{sec:bnli-discrepancy}

For an arbitrary native formal Lie groupoid there need not be a direct map
\(\mathfrak G\to\mathfrak G^{\mathrm{Lie}}\).  The canonical object available
without inserting a reverse materialization arrow is the comparison span of
Theorem~\ref{thm:bnli-completion-map}.  We therefore define discrepancies for
each leg of that span.

Write
\[
    Q_{\mathfrak G}
    \longrightarrow
    Q_{\mathfrak K_{\mathfrak G}^{\mathrm{Lie}}}
    \longleftarrow
    Q_{\mathfrak G^{\mathrm{Lie}}}
\]
for the induced span of effective quotients.

\begin{definition}[Quotient comparison span]
\label{def:bnli-quotient-discrepancy}
Define
\[
\boxed{
    \Delta^Q_{\mathfrak G}
    :=
    \left(
      Q_{\mathfrak G}
      \longrightarrow
      Q_{\mathfrak K_{\mathfrak G}^{\mathrm{Lie}}}
      \longleftarrow
      Q_{\mathfrak G^{\mathrm{Lie}}}
    \right).}
\]
This is a diagram, not a single fibre or cofiber: the quotient objects live in
an unstable topos and no additive discrepancy is manufactured there.
\end{definition}

\begin{definition}[Relation comparison span]
\label{def:bnli-relation-discrepancy}
The relation comparison is the simplicial span
\[
\boxed{
    \Delta^R_{\mathfrak G,\bullet}
    :=
    \left(
      \mathfrak G_\bullet
      \longrightarrow
      \mathfrak K_{\mathfrak G,\bullet}^{\mathrm{Lie}}
      \longleftarrow
      \mathfrak G^{\mathrm{Lie}}_\bullet
    \right).}
\]
Its degree-one terms induce the stable and isotropy comparisons below.
\end{definition}

\subsection{Stable arrow discrepancies}

Intrinsic reduced stabilization is functorial on pointed arrow objects.  Thus
the two legs give
\[
    \mathbb A^{\mathrm{st}}_{\mathfrak G}
    \longrightarrow
    \mathbb A^{\mathrm{st}}_{\mathfrak K_{\mathfrak G}^{\mathrm{Lie}}}
    \longleftarrow
    \mathbb A^{\mathrm{st}}_{\mathfrak G^{\mathrm{Lie}}}.
\]

\begin{definition}[Stable Lie-comparison discrepancies]
\label{def:bnli-stable-discrepancy}
Define the two stable cofibres
\[
\boxed{
\begin{aligned}
    \mathbb C^{\mathrm{Lie,st},L}_{\mathfrak G}
    &:={\rm cofib}\!\left(
      \mathbb A^{\mathrm{st}}_{\mathfrak G}
      \to
      \mathbb A^{\mathrm{st}}_{\mathfrak K_{\mathfrak G}^{\mathrm{Lie}}}
    \right),\\
    \mathbb C^{\mathrm{Lie,st},R}_{\mathfrak G}
    &:={\rm cofib}\!\left(
      \mathbb A^{\mathrm{st}}_{\mathfrak G^{\mathrm{Lie}}}
      \to
      \mathbb A^{\mathrm{st}}_{\mathfrak K_{\mathfrak G}^{\mathrm{Lie}}}
    \right).
\end{aligned}}
\]
No comparison between these two cofibres is imposed by definition.
\end{definition}

\subsection{Actual-isotropy and adjoint discrepancies}

Functoriality of formal isotropy supplies a span of formal group objects
\[
    I^{\mathrm{for}}_{\mathfrak G}
    \longrightarrow
    I^{\mathrm{for}}_{\mathfrak K_{\mathfrak G}^{\mathrm{Lie}}}
    \longleftarrow
    I^{\mathrm{for}}_{\mathfrak G^{\mathrm{Lie}}}.
\]

\begin{definition}[Actual-isotropy Lie-comparison discrepancies]
\label{def:bnli-isotropy-discrepancy}
In the parameterized stable category over \(X\), define
\[
\boxed{
\begin{aligned}
    \mathbb C^{\mathrm{Lie,iso},L}_{\mathfrak G}
    &:={\rm cofib}\!\left(
      \widetilde\Sigma_X^\infty I^{\mathrm{for}}_{\mathfrak G}
      \to
      \widetilde\Sigma_X^\infty
      I^{\mathrm{for}}_{\mathfrak K_{\mathfrak G}^{\mathrm{Lie}}}
    \right),\\
    \mathbb C^{\mathrm{Lie,iso},R}_{\mathfrak G}
    &:={\rm cofib}\!\left(
      \widetilde\Sigma_X^\infty I^{\mathrm{for}}_{\mathfrak G^{\mathrm{Lie}}}
      \to
      \widetilde\Sigma_X^\infty
      I^{\mathrm{for}}_{\mathfrak K_{\mathfrak G}^{\mathrm{Lie}}}
    \right).
\end{aligned}}
\]
\end{definition}

The descended formal adjoint groups give the analogous two-legged comparison
after passage to the common formalized quotient.  We retain those morphisms as
the \emph{adjoint Lie-comparison span}.  No equivalence with actual isotropy is
inserted.

\begin{definition}[Support-marked Spencer comparison spans]
\label{def:bnli-support-marked-discrepancy}
Apply the functorial raw and formalized support-marked Cech constructions of
the Atiyah--Spencer chapter separately to the two arrows
\[
    \mathfrak G\longrightarrow\mathfrak K_{\mathfrak G}^{\mathrm{Lie}},
    \qquad
    \mathfrak G^{\mathrm{Lie}}\longrightarrow
    \mathfrak K_{\mathfrak G}^{\mathrm{Lie}}.
\]
For every finite support set $I$, retain the resulting diagrams and their
stable cofibres as
\[
    \Delta^{\mathrm{Sp,mk},L}_{\mathfrak G}(I),
    \qquad
    \Delta^{\mathrm{Sp,mk},R}_{\mathfrak G}(I).
\]
These are comparison systems in the raw/formalized support-marked branch of
the preceding chapter.  They are not identified with the independent
directional Artinian cotangent-primitive operation system.
\end{definition}

\begin{theorem}[Functoriality of the discrepancy system]
\label{thm:bnli-discrepancy-compatibility}
The quotient, relation, stable-arrow, actual-isotropy, adjoint, and
support-marked Spencer comparison systems above are functorial shadows of the
single canonical span of Theorem~\ref{thm:bnli-completion-map}.

If \(\mathfrak G=\operatorname{Int}^{\mathrm{gpd}}(\mathfrak a)\) is a reconstructed object, the two outer sources are canonically equivalent by Theorem~\ref{thm:bnli-lie-three}.  Under that equivalence the left- and right-leg comparison systems are canonically identified.  The theorem does not assert that either leg into \(\mathfrak K_{\mathfrak G}^{\mathrm{Lie}}\) is an equivalence, so the corresponding discrepancy cofibres need not vanish.
\end{theorem}

\begin{proof}
Every listed construction is functorial for morphisms of fixed-object formal
groupoids.  Apply it separately to the two arrows of
Theorem~\ref{thm:bnli-completion-map}; stable cofibres are then functorial in
maps of arrows.  For a reconstructed source,
Proposition~\ref{prop:bnli-reconstructed-span-compatibility} identifies the
two arrows of the comparison span after the reconstructed-object equivalence.
Applying any of the listed functorial constructions therefore identifies the
left- and right-leg shadows and their stable cofibres.  No equivalence between
either shadow and the middle formalized relation is inferred.
\end{proof}

\section{Operation-theoretic invariance of the integrated controller}
\label{sec:bnli-operation-characterization}

The preceding chapter constructs for every native formal groupoid a canonical
integral equivalence
\[
    \Xi^{\mathrm{Art}}_{\mathfrak G}:
    \mathbb A^{\pi,E_\infty}_{\mathfrak G}
    \xrightarrow{\sim}
    \mathbb A^{\mathrm{Sp,Art}}_{\mathfrak G}.
\]
The native integration theorem supplies a second, independent equivalence on
the controller level.

\begin{theorem}[Artinian Spencer invariance under reconstructed integration]
\label{thm:bnli-spencer-invariance}
For every native formal Lie groupoid \(\mathfrak G\), there is a canonical
equivalence
\[
\boxed{
    \mathbb A^{\mathrm{Sp,Art}}_{\mathfrak G}
    \xrightarrow{\sim}
    \mathbb A^{\mathrm{Sp,Art}}_{\mathfrak G^{\mathrm{Lie}}}.}
\]
It is obtained through the common spectral partition-Lie controller and
preserves the anchor, complete PD operation monad, symmetric actions,
substitution, units, tangent enveloping monads, and formalized Spencer
representation actions.

The theorem does not assert that either raw support-marked geometric system is
equivalent to the other.
\end{theorem}

\begin{proof}
By definition
\[
    \operatorname{Diff}^{\pi,E_\infty}(\mathfrak G)
    =:\mathfrak a_{\mathfrak G}.
\]
The differentiation-recovery equivalence gives
\[
    \operatorname{Diff}^{\pi,E_\infty}
    (\mathfrak G^{\mathrm{Lie}})
    \simeq
    \mathfrak a_{\mathfrak G}.
\]
Apply the integral partition-Lie/formalized-Spencer equivalence of the
Atiyah--Spencer chapter to $\mathfrak G$ and to
$\mathfrak G^{\mathrm{Lie}}$, and compose through the common controller
$\mathfrak a_{\mathfrak G}$.  That theorem identifies the complete Artinian
Spencer monad with the spectral partition-Lie monad, including the anchor,
symmetric actions, PD substitution, unit, and all higher coherence.  Slicing
at the common controller and stabilizing therefore identifies the two tangent
enveloping monads.

The formalized Spencer representation cochains require one additional
functoriality statement rather than only the monad comparison.  The preceding
chapter constructs
$\mathbf{Sp}^{\mathrm{fm},\bullet}_F$ functorially in the formal leaf $F$ and
proves compatibility with tangent differentiation.  The canonical equivalence
\[
    \mathcal L^{\mathrm{fm,sp}}_{\mathfrak G^{\mathrm{Lie}}}
    \simeq
    \mathcal L^{\mathrm{fm,sp}}_{\mathfrak G}
\]
therefore carries the formalized Spencer cochain object and its representation
action for $\mathfrak G^{\mathrm{Lie}}$ to the corresponding object and action
for $\mathfrak G$.  Thus the listed representation actions are preserved as
well.

The raw support-marked branch is independently geometric and is not identified
with the Artinian cotangent-primitive branch by the preceding equivalences, so
no raw marked-cell equivalence follows.
\end{proof}

\begin{corollary}[Operation recovery on reconstructed integration]
\label{cor:bnli-universal-operation-property}
For every spectral partition-Lie algebroid $\mathfrak a$, the complete Artinian
Spencer operation system of
$\operatorname{Int}^{\mathrm{gpd}}(\mathfrak a)$ is canonically the operation
system of $\mathfrak a$.  This recovers the controller and all constructed
Artinian operations after differentiation; it does not assert that raw native
mapping spaces are determined by those operations.
\end{corollary}

\section{Integration of tangent representations}
\label{sec:bnli-representations}

Let
\[
    \mathfrak a\in
    \operatorname{LieAlgd}^{\pi,E_\infty}_{X/S},
    \qquad
    F_{\mathfrak a}:=\operatorname{Int}^{\mathrm{fm}}(\mathfrak a),
    \qquad
    \mathfrak G_{\mathfrak a}
    :=\operatorname{Int}^{\mathrm{gpd}}(\mathfrak a).
\]
The Atiyah--Spencer chapter constructs
\[
    \operatorname{Rep}^{\mathrm{fm}}(F_{\mathfrak a})
    \simeq
    \operatorname{Rep}^{\pi,E_\infty}(\mathfrak a)
\]
by slicing and stabilizing the formal-moduli/partition-Lie equivalence.
Materialization gives an equally canonical stable formal-stack model.  Passage
from that model to raw parameterized Cartan coefficients is a functor, but is
not declared to be an equivalence without a separate stable native-rigidity
theorem.

\subsection{Stable sliced materialization}

\begin{proposition}[Finite-limit stability of the relative formal-stack slice]
\label{prop:bnli-relative-formal-stack-slice}
Let
\[
    Y\in\operatorname{ModuliStk}^{\mathrm{sp}}(X/S).
\]
Then the full inclusion
\[
\boxed{
    \left(
      \operatorname{ModuliStk}^{\mathrm{sp}}(X/S)
    \right)_{/Y}
    \hookrightarrow
    \left(
      (\mathcal X_{/S})_{X/}
    \right)_{/Y}}
\]
preserves finite limits.  Consequently it induces an exact functor on
stabilizations.
\end{proposition}

\begin{proof}
Let
\[
    D:K\longrightarrow
    \left(\operatorname{ModuliStk}^{\mathrm{sp}}(X/S)\right)_{/Y}
\]
be a finite diagram.  Write $D(k)=Z_k\to Y$.  Extend it to the finite
right-cone diagram
\[
    \bar D:K^\triangleright\longrightarrow(\mathcal X_{/S})_{X/}
\]
by putting the cone vertex equal to $Y$ and using the structure maps
$Z_k\to Y$.  The limit in the ambient slice over $Y$ is the object
\[
\boxed{
    Z:=\lim_{K^\triangleright}\bar D}
\]
with its induced centre $X\to Z$ and map $Z\to Y$.

Every vertex of $\bar D$ is $L_S^{\mathrm{Def}}$-local, and the local
subcategory is closed under all limits.  Hence $Z$ is deformation-local.

Apply the left-exact pro-reduction reflector.  For every vertex $W$ of
$\bar D$, including $W=Y$, the centre gives a canonical equivalence
\[
    \mathsf P^{\mathrm{sp}}X\xrightarrow{\sim}\mathsf P^{\mathrm{sp}}W.
\]
Because every arrow of $\bar D$ is under $X$, these equivalences identify the
entire pro-reduced cone coherently with the constant diagram at
$\mathsf P^{\mathrm{sp}}X$.  Since $K^\triangleright$ has a terminal vertex
and is therefore contractible, left exactness gives
\[
\boxed{
    \mathsf P^{\mathrm{sp}}Z
    \simeq
    \lim_{K^\triangleright}\mathsf P^{\mathrm{sp}}\bar D
    \simeq
    \mathsf P^{\mathrm{sp}}X.}
\]
Thus $X\to Z$ is a spectral nil-isomorphism.

It remains to prove spectral laftness of the centre.  Fix $n$ and a filtered
diagram $\{C_\alpha\}$ of $n$-truncated connective affine tests over $Z$, with
colimit $C$.  For a vertex $W$ of $\bar D$, let
\[
    B_W:=W(C),
    \qquad
    B'_W:=\operatorname*{colim}_\alpha W(C_\alpha).
\]
Because the cone is finite and filtered colimits of spaces commute with finite
limits,
\[
    Z(C)\simeq\lim_W B_W,
    \qquad
    \operatorname*{colim}_\alpha Z(C_\alpha)
    \simeq
    \lim_W B'_W.
\]
For every vertex $W$, spectral laftness of $X\to W$ gives
\[
\boxed{
    \operatorname*{colim}_\alpha X(C_\alpha)
    \simeq
    X(C)\times_{B_W}B'_W.}
\]
These equivalences are natural in $W$.  Taking the finite limit over
$K^\triangleright$ and commuting finite limits with finite limits yields
\[
\begin{aligned}
    \operatorname*{colim}_\alpha X(C_\alpha)
    &\simeq
    \lim_W\left(X(C)\times_{B_W}B'_W\right)\\
    &\simeq
    X(C)\times_{\lim_WB_W}\lim_WB'_W\\
    &\simeq
    X(C)\times_{Z(C)}
    \operatorname*{colim}_\alpha Z(C_\alpha).
\end{aligned}
\]
In the first line the limit of the constant left-hand diagram is the same
space because $K^\triangleright$ is contractible.  This is exactly the
relative laft equation for $X\to Z$.

Therefore $Z$ is again a spectral formal-moduli stack.  The inclusion
preserves finite limits.  Stabilization is functorial for finite-limit-
preserving functors and sends such a functor to an exact functor between the
stable $\infty$-categories.
\end{proof}

\begin{theorem}[Stable formal representation materialization]
\label{thm:bnli-representation-lie-three}
There is a canonical equivalence
\[
\boxed{
    \operatorname{Rep}^{\pi,E_\infty}(\mathfrak a)
    \xrightarrow{\sim}
    \operatorname{Sp}
    \left(
      \left(\operatorname{ModuliStk}^{\mathrm{sp}}(X/S)\right)_{/
      \operatorname{Mat}^{\mathrm{sp}}(F_{\mathfrak a})}
    \right).}
\]
Under this equivalence the tangent enveloping monad, universal stable Artinian
sections, and formalized Spencer representation are the materialized forms of
the corresponding tangent formal-moduli constructions.
\end{theorem}

\begin{proof}
The Atiyah--Spencer tangent equivalence identifies
$\operatorname{Rep}^{\pi,E_\infty}(\mathfrak a)$ with the stabilization of the
slice of $\operatorname{FMP}^{\mathrm{sp}}_{X/S}$ over $F_{\mathfrak a}$.
The global materialization equivalence induces an equivalence on slices, and
stabilization preserves equivalences.  The tangent monad and universal stable
Artinian sections are constructed by slicing and stabilizing the same
Schlessinger localization, so their units, multiplication, and representation
actions transport through the equivalence.
\end{proof}

\begin{construction}[Native Cartan realization]
\label{def:bnli-rep-materialization}
Put
\[
    Y_{\mathfrak a}
    :=
    \operatorname{Mat}^{\mathrm{sp}}(F_{\mathfrak a}),
    \qquad
    Q_{\mathfrak a}^{\mathrm{nat}}
    :=
    Q^{\mathrm{nat}}_{F_{\mathfrak a}},
\]
and let
\[
    \kappa_{\mathfrak a}:
    Q_{\mathfrak a}^{\mathrm{nat}}
    \longrightarrow
    Y_{\mathfrak a}
\]
be the native-to-materialized quotient comparison.  Write
\[
    q_{\mathfrak a}:
    X\longrightarrow Q_{\mathfrak a}^{\mathrm{nat}}
\]
for the native effective centre.

By Proposition~\ref{prop:bnli-relative-formal-stack-slice}, the inclusion
\[
    \left(
      \operatorname{ModuliStk}^{\mathrm{sp}}(X/S)
    \right)_{/Y_{\mathfrak a}}
    \hookrightarrow
    \left(
      (\mathcal X_{/S})_{X/}
    \right)_{/Y_{\mathfrak a}}
\]
induces an exact functor after stabilization.  Pullback along
$\kappa_{\mathfrak a}$ is finite-limit preserving and therefore induces
\[
    \kappa_{\mathfrak a}^*:
    \operatorname{Sp}
    \left(
      \left(
        (\mathcal X_{/S})_{X/}
      \right)_{/Y_{\mathfrak a}}
    \right)
    \longrightarrow
    \operatorname{Sp}
    \left(
      \left(
        (\mathcal X_{/S})_{X/}
      \right)_{/Q_{\mathfrak a}^{\mathrm{nat}}}
    \right).
\]

There is a canonical identification
\[
    \left(
      (\mathcal X_{/S})_{X/}
    \right)_{/Q_{\mathfrak a}^{\mathrm{nat}}}
    \simeq
    \left(
      \mathcal X_{/Q_{\mathfrak a}^{\mathrm{nat}}}
    \right)_{q_{\mathfrak a}/}.
\]
Forget the chosen map from $X$:
\[
    \operatorname{oblv}_{q_{\mathfrak a}}:
    \left(
      \mathcal X_{/Q_{\mathfrak a}^{\mathrm{nat}}}
    \right)_{q_{\mathfrak a}/}
    \longrightarrow
    \mathcal X_{/Q_{\mathfrak a}^{\mathrm{nat}}}.
\]
This forgetful functor creates limits and in particular preserves finite
limits.  Stabilization therefore gives a canonical exact functor
\[
\boxed{
    \operatorname{Sp}(\operatorname{oblv}_{q_{\mathfrak a}}):
    \operatorname{Sp}
    \left(
      \left(
        (\mathcal X_{/S})_{X/}
      \right)_{/Q_{\mathfrak a}^{\mathrm{nat}}}
    \right)
    \longrightarrow
    \mathcal E_{Q_{\mathfrak a}^{\mathrm{nat}}},}
\]
where
\[
    \mathcal E_{Q_{\mathfrak a}^{\mathrm{nat}}}
    :=
    \operatorname{Sp}
    \left(
      \mathcal X_{/Q_{\mathfrak a}^{\mathrm{nat}}}
    \right).
\]
For the effective quotient
\[
    X\twoheadrightarrow Q_{\mathfrak a}^{\mathrm{nat}},
\]
the Atiyah--Spencer Cartan coefficient category is intrinsically
\[
    \operatorname{Rep}^{\mathrm{Cart}}(\mathfrak G_{\mathfrak a})
    :=
    \mathcal E_{Q_{\mathfrak a}^{\mathrm{nat}}},
\]
and effective Cartan descent identifies it equivalently with the stable
$J_{\mathfrak G_{\mathfrak a}}$-coalgebra category over $X$.

Composing these exact functors with
Theorem~\ref{thm:bnli-representation-lie-three} defines
\[
\boxed{
\begin{aligned}
    \operatorname{Real}^{\mathrm{Cart}}_{\mathfrak a}:
    \operatorname{Rep}^{\pi,E_\infty}(\mathfrak a)
    &\xrightarrow{\sim}
    \operatorname{Sp}
    \left(
      \operatorname{ModuliStk}^{\mathrm{sp}}(X/S)_{/Y_{\mathfrak a}}
    \right)\\
    &\longrightarrow
    \operatorname{Rep}^{\mathrm{Cart}}
    (\mathfrak G_{\mathfrak a}).
\end{aligned}}
\]
For $M$ in the source write
\[
    \operatorname{IntRep}_{\mathfrak a}(M)
    :=
    \operatorname{Real}^{\mathrm{Cart}}_{\mathfrak a}(M).
\]
\end{construction}

\begin{proposition}[Compatibility and variance of native Cartan realization]
\label{prop:bnli-native-cartan-realization}
Let
\[
    f:\mathfrak a\longrightarrow\mathfrak b
\]
be a morphism of spectral partition-Lie algebroids over the fixed object
space $X/S$.  Then there is a canonical commutative square of exact functors
\[
\begin{tikzcd}
\operatorname{Rep}^{\pi,E_\infty}(\mathfrak b)
    \arrow[r,"\operatorname{Real}^{\mathrm{Cart}}_{\mathfrak b}"]
    \arrow[d,"f^*"'] &
\operatorname{Rep}^{\mathrm{Cart}}(\mathfrak G_{\mathfrak b})
    \arrow[d,"\operatorname{Int}^{\mathrm{gpd}}(f)^*"]\\
\operatorname{Rep}^{\pi,E_\infty}(\mathfrak a)
    \arrow[r,"\operatorname{Real}^{\mathrm{Cart}}_{\mathfrak a}"'] &
\operatorname{Rep}^{\mathrm{Cart}}(\mathfrak G_{\mathfrak a}).
\end{tikzcd}
\]
The square is coherent for identities and composition.

On the Artinian-formalized branch,
$\operatorname{Real}^{\mathrm{Cart}}_{\mathfrak a}$ carries the universal
stable Artinian-section objects, formalized Spencer cochain objects, the
morphisms expressing their tangent-enveloping actions, and their tangent
obstruction morphisms to their exact pulled-back Cartan realizations.  This
asserts realization of the objects and action diagrams; it does not construct
a raw Cartan monad identified with the tangent enveloping monad.  No raw
first-principal-parts or changing-object comparison is asserted without a
separate native rigidity theorem.
\end{proposition}

\begin{proof}
Formal integration sends $f$ to a morphism
\[
    F_{\mathfrak a}\longrightarrow F_{\mathfrak b},
\]
materialization sends this to
\[
    Y_{\mathfrak a}\longrightarrow Y_{\mathfrak b},
\]
and native integration sends it to
\[
    \mathfrak G_{\mathfrak a}
    \longrightarrow
    \mathfrak G_{\mathfrak b}.
\]
Naturality of
\[
    \kappa^{\mathrm{mat}}:
    Q^{\mathrm{nat}}
    \longrightarrow
    \operatorname{Mat}^{\mathrm{sp}}
\]
gives a commutative square of quotient maps.  Pullback on the relative
factorization slices is contravariantly functorial and commutes with this
square.

For every quotient map $q:X\to Q$, the forgetful functor
\[
    \operatorname{oblv}_{q}:
    (\mathcal X_{/Q})_{q/}
    \longrightarrow
    \mathcal X_{/Q}
\]
commutes strictly with pullback in $Q$.  Hence the stabilized forgetful
functors
\[
    \operatorname{Sp}(\operatorname{oblv}_{q})
\]
also commute with the pullback squares induced by
$\operatorname{Int}^{\mathrm{gpd}}(f)$.  After the intrinsic identification
\[
    \operatorname{Rep}^{\mathrm{Cart}}(\mathfrak G)
    =
    \mathcal E_{Q_{\mathfrak G}},
\]
effective Cartan descent identifies these pullbacks with the usual
contravariant pullback of Cartan coefficients.  Combining the exact
formal-stack slice inclusion, pullback along $\kappa$, the stabilized
forgetful functor, and effective Cartan descent gives the displayed square.
Identity and composition coherences are inherited from the corresponding
coherences of pullback and of the native quotient construction.

The stable Artinian sections, formalized Spencer cochains, the morphisms
expressing the tangent-enveloping action, and their tangent obstruction maps
are functorial constructions in the stabilized formal-moduli slice.  The
realization functor is obtained only by exact functors from that slice---
inclusion, pullback, forgetting the distinguished centre map, and effective
descent.  It therefore carries those objects and morphisms, and the complete
action diagrams among them, to their corresponding Cartan realizations.  No
monad comparison on the raw Cartan category and no identification of raw first
neighbourhoods is used.
\end{proof}

\begin{remark}[No representation-recognition input condition]
No claim is made that
$\operatorname{Real}^{\mathrm{Cart}}_{\mathfrak a}$ is fully faithful or that
its essential image can be recognized by an additional field of structure on a
raw Cartan coefficient.  Establishing such an equivalence would require a
stable analogue of the native mapping-rigidity theorem of
Section~\ref{sec:bnli-native-density}.  No such vanishing condition is inserted
as an additional representation-recognition condition.
\end{remark}

\section{Integration of anchor fibres and actual formal isotropy}
\label{sec:bnli-isotropy-integration}

Let
\[
    \mathbb F^{\pi,E_\infty}_{\mathfrak a}
\]
be the zero-anchor partition-Lie algebroid obtained from \(\mathfrak a\) by
the constructed anchor-fibre right adjoint.  Let
\[
    \mathfrak I_{\mathfrak G_{\mathfrak a},\bullet}
    :=
    B_\bullet I^{\mathrm{for}}_{\mathfrak G_{\mathfrak a}}
\]
be the actual formal-isotropy groupoid and put
\[
    \mathfrak i^{\pi,E_\infty}_{\mathfrak G_{\mathfrak a}}
    :=
    \operatorname{Diff}^{\pi,E_\infty}
    \left(
      \mathfrak I_{\mathfrak G_{\mathfrak a},\bullet}
    \right).
\]

\begin{theorem}[Integrated actual-isotropy comparison]
\label{thm:bnli-integrated-isotropy-comparison}
There is a canonical morphism of zero-anchor spectral partition-Lie
algebroids
\[
\boxed{
    \gamma^{\pi,E_\infty}_{\mathfrak a}:
    \mathfrak i^{\pi,E_\infty}_{\mathfrak G_{\mathfrak a}}
    \longrightarrow
    \mathbb F^{\pi,E_\infty}_{\mathfrak a}.}
\]
Its underlying discrepancy is the functorial cofiber of the underlying
pro-coherent morphism.  No equivalence is asserted in general.
\end{theorem}

\begin{proof}
Apply the actual-isotropy-to-anchor-fibre comparison constructed in the
differentiation chapter to
\(\mathfrak G_{\mathfrak a}\).  By
Corollary~\ref{cor:bnli-diff-int-identity}, its ambient spectral
partition-Lie algebroid is canonically \(\mathfrak a\).  The target of the
constructed comparison is therefore exactly
\(\mathbb F^{\pi,E_\infty}_{\mathfrak a}\).  Functoriality and the discrepancy
cofiber are inherited from that theorem.
\end{proof}

\begin{proposition}[Why actual isotropy cannot be identified with the anchor fibre]
\label{prop:bnli-isotropy-nonidentification}
The canonical actual-isotropy-to-anchor-fibre comparison
\[
    \gamma^{\pi,E_\infty}_{\mathfrak G}:
    \mathfrak i^{\pi,E_\infty}_{\mathfrak G}
    \longrightarrow
    \mathbb F^{\pi,E_\infty}_{\mathfrak G}
\]
need not be an equivalence.  More precisely, if $X\to S$ is $0$-truncated and
$\mathbb T^{\mathrm{pc}}_{X/S}\not\simeq0$, then for the identity formal
groupoid $0_{X,\bullet}$ one has
\[
    U\mathfrak i^{\pi,E_\infty}_{0_X}\simeq0,
    \qquad
    U\mathbb F^{\pi,E_\infty}_{0_X}
    \simeq
    \mathbb T^{\mathrm{pc}}_{X/S},
\]
so $\gamma^{\pi,E_\infty}_{0_X}$ is not invertible.
\end{proposition}

\begin{proof}
Choose a $0$-truncated morphism $X\to S$ with
$\mathbb T^{\mathrm{pc}}_{X/S}\not\simeq0$, for example an ordinary smooth
positive-dimensional scheme over its base, and take the identity formal
groupoid $0_{X,\bullet}$.  For a $0$-truncated object the relative inertia is
the identity group object, so the actual formal isotropy of the identity
groupoid is trivial and its differentiated actual-isotropy controller is zero.
The differentiation chapter computes
\[
    \mathbb A^{\pi,E_\infty}_{0_X}
    \simeq
    \left(
      0\longrightarrow \mathbb T^{\mathrm{pc}}_{X/S}[1]
    \right).
\]
Applying the anchor-fibre right adjoint therefore gives underlying coefficient
\[
    U\mathbb F^{\pi,E_\infty}_{0_X}
    \simeq
    \Omega\mathbb T^{\mathrm{pc}}_{X/S}[1]
    \simeq
    \mathbb T^{\mathrm{pc}}_{X/S},
\]
which is not zero by the chosen hypothesis.  Therefore the underlying map of
$\gamma^{\pi,E_\infty}_{0_X}$ has zero source and nonzero target, so the
canonical comparison is not an equivalence.
\end{proof}

\begin{remark}[No isotropy-exactness input condition]
We do not introduce an ``isotropy-exact'' class as an input category.  When a
specific calculation proves
\(\gamma^{\pi,E_\infty}_{\mathfrak a}\) invertible, that is recorded as the
conclusion of that calculation.
\end{remark}

\section{Integration of the adjoint representation}
\label{sec:bnli-adjoint-integration}

The Atiyah--Spencer chapter constructs three different adjoint-level objects:
the nonlinear formalized adjoint group, its reduced stable coefficient, and
the Artinian-linearized formalized adjoint representation
\[
    \operatorname{Ad}^{\pi,\mathrm{fm}}_{\mathfrak G}:
    \mathbb A_{\mathfrak G}^{\pi,E_\infty}
    \curvearrowright
    \mathbf{Ad}_{\mathfrak G}^{\pi,\mathrm{fm}}.
\]
It also constructs the nonlinear counit
\[
    I_{\mathfrak G}^{\mathrm{Ad,fm}}
    \longrightarrow
    I_{\mathfrak G}^{\mathrm{for}}
\]
and its reduced stable shadow.  The Artinian representation coefficient and
the nonlinear stable coefficient are not identified by definition.  Native
integration therefore realizes the tangent representation and retains the
nonlinear stable counit as a separate comparison.

\begin{theorem}[Native realization of the formalized adjoint tangent representation]
\label{thm:bnli-adjoint-materialization}
Let
\[
    \mathfrak G_{\mathfrak a}
    =
    \operatorname{Int}^{\mathrm{gpd}}(\mathfrak a).
\]
Under the recovery equivalence
$\mathbb A_{\mathfrak G_{\mathfrak a}}^{\pi,E_\infty}\simeq\mathfrak a$, the
formalized adjoint representation
\[
    \operatorname{Ad}^{\pi,\mathrm{fm}}_{\mathfrak G_{\mathfrak a}}:
    \mathfrak a
    \curvearrowright
    \mathbf{Ad}_{\mathfrak G_{\mathfrak a}}^{\pi,\mathrm{fm}}
\]
has a canonical native Cartan realization
\[
\boxed{
    \operatorname{Ad}^{\mathrm{Cart,mat}}_{\mathfrak a}
    :=
    \operatorname{Real}^{\mathrm{Cart}}_{\mathfrak a}
    \left(
      \mathbf{Ad}_{\mathfrak G_{\mathfrak a}}^{\pi,\mathrm{fm}}
    \right)
    \in
    \operatorname{Rep}^{\mathrm{Cart}}(\mathfrak G_{\mathfrak a}).}
\]
This realizes only the Artinian-linearized formalized adjoint representation.
No morphism from this Cartan realization to the independently differentiated
actual-isotropy controller, and no equivalence with the raw stable isotropy
coefficient, is asserted.
\end{theorem}

\begin{proof}
The formal-leaf recovery theorem identifies the controller of
$\mathfrak G_{\mathfrak a}$ with $\mathfrak a$.  The Atiyah--Spencer
formalized adjoint coefficient is therefore an object of
$\operatorname{Rep}^{\pi,E_\infty}(\mathfrak a)$.  Apply the native Cartan
realization functor of Construction~\ref{def:bnli-rep-materialization}; this
gives the displayed Cartan coefficient functorially.

The Atiyah--Spencer construction of the formalized adjoint representation is
obtained by stable Artinian sections, Schlessinger reflection, and tangent
differentiation of the formalized nonlinear adjoint coefficient.
Construction~\ref{def:bnli-rep-materialization} is functorial on that
representation object and therefore gives its native Cartan realization.
The independently differentiated actual-isotropy controller is a different
construction, and the preceding chapter proves no comparison theorem from the
present formalized adjoint representation to that controller.  Nothing in
this argument identifies the resulting Cartan realization with the raw stable
coefficient of the nonlinear isotropy group either.
\end{proof}

\begin{proposition}[Separate nonlinear stable adjoint counit]
\label{prop:bnli-adjoint-stable-counit}
For $\mathfrak G_{\mathfrak a}$ there remains the canonical nonlinear stable
comparison
\[
\boxed{
    \mathbb I^{\mathrm{Ad,fm,st}}_{\mathfrak G_{\mathfrak a}}
    \longrightarrow
    \mathbb I^{\mathrm{st}}_{\mathfrak G_{\mathfrak a}}.}
\]
The source of this morphism is the nonlinear stable coefficient from which the
formalized adjoint tangent representation is constructed by stable Artinian
sections and tangent differentiation.  The target is stable actual formal
isotropy.  No theorem identifies the Artinian linearization of that target with
the independently differentiated actual-isotropy controller, and no canonical
morphism from $\operatorname{Ad}^{\mathrm{Cart,mat}}_{\mathfrak a}$ to
$\mathbb I^{\mathrm{st}}_{\mathfrak G_{\mathfrak a}}$ is inferred without a
separate stable native-recognition theorem.
\end{proposition}

\begin{proof}
The Atiyah--Spencer chapter constructs
$K_{\mathfrak G}^{\mathrm{fm}}=\Pi_{\ell_{\mathfrak G}}K_{\mathfrak G}$ and,
after pullback to $X$, the counit
$I_{\mathfrak G}^{\mathrm{Ad,fm}}\to I_{\mathfrak G}^{\mathrm{for}}$.
Reduced stabilization is functorial, giving the displayed stable morphism.
The same chapter defines the formalized adjoint tangent representation by
stable Artinian sections and tangent differentiation of the source coefficient.
It does not identify the corresponding construction on the target stable
coefficient with the separately differentiated actual-isotropy formal-moduli
controller.

The source of the displayed stable map is a nonlinear stable coefficient,
whereas
$\operatorname{Ad}^{\mathrm{Cart,mat}}_{\mathfrak a}$ is the native Cartan
realization of its Artinian-linearized tangent representation.  No theorem in
the present or preceding chapters identifies those two levels, so no further
map is manufactured.
\end{proof}

\section{One-object Lie integration}
\label{sec:bnli-one-object}

Set $X=S$.  Then
\[
    \mathbb T^{\mathrm{pc}}_{S/S}[1]\simeq0,
\]
so spectral partition-Lie algebroids over $S/S$ are spectral partition-Lie
algebras.

\begin{definition}[Formal group of a spectral partition-Lie algebra]
\label{def:bnli-formal-group}
For
\[
    \mathfrak g
    \in
    \operatorname{Alg}_{\operatorname{Lie}^{\pi}_{S,E_\infty}}
    (\operatorname{QC}^{\vee}_S),
\]
define
\[
    B_\bullet\mu^{\mathrm{gpd}}_{\mathfrak g}
    :=
    \operatorname{Int}^{\mathrm{gpd}}_{S/S}(\mathfrak g).
\]
The degree-one arrow group
\[
\boxed{
    \mu^{\mathrm{gpd}}_{\mathfrak g}}
\]
is the \emph{brave-new formal Lie group} integrated from $\mathfrak g$.
\end{definition}

\begin{theorem}[One-object native Lie III]
\label{thm:bnli-one-object-lie-three}
For every spectral partition-Lie algebra $\mathfrak g$,
\[
\boxed{
    \operatorname{Diff}^{\pi,E_\infty}
    \left(
      B_\bullet\mu^{\mathrm{gpd}}_{\mathfrak g}
    \right)
    \xrightarrow{\sim}
    \mathfrak g.}
\]
The construction is functorial in $\mathfrak g$.  No equivalence between the
entire raw category of native formal groups and the partition-Lie algebra
category is asserted without the native rigidity theorem of
Section~\ref{sec:bnli-native-density}.
\end{theorem}

\begin{proof}
Specialize Theorem~\ref{thm:bnli-lie-three} to $X=S$.  The ambient anchor
target vanishes, so the many-object partition-Lie category reduces to the
one-object spectral partition-Lie algebra category.
\end{proof}

\section{Parallel Hopf-derived-primitives and partition-Lie outputs}
\label{sec:bnli-hopf-comparison}

For the integrated formal group $\mu^{\mathrm{gpd}}_{\mathfrak g}$, form the
suspension Hopf object
\[
    H^{\mathrm{gpd}}_{\mathfrak g}
    :=
    \Sigma^\infty_{+,S}\mu^{\mathrm{gpd}}_{\mathfrak g}
\]
and the independent derived-primitives spectral Lie algebra
\[
    \mathfrak l^{\mathrm{sp}}_{\mathfrak g}
    :=
    \operatorname{Prim}^{\mathrm{der}}
    \left(
      H^{\mathrm{gpd}}_{\mathfrak g}
    \right).
\]
The objects $\mathfrak g$ and
$\mathfrak l^{\mathrm{sp}}_{\mathfrak g}$ remain distinct.

\begin{theorem}[One-object parallel Lie outputs]
\label{thm:bnli-hopf-comparison}
The integrated one-object formal group canonically determines, functorially in
$\mathfrak g$:
\begin{enumerate}
    \item the recovered spectral partition-Lie algebra
    \[
        \operatorname{Diff}^{\pi,E_\infty}
        (B_\bullet\mu^{\mathrm{gpd}}_{\mathfrak g})
        \simeq \mathfrak g;
    \]
    \item the independent Hopf-derived-primitives spectral Lie algebra
    $\mathfrak l^{\mathrm{sp}}_{\mathfrak g}$;
    \item the Atiyah--Spencer identification of the mixed Artinian derivative
    of the right-normalized nonlinear commutator with the binary operation of
    the independently differentiated actual-isotropy controller, which in the
    one-object case is the binary operation of $\mathfrak g$.
\end{enumerate}
No morphism
\[
    \mathfrak l^{\mathrm{sp}}_{\mathfrak g}
    \longrightarrow \mathfrak g,
    \qquad
    \mathfrak g\longrightarrow\mathfrak l^{\mathrm{sp}}_{\mathfrak g},
\]
and no aritywise Hopf/partition comparison in a common operation target is
asserted unless a separate theorem constructs such a morphism.
\end{theorem}

\begin{proof}
The first item is Theorem~\ref{thm:bnli-one-object-lie-three}.  The second is
the derived-primitives construction applied to the suspension Hopf object of
the same formal group.  For the third, the one-object formal isotropy group of
$B_\bullet\mu^{\mathrm{gpd}}_{\mathfrak g}$ is the arrow group itself.  The
Atiyah--Spencer arity-two theorem identifies the directional bi-square-zero
derivative of its right-normalized commutator with the binary operation of the
differentiated actual-isotropy formal-moduli problem.  In the one-object case
that formal-moduli problem is the formal group itself, so differentiation
recovery identifies this controller with $\mathfrak g$.

The Hopf-derived-primitives functor is a different construction on the same
nonlinear group object.  The preceding chapters do not construct a morphism
from its spectral Lie operations to the Artinian cotangent-operation target.
Thus common origin in the nonlinear group does not supply an operation-level
comparison, and none is inserted here.
\end{proof}

\begin{remark}[No Hopf--partition discrepancy without a comparison map]
A fibre or cofiber discrepancy requires an actual morphism in a common stable
category.  Since no controller-level or aritywise Hopf/partition morphism has
been constructed, this chapter introduces no Hopf--partition discrepancy
object.  Such discrepancies may be added after a separate comparison theorem.
\end{remark}

\section{Materialized formal spectra and represented comparison}
\label{sec:bnli-represented-materialization}

Materialization is presentation-free.  A complete augmented algebra may
supply an Artinian functor, but no geometric formal-spectrum object is assumed
in order to construct its spectral formal-moduli stack.

\begin{proposition}[Complete almost-finite arrows define formal-moduli functors]
\label{prop:bnli-complete-arrow-fmp}
Let
\[
    A^\wedge\longrightarrow B
    \in
    \operatorname{CAlg}^{\mathrm{sp},\wedge\mathrm{aft}}_{R/B}.
\]
On spectral Artinian tests define
\[
\boxed{
    F_{A^\wedge}(C\to B)
    :=
    \operatorname{Map}_{\operatorname{CAlg}(\operatorname{Sp})_{R//B}}
    (A^\wedge,C).}
\]
Then $F_{A^\wedge}$ belongs canonically to
$\operatorname{FMP}^{\mathrm{sp}}_{B/R}$.
\end{proposition}

\begin{proof}
Let
\[
    \mathfrak a_{A^\wedge}
    :=
    D^{\mathrm{sp}}_{B/R}(A^\wedge)
\]
be the spectral Koszul-dual algebroid of the complete almost-finite arrow.
For every spectral Artinian $C\to B$, full faithfulness of
$D^{\mathrm{sp}}_{B/R}$ gives
\[
\begin{aligned}
    F_{A^\wedge}(C)
    &\simeq
    \operatorname{Map}_{\operatorname{LieAlgd}^{\pi,E_\infty}_{B/R}}
    \left(
      D^{\mathrm{sp}}_{B/R}(C),
      D^{\mathrm{sp}}_{B/R}(A^\wedge)
    \right)\\
    &=
    \operatorname{Int}^{\pi,E_\infty}_{B/R}
    (\mathfrak a_{A^\wedge})(C).
\end{aligned}
\]
The right-hand side is the already constructed formal integration of a
spectral partition-Lie algebroid, hence is a spectral formal-moduli problem.
Thus represented Artinian functoriality is a consequence of complete
Koszul-dual recognition, not an additional materializability hypothesis.
\end{proof}

\begin{definition}[Materialized formal spectrum]
\label{def:bnli-materialized-formal-spectrum}
For
$A^\wedge\to B\in
\operatorname{CAlg}^{\mathrm{sp},\wedge\mathrm{aft}}_{R/B}$ define
\[
\boxed{
    \operatorname{Spf}^{\mathrm{mat}}_{B/R}(A^\wedge)
    :=
    \operatorname{Mat}^{\mathrm{sp}}_{B/R}(F_{A^\wedge}).}
\]
No geometric formal-spectrum object is required in this definition.
\end{definition}

\begin{proposition}[Recognition and presentation independence of materialized formal spectra]
\label{prop:bnli-materialized-spf-recognition}
There is a canonical equivalence
\[
\boxed{
    \operatorname{Leaf}^{\mathrm{Art}}_{B/R}
    \operatorname{Spf}^{\mathrm{mat}}_{B/R}(A^\wedge)
    \xrightarrow{\sim}
    F_{A^\wedge}.}
\]
If $A^\wedge$ and $\widetilde A^\wedge$ determine equivalent Artinian
formal-moduli functors, then they determine canonically equivalent
materialized formal spectra.  A morphism in the complete almost-finite
augmented sector induces the corresponding contravariant morphism of
materialized formal spectra.
\end{proposition}

\begin{proof}
The first assertion is the unit equivalence of affine materialization.  If
$F_{A^\wedge}\simeq F_{\widetilde A^\wedge}$, applying the functor
$\operatorname{Mat}^{\mathrm{sp}}_{B/R}$ gives the second assertion.  A map
$A^\wedge\to\widetilde A^\wedge$ induces by precomposition
\[
    F_{\widetilde A^\wedge}\longrightarrow F_{A^\wedge};
\]
functoriality of materialization gives the contravariant map of materialized
formal spectra.  No complete-algebra presentation is retained by the output.
\end{proof}

\begin{theorem}[Comparison with an independently constructed adic formal spectrum]
\label{thm:bnli-represented-materialization}
Suppose a separate geometric construction assigns to $A^\wedge$ a centred
spectral formal-moduli stack
\[
    \operatorname{Spf}^{\mathrm{adic}}(A^\wedge)
\]
and suppose its Artinian restriction is canonically the represented functor
$F_{A^\wedge}$.  Then there is a unique canonical equivalence of centred
spectral formal-moduli stacks
\[
\boxed{
    \operatorname{Spf}^{\mathrm{mat}}_{B/R}(A^\wedge)
    \xrightarrow{\sim}
    \operatorname{Spf}^{\mathrm{adic}}(A^\wedge).}
\]
Thus geometric adic representability, when independently available, is a
comparison theorem and is not an input to materialization.
\end{theorem}

\begin{proof}
Both sides are spectral formal-moduli stacks and both have Artinian leaf
$F_{A^\wedge}$.  Full faithfulness of affine Artinian restriction therefore
identifies them through a contractible space of equivalences.  The comparison
is the unique morphism whose Artinian restriction is the identity of
$F_{A^\wedge}$.
\end{proof}

\begin{remark}[No formal-spectrum materializability hypothesis]
The chapter never requires the assertion ``$A^\wedge$ has a materializable
formal spectrum'' as a field of structure.  The object
$\operatorname{Spf}^{\mathrm{mat}}(A^\wedge)$ is constructed for every
complete almost-finite augmented arrow by
Proposition~\ref{prop:bnli-complete-arrow-fmp}.  Any independently constructed
adic formal spectrum is compared to it afterward by
Theorem~\ref{thm:bnli-represented-materialization}.
\end{remark}

\section{Rational CE and BCH integration}
\label{sec:bnli-rational-integration}

Assume rational locally coherent coefficients in the dg rectification domain
constructed in the preceding chapters, and let $\mathfrak g_{\mathrm{dg}}$ be
the dg Lie algebra rectifying a spectral partition-Lie algebra $\mathfrak g$.

\begin{definition}[Materialized CE classifying target]
\label{def:bnli-materialized-ce-target}
Put
\[
    F_{\mathfrak g}
    :=
    \operatorname{Int}^{\mathrm{fm}}(\mathfrak g),
    \qquad
    Y^{\mathrm{CE,mat}}_{\mathfrak g}
    :=
    \operatorname{Mat}^{\mathrm{sp}}(F_{\mathfrak g}).
\]
The notation records the rational CE comparison proved below; the object is
constructed spectrally by materialization and does not presuppose a complete
spectral Chevalley--Eilenberg algebra.
\end{definition}

\begin{theorem}[Rational materialized CE comparison]
\label{thm:bnli-rational-comparison}
Rational Artinian rectification identifies the spectral formal-moduli problem
$F_{\mathfrak g}$ with the dg formal-moduli problem controlled by
$\mathfrak g_{\mathrm{dg}}$.  Consequently
\[
\boxed{
    L_S^{\mathrm{Def}}Q^{\mathrm{nat}}_{F_{\mathfrak g}}
    \xrightarrow{\sim}
    Y^{\mathrm{CE,mat}}_{\mathfrak g}.}
\]

Suppose, on a separately stated finite-type dg CE domain, that the complete dg
Chevalley--Eilenberg algebra
\[
    C^*_{\mathrm{CE}}(\mathfrak g_{\mathrm{dg}})
\]
represents the dg Artinian formal-moduli functor of
$\mathfrak g_{\mathrm{dg}}$.  Transport that \emph{Artinian functor}, rather
than an unconstructed category of arbitrary complete algebras, through the
proved rational Artinian rectification equivalence.  Denote the resulting
spectral Artinian formal-moduli problem by
$F^{\mathrm{sp}}_{\mathrm{CE}}(\mathfrak g_{\mathrm{dg}})$.  Then there are
canonical equivalences
\[
\boxed{
    F^{\mathrm{sp}}_{\mathrm{CE}}(\mathfrak g_{\mathrm{dg}})
    \simeq
    F_{\mathfrak g}}
\]
and
\[
\boxed{
    \operatorname{Spf}^{\mathrm{CE,mat}}
    (\mathfrak g_{\mathrm{dg}})
    :=
    \operatorname{Mat}^{\mathrm{sp}}
    \left(F^{\mathrm{sp}}_{\mathrm{CE}}(\mathfrak g_{\mathrm{dg}})\right)
    \xrightarrow{\sim}
    Y^{\mathrm{CE,mat}}_{\mathfrak g}.}
\]
Define the materialized CE formal group by the based loop object
\[
\boxed{
    \mu^{\mathrm{CE,mat}}_{\mathfrak g_{\mathrm{dg}}}
    :=
    S\times_{
      \operatorname{Spf}^{\mathrm{CE,mat}}(\mathfrak g_{\mathrm{dg}})}S.}
\]
No raw equivalence between
$\mu^{\mathrm{gpd}}_{\mathfrak g}$ and
$\mu^{\mathrm{CE,mat}}_{\mathfrak g_{\mathrm{dg}}}$ is asserted without the
separate native rigidity theorem.
\end{theorem}

\begin{proof}
The rational rectification theorem of the differentiation chapter identifies
the spectral and dg Artinian categories, structured square-zero extensions,
Schlessinger generators, and hence the corresponding formal-moduli
categories.  It therefore carries $F_{\mathfrak g}$ to the dg formal-moduli
problem controlled by $\mathfrak g_{\mathrm{dg}}$.  The first displayed
geometric equivalence is
Theorem~\ref{thm:bnli-native-quotient-materialization}.

On the finite-type CE domain, representedness of the dg formal-moduli problem
by $C^*_{\mathrm{CE}}(\mathfrak g_{\mathrm{dg}})$ identifies its dg Artinian
functor.  Applying the proved Artinian rectification equivalence produces
$F^{\mathrm{sp}}_{\mathrm{CE}}(\mathfrak g_{\mathrm{dg}})$ and identifies it
with $F_{\mathfrak g}$.  Materialization is fully faithful and functorial, so
materializing that equivalence gives the displayed equivalence of
classifying targets.  The correctly typed one-object group associated to a
classifying target is its based loop, which gives the final definition.
\end{proof}

\begin{proposition}[Comparison with a separately constructed spectral CE formal spectrum]
\label{prop:bnli-adic-ce-comparison}
If, on a further represented domain, an adic spectral formal spectrum
\[
    \operatorname{Spf}^{\mathrm{adic}}_{
      \mathrm{CE}}(\mathfrak g_{\mathrm{dg}})
\]
is independently constructed and its Artinian restriction is
$F^{\mathrm{sp}}_{\mathrm{CE}}(\mathfrak g_{\mathrm{dg}})$, then
Theorem~\ref{thm:bnli-represented-materialization} gives
\[
\boxed{
    \operatorname{Spf}^{\mathrm{CE,mat}}(\mathfrak g_{\mathrm{dg}})
    \xrightarrow{\sim}
    \operatorname{Spf}^{\mathrm{adic}}_{\mathrm{CE}}
    (\mathfrak g_{\mathrm{dg}}).}
\]
Taking based loops gives the corresponding equivalence of represented CE
formal groups.
\end{proposition}

\begin{corollary}[Classical BCH comparison on its external theorem domain]
\label{cor:bnli-bch-comparison}
On the classical degree-zero finite-dimensional nilpotent domain for which the
ordinary CE/BCH representation theorem is available, applying that classical
theorem to the dg CE model identifies the represented CE based loop with the
classical BCH formal group.  This classical identification is a separate
comparison input and is not used in the proof of native Brave-New Lie III.
No raw native/BCH equivalence is inferred from it.
\end{corollary}

\begin{remark}[No complete-algebra or Sullivan input]
The rational argument transports the already constructed Artinian
formal-moduli functor.  It does not assume an equivalence between categories
of arbitrary complete dg commutative algebras and arbitrary complete rational
spectral commutative algebras.  Likewise, no Sullivan integration functor has
been constructed in the preceding chapters, so no Sullivan comparison is
asserted here.
\end{remark}

\section{Integration of action algebroids}
\label{sec:bnli-action-integration}

Let $\mathfrak g$ be a one-object spectral partition-Lie algebra and let
\[
    \mathfrak g\ltimes X
    \longrightarrow
    \mathbb T^{\mathrm{pc}}_{X/S}[1]
\]
be a constructed spectral partition action algebroid.

\begin{theorem}[Native integration of an action algebroid]
\label{thm:bnli-action-integration}
The native groupoid
\[
\boxed{
    \operatorname{Int}^{\mathrm{gpd}}
    \left(
      \mathfrak g\ltimes X
    \right)}
\]
is canonically defined and differentiates to the given action algebroid:
\[
    \operatorname{Diff}^{\pi,E_\infty}
    \operatorname{Int}^{\mathrm{gpd}}
    (\mathfrak g\ltimes X)
    \xrightarrow{\sim}
    \mathfrak g\ltimes X.
\]

Suppose in addition that a nonlinear action of the integrated formal group
$\mu^{\mathrm{gpd}}_{\mathfrak g}$ on $X$ has been separately constructed and
that differentiation of its formalized action groupoid is identified with the
given action algebroid.  Then there is a canonical Lie-comparison span
\[
\boxed{
    \mu^{\mathrm{gpd}}_{\mathfrak g}\ltimes X
    \longrightarrow
    \mathfrak K^{\mathrm{Lie}}_{\mu^{\mathrm{gpd}}_{\mathfrak g}\ltimes X}
    \longleftarrow
    \operatorname{Int}^{\mathrm{gpd}}
    (\mathfrak g\ltimes X).}
\]
The conditional span is the raw comparison actually supplied by the
pre-existing nonlinear action; no nonlinear action is manufactured from the
infinitesimal action algebroid by this theorem.
\end{theorem}

\begin{proof}
The first assertion is the Native Brave-New Lie III theorem applied to the
action algebroid.  Under the additional nonlinear-action hypothesis, put
\[
    G:=\mu^{\mathrm{gpd}}_{\mathfrak g}.
\]
By Definition~\ref{def:bnli-formal-group}, the one-object groupoid
$B_\bullet G$ is a native formal Lie groupoid.  Degree-one recognition of
unitwise formality from Chapter~13 therefore makes $G$ complete along its
identity section.  Hence the canonical map from the constructed formal
identity group
\[
    \mu^{\mathrm{mod}}_{G/S}\longrightarrow G
\]
is an equivalence.  The action-groupoid completion theorem of Chapter~13 now
gives
\[
    (G\ltimes X)^{\mathrm{dR}}
    \xrightarrow{\sim}
    \mu^{\mathrm{mod}}_{G/S}\ltimes X
    \xrightarrow{\sim}
    G\ltimes X.
\]
Thus the separately constructed nonlinear action groupoid $G\ltimes X$ is
itself a native formal Lie groupoid.  Its controller is
$\mathfrak g\ltimes X$ by the additional hypothesis and the action
differentiation theorem of the preceding chapters.  Its reconstructed
groupoid is therefore the right-hand term above.  Apply the general comparison
span of Theorem~\ref{thm:bnli-completion-map} to $G\ltimes X$.  No direct raw
equivalence of the outer terms is inferred.
\end{proof}

\section{Integration of gauge and Atiyah algebroids}
\label{sec:bnli-gauge-integration}

Let \(P_\chi\to X\) be a principal right \(G\)-object, let
\(\operatorname{Gauge}^{\mathrm{for}}(P_\chi)\) be its formalized gauge
groupoid, and let
\[
    \operatorname{At}^{\pi,E_\infty}(P_\chi/X)
\]
be the spectral Atiyah algebroid constructed in the differentiation chapter.

\begin{theorem}[Groupoid integration of the spectral Atiyah algebroid]
\label{thm:bnli-atiyah-integration}
There is a canonical comparison span
\[
\boxed{
    \operatorname{Gauge}^{\mathrm{for}}(P_\chi)
    \longrightarrow
    \mathfrak K^{\mathrm{Lie}}_{\operatorname{Gauge}^{\mathrm{for}}(P_\chi)}
    \longleftarrow
    \operatorname{Int}^{\mathrm{gpd}}
    \left(
      \operatorname{At}^{\pi,E_\infty}(P_\chi/X)
    \right).}
\]
The source and the integrated target have canonically equivalent spectral
Artinian Atiyah--Spencer controllers and operation systems.

The span is the unconditional raw comparison.  Equality of the spectral controller and Artinian operation system does not by itself imply a direct raw equivalence of the outer groupoids.
\end{theorem}

\begin{proof}
The differentiation chapter computes
\[
    \operatorname{Diff}^{\pi,E_\infty}
    \operatorname{Gauge}^{\mathrm{for}}(P_\chi)
    \simeq
    \operatorname{At}^{\pi,E_\infty}(P_\chi/X).
\]
Apply Theorem~\ref{thm:bnli-completion-map}.  The operation-system statement is
Theorem~\ref{thm:bnli-spencer-invariance}.  No direct raw equivalence is inferred from controller equality alone.
\end{proof}

\begin{corollary}[Atiyah controller recovery]
\label{cor:bnli-atiyah-sequence-integrates}
The native integration of
$\operatorname{At}^{\pi,E_\infty}(P_\chi/X)$ differentiates back to that
spectral Atiyah controller.  The actual formal isotropy and the homotopy anchor
fibre are related by the canonical comparison
$\gamma^{\pi,E_\infty}$ and are not identified unless that comparison is
separately proved invertible in the gauge calculation at hand.
\end{corollary}

\section{Ordinary smooth Lie algebroids and canonical formal integration}
\label{sec:bnli-classical-formal-integration}

Let
\[
    A\longrightarrow T_{X/S}
\]
be an ordinary smooth characteristic-zero Lie algebroid on a smooth classical
base in the locally coherent rational coefficient range.  Regard $A$ as a dg
Lie algebroid $A_{\mathrm{dg}}$ concentrated in degree zero.  Let
\[
    \mathsf R_{\mathbb Q}:
    \operatorname{LieAlgd}^{\pi,E_\infty}_{X/S}
    \xrightarrow{\sim}
    \operatorname{LieAlgd}^{\mathrm{pc}}_{\mathrm{dg}}(X/S)
\]
denote the global rational rectification/desuspension equivalence constructed
in the differentiation chapter.  Define the spectral enhancement of $A$ by
\[
\boxed{
    A^{\pi,E_\infty}
    :=
    \mathsf R_{\mathbb Q}^{-1}(A_{\mathrm{dg}}).}
\]

\begin{definition}[Canonical formal integration of a classical Lie algebroid]
\label{def:bnli-classical-formal-integration}
Define
\[
\boxed{
    \operatorname{Int}^{\mathrm{for}}(A)
    :=
    \operatorname{Int}^{\mathrm{gpd}}
    \left(
      A^{\pi,E_\infty}
    \right).}
\]
\end{definition}

\begin{theorem}[Formal integration without global smooth integrability]
\label{thm:bnli-classical-formal-integration}
Every ordinary smooth characteristic-zero Lie algebroid in the stated
coefficient range admits the canonical native formal groupoid integration
\[
    \operatorname{Int}^{\mathrm{for}}(A)
    \rightrightarrows X.
\]
Its spectral differentiation is canonically the spectral enhancement of
\(A\), and its rational dg differentiation is canonically the original dg
Lie-algebroid controller.

On the local finite-type CE range, the deformation-formalized native quotient,
equivalently the materialized formal-leaf target, is canonically the
materialized CE classifying target of
Theorem~\ref{thm:bnli-rational-comparison}.  If an adic spectral CE formal
spectrum is independently available, Proposition~\ref{prop:bnli-adic-ce-comparison}
identifies it with that materialized target.  The represented first-order
comparison on the formalized target recovers the classical anchor and
Lie-algebroid bracket of \(A\).  No raw representability of the native
Artinian-cell quotient is asserted.

No global smooth Lie-groupoid integration is used in this construction.
\end{theorem}

\begin{proof}
Apply native groupoid integration to \(A^{\pi,E_\infty}\).  The equality of
spectral controllers is Corollary~\ref{cor:bnli-diff-int-identity}, and
rational rectification gives the dg controller.

Since an ordinary smooth Lie algebroid is locally finite locally free in the
classical geometric range, the ordinary dg CE construction represents its dg
Artinian formal-moduli functor on the corresponding local finite-type CE
charts.  Rational Artinian rectification followed by
Theorem~\ref{thm:bnli-rational-comparison} identifies the resulting
materialized CE target with
\(\operatorname{Mat}^{\mathrm{sp}}(F_A)\), and
Theorem~\ref{thm:bnli-native-quotient-materialization} identifies the latter
with the deformation formalization of the native quotient.  The represented
first-order differentiation theorem of the differentiation chapter then
identifies the cotangent-transitivity anchor and the differentiated
right-invariant commutator with the classical anchor and bracket.  This proof never identifies the undeformation-localized native quotient with an
adic or materialized CE classifying target.
\end{proof}

\section{Comparison with a pre-existing smooth integration}
\label{sec:bnli-classical-smooth-comparison}

Suppose an ordinary smooth Lie groupoid
\[
    \mathcal H_1\rightrightarrows X
\]
integrates the ordinary Lie algebroid \(A\).  Form its canonical unitwise
relative de Rham completion
\[
    \widehat{\mathcal H}^{\,\mathrm{dR}}_\bullet.
\]
The classical differentiation theorem of the differentiation chapter
identifies its spectral controller with \(A^{\pi,E_\infty}\).

\begin{theorem}[Formal-germ comparison with a smooth integration]
\label{thm:bnli-smooth-integration-comparison}
There is a canonical comparison span
\[
\boxed{
    \widehat{\mathcal H}^{\,\mathrm{dR}}_\bullet
    \longrightarrow
    \mathfrak K^{\mathrm{Lie}}_{\widehat{\mathcal H}^{\,\mathrm{dR}}}
    \longleftarrow
    \operatorname{Int}^{\mathrm{for}}(A).}
\]
Both outer terms have canonically equivalent spectral partition-Lie
controllers.  The displayed span is the unconditional raw comparison; no
direct raw equivalence is inferred solely from equality of controllers.
\end{theorem}

\begin{proof}
Apply Theorem~\ref{thm:bnli-completion-map} to
\(\widehat{\mathcal H}^{\,\mathrm{dR}}_\bullet\).  Its controller is
\(A^{\pi,E_\infty}\), so its reconstructed groupoid is
\(\operatorname{Int}^{\mathrm{for}}(A)\).  The comparison span is canonical.  A direct raw equivalence would require an additional native comparison theorem and is not inferred from controller equality.
\end{proof}

\begin{corollary}[Source-simply-connected smooth comparison]
\label{cor:bnli-ssc-formal-germ}
Whenever a source-simply-connected smooth integration of $A$ exists, its
unitwise relative de Rham completion and
$\operatorname{Int}^{\mathrm{for}}(A)$ have canonically equivalent spectral
partition-Lie controllers and are related by the comparison span of
Theorem~\ref{thm:bnli-smooth-integration-comparison}.  No raw equivalence is
asserted without the separate native rigidity comparison.
\end{corollary}

\section{Global smooth integrability remains separate}
\label{sec:bnli-global-integrability}

\begin{remark}[Formal Lie III versus global Lie III]
The integration theorem constructs a formal groupoid in the native big
spectral Nisnevich \(\infty\)-topos.  It does not construct a global smooth
Hausdorff Lie groupoid from every ordinary Lie algebroid.

Monodromy, source-simply-connected global integration, smooth
representability, Hausdorffness, properness, and other global analytic
questions remain questions of ordinary global Lie-algebroid integrability.
None is inserted as an input to the formal integration functor.
\end{remark}

\section{Functoriality of groupoid integration}
\label{sec:bnli-functoriality}

\subsection{Fixed-object functoriality}

\begin{proposition}[Fixed-object functoriality]
\label{prop:bnli-fixed-object-functoriality}
For fixed \(X/S\), every morphism
\[
    f:\mathfrak a\longrightarrow\mathfrak b
\]
of spectral partition-Lie algebroids induces canonically
\[
    \operatorname{Int}^{\mathrm{gpd}}(f):
    \operatorname{Int}^{\mathrm{gpd}}(\mathfrak a)
    \longrightarrow
    \operatorname{Int}^{\mathrm{gpd}}(\mathfrak b).
\]
The construction preserves identities, composition, and all higher coherence.
\end{proposition}

\begin{proof}
This is the fixed-object part of
Theorem~\ref{thm:bnli-integration-functor}.
\end{proof}

\subsection{Formalized and coefficient-level spectral-etale transport}

\begin{proposition}[Etale transport after deformation formalization]
\label{prop:bnli-etale-formalized-transport}
Let
\[
    c:S'\longrightarrow S
\]
be spectral etale in the coherent base-change range of the differentiation
chapter, and let
\[
    u:X'\longrightarrow X\times_SS'
\]
be spectral etale in its fixed-base varying-object range.  Let $f$ denote the
composite change of coefficient pair obtained by first applying $c$ and then
$u$.  For
\[
    \mathfrak a\in
    \operatorname{LieAlgd}^{\pi,E_\infty}_{X/S}
\]
write
\[
    f_{\mathrm{FMP}}^*
    :=u_{\mathrm{FMP}}^*c_{\mathrm{FMP}}^*,
    \qquad
    f_{\mathrm{Lie}}^*
    :=u_{\mathrm{Lie}}^*c_{\mathrm{Lie}}^*
\]
for the two already constructed restriction functors.  Then formal integration
satisfies the canonical exchange equivalence
\[
\boxed{
    f_{\mathrm{FMP}}^*F_{\mathfrak a}
    \xrightarrow{\sim}
    F_{f_{\mathrm{Lie}}^*\mathfrak a}.}
\]
The materialized targets obey the corresponding formal-completion comparison
\[
\boxed{
    \operatorname{Mat}^{\mathrm{sp}}_{X'/S'}
    \left(f_{\mathrm{FMP}}^*F_{\mathfrak a}\right)
    \simeq
    \widehat{
      \left(
        \operatorname{Mat}^{\mathrm{sp}}_{X/S}(F_{\mathfrak a})
        \times_SS'
      \right)
    }^{\,\mathrm{sp}}_{X'}.}
\]
Equivalently, after deformation formalization, the corresponding native
quotients have the same formal-moduli-stack images.

No raw equivalence
\[
    f^*\operatorname{Int}^{\mathrm{gpd}}_{X/S}(\mathfrak a)
    \simeq
    \operatorname{Int}^{\mathrm{gpd}}_{X'/S'}
    (f_{\mathrm{Lie}}^*\mathfrak a)
\]
is asserted solely from this changing-object spectral-etale coefficient
transport.
\end{proposition}

\begin{proof}
The differentiation chapter constructs spectral-etale base change in the
base variable and spectral-etale restriction in the object variable, proves
coherent compatibility of the formal-moduli/partition-Lie equivalence with
both, and proves their two-variable pasting compatibility.  The construction
of $\operatorname{Int}^{\mathrm{fm}}$ in
Definition~\ref{def:bnli-formal-integration} is assembled from the inverse
affine equivalences using precisely those coherences.  Applying first the base
change and then the object restriction therefore gives
\[
    f_{\mathrm{FMP}}^*F_{\mathfrak a}
    \simeq
    F_{f_{\mathrm{Lie}}^*\mathfrak a}.
\]

Spectral-etale base change of the deformation reflector handles the first
stage geometrically.  Proposition~\ref{prop:bnli-affine-materialization-etale}
and the global relative-reduced formal-completion descent of
Section~\ref{sec:bnli-materialization-descent} handle the subsequent change of
centre.  Together they give the displayed materialization comparison.
Theorem~\ref{thm:bnli-native-quotient-materialization} identifies the two
sides with the deformation localizations of the corresponding native
quotients.

The final sentence records the exact scope of the nonlinear theory.  The
constructed coefficient variance is the composite of the two spectral-etale
variances just used; it is not an arbitrary varying-object functor.  Moreover
Chapter~13 shows that replacing the fixed object space by an arbitrary atlas
need not preserve unitwise formality.  Hence no raw atlas-invariance statement
follows.
\end{proof}

\begin{remark}[Base change of an already constructed native groupoid]
This restriction does not contradict ordinary base change of a fixed native
formal groupoid along a Cartesian change of spectral base: Chapter~13 proves
that the nonlinear quotient, relation, isotropy, and unitwise formalization
commute with such base change.  The present point is narrower: the integration
construction is not declared invariant under replacing its fixed object space
by an arbitrary spectral-etale atlas.
\end{remark}

\section{Deformation-formalized zero recovery and finite limits}
\label{sec:bnli-products}

\begin{proposition}[Deformation-formalized zero recovery]
\label{prop:bnli-zero-integration}
Let $0_{X/S}$ denote the zero spectral partition-Lie algebroid.  Then
\[
\boxed{
    L_S^{\mathrm{Def}}Q^{\mathrm{nat}}_{F_0}
    \xrightarrow{\sim} X,}
\]
where $F_0=\operatorname{Int}^{\mathrm{fm}}(0_{X/S})$.  Equivalently, the
formal leaf of the native integrated groupoid is the formal integration of the
zero controller and
\[
\boxed{
    \operatorname{Diff}^{\pi,E_\infty}
    \operatorname{Int}^{\mathrm{gpd}}(0_{X/S})
    \simeq0_{X/S}.}
\]
No raw equivalence
$\operatorname{Int}^{\mathrm{gpd}}(0_{X/S})\simeq0_{X,\bullet}$ is asserted.
\end{proposition}

\begin{proof}
The identity formal groupoid $0_{X,\bullet}$ has effective quotient
$X\xrightarrow{\mathrm{id}}X$.  Its global formal leaf is the Artinian lifting
functor of the represented centre $X\to X$, and the differentiation chapter
computes its controller as the zero algebroid
\[
    0\longrightarrow\mathbb T^{\mathrm{pc}}_{X/S}[1].
\]
By the formal-integration equivalence, that Artinian lifting functor is
canonically $F_0$.  The global materialization theorem therefore identifies
\[
    \operatorname{Mat}^{\mathrm{sp}}(F_0)\simeq X
\]
as centred deformation-local targets.  Applying
Theorem~\ref{thm:bnli-native-quotient-materialization} gives
\[
    L_S^{\mathrm{Def}}Q^{\mathrm{nat}}_{F_0}
    \simeq X.
\]
Differentiation recovery is the specialization of
Corollary~\ref{cor:bnli-diff-int-identity}.

The formal-moduli problem $F_0$ is not supported only at the identity Artinian
augmentation: split zero-derivation square-zero tests have canonical points
mapping their free zero-anchor Koszul-dual cells to the zero algebroid.  Hence
the native Artinian-cell colimit need not collapse to $X$ before deformation
localization, and no raw zero formula follows.
\end{proof}

\begin{proposition}[Finite limits after deformation formalization]
\label{prop:bnli-products}
Let $K$ be a finite diagram of spectral partition-Lie algebroids and let
\[
    \mathfrak a
    =
    \operatorname*{lim}_{k\in K}
    \mathfrak a_k
\]
be its limit in
$\operatorname{LieAlgd}^{\pi,E_\infty}_{X/S}$.  Then
\[
\boxed{
    L_S^{\mathrm{Def}}Q^{\mathrm{nat}}_{F_{\mathfrak a}}
    \simeq
    \operatorname*{lim}^{\operatorname{ModuliStk}^{\mathrm{sp}}(X/S)}_{k\in K}
    \operatorname{Mat}^{\mathrm{sp}}(F_{\mathfrak a_k}).}
\]
Thus the limit on the right is the categorical limit in the
formal-moduli-stack category, not an ambient limit in
$(\mathcal X_{/S})_{X/}$.

In particular,
\[
    L_S^{\mathrm{Def}}Q^{\mathrm{nat}}_{F_{\mathfrak a\times\mathfrak b}}
    \simeq
    \operatorname{Mat}^{\mathrm{sp}}(F_{\mathfrak a})
    \mathop{\times}^{\operatorname{ModuliStk}^{\mathrm{sp}}(X/S)}
    \operatorname{Mat}^{\mathrm{sp}}(F_{\mathfrak b}),
\]
where the displayed notation means the binary product formed in
$\operatorname{ModuliStk}^{\mathrm{sp}}(X/S)$.  No raw finite-limit formula
for $\operatorname{Int}^{\mathrm{gpd}}$ is asserted.
\end{proposition}

\begin{proof}
The formal-moduli/partition-Lie equivalence is an equivalence of
$\infty$-categories and therefore preserves all existing limits; so does its
inverse formal integration.  The global materialization equivalence
\[
    \operatorname{FMP}^{\mathrm{sp}}_{X/S}
    \simeq
    \operatorname{ModuliStk}^{\mathrm{sp}}(X/S)
\]
likewise preserves categorical limits.  Hence
\[
    \operatorname{Mat}^{\mathrm{sp}}(F_{\mathfrak a})
    \simeq
    \operatorname*{lim}^{\operatorname{ModuliStk}^{\mathrm{sp}}(X/S)}_{k\in K}
    \operatorname{Mat}^{\mathrm{sp}}(F_{\mathfrak a_k}).
\]
Apply Theorem~\ref{thm:bnli-native-quotient-materialization} to the left-hand
side.  This gives the asserted equivalence.  No finite-limit preservation of
the raw native Artinian-cell colimit is used or implied.
\end{proof}

\begin{remark}[Ambient groupoid products remain separate]
Chapter~13 proves that the identity groupoid is initial and the maximal formal
relation $R_{X/S,\bullet}$ is terminal in the native fixed-object formal
groupoid category.  Thus an ambient raw product, when independently known to
compute a desired comparison, is a pullback over $R_{X/S,\bullet}$, not over
the identity groupoid.  The proposition above does not assert that such an
ambient product computes integration of a product algebroid.
\end{remark}

\section{Rational representation, PBW, and CE exchange}
\label{sec:bnli-rational-representation-exchange}

Assume the rational dg representation domain of the Atiyah--Spencer chapter.

\begin{theorem}[Rational stable representation materialization]
\label{thm:bnli-rational-rep-comparison}
Let $\mathfrak a$ be a rational spectral partition-Lie algebroid and
$\mathfrak a_{\mathrm{dg}}$ its dg rectification.  Under
Theorem~\ref{thm:bnli-representation-lie-three}, the stabilized formal-stack
representation category rectifies to the dg tangent formal-moduli
representation category of $\mathfrak a_{\mathrm{dg}}$.

On the ordinary PBW comparison domain, the tangent enveloping monad rectifies
to the ordinary bimodule enveloping-module monad.  On the CE representation
domain, the Artinian-formalized Spencer object and its stabilized
materialization rectify to the CE--Spencer resolution constructed in the
preceding chapter.

The source formalized stages entering
$\operatorname{Real}^{\mathrm{Cart}}_{\mathfrak a}$ therefore admit the stated
dg rectification.  The final pullback to
$Q^{\mathrm{nat}}_{F_{\mathfrak a}}$ and passage to the raw parameterized
Cartan category remain spectral native constructions.  No dg equivalence of
that raw target category, and no equivalence with all raw Cartan
coefficients, is asserted.
\end{theorem}

\begin{proof}
The stable formal representation materialization is obtained from the
stabilized slice of the same Artinian formal-moduli category which rationally
rectifies to the dg tangent formal-moduli category.  Slicing and stabilization
commute with the constructed rational equivalence, so the formalized
representation category and its tangent monad rectify on the stated dg
domain.  The PBW and CE assertions are precisely the separately scoped
comparison theorems of the Atiyah--Spencer chapter transported through this
equivalence.

Construction~\ref{def:bnli-rep-materialization} then applies exact spectral
functors: inclusion into the ambient stabilized slice, pullback along the
native-to-materialized quotient map, forgetting the distinguished centre map,
and effective Cartan descent.  The first formalized stages have rational dg
counterparts by the preceding paragraph.  The target
$Q^{\mathrm{nat}}_{F_{\mathfrak a}}$ and its parameterized stable category,
however, are raw native spectral objects; the preceding chapters construct no
rational rectification equivalence for that raw category.  The theorem
therefore stops at the formalized dg comparison and retains the final native
Cartan realization as a spectral construction.
\end{proof}

\section{Operator-valued Maurer--Cartan integration}
\label{sec:bnli-operator-mc}

\begin{theorem}[Integration of the operator-valued Maurer--Cartan system]
\label{thm:bnli-operator-mc}
For
\[
    \mathfrak a
    \in
    \operatorname{LieAlgd}^{\pi,E_\infty}_{X/S},
\]
the native integrated groupoid
\[
    \mathfrak G_{\mathfrak a}
    =
    \operatorname{Int}^{\mathrm{gpd}}(\mathfrak a)
\]
carries its complete Cartan transport.  The Artinian Spencer linearization of
that transport is the canonical tangent-monadic twisting system of
\(\mathfrak a\).

In the rational dg domain, this transport differentiates to the convolution
\(L_\infty\) Maurer--Cartan element constructed in the Atiyah--Spencer chapter.
The transferred infinite Maurer--Cartan series is formed in the same
canonically complete convolution mapping object; no pronilpotence condition is
introduced.
\end{theorem}

\begin{proof}
The native quotient of \(\mathfrak G_{\mathfrak a}\) is effective, so its Cech
relation carries the canonical Cartan descent cocycle.  Formal-leaf recovery
identifies its Artinian formalized controller with \(\mathfrak a\), and
Theorem~\ref{thm:bnli-spencer-invariance} identifies the normalized Artinian
Spencer system with the tangent-monadic system.  The rational conclusion is
then the rooted-tree/CE--Spencer comparison of the preceding chapter.  The
convolution target is complete because it is constructed there as the inverse
limit obtained by mapping out of the exhaustive one-marked bar.
\end{proof}

\section{Audit of constructions and theorem scopes}
\label{sec:bnli-audit}

The chapter contains three logically separate layers:
\begin{enumerate}
    \item spectral formal integration and formal-moduli materialization;
    \item native fixed-object reconstruction from induced Artinian cells, with
    one-sided differentiation recovery;
    \item comparisons with arbitrary pre-existing nonlinear groupoids and with
    rational or represented models.
\end{enumerate}
Their separation records theorem strength and prevents a property of a
canonical comparison from being turned into input structure.

\subsection{Unconditional constructions in the global spectral coefficient domain}

For every
\[
    \mathfrak a
    \in
    \operatorname{LieAlgd}^{\pi,E_\infty}_{X/S},
\]
the chapter constructs, without an additional integrability hypothesis:
\begin{enumerate}
    \item its spectral formal integration
    $F_{\mathfrak a}=\operatorname{Int}^{\mathrm{fm}}(\mathfrak a)$;
    \item spectral pro-reduction, the left-exact idempotent reflector
    $\mathsf P^{\mathrm{sp}}$, and the intrinsic infinitesimal prestacks
    $U^{\mathrm{inf,sp}}=\mathsf P^{\mathrm{sp}}U$;
    \item the left-exact global spectral deformation localization
    $L_S^{\mathrm{Def}}$;
    \item finite square-zero generation without a finite-presentation input,
    with coefficients chosen connective and eventually coconnective over the
    centre;
    \item spectral Artinian effectivity and Cech descent in the deformation-local
    $\infty$-topos;
    \item bounded coherent approximation of unrestricted square-zero
    coefficients, correctly varianced relative obstruction theory for backward
    Artinian edges, filtered contractible coherent-lift categories, finite coherent
    localization-resolution propagation, unrestricted square-zero excision, relative
    Postnikov reconstruction, explicit Artinian factorization categories,
    Cartesian colimit fibres, relative split-tangent continuity for nil-laft
    centres, the first nilpotent left Kan extension, and first-stage
    recognition;
    \item the bounded-test model of convergent prestacks, filtered representative
    categories, the pointed filtered-fibre and pointed pro-reduction fibre
    formulae, simultaneous finite-square representative refinement, bounded
    pro-pushout and pointed pro-pushout systems, relative Postnikov recovery
    and Postnikov comparison for nilpotent pullbacks, the second-stage pro-reduced
    dotted-lift materialization on unrestricted nilpotent neighbourhoods,
    pro-infinitesimal completion of arbitrary centred targets, and the bridge
    from ambient deformation locality plus spectral laftness to nil-laftness on bounded nilpotent diagrams;
    \item the raw split square-zero tangent spectrum constructed already for
    convergent nilpotent-excisive prestacks, its exact extension by
    stabilization, and the resulting spectral Nisnevich descent theorem for
    laft nil-isomorphic deformation prestacks, together with spectral-etale
    cocartesian Artinian restriction and the affine
    restriction/materialization equivalence;
    \item reduced-centred completion stability, the native relative reduced
    target system $W_{U/V}\to W_{X/S}$, representable open chart morphisms on
    finite affine open covers, formal-completion rather than raw
    target-base-change transition maps under coherent spectral-etale changes
    of centre, and global materialization by ordinary effective Cech/basis
    descent in the left-exact deformation-local subtopos of the
    non-hypercomplete native spectral Nisnevich topos in the locally coherent
    qcqs range;
    \item for every local Artinian test, its native effective image and formal
    Cech groupoid;
    \item the minimal affine native Artinian-cell quotient and its local raw
    mapping discrepancy system;
    \item the global Artinian Grothendieck category and the contravariant
    induced-cell functors
    \[
        c\longmapsto
        \overline I_c=X\amalg_{U_c}I_c,
        \qquad
        c\longmapsto
        \overline T_c=X\amalg_{U_c}T_c;
    \]
    \item the global native quotient
    \[
        X\twoheadrightarrow Q^{\mathrm{nat}}_{F_{\mathfrak a}};
    \]
    \item the native formal groupoid
    $\operatorname{Int}^{\mathrm{gpd}}(\mathfrak a)$;
    \item the deformation-formalization comparison
    \[
        L_S^{\mathrm{Def}}Q^{\mathrm{nat}}_{F_{\mathfrak a}}
        \simeq
        \operatorname{Mat}^{\mathrm{sp}}(F_{\mathfrak a});
    \]
    \item formal-leaf recovery and
    \[
        \operatorname{Diff}^{\pi,E_\infty}
        \operatorname{Int}^{\mathrm{gpd}}(\mathfrak a)
        \simeq
        \mathfrak a;
    \]
    \item fixed-object functoriality of native integration;
    \item formalized and coefficient-level coherent spectral-etale transport,
    without a raw changing-object atlas-invariance theorem;
    \item local and global native-cell mapping-evaluation maps, their canonical
    sections, and the corresponding mapping discrepancy formulae;
    \item the arbitrary-groupoid Lie-comparison span and the theorem identifying
    its two legs after the reconstructed-object equivalence;
    \item finite-limit stability of the relative formal-stack slice, stable
    formal representation materialization, and the contravariantly functorial
    native Cartan realization on the Artinian-formalized branch;
    \item actual-isotropy-to-anchor-fibre comparison and its discrepancy;
    \item native realization of the formalized adjoint tangent representation,
    together with the separate nonlinear stable formalized-adjoint-to-actual
    isotropy counit;
    \item one-object formal group integration and the independent
    Hopf-derived-primitives spectral Lie algebra attached to that group;
    \item materialized formal-spectrum construction and its conditional adic comparison,
    together with rational Artinian CE materialization, the correctly typed
    based-loop CE group, and the separately scoped classical BCH comparison;
    \item unconditional integration of action algebroids and, only when a
    nonlinear action has separately been constructed, its comparison span;
    \item gauge, classical, and pre-existing smooth comparison spans;
    \item deformation-formalized zero recovery and finite-limit comparison
    after deformation formalization.
\end{enumerate}

\subsection{The unresolved native rigidity is constructed as a discrepancy}

Affinely, for every Artinian $A\to B$ and formal-moduli problem $G$, the
chapter constructs
\[
    \operatorname{ev}^{\mathrm{nat,aff}}_{A,G}:
    \operatorname{Map}_{U/}(I_A,Q_G^{\mathrm{nat,aff}})
    \longrightarrow
    G(A)
\]
and its canonical section.  The affine mapping discrepancy theorem expresses
the fibre of
\[
    \operatorname{Map}
    (\mathfrak G_F^{\mathrm{nat,aff}},\mathfrak G_G^{\mathrm{nat,aff}})
    \longrightarrow
    \operatorname{Map}(F,G)
\]
as the limit of the corresponding cellwise fibres and gives the exact local
if-and-only-if criterion for full faithfulness.

Globally, for every object
$c=((U,V),A\to B)$ of the global Artinian test category and every global
formal-moduli problem $G$, the chapter constructs
\[
    \operatorname{ev}^{\mathrm{nat}}_{c,G}:
    \operatorname{Map}_{X/}(\overline I_c,Q_G^{\mathrm{nat}})
    \longrightarrow
    G_c(A)
\]
and its canonical section.  The global mapping discrepancy theorem expresses
the fibre of
\[
    \operatorname{Map}
    (\mathfrak G_F^{\mathrm{nat}},\mathfrak G_G^{\mathrm{nat}})
    \longrightarrow
    \operatorname{Map}(F,G)
\]
as the limit of these global cellwise discrepancy spaces.  Vanishing of all
such discrepancies therefore implies raw full faithfulness.  The chapter does
not impose this vanishing as an input condition.

No class of ``native-rigid'', ``Artinian-dense'', ``materializable'', or
``admissible'' groupoids is defined by imposing the conclusion of the
unresolved raw rigidity theorem.  The word ``reconstructed'' is used only for
objects already produced as $\operatorname{Int}^{\mathrm{gpd}}(\mathfrak a)$;
it is not a rigidity condition.

\subsection{Canonical comparisons retained without forced equivalence}

The chapter retains, as primary constructions:
\begin{enumerate}
    \item the two-legged quotient and relation comparison spans for an arbitrary
    native formal groupoid;
    \item stable-arrow and actual-isotropy discrepancy cofibres for the two
    legs;
    \item the support-marked Spencer comparison spans obtained by applying the
    preceding chapter's raw/formalized support-marked functors to those legs;
    \item the nonlinear formalized-adjoint-to-stable-actual-isotropy counit and its
    stable shadow;
    \item the Artinian-linearized formalized adjoint representation and its
    native Cartan realization, retained independently of the differentiated
    actual-isotropy controller;
    \item the actual-isotropy-to-homotopy-anchor-fibre comparison;
    \item local and global native cell mapping discrepancies;
    \item the parallel one-object partition-Lie controller and
    Hopf-derived-primitives spectral Lie algebra, without an unconstructed
    morphism between them;
    \item the canonical construction
    $\operatorname{Spf}^{\mathrm{mat}}(A^\wedge)=\operatorname{Mat}(F_{A^\wedge})$
    and its comparison with an independently constructed adic formal spectrum;
    \item on the rational finite-type CE range, the materialized based-loop CE
    formal group
    \[
        S\times_{\operatorname{Spf}^{\mathrm{CE,mat}}(\mathfrak g_{\mathrm{dg}})}S,
    \]
    together with a conditional comparison to any independently constructed
    adic CE formal spectrum.
\end{enumerate}
The $0$-truncated identity-groupoid calculation proves that actual formal
isotropy and the homotopy anchor fibre cannot be identified unconditionally.
The zero-controller calculation separately shows why deformation-formalized
zero recovery does not imply a raw identity-groupoid formula.

\subsection{No hidden input conditions}

None of the following is an input condition or retained auxiliary choice for a spectral partition-Lie algebroid or
of its native integration:
\begin{enumerate}
    \item a prorepresenting complete algebra or formal-spectrum presentation;
    \item an Artinian square-zero chain or chosen nilpotent filtration;
    \item a coherent approximation, obstruction representative, or chosen structured
    derivation factorization;
    \item an affine atlas or independently retained descent field;
    \item a source-simply-connected smooth integration or monodromy choice;
    \item a pronilpotence condition;
    \item a materializability, Artinian-density, native-rigidity,
    deformation-completeness, admissibility, isotropy-exactness,
    primitive-exactness, or representation-recognition structure;
    \item a connection, splitting, curvature, Bianchi, or flatness structure;
    \item a Hopf algebra presentation or an identification of Hopf-derived
    primitives with the partition-Lie controller;
    \item an identification of actual formal isotropy with the homotopy anchor
    fibre;
    \item an additional PBW, CE, finite-duality, BCH, or Sullivan input structure;
    \item a chosen finite-order jet tower;
    \item a chosen completion topology or adic generator;
    \item a chosen changing-object atlas for native integration;
    \item an arbitrary changing-object spectral partition-Lie pullback.
\end{enumerate}

The replacements used in the chapter are constructions: the left-exact
pro-reduction reflector $\mathsf P^{\mathrm{sp}}$; left-exact deformation
localization; finite square-zero generation; bounded coherent approximation
categories; correctly varianced relative obstruction maps; filtered
contractible coherent-lift categories; finite coherent localization resolution in groupoid
completions; explicit Artinian factorization Grothendieck constructions;
unrestricted first-stage square-zero excision; relative Postnikov
reconstruction; first-stage nilpotent left Kan extension; relative
split-tangent continuity for nil-laft centres; the bounded-test model of
convergent prestacks; filtered representative categories; pointed
filtered-fibre and pointed pro-reduction fibre formulae; simultaneous
nilpotent-square representative refinement; bounded pro-pushout and pointed
pro-pushout systems; Postnikov comparison for nilpotent pullbacks;
second-stage dotted-lift
materialization on unrestricted nilpotent neighbourhoods; pro-infinitesimal
completion of arbitrary targets; the bridge from spectral laftness to
nil-laftness on bounded nilpotent diagrams; the raw split square-zero tangent
spectrum and its stabilization extension; spectral Nisnevich descent;
reduced-centred completion stability and formal-completion transition maps;
the formal-completion Cartesian-section category; finite-subobject descent;
recovery of arbitrary coherent affine-etale pairs from finite reduced-open
bases; ordinary effective Cech/basis descent in the non-hypercomplete native
spectral Nisnevich topos; materialized formal spectra obtained from complete
almost-finite Koszul-dual functors; global Artinian restriction and its spectral-etale
cocartesianity; deformation-effective Artinian centres; native effective images;
local native Artinian cells; the global Artinian Grothendieck category;
functorial induction of local cells by pushout to the fixed object space $X$;
global native quotient colimits; reconstructed-span compatibility; stabilized
relative-slice materialization; native Cartan realization; groupoid
effectivity; unitwise formalization coreflection; functorial comparison maps;
and native mapping discrepancies.

\subsection{Genuine theorem domains}

The following remain domains of constructions rather than fields carried by an
object:
\begin{enumerate}
    \item coherence of spectral coefficient rings for pro-coherent
    partition-Lie theory;
    \item local coherence and the qcqs scope used for the coherent global
    spectral-etale coefficient system and global materialization;
    \item spectral etaleness for the presently constructed coefficient-level
    varying-object variance and formal-completion transport;
    \item representability when geometric quasi-coherent cotangent complexes,
    independently constructed adic formal spectra, principal parts, or
    ordinary tangent bundles are invoked;
    \item the complete almost-finite augmented sector when a complete algebra
    is converted to its materialized formal spectrum by Koszul duality;
    \item rationality and the boundedness/model range used for dg
    rectification and dg representation theory;
    \item the finite-type dg CE representability range when a literal complete
    Chevalley--Eilenberg algebra is used as a representing model;
    \item perfectness only where an ordinary symmetric-power dual formula is
    used in PBW comparison;
    \item finite-dimensionality and nilpotence where the classical BCH
    representation theorem is invoked;
    \item smooth representability for comparisons with ordinary smooth Lie
    groupoids.
\end{enumerate}

\subsection{Still not asserted}

The chapter does not assert:
\begin{enumerate}
    \item raw full faithfulness of
    $\operatorname{Int}^{\mathrm{gpd}}$ before the native cell rigidity theorem
    is proved;
    \item an equivalence between a full raw subgroupoid category and
    $\operatorname{LieAlgd}^{\pi,E_\infty}_{X/S}$;
    \item an adjunction
    $\operatorname{Diff}\dashv\operatorname{Int}^{\mathrm{gpd}}$
    on the entire unrestricted native formal-groupoid category;
    \item a direct map
    $\mathfrak G\to\operatorname{Int}^{\mathrm{gpd}}
      \operatorname{Diff}(\mathfrak G)$
    for an arbitrary native formal groupoid;
    \item an equivalence between all raw Cartan coefficients on an integrated
    groupoid and tangent representations of its controller;
    \item a geometric dg materialization functor on arbitrary materialized
    spectral prestacks, or a rational rectification equivalence for the raw
    native parameterized Cartan target category;
    \item preservation of the raw represented first-principal-parts object by
    native Cartan realization;
    \item a canonical morphism from the native Cartan realization of the
    formalized adjoint tangent representation to raw stable actual isotropy;
    \item raw finite-limit preservation by native integration;
    \item a raw equivalence
    $\operatorname{Int}^{\mathrm{gpd}}(0_{X/S})\simeq0_{X,\bullet}$;
    \item a raw changing-object spectral-etale exchange theorem for native
    integration, or Morita/atlas invariance of fixed-object formality;
    \item arbitrary global smooth Lie-groupoid integrability of every ordinary
    Lie algebroid;
    \item arbitrary changing-object functoriality beyond the constructed
    coefficient-level spectral-etale variance;
    \item an integral spectral-to-animated Artinian equivalence;
    \item any controller-level or aritywise comparison morphism between
    $\mathfrak l^{\mathrm{sp}}$ and the spectral partition-Lie controller;
    \item unconditional identification of actual formal isotropy with the
    homotopy anchor fibre;
    \item unconditional representability of the raw native quotient by a
    formal spectral scheme;
    \item raw equivalence of the native one-object group with the represented
    CE or BCH group without native rigidity;
    \item integration of an arbitrary infinitesimal action algebroid to an
    action presentation by $\mu^{\mathrm{gpd}}_{\mathfrak g}$ unless the
    nonlinear action has separately been constructed;
    \item a Sullivan comparison before a Sullivan realization functor has been
    constructed.
\end{enumerate}

\section{Brave-New Lie Integration Theorem Summary}
\label{sec:bnli-package}

\begin{theorem}[Brave-New Lie Integration Summary]
\label{thm:bnli-package}
Let
\[
    \mathfrak a
    \in
    \operatorname{LieAlgd}^{\pi,E_\infty}_{X/S}.
\]
The constructions of this chapter canonically provide:
\begin{enumerate}
    \item the spectral formal leaf
    $F_{\mathfrak a}=\operatorname{Int}^{\mathrm{fm}}(\mathfrak a)$;
    \item the left-exact pro-reduction reflector
    $\mathsf P^{\mathrm{sp}}$, intrinsic infinitesimal prestacks, and the
    presentation-free formal-moduli-stack materialization
    $\operatorname{Mat}^{\mathrm{sp}}(F_{\mathfrak a})$;
    \item left exactness of $L_S^{\mathrm{Def}}$, finite square-zero
    generation without finite-presentation input with bounded coefficients,
    bounded coherent approximation, the relative formal-lifting obstruction
    $L_{A_1/A}\to M[1]$, filtered contractible coherent obstruction
    representatives, finite coherent localization-resolution propagation, unrestricted square-zero
    excision for the first Kan extension, relative Postnikov reconstruction,
    Cartesian colimit fibres, relative split-tangent continuity, and Artinian
    Cech effectivity in the deformation-local topos;
    \item explicit first-stage Artinian factorization categories, the
    bounded-test model for convergent prestacks, filtered representative
    categories, pointed filtered-fibre and pro-reduction fibre formulae,
    simultaneous finite-square representative refinement, bounded
    pro-pushouts and pointed pro-pushouts, Postnikov extension of nilpotent
    excision, relative
    Postnikov recovery of unrestricted nilpotent neighbourhoods, the
    ambient-deformation-local spectral-laft-to-nil-laft bridge, pro-infinitesimal completion, affine
    two-stage materialization, the stabilized raw tangent-spectrum proof of
    spectral Nisnevich descent, spectral-etale cocartesian Artinian
    restriction, the relative reduced chart system, finite-subobject descent, recovery of every
    coherent affine-etale pair from a finite reduced-open basis, and ordinary
    effective Cech/basis global materialization in the left-exact deformation-local
    subtopos of the non-hypercomplete native topos with reduced-centred
    formal-completion transition maps;
    \item the native Artinian effective-image cells $I_A$ and their formal Cech
    groupoids $\mathfrak G_A$;
    \item the minimal affine native quotient and its local mapping discrepancy
    theorem;
    \item the global Artinian Grothendieck category and the functorial induced
    fixed-object cells
    \[
        \overline I_c=X\amalg_{U_c}I_c,
        \qquad
        \overline T_c=X\amalg_{U_c}T_c;
    \]
    \item the canonical global native quotient
    \[
        X\twoheadrightarrow Q^{\mathrm{nat}}_{\mathfrak a};
    \]
    \item the native formal groupoid integration
    \[
\boxed{
        \operatorname{Int}^{\mathrm{gpd}}(\mathfrak a)
        =
        \check C_\bullet
        \left(
          X\twoheadrightarrow Q^{\mathrm{nat}}_{\mathfrak a}
        \right);}
    \]
    \item the deformation-formalization equivalence
    \[
        L_S^{\mathrm{Def}}Q^{\mathrm{nat}}_{\mathfrak a}
        \xrightarrow{\sim}
        \operatorname{Mat}^{\mathrm{sp}}(F_{\mathfrak a});
    \]
    \item the formal-leaf recovery equivalence
    \[
        \mathcal L^{\mathrm{fm,sp}}_{
        \operatorname{Int}^{\mathrm{gpd}}(\mathfrak a)}
        \simeq
        F_{\mathfrak a};
    \]
    \item the differentiation-after-integration equivalence
    \[
        \operatorname{Diff}^{\pi,E_\infty}
        \operatorname{Int}^{\mathrm{gpd}}(\mathfrak a)
        \xrightarrow{\sim}
        \mathfrak a;
    \]
    \item local and global native cell evaluation maps, their canonical
    sections, and mapping discrepancy formulae isolating the remaining raw
    rigidity problem;
    \item for an arbitrary native formal groupoid, the canonical Lie-comparison
    span
    \[
        \mathfrak G
        \longrightarrow
        \mathfrak K^{\mathrm{Lie}}_{\mathfrak G}
        \longleftarrow
        \mathfrak G^{\mathrm{Lie}},
    \]
    together with the theorem that its two legs coincide after the canonical
    reconstructed-object equivalence when $\mathfrak G$ is reconstructed;
    \item quotient, relation, stable, actual-isotropy, adjoint,
    support-marked, and native mapping discrepancy systems;
    \item invariance of the Artinian Spencer/partition-Lie operation system
    under reconstructed integration, together with functorial recovery of the
    formalized Spencer representation cochains;
    \item finite-limit stability of the relative formal-stack slice, the
    stable formal representation materialization equivalence, and the
    contravariantly functorial native Cartan realization
    \[
        \operatorname{Rep}^{\pi,E_\infty}(\mathfrak a)
        \longrightarrow
        \operatorname{Rep}^{\mathrm{Cart}}
        (\operatorname{Int}^{\mathrm{gpd}}(\mathfrak a));
    \]
    \item the integrated actual-isotropy-to-anchor-fibre comparison and its
    discrepancy;
    \item native realization of the formalized adjoint tangent representation
    and, separately, the nonlinear stable formalized-adjoint-to-stable-actual-isotropy
    counit;
    \item the one-object formal group
    $\mu^{\mathrm{gpd}}_{\mathfrak g}$, its differentiation recovery, and the
    independent Hopf-derived-primitives spectral Lie algebra;
    \item materialized formal-spectrum construction, its conditional adic comparison,
    rational Artinian CE materialization, and the based-loop CE/BCH comparison
    on the separately stated classical range, without a raw native/represented
    identification;
    \item unconditional native integration of action algebroids and conditional
    comparison with an action-groupoid presentation when the nonlinear action
    is separately available;
    \item gauge and Atiyah comparison spans;
    \item canonical formal integration of ordinary smooth characteristic-zero
    Lie algebroids in the rational coefficient range;
    \item comparison spans with pre-existing smooth integrations after
    unitwise formal completion;
    \item fixed-object functoriality, coefficient-level formalized
    spectral-etale transport, deformation-formalized zero recovery, and
    finite-limit comparison after deformation formalization.
\end{enumerate}

Every item is obtained from an equivalence, left-exact localization,
structured factorization, bounded coherent approximation, relative
split-tangent continuity, Artinian factorization Grothendieck construction,
bounded-test right Kan extension, bounded pro-pushout or pointed pro-pushout,
relative Postnikov reconstruction, effective-image factorization, Cech nerve,
category-of-elements colimit, pushout to the fixed object space, native
spectral-Nisnevich descent, ordinary effective Cech/basis descent, stabilized
relative-slice materialization, the exact forgetful functor to parameterized
Cartan coefficients, formal-completion descent, coreflection, or
comparison morphism constructed in this chapter or the preceding three
chapters.  No global smooth integrability, prorepresenting
algebra, connection, splitting, chosen Artinian presentation, admissibility
structure, native-rigidity condition, or Hopf presentation is inserted into the
input.
\end{theorem}

\begin{proof}
Formal integration is Theorem~\ref{thm:bnli-formal-lie-three}.  Spectral
pro-reduction and its localization theorem are
Theorem~\ref{thm:bnli-prored-localization}; left exactness of deformation
localization and Artinian Cech effectivity are
Theorem~\ref{thm:bnli-def-left-exact} and
Theorem~\ref{thm:bnli-artinian-cech-descent}.  Unrestricted finite
square-zero generation is
Proposition~\ref{prop:bnli-unrestricted-finite-square-zero}.  The corrected
backward-edge construction and unrestricted square-zero excision are
Lemma~\ref{lem:bnli-coherent-backward-lift} and
Theorem~\ref{thm:bnli-unrestricted-square-zero-kan}, with finite coherent localization
resolution supplied by Lemma~\ref{lem:bnli-finite-zigzag-resolution}.
Relative Postnikov reconstruction is
Lemma~\ref{lem:bnli-relative-postnikov-reconstruction}; bounded coherent
approximation and relative split-tangent continuity are
Lemma~\ref{lem:bnli-bounded-coherent-approximation} and
Lemma~\ref{lem:bnli-relative-split-tangent-continuity}.  Together with
Theorem~\ref{thm:bnli-first-stage-recognition} these give first-stage
recognition.

The bounded-test model for convergent prestacks is
Proposition~\ref{prop:bnli-convergent-bounded-tests}.  Filtered
representatives, pointed filtered fibres, pointed pro-reduction fibres, and
simultaneous representatives for bounded nilpotent squares are
Lemmas~\ref{lem:bnli-filtered-representatives},
~\ref{lem:bnli-pointed-filtered-fibre},
~\ref{lem:bnli-pointed-prored-fibre}, and
~\ref{lem:bnli-simultaneous-square-representatives}.  Bounded pro-pushouts
and pointed pro-pushouts are Definitions~\ref{def:bnli-pro-pushout}
and~\ref{def:bnli-pointed-pro-pushout}, and their evaluation formula is
Proposition~\ref{prop:bnli-dotted-lift-formula}.  Extension from bounded
nilpotent excision to arbitrary connective tests is
Lemma~\ref{lem:bnli-convergent-excision-extension}, based on
Lemma~\ref{lem:bnli-postnikov-nilpotent-pullback}.  The second-stage
equivalence and pro-infinitesimal completion are
Theorem~\ref{thm:bnli-second-stage-equivalence} and
Proposition~\ref{prop:bnli-pro-infinitesimal-universal}; the
ambient-deformation-local spectral-laft-to-nil-laft bridge is
Proposition~\ref{prop:bnli-spectral-laft-implies-nil-laft}, applied on the ambient deformation-local locus.  Spectral
Nisnevich
descent is Theorem~\ref{thm:bnli-formal-stack-nisnevich}, based on the raw
tangent spectrum and its stabilization extension in
Lemma~\ref{lem:bnli-raw-square-zero-tangent}.  Spectral-etale cocartesianity
of centred Artinian restriction is
Proposition~\ref{prop:bnli-artinian-leaf-etale}.  Affine and global
materialization are
Theorems~\ref{thm:bnli-affine-materialization-equivalence}
and~\ref{thm:bnli-global-materialization}; reduced-centred completion
stability, finite-subobject descent, finite-open locality of laftness, basis
recovery, and ordinary effective Cech/basis descent are
Proposition~\ref{prop:bnli-completion-stability},
Lemmas~\ref{lem:bnli-finite-subobject-descent},
~\ref{lem:bnli-finite-open-laft}, and~\ref{lem:bnli-basis-recovery}, and
Proposition~\ref{prop:bnli-formal-completion-descent}.

Native Artinian cells are
Proposition~\ref{prop:bnli-artinian-cell-formal}, and their global
functoriality is Proposition~\ref{prop:bnli-induced-cell-functoriality}.
The global native quotient and its comparison with materialization are
Theorems~\ref{thm:bnli-native-quotient-formal} and
Theorem~\ref{thm:bnli-native-quotient-materialization}.  Formal-leaf recovery
and differentiation-after-integration are
Theorem~\ref{thm:bnli-leaf-recovery} and
Corollary~\ref{cor:bnli-diff-int-identity}.

The local and global mapping discrepancies are
Theorems~\ref{thm:bnli-affine-mapping-discrepancy} and
Theorem~\ref{thm:bnli-mapping-discrepancy}.  The arbitrary-groupoid
comparison and its reconstructed compatibility are
Theorem~\ref{thm:bnli-completion-map} and
Proposition~\ref{prop:bnli-reconstructed-span-compatibility}.  Operation and
representation recovery are
Theorem~\ref{thm:bnli-spencer-invariance},
Theorem~\ref{thm:bnli-representation-lie-three},
Proposition~\ref{prop:bnli-relative-formal-stack-slice}, and
Proposition~\ref{prop:bnli-native-cartan-realization}.  The isotropy and
adjoint statements are
Theorem~\ref{thm:bnli-integrated-isotropy-comparison},
Theorem~\ref{thm:bnli-adjoint-materialization}, and
Proposition~\ref{prop:bnli-adjoint-stable-counit}.  The remaining items are
the corresponding one-object, represented, rational, action, gauge,
classical, fixed-object functoriality, zero-recovery, and
formation-after-localization results of the later sections.
\end{proof}

\section{Final synthesis}
\label{sec:bnli-final-synthesis}

The four brave-new Lie chapters now form the sequence
\[
\boxed{
\text{formal groupoids}
\longrightarrow
\text{spectral differentiation}
\longrightarrow
\text{Atiyah--Spencer operations}
\longrightarrow
\text{native integration and comparison}.}
\]

The nonlinear category contains fixed-object geometric information which need
not be detected by spectral Artinian deformation theory.  Spectral
differentiation extracts the global formal-moduli controller, while the
Atiyah--Spencer chapter identifies that controller with the independently
constructed Artinian cotangent-primitive Spencer operation system.  The
present chapter solves the reverse geometric problem without inserting a
presentation or integrability condition into the algebroid.

The first new ingredient is the intrinsic spectral pro-reduction reflector
\[
    \mathsf P^{\mathrm{sp}}.
\]
For a spectral prestack $Y$ it is defined by
\[
    (\mathsf P^{\mathrm{sp}}Y)(C)
    =
    \operatorname*{colim}_{I\subseteq\pi_0C,\ I\ {\rm nil}}
    Y(H(\pi_0C/I)).
\]
Filtered colimits of spaces commute with finite limits, and a cofinality
calculation on pairs of nested nilpotent ideals proves
\[
\boxed{
    (\mathsf P^{\mathrm{sp}})^2
    \simeq
    \mathsf P^{\mathrm{sp}}}
\]
and left exactness.  Spectral nil-isomorphisms are therefore precisely the
$\mathsf P^{\mathrm{sp}}$-equivalences.  In particular
\[
    U^{\mathrm{inf,sp}}
    =
    \mathsf P^{\mathrm{sp}}U
\]
is not merely notation for a chosen pro-system: it is the value of an
idempotent prestack localization, and a nil-isomorphism $U\to Y$ determines a
contractibly unique map
\[
    Y\longrightarrow U^{\mathrm{inf,sp}}
\]
under $U$.

Independently, the global deformation reflector
\[
    L_S^{\mathrm{Def}}
\]
is proved left exact.  Its local objects consequently form an
$\infty$-topos.  The proof checks pullback stability of the Postnikov and
nilpotent-excision generators, first on affine pullbacks and then on arbitrary
pullbacks by affine density and universality of colimits.

The finite nilpotent generation theorem is strengthened beyond the Artinian
finite-presentation range.  If
\[
    A\longrightarrow B
\]
is connective, surjective on $\pi_0$ with nilpotent kernel, and has eventually
coconnective fibre, then it admits a finite structured square-zero
factorization with coefficients restricted from connective eventually
coconnective $B$-modules.  The
proof separates the discrete nilpotent tower and then kills the positive
relative Postnikov groups one at a time.  The key boundedness input is the
cofiber of
\[
    A\longrightarrow
    H\pi_0(A)\times_{H\pi_0(B)}B,
\]
whose upper Postnikov bound follows from fibre transitivity.  Relative spectral
Hurewicz identifies the first unresolved cofiber group with the first
cotangent group, and octahedrality gives
\[
    C_{i+1}
    \simeq
    \tau_{\geq i+1}C_i.
\]
No finite-presentation hypothesis is used in this Postnikov-killing argument.

Left exactness then makes every structured square-zero centre an effective
epimorphism in the deformation-local topos.  Finite square-zero generation
gives, for every spectral Artinian arrow $A\to B$,
\[
    \operatorname{Spec}B
    \twoheadrightarrow
    \operatorname{Spec}A
\]
after deformation localization, and therefore
\[
\boxed{
    \operatorname{Map}(\operatorname{Spec}A,Z)
    \simeq
    \operatorname*{Tot}_{[n]\in\Delta}
    \operatorname{Map}
    \left(
      \operatorname{Spec}(B^{\otimes_A(n+1)}),
      Z
    \right)}
\]
for every deformation-local target $Z$.  The resulting Cech descent is
geometric and requires no separate adic bar-length convergence argument.

Affine materialization is then constructed in two stages.  The first stage is
left Kan extension from spectral Artinian neighbourhoods to all nilpotent
neighbourhoods whose relative fibre over the coherent centre is eventually
coconnective.  The square-zero proof uses the fixed-extension obstruction in
its correct variance.  For a backward Artinian edge
\[
    A\longrightarrow A_1\longrightarrow C_0
\]
against a square-zero extension $C\to C_0$ with coefficient $M$, the
obstruction is
\[
\boxed{
    L_{A_1/A}\otimes_{A_1}C_0
    \longrightarrow L_{C_0/A}
    \longrightarrow M[1],}
\]
equivalently a map $L_{A_1/A}\to\operatorname{Res}_{A_1}^{C_0}M[1]$.
Almost finite presentation of $A\to A_1$ makes this cotangent almost perfect;
bounded coherent approximation then produces a filtered contractible
category of coherent obstruction representatives.  Each representative gives
an actual coherent Artinian square-zero refinement of the backward edge.
Finite coherent localization resolution in the groupoid completion of the Artinian
factorization category propagates these refinements through an arbitrary path
of Artinian representatives.  This proves unrestricted square-zero excision
directly, without an auxiliary density hypothesis, and hence gives
deformation theory for the first left Kan extension.

A Cartesian-colimit fibre lemma supplies the missing global fibre calculation
for the first-stage left Kan extension: a cartesian natural transformation of
space-valued diagrams remains cartesian after passage from any chosen cell to
the two colimits.  Applied to the Artinian identity cell, it turns the
cellwise almost-finite-presentation calculation into the nil-laft square for
the centre.

The converse recognition argument uses a different mechanism.  For a
nil-laft centre $U\to G$, the relative split tangent spectrum
\[
    \Theta_{U/G}^{\mathrm{sp}}(M)
    =
    \operatorname{fib}
    \left(
      \Theta_U^{\mathrm{sp}}(M)
      \to
      \Theta_G^{\mathrm{sp}}(M)
    \right)
\]
commutes with bounded coherent approximation of every connective eventually
coconnective $B$-module $M$.  The Artinian counit is an equivalence on the
coherent split coefficients; relative tangent continuity and the common
ambient tangent therefore make it an equivalence on every bounded split
coefficient.  Split coefficient base change propagates this comparison over
an arbitrary intermediate nilpotent neighbourhood, and the bounded
finite-square-zero factorization theorem propagates it through every
structured stage.  This proves the recognition equivalence
\[
\operatorname{FMP}^{\mathrm{sp}}_{B/R}
\simeq
\left\{
\begin{array}{c}
\text{nilpotent-neighbourhood functors with deformation theory}\\
\text{and nil-laft centre}
\end{array}
\right\}
\]
without an object-level Artinian-density hypothesis.

The second stage is performed in the category of convergent prestacks.  This
category is itself constructed rather than assumed: restriction to eventually
coconnective affine tests has fully faithful right Kan extension
\[
    H\longmapsto
    \left(
      C\longmapsto
      \lim_nH(\tau_{\leq n}C)
    \right),
\]
so convergent prestacks are equivalent to presheaves on bounded tests and
their colimits are computed pointwise there.  A bounded point
\[
    x:T\longrightarrow U^{\mathrm{inf,sp}}
\]
has a filtered, hence weakly contractible, category of representatives on
nilpotent discrete quotients of $\pi_0T$; finite coherence diagrams admit a
common refinement.  A pointed filtered-fibre theorem identifies colimits of
the varying stagewise fibres with the fibre over the resulting pro-point, and
a flag-category refinement gives the pointed pro-reduction fibre formula.
For a finite nilpotent square, the category of simultaneous four-vertex
representatives is filtered and cofinal in every individual representative
system.  Pushing out each representative against the common centre $U$ gives
the bounded pointed pro-pushout system $P_x(T)$.  The
unpointed pro-pushout notation is likewise used only when the affine test
itself is eventually coconnective; no generic boundedness is asserted for the
pullback
$B\times_{H(\pi_0C/I)}C$
when $C$ is unbounded.  The dotted-lift adjunction satisfies
\[
\boxed{
    \operatorname{Map}_{/U^{\mathrm{inf,sp}}}
    (T,\widetilde v_!G)_x
    \simeq
    G(P_x(T)).}
\]
For a bounded nilpotent neighbourhood this specializes to
\[
    G(U\amalg_{T^{\mathrm{prored}}}T).
\]
Bounded nilpotent excision is promoted to arbitrary connective tests by
comparing the actual pullback with the pullbacks of the termwise Postnikov
truncations.  The unit at an unrestricted nilpotent neighbourhood is proved
separately by the relative Postnikov theorem: if
$K=\operatorname{fib}(C\to B)$ is $(n-1)$-coconnective, then
\[
    K
    \simeq
    \operatorname{fib}(\tau_{\leq n}C\to\tau_{\leq n}B)
\]
by the long exact homotopy sequence, and hence
\[
    T\simeq U\amalg_{U_n}T_n.
\]
No exactness of $\tau_{\leq n}$ is asserted and no absolute coconnective
bound on the coherent centre is introduced.  The
counit is checked on bounded pointed tests and then extended by convergence.
Thus the unit and counit are equivalences exactly on the nil-isomorphic
deformation diagrams.

For an arbitrary centred target $U\to Y$, the construction
\[
    Y_U^{\mathrm{pinf}}
    :=
    Y\times_{\mathsf P^{\mathrm{sp}}Y}
    \mathsf P^{\mathrm{sp}}U
\]
supplies the missing universal target for the affine materialization
adjunction.  Any nil-isomorphic source under $U$ maps to
$Y_U^{\mathrm{pinf}}$ exactly as it maps to $Y$.  Consequently
\[
\boxed{
    \operatorname{Mat}^{\mathrm{sp}}_{B/R}
    \dashv
    \operatorname{Leaf}^{\mathrm{Art}}_{B/R}}
\]
on all centred deformation-local targets.  A separate relative Postnikov
argument proves that a spectrally laft deformation-local centre restricts to a
nil-laft centre on every uniformly bounded nilpotent diagram.  The first-stage
recognition theorem is therefore applicable to every object on the claimed
essential-image locus, and that locus consists precisely of the spectrally
laft nil-centred objects.

The fact that these materializations are native spectral Nisnevich sheaves is
also derived rather than assumed.  Before any sheaf condition is known, a
convergent nilpotent-excisive prestack $Y$ has split square-zero lifting fibres
which over an arbitrary connective centre assemble into
an exact spectrum-valued functor
\[
    \Theta^{\mathrm{sp}}_{Y,D,x}:
    \operatorname{Mod}_D
    \longrightarrow
    \operatorname{Sp},
\]
before any coherent or pro-coherent approximation.  Nilpotent excision first
makes the connective lifting functor reduced and excisive; the universal
property of
\[
    \operatorname{Sp}(\operatorname{Mod}_{D,\geq0})
    \simeq
    \operatorname{Mod}_D
\]
then gives its unique exact extension to all $D$-modules.  Spectral-etale
restriction-of-scalars compatibility extends through the same stabilization.
After a finite square-zero factorization of a bounded affine test, ordinary
quasi-coherent Nisnevich descent for the coefficient modules and exactness of
$\Theta^{\mathrm{sp}}$ inductively prove the fibrewise Nisnevich square.
Thus
\[
\boxed{
\text{laft}+\text{nil-isomorphism}+\text{convergence}+\text{nilpotent excision}
\Longrightarrow
\text{spectral Nisnevich descent}.}
\]
Once the prestack is a native sheaf, the already verified convergence and
nilpotent excision identify it as $L_S^{\mathrm{Def}}$-local.  This ordering
removes any circular use of deformation locality in the proof of sheafness.

Reduced-centred completion
\[
    \widehat Y^{\,\mathrm{sp}}_U
    =
    Y\times_{D_{\mathrm{red}}Y}D_{\mathrm{red}}U
\]
is shown to preserve deformation locality and laftness and to leave Artinian
leaves unchanged.  The correct global base for these completions is the
relative reduced target
\[
    W_{X/S}
    =
    S\times_{D_{\mathrm{red}}S}D_{\mathrm{red}}X.
\]
For every coherent affine-etale pair $(U,V)$ the relative reduced chart
\[
    W_{U/V}
    =
    V\times_{D_{\mathrm{red}}V}D_{\mathrm{red}}U
\]
satisfies the Cartesian reconstruction
\[
    U
    \simeq
    X\times_{W_{X/S}}W_{U/V}.
\]
Moreover formal-completion transport becomes ordinary pullback on these
charts:
\[
    \widehat{(Y\times_VV')}^{\,\mathrm{sp}}_{U'}
    \simeq
    Y\times_{W_{U/V}}W_{U'/V'}.
\]

Global materialization is therefore obtained from the constructed category of
formal-completion Cartesian sections by ordinary effective Cech descent along
a finite family of representable quasi-compact open morphisms
$W_{U/V}\to W_{X/S}$; the native targets $W_{U/V}$ themselves are not
asserted to be represented.  Finite-subobject descent avoids commuting a
filtered colimit with an infinite Cech totalization, and the basis-recovery
theorem proves that every coherent affine-etale pair is recovered from the
finite reduced-open basis.  The descent is performed
inside the left-exact deformation-local subtopos
$\mathcal X^{\mathrm{Def}}_{/S}$, so deformation locality of the descended
target is automatic rather than re-proved after gluing.  Finite-open locality
then globalizes spectral laftness.  Reduced-classical descent identifies the
global centre as a reduced-classical equivalence, and reduced-centre detection
upgrades it to a spectral nil-isomorphism.  This gives
\[
\boxed{
    \operatorname{ModuliStk}^{\mathrm{sp}}(X/S)
    \simeq
    \operatorname{FMP}^{\mathrm{sp}}_{X/S}.}
\]
The transition law remains formal completion along the new centre; it becomes
ordinary Cartesian transport only on the relative reduced chart system.

The native groupoid is constructed by a second realization which is kept
separate from materialization.  Every local Artinian test
\[
    c=((U,V),A\to B)
\]
has an effective-image factorization
\[
    U\twoheadrightarrow I_c\hookrightarrow T_c,
    \qquad T_c=\operatorname{Spec}A.
\]
Artinian Cech effectivity gives
\[
    L_S^{\mathrm{Def}}I_c
    \simeq
    T_c,
\]
while de Rham localization shows that $U\to I_c$ is already a formal
effective quotient.  Coherent Artinian spectral-etale transport makes
\[
    \overline I_c
    =
    X\amalg_UI_c,
    \qquad
    \overline T_c
    =
    X\amalg_UT_c
\]
contravariantly functorial on the global Artinian test category.  For a global
formal-moduli problem $F$, the category-of-elements colimit
\[
\boxed{
    Q_F^{\mathrm{nat}}
    :=
    \operatorname*{colim}_{(c,\xi)\in(\int_XF)^{\mathrm{op}}}
    \overline I_c}
\]
is an effective formal quotient of the fixed object space $X$.

The materialized and native realizations meet only after a theorem:
\[
\boxed{
    L_S^{\mathrm{Def}}Q_F^{\mathrm{nat}}
    \xrightarrow{\sim}
    \operatorname{Mat}^{\mathrm{sp}}_{X/S}(F).}
\]
The proof uses the termwise deformation equivalences
$\overline I_c\to\overline T_c$, the global Artinian co-Yoneda formula, and
the materialization adjunction.  Therefore, for
\[
    F_{\mathfrak a}
    =
    \operatorname{Int}^{\mathrm{fm}}(\mathfrak a),
\]
one may define
\[
\boxed{
    \operatorname{Int}^{\mathrm{gpd}}(\mathfrak a)
    :=
    \check C_\bullet
    \left(
      X\twoheadrightarrow Q_{F_{\mathfrak a}}^{\mathrm{nat}}
    \right),}
\]
and obtain
\[
\boxed{
\begin{aligned}
    \mathcal L^{\mathrm{fm,sp}}_{
      \operatorname{Int}^{\mathrm{gpd}}(\mathfrak a)}
    &\simeq
    F_{\mathfrak a},\\
    \operatorname{Diff}^{\pi,E_\infty}
    \operatorname{Int}^{\mathrm{gpd}}(\mathfrak a)
    &\simeq
    \mathfrak a.
\end{aligned}}
\]
This is the Native Brave-New Lie III theorem proved in the chapter.

The theorem is deliberately one-sided at the raw nonlinear level.  The local
and global native Artinian cells carry evaluation maps with canonical
sections, and the fibre of a raw integrated-groupoid mapping space over a
formal-moduli morphism is the limit of the corresponding cellwise discrepancy
spaces.  Raw full faithfulness is therefore isolated as a concrete native
rigidity theorem rather than imposed as an admissibility condition.

For an arbitrary formal groupoid $\mathfrak G$, the variance of the
materialization counit gives the canonical comparison span
\[
\boxed{
    \mathfrak G
    \longrightarrow
    \mathfrak K^{\mathrm{Lie}}_{\mathfrak G}
    \longleftarrow
    \operatorname{Int}^{\mathrm{gpd}}
    \operatorname{Diff}^{\pi,E_\infty}(\mathfrak G).}
\]
For a reconstructed groupoid, the two outer sources are canonically
equivalent.  The reconstructed-span compatibility theorem further proves that
the two arrows into the middle comparison object agree after that
identification; consequently all functorial stable, isotropy, adjoint, and
support-marked comparison shadows agree on the two legs.  Neither leg is
thereby forced to be an equivalence.

On the representation side, the relevant relative slice
\[
    \operatorname{ModuliStk}^{\mathrm{sp}}(X/S)_{/Y}
\]
is proved finite-limit closed in the ambient factorization slice.  Hence its
stabilization maps canonically to the ambient stabilized slice.  Pullback along
\[
    Q^{\mathrm{nat}}_{F_{\mathfrak a}}
    \longrightarrow
    \operatorname{Mat}^{\mathrm{sp}}(F_{\mathfrak a})
\]
is followed by the exact stabilized forgetful functor from the factorization
slice under $X$ to the parameterized stable category
$\mathcal E_{Q^{\mathrm{nat}}_{F_{\mathfrak a}}}$; effective Cartan descent
then gives a correctly typed functor
\[
    \operatorname{Real}^{\mathrm{Cart}}_{\mathfrak a}:
    \operatorname{Rep}^{\pi,E_\infty}(\mathfrak a)
    \longrightarrow
    \operatorname{Rep}^{\mathrm{Cart}}
    (\operatorname{Int}^{\mathrm{gpd}}(\mathfrak a)).
\]
Its functoriality is contravariant on representation categories, as required
by pullback.

The operation-theoretic conclusions retain the separations established in the
preceding chapter.  Reconstructed integration preserves the Artinian
Spencer/partition-Lie controller, its complete PD operation monad, tangent
enveloping monad, and the formalized Spencer representation cochains.  The
raw support-marked geometric branch remains independent.  Actual formal
isotropy maps canonically to the homotopy anchor fibre but is not identified
with it; the $0$-truncated identity-groupoid calculation shows that such an
identification is false in general.  The formalized adjoint tangent
representation has a native Cartan realization, while the nonlinear stable
formalized-adjoint-to-actual-isotropy counit remains a separate comparison.

In the one-object case every spectral partition-Lie algebra $\mathfrak g$
therefore has a canonical brave-new formal Lie group
\[
    \mu^{\mathrm{gpd}}_{\mathfrak g}
\]
with
\[
    \operatorname{Diff}^{\pi,E_\infty}
    (B_\bullet\mu^{\mathrm{gpd}}_{\mathfrak g})
    \simeq
    \mathfrak g.
\]
The same nonlinear group independently determines the Hopf-derived-primitives
spectral Lie algebra
$\mathfrak l^{\mathrm{sp}}_{\mathfrak g}$.  No controller-level or aritywise
comparison map between these two outputs is manufactured.  On the rational finite-type CE range, rational Artinian rectification first
identifies the formal-moduli functor controlled by
$\mathfrak g_{\mathrm{dg}}$.  Materialization then constructs the classifying
target
\[
    \operatorname{Spf}^{\mathrm{CE,mat}}(\mathfrak g_{\mathrm{dg}}),
\]
and the correctly typed CE formal group is its based loop
\[
    S\times_{\operatorname{Spf}^{\mathrm{CE,mat}}(\mathfrak g_{\mathrm{dg}})}S.
\]
Any independently constructed adic CE formal spectrum is compared with this
object only after its Artinian restriction has been identified.  On the
separately stated finite-dimensional nilpotent classical range, the external
classical CE/BCH theorem identifies the represented CE loop with the BCH
formal group.  No raw native/CE/BCH equivalence is inferred without native
rigidity.

For ordinary smooth characteristic-zero Lie algebroids, native integration
supplies a canonical formal groupoid without solving or presupposing the
separate global smooth integrability problem.  For action algebroids, native
integration is unconditional, while comparison with an action-groupoid
presentation is made only when the nonlinear action has independently been
constructed.

The final scope can therefore be recorded as
\[
\boxed{
\begin{aligned}
\text{spectral partition-Lie algebroids}
&\text{ admit canonical native fixed-object formal integrations},\\
\text{native integration}
&\text{ differentiates back identically},\\
\text{pro-reduction and deformation localization}
&\text{ are left exact at the levels used in materialization},\\
\text{Artinian centres}
&\text{ are Cech-effective after deformation localization},\\
\text{formal-moduli stacks}
&\text{ are presentation-free two-stage materializations},\\
\text{formalized spectral-etale transport}
&\text{ uses completion without raw atlas invariance},\\
\text{arbitrary native groupoids}
&\text{ compare to their reconstructed controllers through a canonical span},\\
\text{native mapping discrepancies}
&\text{ isolate the remaining raw rigidity problem},\\
\text{no admissibility condition}
&\text{ is introduced to force those discrepancies to vanish}.
\end{aligned}}
\]
